% Vorlage fuer Abschlussarbeiten am Fachbereich Mathematik der TU Darmstadt.
% Geeignet fuer Bachelorarbeiten, Masterarbeiten und Diplomarbeiten.

%\documentclass[11pt,a4paper,twoside,openright]{report}
\documentclass[12pt,a4paper,twoside]{report}
% Benutzt man die Option draft, so kann man die Umbrueche ueberpruefen.

% Hier werden alle benoetigten Pakete und Einstellungen geladen. Auch
% hier sind Sie frei diese direkt in der Praeambel zu laden.


\usepackage[ngerman]{babel}

% Kodierung fuer unixoide Systeme und Windowssysteme.

\usepackage[utf8]{inputenc}

% Schriftart Times New Roman.

%\usepackage{lmodern}
%\usepackage{times}
%\usepackage{newtxtext,newtxmath}
\usepackage{mathptmx}

% Erweitert den Zeichenvorrat, so dass z.B. auch Umlaute im PDF-Dokument
% gefunden werden.

\usepackage[T1]{fontenc}

%Bibliography Fehlerbehebung
\usepackage{bibgerm}

% Zur Einbindung von Bildern.

\usepackage{graphicx}

% Erweiterte enumerate-Umgebung.

\usepackage{enumerate}


% Verschiedene Pakete, die nuetzlich sind um Mathematik in LaTeX zu setzen.

\usepackage{mathtools}
%\mathtoolsset{showonlyrefs=false}
\usepackage{amsmath, amssymb, amsthm, dsfont}
\usepackage{mathrsfs}
\usepackage{listings}
\usepackage{xcolor}
\definecolor{dunkelgruen}{HTML}{0B6121}
\lstset{
language=R,                     % the language of the code
  basicstyle=\tiny, % the size of the fonts that are used for the code
  numbers=left,                   % where to put the line-numbers
  numberstyle=\tiny\color{gray},  % the style that is used for the line-numbers
  stepnumber=1,                   % the step between two line-numbers. If it's 1, each line
                                  % will be numbered
  numbersep=5pt,                  % how far the line-numbers are from the code
 % backgroundcolor=\color{white},  % choose the background color. You must add \usepackage{color}
  showspaces=false,               % show spaces adding particular underscores
  showstringspaces=false,         % underline spaces within strings
  showtabs=false,                 % show tabs within strings adding particular underscores
%  frame=single,                   % adds a frame around the code
%  rulecolor=\color{black},        % if not set, the frame-color may be changed on line-breaks within not-black text (e.g. commens (green here))
  tabsize=2,                      % sets default tabsize to 2 spaces
  captionpos=t,                   % sets the caption-position to bottom
    breaklines=true,                % sets automatic line breaking
  breakatwhitespace=false,        % sets if automatic breaks should only happen at whitespace
%  title=\lstname,                 % show the filename of files included with \lstinputlisting;
  otherkeywords={},                                % also try caption instead of title
  deletekeywords={off,Q,t,q,file,data},
  keywordstyle=\color{blue},      % keyword style
  commentstyle=\color{dunkelgruen},   % comment style
 stringstyle=\color{violet},      % string literal style
 %  escapeinside={\%*}{*)},         % if you want to add a comment within your code
%  morekeywords={*,...},           % if you want to add more keywords to the set
%alsoletter={!},  
  literate={ö}{{\"o}}1
           {ä}{{\"a}}1
           {ü}{{\"u}}1
           {ß}{{\ss}}1
}

% Doppelte 1 (\mathbbm{1})

\usepackage{bbm}

% Indentation

%\usepackage{scrextend}

% Spaces for newcommand

%\usepackage{xspace}

% tikz
\usepackage{tikz}
%\usetikzlibrary{external}
%\tikzexternalize[prefix=figures/]
%\usetikzlibrary{arrows.meta}% arrow tip library
%\usetikzlibrary{bending}% better arrow head for bended lines

% Hyphenation of words that already contain a hyphen
% http://tex.stackexchange.com/questions/2706/adequate-hyphenation-of-words-already-containing-a-hyphen
% Example:
% \textsc{Alexandroff}\hyp{}\textsc{Urysohn}\hyp{}Kompaktheit

%\usepackage{hyphenat}


% Hyperlinks

%\usepackage{hyperref}
\usepackage[pdfauthor={Adrian Gabel},%
pdftitle={Bachelorarbeit - Bewertung Amerikanischer Optionen durch Lösen eines optimalen Stoppproblems},%
%pagebackref=true,%
%pdftex
]{hyperref}
%\hypersetup{ colorlinks   =  false}
\usepackage[a-1b]{pdfx}
%\usepackage{url}
%\def\UrlBreaks{\do\/\do-}

%\usepackage[
%%nonumberlist, %keine Seitenzahlen anzeigen
%toc,          %Einträge im Inhaltsverzeichnis
%section]      %im Inhaltsverzeichnis auf section-Ebene erscheinen
%{glossaries}

%%%%%%%%%%%%%%%%
% Seitenlayout %
%%%%%%%%%%%%%%%%

% DIV# gibt den Divisor für die Layoutberechnung an.
% Vergrößern des Divisors vergrößert den Textbereich.
% BCOR#cm gibt die Breite des Bundstegs an.
\usepackage[DIV14,BCOR2cm]{typearea}

% Abstand obere Blattkante zur Kopfzeile ist 2.54cm - 15mm
\setlength{\topmargin}{-15mm}

% Keine Einrueckung nach einem Absatz.

\parindent 0pt

% Abstand zwischen zwei Abs\"atzen.

%\parskip 12pt

% Zeilenabstand.

\linespread{1.25}

% Inhaltsverzeichnis erstellen.

\usepackage{pgfplots}

\pgfplotsset{compat=1.12}
\pgfplotsset{samples=50}

\usepgfplotslibrary{statistics}


\usepackage{multicol}

\usepackage{subcaption}
\usepackage{microtype}
\usepackage{parcolumns}
\usepackage{changepage}

\usepackage{fancyhdr}
\usepackage{booktabs}
\usepackage{caption}

\usepackage[utf8]{inputenc}
\usepackage{pgfplots}
\pgfplotsset{compat=newest}
\usepgfplotslibrary{groupplots}
\usepgfplotslibrary{dateplot}
\usepgfplotslibrary{patchplots}

\usepackage{filecontents}

\usetikzlibrary{matrix,chains,positioning,decorations.pathreplacing,arrows}
\usetikzlibrary{positioning,calc}
\tikzset{%
  every neuron/.style={circle,draw,minimum size=1cm},
  neuron missing/.style={draw=none,scale=4,text height=0.333cm,execute at begin node=\color{black}$\vdots$},
}

\usepackage{todonotes}

\usepackage{float}

% Hier werden Makros und Umgebungen eingebunden. Diese werden separat
% in der Datei befehle.tex definiert. Sie sind frei diese Befehle auch
% direkt in der Praeambel zu definieren.
% Makros

\newcommand{\F}{\mathscr{F}}
\newcommand{\E}{\mathbb{E}\,}
\newcommand{\e}{\mathrm{e}\,}
\newcommand{\T}{\mathscr{T}}


\newcommand{\N}{\mathbb{N}} % natuerliche Zahlen
\newcommand{\Z}{\mathbb{Z}} % ganze Zahlen
\newcommand{\Q}{\mathbb{Q}} % rationale Zahlen
\newcommand{\R}{\mathbb{R}} % reelle Zahlen
\newcommand{\K}{\mathbb{K}} % Körper
\newcommand{\co}{\mathrm{co}} % compact open
\newcommand{\W}{\mathcal{W}} % Weyl-Gruppe
\newcommand{\1}{\mathbbm{1}} % Körper
\newcommand{\id}{\mathrm{id}}
\newcommand{\sq}{\mathrm{sq}}
\newcommand{\mult}{\mathrm{mult}}
\newcommand{\ReLU}{\mathrm{ReLU}}
\newcommand{\net}{\mathrm{net}}
\newcommand{\bj}{\mathbf{j}}
\newcommand{\bi}{\mathbf{i}}
\newcommand{\bk}{\mathbf{k}}
\newcommand{\bA}{\mathbf{A}}
\newcommand{\bU}{\mathbf{U}}
\newcommand{\bY}{\mathbf{Y}}
\newcommand{\ba}{\mathbf{a}}
\newcommand{\bb}{\mathbf{b}}








%\newcommand{\C}{\mathcal{C}} % stetige Funktionen

\DeclareMathOperator{\Ret}{Re} % Realteil
%\newcommand{\PIT}{\textbf{PIT }} % Prime Ideal Theorem
\newcommand{\AC}{\textbf{AC}\xspace} % Prime Ideal Theorem
\newcommand{\UFT}{\textbf{UFT}\xspace} % Prime Ideal Theorem
\newcommand{\PIT}{\textbf{PIT}\xspace}
\newcommand{\ZF}{\textbf{ZF}\xspace}
\newcommand{\ZFC}{\textbf{ZFC}\xspace}
\newcommand{\CC}{\textbf{CC}\xspace}
\newcommand{\CCR}{\textbf{CC}($\R$)\xspace}
\newcommand{\PCC}{\textbf{PCC}\xspace}
\newcommand{\PCCR}{\textbf{PCC}($\R$)\xspace}

\newcommand{\grad}[1]{\gv{\nabla} #1} % for gradient
\let\divsymb=\div % rename builtin command \div to \divsymb
\renewcommand{\div}[1]{\mathrm{div\,} #1} % for divergence
\renewcommand{\d}[1]{\ensuremath\, {\operatorname{d}\!{#1}}}
\newcommand{\norm}[1]{\lVert #1 \rVert}

\newcommand{\alaoglu}{\textsc{Alaoglu}}
\newcommand{\banach}{\textsc{Banach}}
\newcommand{\cauchy}{\textsc{Cauchy}}
\newcommand{\euklid}{\textsc{Euklid}}
\newcommand{\fubini}{\textsc{Fubini}}
\newcommand{\hahn}{\textsc{Hahn}}
\newcommand{\hausdorff}{\textsc{Hausdorff}}
\newcommand{\helmholtz}{\textsc{Helmholtz}}
\newcommand{\hilbert}{\textsc{Hilbert}}
\newcommand{\hoelder}{\textsc{Hölder}}
\newcommand{\killing}{\textsc{Killing}}
\newcommand{\kondrachov}{\textsc{Kondrachov}}
\newcommand{\laplace}{\textsc{Laplace}}
\newcommand{\lebesgue}{\textsc{Lebesgue}}
\newcommand{\leibniz}{\textsc{Leibniz}}
\newcommand{\lie}{\textsc{Lie}}
\newcommand{\lipschitz}{\textsc{Lipschitz}}
\newcommand{\navier}{\textsc{Navier}}
\newcommand{\neumann}{\textsc{Neumann}}
\newcommand{\newton}{\textsc{Newton}}
\newcommand{\poincare}{\textsc{Poincaré}}
\newcommand{\rellich}{\textsc{Rellich}}
\newcommand{\riesz}{\textsc{Riesz}}
\newcommand{\schwartz}{\textsc{Schwartz}}
\newcommand{\sobolev}{\textsc{Sobolev}}
\newcommand{\stokes}{\textsc{Stokes}}
\newcommand{\stein}{\textsc{Stein}}
\newcommand{\urysohn}{\textsc{Urysohn}}

% Umgebungen für Definitionen, Sätze, usw.

%\newtheorem{defn}{Definition}[section]
%\newtheorem{ex}{Beispiel}[chapter]

\theoremstyle{plain}
\newtheorem{thm}{Satz}[chapter]%[section]
\newtheorem{mthm}{Hauptsatz}[chapter]%[section]
\newtheorem{lem}[thm]{Lemma}
\newtheorem{cor}[thm]{Korollar}
\newtheorem{prop}[thm]{Proposition}

\theoremstyle{definition}
\newtheorem{defn}[thm]{Definition}

\theoremstyle{remark}
\newtheorem*{bem}{Bemerkung}
\newtheorem*{beh}{Behauptung}
\newtheorem{bemnumber}[thm]{Bemerkung}

\def\Satzrefname{Satz}

\DeclareMathOperator{\spn}{span}
\DeclareMathOperator{\sign}{sign}
\DeclareMathOperator{\GL}{GL}
\DeclareMathOperator{\Sym}{Sym}
\DeclareMathOperator{\hgt}{ht}
\DeclareMathOperator{\dist}{dist}
\DeclareMathOperator{\loc}{loc}
\DeclareMathOperator{\supp}{supp}
\DeclareMathOperator{\diam}{diam}

%\renewcommand{\refname}{Literaturverzeichnis}
\addto\captionsngerman{%
\renewcommand*{\bibname}{Literaturverzeichnis}
}

\fancypagestyle{intro}{%
\fancyhf{} % clear all header and footer fields
\fancyfoot[LE,RO]{\thepage} % except the center
\renewcommand{\headrulewidth}{0.0pt}
\renewcommand{\footrulewidth}{0pt}
}

\fancypagestyle{plain}{%
\fancyhf{} % clear all header and footer fields
\fancyhead[LE,RO]{\thepage} % except the center
\renewcommand{\headrulewidth}{1.0pt}
\renewcommand{\footrulewidth}{0pt}
}

\fancypagestyle{MyStyle}{
\fancyhead{} % clear all header fields
\fancyhead[RE]{\normalfont \nouppercase{\leftmark}}
\fancyhead[LO]{\normalfont \nouppercase{\rightmark}}
\fancyfoot{} % clear all footer fields
\fancyhead[LE,RO]{\thepage}
\renewcommand{\headrulewidth}{1.0pt}
\renewcommand{\footrulewidth}{0.0pt}
%\renewcommand{\sectionmark}[1]{\markboth{{\normalfont{##1}}}{}}
\renewcommand{\sectionmark}[1]{\markboth{\thesection\quad {\normalfont{##1}}}{}}
}



\begin{document}

% Auf der Titelseite und im Inhaltsverzeichnis sollen keine
% Seitenzahlen dargestellt werden.
\pagestyle{empty}

% Hier wird die Titelseite eingebunden.
\begin{titlepage}
  \begin{center}
    \vspace{1cm}
    \includegraphics[width=0.5\linewidth]{TU_Darmstadt_Logo.pdf}
    \vspace{1cm}
    
    \large{Fachbereich Mathematik}
    \vspace{2.5cm}
    
    \large{Masterarbeit}
    \vspace{2cm}

    \Large{\textbf{Analyse der Konvergenzgeschwindigkeit eines einfach berechenbaren Neuronale-Netze-Regressionsschätzers}}
    
    \vspace*{3cm}    
    
		\large
                %\href{mailto:gabel@mathematik.tu-darmstadt.de}{Fabian Gabel}
                Adrian Gabel
    \vspace*{1.0cm}

    XX.03.2020 \\
    \vspace*{2cm}

    Betreuer: Prof. Dr. Michael Kohler

    \vspace*{.5cm}

    %Zweiter Gutachter: Name des zweiten Gutachters\\[2ex]
    \vspace*{\fill}
    %\tiny{Version vom \today}
  \end{center}
\end{titlepage}
\vspace*{\fill}

% Eidesstattliche Erklaerung.
\cleardoublepage 
% Keine Kopf-/Fußzeile auf dieser Seite
\thispagestyle{empty}

\vspace*{4cm}
\section*{Erklärung}\index{Erklärung}

Hiermit versichere ich, dass ich die vorliegende Arbeit selbstständig verfasst habe und alle benutzten Quellen einschließlich der Quellen aus dem Internet und alle sonstigen Hilfsmittel angegeben habe.\vspace{20pt}

\noindent
Darmstadt, den XX.03.2020\vspace{60pt}

% Unterschrift (handgeschrieben)

\noindent
Adrian Gabel




% Inhaltsverzeichnis erstellen.
\cleardoublepage
\tableofcontents

% Ab sofort werden Seitenzahlen in der Kopfzeile angezeigt.
\pagestyle{headings}

% Einbinden der Einleitung.
\chapter*{Einleitung}\index{Einleitung}
\addcontentsline{toc}{chapter}{Einleitung}\index{Einleitung}

Erfinder träumen schon lange davon, Maschinen zu schaffen, die denken. Dieser Wunsch geht zumindest auf die Zeit des antiken Griechenlands zurück. Die mythischen Figuren Pygmalion, Daedalus und Hephaestus können alle als legendäre Erfinder interpretiert werden, und Galatea, Talos und Pandora können alle als künstliches Leben betrachtet werden (siehe \cite{2005metamorphoses}, \cite{sparkes1996red}, \cite{hesiod1996works}).

Als programmierbare Computer zum ersten Mal konzipiert wurden, fragten sich die Menschen, ob solche Maschinen intelligent werden könnten, mehr als hundert Jahre bevor man sie baute. Heute ist die künstliche Intelligenz (\emph{KI}) ein blühendes Feld mit vielen praktischen Anwendungen und aktiven Forschungsthemen. 
Künstliche Intelligenz ist längst in unserem Alltag präsent und dringt in immer mehr Bereiche vor. Sprachassistenten etwa sind bereits als Helfer auf dem Smartphone, im Auto oder zu Hause Normalität geworden. Fortschritte im Bereich der KI beruhen vor allem auf der Verwendung \emph{neuronaler Netze}. Vergleichbar mit der Funktionsweise des menschlichen Gehirns verknüpfen sie mathematisch definierte Einheiten miteinander \cite{goodfellow2016deep}.

In vielen Anwendungen der KI geht es darum, aus einer Menge von Daten eine allgemeine Regel abzuleiten (\emph{maschinelles Lernen}). Maschinelles Lernen kann \emph{als Lernen einer Funktion} $f$ zusammengefasst werden, die Eingangsvariablen~$X$ auf Ausgangsvariablen~$Y$ abbildet, sodass die Abbildungsvorschrift~$f(X) = Y$ entsteht.

Ein Algorithmus lernt diese Zielabbildungsfunktion aus Trainingsdaten.
Die Form der Funktion ist unbekannt, sodass es unsere Aufgabe als Praktiker des maschinellen Lernens ist, verschiedene Algorithmen des maschinellen Lernens zu evaluieren und zu sehen, welcher die zugrunde liegende Funktion besser annähert.
Unterschiedliche Algorithmen machen unterschiedliche Annahmen über die Form der Funktion und die Art und Weise, wie sie gelernt werden kann.

Algorithmen, die keine Annahmen über die Form der Abbildungsfunktion treffen, werden als nichtparametrische Algorithmen des maschinellen Lernens bezeichnet. Indem sie keine Annahmen treffen, können sie jede beliebige Funktionsform aus den Trainingsdaten lernen.
Nichtparametrische Methoden sind gut, wenn viele Daten verfügbar sind und man sich nicht allzu sehr um die Auswahl der richtigen Funktionen kümmern will \cite[Seite~757]{AI}. Mathematisch führt dies zu einem Approximationsproblem. 

Im Kontext der KI wurden hierzu unter anderem neuronale Netze vorgeschlagen, die als universale Funktionsapproximatoren (\emph{Schätzer}) eingesetzt werden können, jedoch insbesondere bei vielen verdeckten Schichten schwer zu analysieren sind. 
Dies führt unter anderem dazu, dass eine große Lücke zwischen den Schätzungen, für die schöne Konvergenzergebnisse in der Theorie nachgewiesen wurden, und den Schätzungen, die in der Praxis verwendet werden können, entsteht.

Ziel dieser Arbeit ist es, die folgende Frage genauer zu betrachten: Wenn wir eine Regressionsschätzung des neuronalen Netzes theoretisch genau so definieren, wie sie in der Praxis umgesetzt wird, welches Konvergenzergebnis können wir dann für diese Schätzung vorweisen? 

Als Erstes werden wir in Kapitel~\ref{chap:1} Grundlagen zu neuronalen Netzen für den weiteren Verlauf der Arbeit sammeln.
Anschließend definieren wir in Kapitel~\ref{chap:2} eine neue, leicht zu implementierende Neuronale-Netze-Regressionsschätzung. Die Besonderheit des dem Schätzer zugrunde liegenden neuronalen Netzes besteht darin, dass die meisten Gewichte unabhängig von den Daten gewählt werden.
%die durch einige neuere Approximationsergebnisse für neuronale Netze motiviert sind. 
In Kapitel~\ref{chap:3} zeigen wir, dass wir für diesen Schätzer Konvergenzraten ableiten können, falls die Regressionsfunktion bestimmte Glattheitsvoraussetzungen erfüllt. 
%Mit diesem Hauptergebnis ist es nun möglich, einfach zu implementierende Regressionsverfahren für neuronale Netze zu definieren, die die gleiche Konvergenzrate wie lineare Regressionsschätzungen (z.B. Kernel- oder Spline-Schätzungen) erreichen, d.h. sie erreichen (bis zu einem logarithmischen Faktor) die optimale Minimax-Konvergenzrate $n-2p/(2p+d)$ im Falle einer ($p,C$)-glatten Regressionsfunktion, für jedes $p > 0$.
Abschließend wird in Kapitel~\ref{chap:4} eine Simulation durchgeführt, bei der die Leistung unseres vorgestellten Schätzers mit anderen verbreiteten Schätzern verglichen wird. Diese Arbeit orientiert sich an \cite{kohler19}.

\chapter{Grundlagen und Hilfsresultate}
\chapter{Konstruktion des Neuronal Netze Regresssionsschätzers}
\chapter{Anwendungsbeispiel}
\chapter*{Appendix}

% Einbinden der einzelnen Kapitel (zur Übersicht eine Datei je Kapitel)

%\chapter{Grundlagen und Hilfsresultate}
\label{chap:1}

Der Zweck dieses Kapitels ist es, grundlegende Definitionen zu sammeln, die in den folgenden Kapiteln verwendet werden. Weiterhin werden wir Hilfsresultate darstellen und beweisen welche wir vor allem für das Resultat der Konvergenzgeschwindigkeit des einfach berechenbaren Neuronale Netze Regressionsschätzer benötigt werden.

In dieser Arbeit behandeln wir Neuronale-Netze-Regressionsschätzer im Kontext der nichtparametrischen Regression mit zufälligem Design. Im Gegensatz zur parametrischen Regression ist bei der nichtparametrischen, die Bauart der schätzenden Funktion komplett unbekannt, was von Vorteil hat dass weniger Annahmen getroffen werden müssen, man aber dadurch noch mehr Daten benötigt um eine Funktion zu schätzen.

Bei der nichtparametrischen Regressionsschätzung ist seien $(X, Y), (X_1, Y_1), (X_2, Y_2), \dots$ u.i.v $\R^d \times \R$-wertige Zufallsvariablen mit $\E[Y^2] < \infty$. Zudem sei $m\colon \R^d \to R$ definiert durch $m(x) = \E[Y \mid X = x]$ die zugehörige Regressionsfunktion. Ausgehend von 
$$ (X_1, Y_1),\dots,(X_n, Y_n)$$ 
soll $m$ geschätzt werden.

Das Problem der Regressionsschätzung bei zufälligem Design lässt sich wie gefolgt erläutern. In Anwerndung ist üblicherweise die Verteilung von $(X, Y)$ unbekannt, daher kann $m(x) = \E[Y \mid X = x]$ nicht berechnet werden. Oft ist es aber möglich, Werte von $(X, Y)$ zu beobachten. Ziel ist es dann, daraus die Regressionsfunktion zu schätzen. Im Hinblick auf die Minimierung des $L_2$-Risikos sollte dabei der $L_2$-Fehler der Schätzfunktion möglichst klein sein. 

Für das $L_2$-Risiko einer beliebigen messbaren Funktion $f\colon \R^d \to \R$  gilt$\colon$
$$\E[|f(X) - Y|^2] = \E[|m(X) - Y|^2] + \int_{\R^d}|f(x) - m(x)|^2 \mathds{P}_X (dx),$$
d.h. der mittlere quadratische Vorhersagefehler einer Funktion ist darstellbar als Summe des $L_2$-Risikos der Regressionsfunktion (unvermeidbarer Fehler) und des $L_2$-Fehlers der entsteht aufgrund der Verwendung von $f$ anstelle von $m$ bei der Vorhersage bzw. Approximation des Wertes von Y.

Formal führt das daher auf folgende Problemstellung$\colon$
$(X, Y), (X_1, Y_1), (X_2, Y_2), \dots$ seien unabhängig identisch verteilte $\R^d \times \R$ wertige Zufallsvariablen mit $\E[Y^2] < \infty$ und $m\colon\R^d \to \R$ definiert durch $m(x) = \E[Y \mid X = x]$ sei die zugehörige Regressionsfunktion. Gegeben ist die Datenmenge 
$$ \mathcal{D}_n = \{(X_1, Y_1),\dots,(X_n, Y_n)\}.$$
Gesucht ist eine Schätzung 
$$m_n(\cdot) = m_n(\cdot, \mathcal{D}_n)\colon\R^d \to \R $$
von $m$, für die der $L_2$-Fehler 
$$\int |m_n(x) - m(x)|^2 \mathds{P}_X(dx)$$
möglichst \glqq klein\grqq ist. (Referenz Györfi (2002))

\section{Definitionen}
Es ist bekannt, dass man Glattheitsvoraussetzungen an die Regressionsfunktion haben muss um nichttriviale Konvergenzresultate für nichtparametrische Regressionsschätzer herzuleiten. Dafür verwenden wir die folgende Definition.
\begin{defn}[($p,C$)-Glattheit]
\label{def:pc}
   Sei $p = q + s$ mit $q \in \N_0$ und $s \in (0,1]$ (also $p \in (0, \infty)$ und sei $C > 0$. Eine Funktion $f\colon \R^d \to \R$ heißt ($p, C$)-glatt, falls für alle $\alpha = (\alpha_1,\dots,\alpha_d) \in \N_0^d$ mit $\sum_{j = 1}^{d}\alpha_j = q$ die partielle Ableitung 
   $$ \frac{\partial^qf}{\partial x_1^{\alpha_1}\dots\partial x_d^{\alpha_d}}$$
   existiert und falls für alle $x, z \in \R^d$ gilt
   $$ \bigg|\frac{\partial^qf}{\partial x_1^{\alpha_1}\dots\partial x_d^{\alpha_d}}(x) - \frac{\partial^qf}{\partial x_1^{\alpha_1}\dots\partial x_d^{\alpha_d}}(z) \bigg| \leq C \cdot \|x - z\|^r,$$
   wobei $\|\cdot\|$ die euklidische Norm ist.  
\end{defn}
\begin{bemnumber}
Im Falle von $p \leq 1$ ist keine Funktion ($p, C$)-glatt genau dann, wenn sie Hölder-stetig ist mit Exponent $p$ und Hölder-Konstante $C$.
\end{bemnumber}

Der Ausgangspunkt für die Definition eines neuronalen Netzes ist die Wahl einer Aktivierungsfunktion $\sigma\colon \R \to \R$. Wir haben uns in dieser Arbeit für die sogenannten \glqq squashing functions\grqq entschieden, welche eine monoton wachsend ist und für die $\lim_{x \to -\infty}\sigma(x) = 0$ und $\lim_{x \to \infty}\sigma(x) = 1$ gilt. Ein Beispiel für eine squashing function ist der sogenannte sigmoidal bzw. logisitsche squasher
\begin{align}
\label{logsquasher}
\sigma(x) = \frac{1}{1 + \exp(-x)} \quad (x \in \R).
\end{align}

\begin{defn}
\label{nzulässig}
Sei $n \in \N_0$. Eine Funktion $\sigma\colon \R \to [0, 1]$ wird $N$-zulässig genannt, wenn monoton wachsend und Lipschitz stetig (REFERENZ) ist und wenn zusätzlich die folgenden drei Bedingungen erfüllt sind$\colon$
\begin{itemize}
\item[(i)] Die Funktion $\sigma$ ist $N + 1$ mal stetig differenzierbar mit beschränkten Ableitungen.
\item[(ii)] Es existiert ein Punkt $t_{\sigma} \in \R$, in welchem alle Ableitungen bis hin zur $N$-ten Ableitung von $\sigma$ ungleich Null sind.
\item[(iii)] Wenn $y > 0$ ist, gilt $|\sigma(y) - 1| \leq \frac{1}{y}$. Wenn $y < 0$ ist, gilt $|\sigma(y)| \leq \frac{1}{|y|}$.
\end{itemize}  
\end{defn}  

In Lemma \ref{lem:logsquasher} werden wir zudem zeigen, dass der logistische squasher (\ref{logsquasher}) $N$-zulässig ist für beliebiges $N \in \N.$ 

\section{Hilfsresultate}
\begin{lem}
\label{lem:logsquasher}
Sei $N \in \N$ beliebig, dann erfüllt der logistische squasher $\sigma\colon \R \to [0, 1], \sigma(x) = \frac{1}{1 + \exp(-x)}$ die Bedingungen aus Definition \ref{logsquasher}.  
\end{lem}
\begin{proof}
Sei $N \in \N$ beliebig. Wir wissen, dass $\sigma$ monoton wachsend ist, da für beliebige $s, t \in \R$ mit $s \leq t$ gilt:
$$\sigma(s) = \frac{1}{1 + \exp(-s)} \leq \frac{1}{1 + \exp(-t)} = \sigma(t),$$
wobei wir bei der Ungleichung die Monotonie der Exponentialfunktion verwendet haben und die obige Ungleichung aus
\begin{equation*}
\begin{split}
& \quad \exp(s) \leq \exp(t) \\
 \Leftrightarrow & \quad \exp(-s) \geq \exp(-t) \\
 \Leftrightarrow & \quad 1 + \exp(-s) \geq 1 + \exp(-t) \\
 \Leftrightarrow & \quad \frac{1}{1 + \exp(-s)} \leq \frac{1}{1 + \exp(-t)} \\
\end{split}
\end{equation*}
folgt. Zudem ist $\sigma$ als Komposition $N + 1$ mal stetig differenzierbarer Funktionen selber auch $N + 1$ mal stetig differenzierbar. Die Ableitungen von $\sigma$ haben die Form$\colon$
\begin{equation*}
\begin{split}
\frac{\partial \sigma}{\partial x}(x) &= -\frac{1}{(1 + \exp(-x))^2} \cdot (-\exp(-x)) \\
& = \frac{\exp(-x)}{1 + \exp(-x)} \cdot \frac{1}{1 + \exp(-x)} \\
& = \bigg(1 - \frac{1}{1 + \exp(-x)}\bigg) \cdot \frac{1}{1 + \exp(-x)} \\
& = (1 - \sigma(x)) \cdot \sigma(x).
\end{split}
\end{equation*}
Da wir bei weiterem Ableiten die Produktregel wiederholt anwenden sind alle Ableitungen von $\sigma$, Polynome in $\sigma$. Dadurch folgt Bedingung (i) aus Definition \ref{logsquasher}, da $sigma$ nach Voraussetzung durch $0$ und $1$ beschränkt ist, und die Ableitungen von $\sigma$ als Produkt von beschränkten Faktoren daher auch. Da hiermit auch die erste Ableitung von $\sigma$ beschränkt ist wissen wir nach Satz ... aus (REFERENZ), dass $\sigma$ Lipschitz stetig ist.
Nun kommen wir zum Beweis von Bedingung (ii). Polynome, die nicht das $0$-Polynom sind, haben nach Satz ... (REFERENZ) auf $(0, 1)$ endlich viele Nullstellen und $\sigma$ bildet nach Voraussetzung in das Intervall $[0, 1] \supseteq (0, 1)$ ab. Da die Ableitungen von $\sigma$, als Zusammensetzung von Polynome in $\sigma$, wieder Polynome sind für die die obere Eigenschaft ebenfalls gilt, existiert ein $t_{\sigma} \in \R$ mit$\sigma(t_{\sigma}) \neq 0$ sodass alle Ableitungen bis zum Grad $N$ von $\sigma$, aufgrund ihrer Struktur ungleich $0$ sind. Daher ist Bedingung (ii) ebenfalls erfüllt.
Betrachten wir nun ein beliebiges $x > 0$. Dann wissen wir nach dem Mittelwertsatz (REFERENZ) dass ein $z \in (0, \infty)$ existiert, sodass mit $\exp(z) > 1$, da $z > 0$ ist, gilt$\colon$
\begin{equation*}
\begin{split}
(\exp(x - 1) \cdot (x - 0) = e^z
\end{split}
\end{equation*} 
und da $x \neq 0$ ist, daraus folgt, dass $$\frac{\exp(x) - 1}{x - 0} > 1$$ gilt und damit dann auch durch Multiplikation mit $x$ insbesondere $$ x \leq \exp(x) + 1.$$ Daraus erhalten wir mit Umformungen da $x > 0$ und $1 + \exp(-x) > 0$ ist$\colon$
\begin{equation*}
\begin{split}
& \quad x \leq \exp(x) + 1 \\
 \Leftrightarrow & \quad x \cdot \exp(-x) \leq 1 + \exp(-x) \\
 \Leftrightarrow & \quad \frac{\exp(-x)}{1 + \exp(-x)} \leq \frac{1}{x} \\
 \Leftrightarrow & \quad 1 - \frac{1}{1 + \exp(-x)} \leq \frac{1}{x} \\
 \Leftrightarrow & \quad |\sigma(x) - 1| \leq \frac{1}{x}.
\end{split}
\end{equation*}
Wobei die letzte Ungleichung aus der Eigenschaft des Betrags kommt, da $\frac{1}{1 + \exp(-x)} - 1< 0$ ist, weil $1 + \exp(-x) > 1$, da $\exp(-x) > 0$. Dies zeigt die erste Relation aus Bedingung (iii).
Die zweite Relation folgt durch die gleiche Art und Weise, da wir durch $$\frac{1}{1 + \exp(x)} - \frac{1}{2}= \sigma(0 - x) - \frac{1}{2} = -\sigma(0 + x) + \frac{1}{2} = -\frac{1}{1 + \exp(-x)} + \frac{1}{2}$$ 
wissen, dass $\sigma$ punktsymmetrisch in $(0, \frac{1}{2})$ ist. Die obige Gleichheit folgt aus
\begin{equation*}
\begin{split}
 & \quad \frac{1}{1 + \exp(x)} - \frac{1}{2} = -\frac{1}{1 + \exp(-x)} + \frac{1}{2} \\
 \Leftrightarrow  & \quad \frac{1}{1 + \exp(x)} + \frac{1}{2} -\frac{1}{1 + \exp(-x)} = \frac{1}{2}\\
 \Leftrightarrow & \quad  \frac{1 + \exp(-x) + 1 + \exp(x)}{(1 + \exp(x)) \cdot (1 + \exp(-x))} - \frac{1}{2} = \frac{1}{2} \\
 \Leftrightarrow & \quad \frac{2 + \exp(-x) + \exp(x)}{2 + \exp(-x) + \exp(x)} - \frac{1}{2} = \frac{1}{2} \\
 \Leftrightarrow & \quad 1 - \frac{1}{2} = \frac{1}{2}.
\end{split}
\end{equation*}
Aus dieser Eigenschaft folgt mit $$\sigma(- x) - 1 = \frac{1}{1 + \exp(x)} - 1 = -\frac{1}{1 + \exp(-x)} = -\sigma(x)$$ für $x < 0$ aus der ersten Relation, da $-x > 0 $ ist$\colon$  
$$|\sigma(x)| = |-\sigma(x)| = |\sigma(- x) - 1| \leq \frac{1}{-x} \leq \frac{1}{|x|}.$$
Damit haben wir alle drei Bedingungen aus Definition \ref{logsquasher} gezeigt und unsere Aussage bewiesen.
\end{proof}
\begin{lem}
  \label{lem:1}
  Sei $\sigma \colon \R \to \R$ eine Funktion und $R$, $a > 0$.
  \begin{itemize}
  \item[a)] Angenommen $\sigma$ ist zwei mal stetig differenzierbar und $t_{\sigma,id} \in \R$ so, dass $sigma'(t_{\sigma, id}) \neq 0$ ist. Dann gilt mit
  $$ f_{id}(x) = \frac{R}{\sigma'(t_{\sigma, id})} \cdot \left(\sigma\left(\frac{x}{R} + t_{\sigma, id}\right) - \sigma(t_{\sigma, id})\right)$$
  für beliebige $x \in [-a, a]\colon$ 
  $$ |f_{id}(x) - x| \leq \frac{\|\sigma''\|_{\infty} \cdot a^2}{2 \cdot |\sigma'(t_{\sigma, id})|} \cdot \frac{1}{R}.$$
  \item[b)] Angenommen $\sigma$ ist drei mal stetig differenzierbar und $t_{\sigma,sq} \in \R$ so, dass $sigma''(t_{\sigma, sq}) \neq 0$ ist. Dann gilt mit
  $$ f_{sq}(x) = \frac{R^2}{\sigma''(t_{\sigma, sq})} \cdot \left(\sigma\left(\frac{2 \cdot x}{R} + t_{\sigma, sq}\right) - 2 \cdot \sigma(\frac{x}{R} + t_{\sigma, sq})+ \sigma(t_{\sigma, sq})\right)$$
  für beliebige $x \in [-a, a]\colon$ 
  $$ |f_{sq}(x) - x^2| \leq \frac{5 \cdot \|\sigma'''\|_{\infty} \cdot a^3}{3 \cdot |\sigma'(t_{\sigma, sq})|} \cdot \frac{1}{R}.$$
  \end{itemize}
\end{lem}
\begin{proof}
	\begin{itemize}
  	\item[a)] Sei $u = \frac{c}{R} + t_{\sigma, id}$, $\xi = 0$ und  $x \in [-a, a]$beliebig. Wir wissen, dass $f_{id}$ 1-mal differenzierbar ist, da nach Vorraussetzung $\sigma$ 2-mal stetig differenzierbar ist,  existiert nach der Restgliedformel von Lagrange (REFERENZ) ein $c \in [\xi, x] $, sodass mit Ausklammern von $\frac{R}{\sigma'(t_{\sigma, id})} $ folgt$\colon
$

  	\begin{equation*}
  	\begin{split}
  	 |f_{id}(x) -  & x| \\
  	& \leq \bigg|\frac{R}{\sigma'(t_{\sigma, id})} \cdot \bigg(\sigma\left(\frac{\xi}{R} + t_{\sigma, id}\right) - \sigma(t_{\sigma, id}) + \frac{1}{R} \sigma'\left(\frac{\xi}{R} + t_{\sigma, id}\right) (x - \xi) \\ & \qquad + \frac{1}{2R^2} \sigma''(\frac{c}{R} + t_{\sigma, id}) (x - \xi)^2)\bigg) - x \bigg| \\
  	& = \bigg|\frac{R}{\sigma'(t_{\sigma, id})} \cdot \bigg(\sigma(t_{\sigma, id}) - \sigma(t_{\sigma, id}) + \frac{x}{R} \sigma'(t_{\sigma, id}) + \frac{x^2}{2R^2} \sigma''(\frac{c}{R} + t_{\sigma, id})\bigg) - x\bigg| \\
  	& = \bigg|\frac{R}{\sigma'(t_{\sigma, id})} \cdot \bigg(\frac{x}{R} \sigma'(t_{\sigma, id}) + \frac{x^2}{2R^2}\sigma''(u)\bigg) - x\bigg| \\
  	& = \bigg| \frac{\sigma''(u) \cdot x^2}{2R \cdot \sigma'(t_{\sigma, id})} + x - x\big| \\
  	& \leq \frac{\| \sigma'' \|_{\infty} \cdot a^2}{2 \cdot |\sigma'(t_{\sigma, id})|} \cdot \frac{1}{R},  
  	\end{split}
  	\end{equation*}
  	Wobei sich die letzte Ungleichung aus den Eigenschaften der Supremumsnorm ergibt und zudem aus $x \in [-a,a] \Leftrightarrow -a \leq x \leq a$ durch Quadrieren der Ungleichung folgt, dass $x^2 \leq a^2$ ist.
  	\item[b)] Folgt analog wie in a) durch 2-maliges Anwenden der Restgliedformel von Lagrange (REFeRENZ) auf die Funktion $f$ die hier nun 2-mal differenzierbar ist, da $\sigma$ nach Voraussetzung  3-mal stetig differenzierbar ist.
 	\end{itemize}
\end{proof}

\begin{lem}
  \label{lem:2}
  Sei $\sigma \colon \R \to [0, 1]$ nach Definition 1, 2-zulässig. Zudem sei $R > 0$ und $a > 0$ beliebig. Dann gilt für das neuronale Netz
  \begin{equation*}
  	\begin{split}
  	f_{mult}(x, y) = \frac{R^2}{4 \cdot \sigma''(t_{\sigma})} \cdot & \bigg(\sigma\Big(\frac{2 \cdot (x + y)}{R} + t_{\sigma} \Big) - 2 \cdot \sigma\Big(\frac{x + y}{R} + t_{\sigma}\Big) \\
  	& - \sigma \Big(\frac{2 \cdot (x - y)}{R} + t_{\sigma} \Big) + 2 \cdot \sigma \Big(\frac{x - y}{R} + t_{\sigma} \Big) \bigg)
  	\end{split}
  	\end{equation*}
  	für beliebige $x, y \in [-a, a]$ die folgende Ungleichung$\colon$
  	$$|f_{mult}(x, y) - x \cdot y| \leq \frac{20 \cdot \|\sigma'''\|_{\infty} \cdot a ^3}{3 \cdot |\sigma''(t_{\sigma})|} \cdot \frac{1}{R}.$$
  \end{lem}
  \begin{proof}
  Durch Ausmultiplizieren erhalten wir $$f_{mult}(x,y) = \frac{1}{4}(f_{sq}(x + y) - f_{sq}(x - y))$$ und $$x \cdot y = \frac{1}{4}\big((x + y)^2 - (x - y)^2\big).$$
  Aus diesen beiden Gleichungen folgt durch Ausklammern von $\frac{1}{4}$, der Homogenität des Betrags und der Anwendung der Dreickecksungleichung$\colon$
  \begin{equation*}
  \begin{split}
  |f_{mult}(x, y) - x \cdot y| & = \frac{1}{4} \cdot \big|f_{sq}(x + y) - f_{sq}(x - y) - (x + y)^2 + (x - y)^2\big| \\
  & \leq \frac{1}{4} \cdot \big|f_{sq}(x + y) - (x + y)^2\big| + \frac{1}{4}\cdot\big| (x - y)^2 - f_{sq}(x - y)\big| \\
  & \leq 2 \cdot \frac{1}{4} \cdot \frac{40 \cdot \|\sigma'''\|_{\infty} \cdot a^3}{3 \cdot |\sigma'(t_{\sigma, sq})|} \cdot \frac{1}{R} \\
  & = \frac{20 \cdot \|\sigma'''\|_{\infty} \cdot a ^3}{3 \cdot |\sigma''(t_{\sigma})|} \cdot \frac{1}{R},
  \end{split}
\end{equation*}   
wobei bei bei der letzten Ungleichung verwendet haben, dass $a > 0$ nach Lemma \ref{lem:1}b) beliebig gewählt wurde und daher insbesondere für beliebiges $x \in [-2a,2a]$ $$ |f_{sq}(x) - x^2| \leq \frac{40 \cdot \|\sigma'''\|_{\infty} \cdot a^3}{3 \cdot |\sigma'(t_{\sigma, sq})|} \cdot \frac{1}{R}$$ gilt. 
  \end{proof}
  \begin{lem}
  \label{lem:3}
  Sei $\sigma\colon \R \to [0, 1]$ nach Definition 1, 2-zulässig. Sei $f_{mult}$ das neuronale Netz aus Lemma \ref{lem:2} und $f_{id}$ das neuronale Netz aus Lemma \ref{lem:1}. Angenommen es gilt $$a \geq 1 \text{\quad und \quad} R \geq \frac{\|\sigma''\|_{\infty} \cdot a}{2 \cdot |\sigma'(t_{\sigma, id})|}.$$ 
  Dann erfüllt das neuronale Netz 
  \begin{equation*}
  \begin{split}
  f_{ReLU}(x) & = f_{mult}(f_{id}(x), \sigma(R \cdot x)) 
  \end{split}
  \end{equation*}
 für alle $x \in [-a, a]\colon$
 $$|f_{ReLU}(x) - \max\{x, 0\}| \leq 56 \cdot \frac{\max\{|\sigma''\|_{\infty}, \|\sigma'''\|_{\infty}, 1\}}{\min\{2 \cdot |\sigma'(t_{\sigma, id}), |\sigma''(t_{\sigma})|, 1\}} \cdot a^3 \cdot \frac{1}{R}.$$
  \end{lem}
  \begin{proof}
  Da $\sigma$ nach Voraussetzung 2-zulässig nach Definition 1 ist, gilt für $R \geq 0,$ und $x \in \R\setminus\{0\}\colon$ $$|\sigma(R \cdot x) - 1| \leq \frac{1}{R\cdot x} \text{\quad für \quad} x > 0$$ und $$|\sigma(R \cdot x)| \leq \frac{1}{|R \cdot x|} \text{\quad für \quad} x < 0.$$
Damit folgt aus der Homogenität des Betrags $$|\sigma(R \cdot x) - \1_{[0, \infty)}(x)| \leq \frac{1}{|R \cdot x|} = \frac{1}{R \cdot |x|}.$$ 
Nach Lemma \ref{lem:1} und Lemma \ref{lem:2} gilt$\colon$
$$ |f_{id}(x) - x| \leq \frac{\|\sigma''\|_{\infty} \cdot a^2}{2 \cdot |\sigma'(t_{\sigma, id})|} \cdot \frac{1}{R} \text{\quad für \quad} x \in [-a, a]$$ und 
$$ |f_{mult}(x, y) - x \cdot y| \leq \frac{160 \cdot \|\sigma'''\|_{\infty} \cdot a ^3}{3 \cdot |\sigma''(t_{\sigma})|} \cdot \frac{1}{R} \text{\quad für \quad} x \in [-2a, 2a].$$ 
Da nach Voraussetzung $a \geq 1$ ist gilt insbesondere $[0, 1] \in [-2a, 2a]$ und daher gilt insbesondere $\sigma(x) \in [-2a, 2a].$ 
Zudem erhalten wir durch eine Nulladdition, das Anwenden der Dreiecksungleichung, die Verwendung von Lemma \ref{lem:1} und der Voraussetzung für $R\colon$
\begin{equation*}
\begin{split}
|f_{id}(x)| & = |f_{id}(x) - x + x| \\
& = |f_{id}(x) -x| + |x| \\
& \leq |f_{id}(x) - x| \leq \frac{\|\sigma''\|_{\infty} \cdot a^2}{2 \cdot |\sigma'(t_{\sigma, id})|} \cdot \frac{1}{R} + |x| \\
& \leq |f_{id}(x) - x| \leq \frac{\|\sigma''\|_{\infty} \cdot a^2}{2 \cdot |\sigma'(t_{\sigma, id})|} \cdot \frac{2 \cdot |\sigma'(t_{\sigma, id})|}{\|\sigma''\|_{\infty} \cdot a} + |x| \\
& = a + |x| \\
& = 2 \cdot a
\end{split}
\end{equation*}
wobei $x \in [-a, a]$. Daraus folgt insbesondere $f_{id}(x) \in [-2a, 2a].$
Mithilfe von $\max\{x, 0 \} = x \cdot \1_{[0, \infty)}(x)$, der Voraussetzung, zweier Nulladdition und dem zweifachen Anwenden der Dreiecksungleichung erhalten wir$\colon$
\begin{equation*}
\begin{split}
|f_{ReLU}(x) - & \max\{x, 0\}|  \\
& = \big| f_{mult}(f_{id}(x), \sigma(R \cdot x)) - x \cdot \1_{[0, \infty)}(x)\big| \\
& \leq \big| f_{mult}(f_{id}(x), \sigma(R \cdot x)) - f_{id}(x)\cdot\sigma(R \cdot x)\big| \\
& \qquad + \big| f_{id}(x)\cdot\sigma(R \cdot x) - x \cdot \sigma(R \cdot x)\big| + \big| x \cdot \sigma(R \cdot x) - x \cdot \1_{[0, \infty)}(x)\big|. \\
& \text{Daraus ergibt sich mithilfe der obigen Eigenschaften und $a^3 \geq 1$}\\
& \leq \frac{160 \cdot \|\sigma'''\|_{\infty} \cdot a ^3}{3 \cdot |\sigma''(t_{\sigma})|} \cdot \frac{1}{R} + \frac{\|\sigma''\|_{\infty} \cdot a^3}{2 \cdot |\sigma'(t_{\sigma, id})|} \cdot \frac{1}{R} \cdot 1+ \frac{1}{R} \\
& \leq \bigg(\frac{160}{3} \cdot \frac{\|\sigma'''\|_{\infty} \cdot a ^3}{|\sigma''(t_{\sigma})|} + \frac{\|\sigma''\|_{\infty} \cdot a^3}{2 \cdot |\sigma'(t_{\sigma, id})|} + \frac{a^3}{a^3} \bigg) \cdot \frac{1}{R} \\ 
& \leq \bigg(\frac{160 \cdot\|\sigma'''\|_{\infty} \cdot a ^3 + 3 \cdot \|\sigma''\|_{\infty} \cdot a^3 + 3 \cdot a^3}{3 \cdot \min\{ 2 \cdot \sigma'(t_{\sigma, id})|, |\sigma''(t_{\sigma})|, 1\}}\bigg) \cdot \frac{1}{R}\\
& \leq \frac{166}{3} \cdot \bigg(\frac{\max\{\|\sigma'''\|_{\infty}, \|\sigma''\|_{\infty} , 1\}}{\min\{ 2 \cdot \sigma'(t_{\sigma, id})|, |\sigma''(t_{\sigma})|, 1\}}\bigg) \cdot a^3 \cdot  \frac{1}{R} \\
& \leq 56 \cdot \frac{\max\{|\sigma''\|_{\infty}, \|\sigma'''\|_{\infty}, 1\}}{\min\{2 \cdot |\sigma'(t_{\sigma, id}), |\sigma''(t_{\sigma})|, 1\}} \cdot a^3 \cdot \frac{1}{R}.
\end{split}
\end{equation*}
  \end{proof}
  \begin{lem}
  \label{lem:4}
  Sei $M \in \N$ und sei $\sigma\colon \R \to [0, 1]$ 2-zulässig nahc Definition .... Sei $a > 0$ und $$R \geq \frac{\|\sigma''\|_{}\infty \cdot (M + 1)}{2 \cdot |\sigma'(t_{\sigma, id})|},$$ sei $y \in [-a, a]$ und $f_{ReLU}$ das neuronale Netz aus Lemma \ref{lem:3}. Dann erfüllt das neuronale Netz 
  \begin{equation*}
  \begin{split}
  f_{hat,y}(x) = f_{ReLU}& \bigg(\frac{M}{2a} \cdot (x - y) + 1\bigg) - 2 \cdot f_{ReLU}\bigg(\frac{M}{2a} \cdot (x - y)\bigg) \\ &+ f_{ReLU}\bigg(\frac{M}{2a} \cdot (x - y) - 1\bigg)
  \end{split}
  \end{equation*}
  für alle $x \in [-a ,a]\colon$ 
  $$\bigg|f_{hat,y}(x) - \bigg(1 - \frac{M}{2a} \cdot |x - y|\bigg)_+\bigg| \leq 1792 \cdot \frac{\max\{|\sigma''\|_{\infty}, \|\sigma'''\|_{\infty}, 1\}}{\min\{2 \cdot |\sigma'(t_{\sigma, id}), |\sigma''(t_{\sigma})|, 1\}} \cdot M^3 \cdot \frac{1}{R}.$$
  \end{lem}
  \begin{proof}
  Als erstes zeigen wir $$(1 - \frac{M}{2a} \cdot |x|)_+ = \max\{\frac{M}{2a} \cdot x + 1, 0\} - 2 \cdot \max\{\frac{M}{2a} \cdot x, 0\} + \max\{\frac{M}{2a} \cdot x - 1, 0\}, \qquad (x \in \R)$$ damit wir das Resultat mithilfe von Lemma \ref{lem:3} beweisen können.
  Um die obige Gleichung zu zeigen unterscheiden wir vier Fälle.
  \begin{itemize}
  \item[Fall 1] ($x < - \frac{M}{2a}$) In diesem Fall hat die linke Seite durch $z_+ = \max\{z, 0\} \quad (z \in \R)$ und nach der Definition des Betrags die Gestalt $$\max\{1 + \frac{M}{2a} \cdot x, 0\}$$ und die rechte Seite die Form $$\max\{\frac{M}{2a} \cdot x + 1, 0\} - 2 \cdot 0 + 0,$$ da $x < 0$ und damit die letzten zwei Summanden 0 sind. Damit sind stimmt die rechte Seite mit der linken überein. $\hfill(\square)$
  \item[Fall 2] ($ - \frac{M}{2a} \leq x \leq 0$) Dieser Fall liefert aufgrund der nicht Positivität von $x$ analog das selbe Resultat wie Fall 1. $\hfill(\square)$
  \item[Fall 3] ($0 < x \leq \frac{M}{2a}$) In diesem Fall hat die linke Seite nach der Definition des Betrags die Gestalt $$\max\{1 - \frac{M}{2a} \cdot x, 0\}$$ und die rechte Seite die Form $$\max\{\frac{M}{2a} \cdot x + 1, 0\} - 2 \cdot \max\{\frac{M}{2a} \cdot x, 0\} + \max\{\frac{M}{2a} \cdot x - 1, 0\},$$ und erfordert daher eine weitere Fallunterscheidung.
  \item[Fall 3.1] ($x \cdot \frac{M}{2a} \leq 1$) In diesem Fall gilt für die linke Seite$\colon$ 
  $$\max\{1 - \frac{M}{2a} \cdot x, 0\} = 1 - \frac{M}{2a} \cdot x$$ und für die rechte Seite$\colon$ 
  $$\max\{\frac{M}{2a} \cdot x + 1, 0\} - 2 \cdot \max\{\frac{M}{2a} \cdot x, 0\} + 0 = \frac{M}{2a} \cdot x + 1 - 2 \cdot \frac{M}{2a} \cdot x = 1 - \frac{M}{2a} \cdot x,$$ und stimmt daher mit der linken Seite überein.
  \item[Fall 3.2] ($x \cdot \frac{M}{2a} > 1$) In diesem Fall gilt für die linke Seite$\colon$ 
  $$\max\{1 - \frac{M}{2a} \cdot x, 0\} = 0$$ und für die rechte Seite$\colon$ 
  $$\frac{M}{2a} \cdot x + 1 - 2 \cdot \frac{M}{2a} \cdot x + \frac{M}{2a} \cdot x - 1 = 0$$ und stimmt daher mit der linken Seite überein. Damit ist Fall 3 gezeigt. $\hfill(\square)$
  \item[Fall 4] ($\frac{M}{2a} < x$)  In diesem Fall hat die linke Seite nach der Definition des Betrags die Gestalt $$\max\{1 - \frac{M}{2a} \cdot x, 0\}$$ und die rechte Seite die Form $$\max\{\frac{M}{2a} \cdot x + 1, 0\} - 2 \cdot \max\{\frac{M}{2a} \cdot x, 0\} + \max\{\frac{M}{2a} \cdot x - 1, 0\},$$ und erfordert daher eine weitere Fallunterscheidung.
  \item[Fall 4.1] ($\frac{M}{2a} \cdot x \leq 1$) In diesem Fall gilt für die linke Seite$\colon$ 
  $$\max\{1 - \frac{M}{2a} \cdot x, 0\} = 1 - \frac{M}{2a} \cdot x$$ und für die rechte Seite$\colon$ 
  $$\frac{M}{2a} \cdot x + 1 - 2 \cdot \frac{M}{2a} \cdot x + 0 = 1 - \frac{M}{2a}$$ und stimmt daher mit der linken Seite überein.
  \item[Fall 4.2] ($\frac{M}{2a} \cdot x > 1$) In diesem Fall gilt für die linke Seite$\colon$ 
  $$\max\{1 - \frac{M}{2a} \cdot x, 0\} = 0$$ und für die rechte Seite$\colon$ 
  $$\frac{M}{2a} \cdot x + 1 - 2 \cdot \frac{M}{2a} \cdot x + \frac{M}{2a} \cdot x - 1 = 0$$ und stimmt daher mit der linken Seite überein. Damit ist Fall 4 gezeigt. $\hfill(\square)$
\end{itemize} 
Durch diese Fallunterscheidung wurde die obige Gleichung (REFERENZ) bewiesen. Daraus folgt mit der Definition von $f_{hat,y}(x)$ und zwei mal der Dreiecksungleichung
\begin{equation*}
\begin{split}
\bigg|f_{hat,y}(x) - \bigg(1 - \frac{M}{2a} \cdot |x - y|\bigg)_+\bigg| \leq \bigg|&f_{ReLU} \bigg(\frac{M}{2a} \cdot (x - y) + 1\bigg) - \max\{\frac{M}{2a} \cdot (x - y) + 1, 0\}\bigg| \\ 
& + 2 \cdot \bigg|f_{ReLU}\bigg(\frac{M}{2a} \cdot (x - y)\bigg) - \max\{\frac{M}{2a} \cdot (x - y), 0\}\bigg| \\
& + \bigg|f_{ReLU}\bigg(\frac{M}{2a} \cdot (x - y) - 1\bigg) - \max\{\frac{M}{2a} \cdot (x - y) - 1, 0\}\bigg|\\ 
& \leq 1792 \cdot \frac{\max\{|\sigma''\|_{\infty}, \|\sigma'''\|_{\infty}, 1\}}{\min\{2 \cdot |\sigma'(t_{\sigma, id}), |\sigma''(t_{\sigma})|, 1\}} \cdot M^3 \cdot \frac{1}{R},
\end{split}
\end{equation*} 
wobei die letzte Ungleichung daraus folgt, dass wir auf jeden Summanden mit $1 \leq a = M + 1$ Lemma \ref{lem:3} angewendet haben und die Abschätzung 
$$ (M + 1)^3 = M^3 + 3 \cdot M^2 + 3 \cdot M + 1 \leq M^3 + 3 \cdot M^3 + 3 \cdot M ^3 + M^3 = 8 \cdot M^3 \quad (M \in \N)$$ verwendet haben.
  \end{proof}
Die nächsten Lemmata benötigen wir für den Beweis unseres Hauptresultats, einer Aussage über die Konvergenzgeschwindigkeit unseres neuronale Netze Schätzers.. Diese Lemmata werden hier nur der Vollständigkeit halber und ohne Beweis aufgeführt.
\begin{lem}
\label{lem:5}
Sei $M \in \N$ und $\sigma\colon \R \to [0, 1]$ $2$-zulässig nach Definition \ref{nzulässig}. (AUF SCHREIBWEISE VON FNET AUFPASSEN)
Sei $a \geq 1$ und 
\begin{equation*}
\begin{split}
R & \geq \max\biggl\{\frac{\|\sigma''\|_{\infty} \cdot (M + 1)}{2 \cdot |\sigma'(t_{\sigma, id})|}, \frac{9 \cdot \|\sigma''\|_{\infty} \cdot a}{|\sigma'(t_{\sigma, id}|}, \\
& \quad \frac{20 \cdot \|\sigma'''\|_{\infty}}{3 \cdot |\sigma''(t_{\sigma})|} \cdot 3^{3 \cdot 3^s} \cdot a^{3 \cdot 2^s}, 1792 \cdot \frac{\max\{\|\sigma''\|_{\infty},\|\sigma'''\|_{\infty}, 1\}}{\min\{2 \cdot |\sigma'(t_{\sigma, id})|, |\sigma''(t_{\sigma})|, 1\}} \cdot M^3 \biggr\}
\end{split}
\end{equation*}
und sei $y \in [-a, a]^d.$ Sei $N \in \N$ und $j_1,\dots,j_d \in \N_0$ so, dass $j_1 + \cdots + j_d \leq N$ gilt und wir setzen $s = \lceil\log_2(N + d)\rceil$. Sei $f_{id}, f_{mult}$ und $f_{hat, z}$ (für $z \in \R$) die neuronalen Netze wir in Lemma \ref{lem:1}, Lemma \ref{lem:2} und Lemma \ref{lem:4}. Wir definieren das Netz $f_{net,j_1,\dots,j_d,y}$ durch:
\begin{align*}
\label{fnet}
& f_{net,j_1,\dots,j_d,y}(x) = f_1^{(0)}(x). \\
\intertext{wobei} \\
& f_k^{(l)}(x) = f_{mult}\Big(f_{2k - 1}^{(l + 1)}(x),f_{2k}^{(l + 1)}(x)\Big) \\
\intertext{für $k \in \{1, 2, \dots, 2^l \}$ und $l \in\{0,\dots,s - 1\}, $ und} \\
& f_k^{(s)}(x) = f_{id}(f_{id}(x^{(l)} - y^{(l)}))  \\
\intertext{für $j_1 + j_2 + \dots + j_{l-1} + 1 \leq k \leq j_1 + j_2 + \dots + j_l$ und $l = 1,\dots,d$ und} \\ 
& f_{j_1 + j_2 + \dots + j_d + k}^{(s)}(x) = f_{hat,y^{(k)}}(x^{(k)}) \\
\intertext{für $k = 1,\dots,d$ und} \\
& f_k^{(s)}(x) = 1 \\
\intertext{für $k = j_1 + j_2 + \dots + j_d + d + 1,j_1 + j_2 + \dots + j_d + d + 2,\dots,2^s.$}
\end{align*} 
Dann erhalten wir für $x \in [-a, a]^d$:
\begin{equation*}
\begin{split}
& \bigg|f_{net, y}(x) - (x^{(1)} - y^{(1)})^{j_1} \cdots (x^{(d)} - y^{(d)})^{j_d} \prod_{j = 1}^d (1 - \frac{M}{2a} \cdot |x^{(j)} - y^{(j)}|)_+\bigg| \\
& \leq c_{12} \cdot 3^{3 \cdot 3^s} \cdot a^{3 \cdot 2^s} \cdot M^3 \cdot \frac{1}{R}.
\end{split}
\end{equation*}
\end{lem}
  \begin{lem}
  \label{lem:8}
Sei $\beta_n = c_6 \cdot \log(n)$ für eine hinreichend große Konstante $c_6 > 0$. Angenommen die Verteilung von $(X, Y)$ erfüllt 
$$ \E\Big(\mathrm{e}^{c_4 \cdot |Y|^2}\Big) < \infty$$
für eine Konstante $c_4 > 0$ und dass der Betrag der Regressionsfunktion $m$ beschränkt ist. Sei $\mathcal{F}_n$ eine Menge von Funktionen $f\colon \R^d \to \R$ und wir nehmen an, dass der Schätzer $m_n$ 
$$m_n = T_{\beta_n}\tilde{m}_n$$ 
erfüllt, mit 
$$\tilde{m}_n(\cdot) = \tilde{m}_n(\cdot,(X_1, Y_1),\dots,(X_n, Y_n)) \in \mathcal{F}_n$$
und 
$$\frac{1}{n} \sum_{i = 1}^n |Y_i - \tilde{m}_n(X_i)|^2 \leq \min_{l \in \Theta_n}\bigg(\frac{1}{n}\sum_{i = 1}^n |Y_i - g_{n,l}(X_i)|^2 + pen_n(g_n,l)\bigg)$$
mit einer nichtleeren Parametermenge $\Theta_n$, zufällige Funktionen $g_{n,l}\colon \R^d \to \R$ und deterministischen penalty Termen $pen_n(g_{n,l}) \geq 0$, wobei die zufälligen Funktionen $g_{n,l}\colon \R^d \to \R$ nur von den Zufallsvariablen
$$\mathbf{b}_1^{(1)},\dots,\mathbf{b}_r^{(1)},\dots,\mathbf{b}_1^{(I_n)},\dots,\mathbf{b}_r^{(I_n)},$$
abhängen, die unabhängig von $(X_1, Y_1), (X_2, Y_2),\dots$ sind.
Dann erfüllt $m_n\colon$
\begin{equation*}
\begin{split}
& \E \int |m_n(x) - m(x)|^2 \mathds{P}_X(dx) \\
& \leq \frac{c_{13} \cdot \log(n)^2 \cdot \big(\log\big(\sup_{x_1^n \in (supp(X))^n}\mathcal{N}_1\big(\frac{1}{n \cdot \beta_n},\mathcal{F}_n,x_1^n\big)\big) + 1\big)}{n} \\
& \quad + 2 \cdot \E\bigg(\min_{l \in \Theta_n} \int |g_{n,l}(x) - m(x)|^2 \mathds{P}_X(dx) + pen_n(g_{n,l})\bigg),
\end{split}
\end{equation*}
für $n > 1$ und einer Konstante $c_{13} > 0$ welche nicht von $n$ abhängt. (DEFINITION VON LP-e-ÜBERDECKUNGSZAHLEN)
  \end{lem}
Das nächste Lemma benötigen wir um eine Schranke für die Überdeckungszahl $\mathcal{N}_1\big(\frac{1}{n \cdot \beta_n},\mathcal{F}_n,x_1^n\big)$ zu finden.
\begin{lem}
\label{lem:9}
Sei $a > 0$ und $d, N, J_n \in \N$ so, dass $J_n \leq n^{c_{14}}$ und setze $\beta_n = c_6 \cdot \log(n).$
Sei $\sigma$ 2-zulässig nach Definition \ref{nzulässig}. Sei $\mathcal{F}$ die Menge aller Funktionen die durch (\ref{networkarch}) definiert sind mit $k_1 = k_2 = \cdots = k_L = 24 \cdot (N + d)$ und dass der Betrag der Gewichte durch $c_{15} \cdot n^{c_{16}}$ beschränkt ist. Sei
$$ \mathcal{F}^{(J_n)} = \biggl\{\sum_{j = 1}^{J_n} a_j \cdot f_j : f_j \in \mathcal{F} \quad \text{und} \quad \sum_{j = 1}^{J_n} a_j^2 \leq c_{17} \cdot n^{c_{18}}\biggr\}.$$
Dann gilt für $n > 1:$
$$\log\bigg(supp_{x_1^n\in[-a,a]^{d \cdot n}} \mathcal{N}_1\bigg(\frac{1}{n \cdot \beta_n}, \mathcal{F}^{(J_n)},x_1^n\bigg)\bigg) \leq c_{19} \cdot \log(n) \cdot J_n,$$
für eine Konstante $c_{19}$ die nur von $L, N, a$ und $d$ abhängt.
\end{lem}
%\chapter{Konstruktion des Neuronale-Netze-Schätzers}
\label{chap:2}

In diesem Kapitel werden wir mithilfe der uns gegebenen Datenmenge 
\begin{equation}
\label{dataset}
\mathcal{D}_n = \{(X_1, Y_1),\dots,(X_n, Y_n)\},
\end{equation}
unseren Regressionsschätzer $\tilde{m}_n$ konstruieren. 
In Kapitel~\ref{chap:1} haben wir bereits in Definition~\ref{def:nn} vorgestellt, was wir unter einem mehrschichtigen feedforward neuronalen Netz verstehen.

Für die Konstruktion unseres Neuronale-Netze-Schätzers wählen wir den logistischen Squasher~(\ref{logsquasher}) als Aktivierungsfunktion, verwenden die gegebene Datenmenge $\mathcal{D}_n$ und wählen die Gewichte des neuronalen Netzes so, dass die resultierende Funktion aus Gleichung~(\ref{networkarch}) eine gute Schätzung für die Regressionsfunktion $m$ ist. Dafür wählen wir die Gewichte bis auf die in der Ausgabeschicht fest und schätzen die Gewichte in der Ausgabeschicht, indem wir mit unserer Datenmenge~(\ref{dataset}) ein regularisiertes Kleinste-Quadrate-Problem (REFERENZ) lösen.

\section{Definition der Netzwerkarchitektur}
Sei $a > 0$ fest und $M \in \N$. Zunächst fixieren wir die Multiindexnotation, die wir aufgrund der Übersichtlichkeit im weiteren Verlauf dieser Arbeit verwenden werden.

Sei $[M]^d \coloneqq\{0, 1, \dots, M\}^d.$ 
Für $(\mathbf{i}^{(1)},\dots,\mathbf{i}^{(d)}) = \bi \in [M]^d$ und $x \in \R^d$ definieren wir
$$|\mathbf{i}|_1 \coloneqq \sum_{k= 1}^d \mathbf{i}^{(k)} \quad \mathbf{i}! \coloneqq \mathbf{i}^{(1)}! \cdots \, \mathbf{i}^{(d)}! \quad x^{\mathbf{i}} \coloneqq x_1^{\mathbf{i}^{(1)}} \cdots \,    x_d^{\mathbf{i}^{(d)}}.$$
Für $f\colon \R^d \to \R$ ausreichend oft differenzierbar definieren wir 
$$\partial^{\mathbf{i}}f(x) \coloneqq \frac{\partial^{|\mathbf{i}|_1}f}{\partial^{\mathbf{i}^{(1)}} x_1 \cdots \, \partial^{\mathbf{i}^{(d)}} x_d} (x).$$
Wir betrachten im Folgenden ein $d$-dimensionales äquidistantes Gitter im Würfel $[-a, a]^d$ mit Schrittweite $\frac{2a}{M}.$ 
%Sei $\bi_1,\dots,\bi_{(M + 1)^d}$ eine Aufzählung der Elemente von $[M]^d$.
Dann ordnen wir jedem Multiindex $\bi \in [M]^d$ einen Gitterpunkt $x_{\bi}$ mit:
\begin{equation}
\label{eq:gitter}
x_{\bi} = \bigg( -a + \bi^{(1)} \cdot \frac{2a}{M},\dots, -a + \bi^{(d)} \cdot \frac{2a}{M}\bigg) = -\mathbf{a} + \frac{2a}{M} \cdot \bi,
\end{equation}
zu, mit $\mathbf{a} = (a, a, \dots, a) \in \R^d$.

Hiermit lässt sich das zu $m$ gehörige Taylorpolynom der Ordnung $q$ mit Entwicklungspunkt $x_{\bi}$ schreiben als
$$p_{\bi}^m(x) = \sum_{\substack{\mathbf{j} \in [q]^d \\|\mathbf{j}|_1 \leq q}} \partial^{\mathbf{j}}m(x_{\mathbf{i}}) \cdot \frac{(x - x_{\mathbf{i}})^{\mathbf{j}}}{\mathbf{j}!}.$$
Zudem betrachten wir eine Funktion
\begin{equation}
\label{konvexkomb}
P_m(x) = \sum_{\bi \in [M]^d} p_{\bi}^m(x) \prod_{j = 1}^{d} \bigg(1 - \frac{M}{2a} \cdot |x^{(j)} - x_{\mathbf{i}}^{(j)}|\bigg)_+,
\end{equation}
für die wir im weiteren Verlauf dieser Arbeit zeigen werden, dass diese die Regeressionsfunktion $m$ approximiert.

Wir zeigen mit folgenden Lemma dass $P_m(x)$ eine lokale Spline Interpolation von Taylorpolynomen von $m$ ist.
\begin{lem}
\label{lem:loccon}
Sei $a >0$ und $M \in \N$. Dann sind für $\bi \in [M]^d$ und $x \in \R^d$ die Funktionen 
$$B_{\bi}(x) = \prod_{j = 1}^{d} \bigg(1 - \frac{M}{2a} \cdot |x^{(j)} - x_{\mathbf{i}}^{(j)}|\bigg)_+,$$
B-Splines im $R^d$, wobei wir unter $x_{\bi}$ die Gitterpunkte aus Gleichung~(\ref{eq:gitter}) verstehen. Dafür sind folgende drei Bedingungen für $x \in \R^d$ zu überprüfen:
\begin{itemize}
\item[i)] Zerlegung der Eins: $\sum_{\bi \in [M]^d} B_{\bi}(x) = 1$.
\item[ii)] Nicht-Negativität: $B_{\bi}(x) \geq 0$ für alle $\bi \in [M]^d$.
\item[iii)] Lokaler Träger: Für $\bi \in [M]^d$ ist $B_{\bi}(x) > 0$ falls $|x^{(j)} - x_{\bi}^{(j)}| < \frac{2a}{M}$ für alle $j \in \{1,\dots,d\}$ gilt und andernfalls $B_{\bi}(x) = 0$.
\end{itemize}
\end{lem}
\begin{proof}
Als erstes möchten wir für $d = 2$ und $M = 3$ eine Skizze angeben um die Idee des Beweises zu veranschaulichen. 
\begin{figure}[htp]
\centering
\begin{tikzpicture} 
   %Raster zeichnen 
   %\draw [color=gray!50]  [step=20mm] (-3,-3) grid (4,4); 
   \draw [color=gray!50] (0,0) -- (6,0) -- (6,6) -- (0,6) -- (0,0);			
   \draw [color=gray!50] (0,2) -- (6,2);
   \draw [color=gray!50] (0,4) -- (6,4);
        \draw [color=gray!50] (2,6) -- (2,0);
                \draw [color=gray!50] (4,6) -- (4,0);

	\fill[black] (0,0) circle (0.08cm) node[label=below:{$(-a, -a)$}]{};
		\fill[black] (6,0) circle (0.08cm) node[label=below:{$(a, -a)$}]{};
				\fill[black] (0,6) circle (0.08cm) node[label=above:{$(-a, a)$}]{};
						\fill[black] (6,6) circle (0.08cm) node[label=above:{$(a, a)$}]{};
						\fill[black] (0,0) circle (0.00cm) node[label=left:{$x_{\mathbf{i}_1}$}]{};
						\fill[black] (2,0) circle (0.00cm) node[label=below:{$x_{\mathbf{i}_2}$}]{};
						\fill[black] (4,0) circle (0.00cm) node[label=below:{$x_{\mathbf{i}_3}$}]{};
						\fill[black] (6,0) circle (0.00cm) node[label=right:{$x_{\mathbf{i}_4}$}]{};

						\fill[black] (9,6) circle (0.00cm) node[label=right:{$\mathbf{i}_1 = (0, 0)$}]{};			
						\fill[black] (9,5.5) circle (0.00cm) node[label=right:{$\mathbf{i}_2 = (1, 0)$}]{};					
						\fill[black] (9,5) circle (0.00cm) node[label=right:{$\mathbf{i}_3 = (2, 0)$}]{};	
						\fill[black] (9,4.5) circle (0.00cm) node[label=right:{$\mathbf{i}_4 = (3, 0)$}]{};	
						\draw [dotted, ultra thick] (10,4) -- (10,3.65);
						\fill[black] (9,3.15) circle (0.00cm) node[label=right:{$\mathbf{i}_{16} = (3, 3)$}]{};	
						
						\fill[black] (0,2) circle (0.08cm);
						\fill[black] (0,4) circle (0.08cm); 
						\fill[black] (2,2) circle (0.08cm);
						\fill[black] (2,4) circle (0.08cm);
						\fill[black] (2,6) circle (0.08cm);
						\fill[black] (2,0) circle (0.08cm);
						\fill[black] (0,4) circle (0.08cm);
						\fill[black] (2,4) circle (0.08cm);
						\fill[black] (4,4) circle (0.08cm);
						\fill[black] (4,6) circle (0.08cm);																\fill[black] (4,0) circle (0.08cm);
						\fill[black] (4,2) circle (0.08cm);						
						\fill[black] (6,2) circle (0.08cm);
						\fill[black] (6,4) circle (0.08cm);
\end{tikzpicture}
\caption{Beispielhafte Darstellung der $x_{\mathbf{i}_k}$ für $d = 2$ und $M = 3$.}
\label{fig:gitter}
\end{figure}
Es ist ein Gitter mit $(M + 1)^d$ Gitterpunkten die den $x_{\mathbf{i}_k}$ entsprechen. Der Abstand zwischen zwei Gitterpunkten beträgt $\frac{2a}{M}.$ Man betrachtet immer den Abstand zu den nächsten $2^d$ Gitterpunkten, da  $(1 - \frac{M}{2a} \cdot |x^{(j)} - x_{\mathbf{i}}^{(j)}|)_+ = 0$ immer dann gilt, wenn der Abstand zwischen $x^{(j)}$ und $x_{\mathbf{i}}^{(j)}$ größer als $\frac{2a}{M}$ ist.     

\textbf{i)} Im Folgenden wollen wir 
\begin{equation}
\label{induktiongl1}
\sum_{\bi \in [M]^d} B_{\bi}(x) = 1 \quad (x \in \R^d)
\end{equation}
per Induktion über $d$ zeigen. 
	
	 \emph{Induktionsanfang} (IA): Für $d = 1$ kann $x$ nur zwischen zwei Gitterpunkten $x_{\mathbf{i}_1} \neq x_{\mathbf{i}_2}$ liegen. Sei oBdA $x_{\mathbf{i}_1}\leq x \leq x_{\mathbf{i}_2}$, dann gilt mit der gleichen Begründung wie im einleitenden Beispiel:
	\begin{equation*}
	\begin{split}
	\sum_{\bi \in [M]^d} (1 - \frac{M}{2a} \cdot |x - x_{\mathbf{i}}|)_+ & = (1 - \frac{M}{2a} \cdot |x - x_{\mathbf{i}_1}|)_+ + (1 - \frac{M}{2a} \cdot |x - x_{\mathbf{i}_2}|)_+ \\
	& = 1 + 1 - \frac{M}{2a} \cdot (x - x_{\mathbf{i}_1} + x_{\mathbf{i}_2} - x) \\
	& = 1 + 1 - \frac{M}{2a} \cdot \frac{2a}{M} \\
	& = 1,
	\end{split}
	\end{equation*} wobei wir unter anderem verwendet haben, dass beide Summanden unabhängig von dem Positivteil nichtnegativ sind, da der Abstand von $x$ zu den beiden Gitterpunkten $x_{\mathbf{i}_1}$ und$ x_{\mathbf{i}_2}$ kleiner gleich $\frac{2a}{M}$ ist. Zudem haben wir verwendet, dass $x_{\mathbf{i}_2} - x_{\mathbf{i}_1} = \frac{2a}{M}$ gilt, da beides Gitterpunkte sind.     
	
\emph{Induktionshypothese} (IH): Aussage~(\ref{induktiongl1}) gelte für ein beliebiges aber festes $d \in \N.$

\emph{Induktionsschritt} (IS): Wir nehmen oBdA an, dass $x_{(0,\dots,0)} \leq x \leq x_{(1,\dots,1)}$ gilt bzgl.\@ $(d+1)$-stelliger Relation. Das heißt also, dass $x \in [-a, -a + \frac{2a}{M}]^{d + 1}$ gilt. Im Folgenden zeigen wir $$\sum_{\bi \in [M]^{(d + 1)}} B_{\bi}(x) = \sum_{\bi \in [M]^{(d + 1)}}\prod_{j = 1}^{d + 1} \bigg(1 - \frac{M}{2a} \cdot \big|x^{(j)} - x_{\mathbf{i}}^{(j)}\big|\bigg)_+ = 1.$$
Ein Summand ist Null, wenn ein $j \in \{1,\dots,d+1\}$ existiert mit $\big|x^{(j)} - x_{\mathbf{i}}^{(j)}\big| \geq \frac{2a}{M}$. Zudem haben wir oBdA angenommen dass $x \in [-a, -a + \frac{2a}{M}]^{d + 1}$ gilt, damit haben wir also nur noch $2^{d + 1}$ Summanden, was der Anzahl der Gitterpunkte die am nächsten bei $x$ liegen entspricht. Zudem wissen wir, dass alle Gitterpunkte, die in der $(d + 1)$-ten Komponente den selben Wert haben, in dieser Dimension gleich weit von $x^{(d + 1)}$ entfernt sind. Das heißt, in jedem Summanden kommt der Faktor $(1 - \frac{M}{2a} \cdot \big|x^{(d + 1)} - x_{(0,\dots,0)}^{(d + 1)}\big|)$ bzw. $(1 - \frac{M}{2a} \cdot \big|x^{(d + 1)} - x_{(1,\dots,1)}^{(d + 1)}\big|)$ vor, da 
\begin{equation*}
\bigg(1 - \frac{M}{2a} \cdot \Big|x^{(d + 1)} - x_\mathbf{i}^{(d + 1)}\Big|\bigg) = \begin{cases}
(1 - \frac{M}{2a} \cdot \big|x^{(d + 1)} - x_{(0,\dots,0)}^{(d + 1)}\big|) &\text{$\mathbf{i} \in \{0, 1\}^d \times \{0\}$}\\
(1 - \frac{M}{2a} \cdot |x^{(d + 1)} - x_{(1,\dots,1)}^{(d + 1)}|) &\text{$\mathbf{i} \in \{0, 1\}^d \times \{1\}$}
\end{cases}
\end{equation*}
gilt. Daraus ergibt sich:
\begin{equation*}
\begin{split}
\sum_{\bi \in [M]^{(d + 1)}} & \prod_{j = 1}^{d + 1} \bigg(1 - \frac{M}{2a} \cdot |x^{(j)} - x_{\mathbf{i}}^{(j)}|\bigg)_+ \\
& = \sum_{\mathbf{i} \in \{0, 1\}^{d + 1}} \prod_{j = 1}^{d + 1} \bigg(1 - \frac{M}{2a} \cdot |x^{(j)} - x_{\bi}^{(j)}|\bigg) \\
& = \Bigg(\sum_{\mathbf{i} \in \{0, 1\}^d \times \{0\}} \prod_{j = 1}^{d} \bigg(1 - \frac{M}{2a} \cdot |x^{(j)} - x_{\bi}^{(j)}|\bigg)\Bigg) \cdot \bigg(1 - \frac{M}{2a} \cdot |x^{(d + 1)} - x_{(0,\dots,0)}^{(d + 1)}|\bigg) \\
& \qquad + \Bigg(\sum_{\mathbf{i} \in \{0, 1\}^d \times \{1\}} \prod_{j = 1}^{d} \bigg(1 - \frac{M}{2a} \cdot |x^{(j)} - x_{\bi}^{(j)}|\bigg)\Bigg) \cdot \bigg(1 - \frac{M}{2a} \cdot |x^{(d + 1)} - x_{(1,\dots,1)}^{(d + 1)}|\bigg) \\
& \stackrel{(\text{IV})}{=} 1 \cdot \bigg(1 - \frac{M}{2a} \cdot |x^{(d + 1)} - x_{(0,\dots,0)}^{(d + 1)}|\bigg) + 1 \cdot \bigg(1 - \frac{M}{2a} \cdot |x^{(d + 1)} - x_{(1,\dots,1)}^{(d + 1)}|\bigg) \\
& = 1 + 1 - \frac{M}{2a} \cdot |x^{(d + 1)} - x_{(0,\dots,0)}^{(d + 1)} + x_{(1,\dots,1)}^{(d + 1)} - x^{(d + 1)}| \\
& = 1 + 1 - 1 \\
& = 1,
\end{split}
\end{equation*}
wobei wir bei der vorletzten Gleichung angewendet haben, dass $x_{(1,\dots,1)}^{(d + 1)} - x_{(0,\dots,0)}^{(d + 1)} = \frac{2a}{M}$ ist, da beides Gitterpunkte sind.  $\hfill(\square)$ 		

\textbf{ii)} Es folgt $\prod_{j = 1}^d (1 - \frac{M}{2a} \cdot |x^{(j)} - x_{\mathbf{i}}^{(j)}|)_+ \geq 0$ für alle $\bi \in [M]^d$, da $$z_+= \max\{z, 0\} \geq 0 \quad (z \in \R)$$ gilt. Damit wäre die Nichtnegativität der Koeffizienten der Linearkombination gezeigt. Damit ist jeder Summand in $$\sum_{\bi \in [M]^d} \prod_{j = 1}^d \bigg(1 - \frac{M}{2a} \cdot |x^{(j)} - x_{\mathbf{i}}^{(j)}|\bigg)_+$$ größer gleich Null und wegen (\ref{induktiongl1}) muss dann auch 
\begin{equation}
\label{induktiongl2}
\prod_{j = 1}^d \bigg(1 - \frac{M}{2a} \cdot |x^{(j)} - x_{\mathbf{i}}^{(j)}|\bigg)_+ \leq 1
\end{equation}
gelten.

\textbf{iii)} Es handelt sich hierbei um einen lokale Träger, da nach der Konstruktion von $B_{\bi}(x)$ der Funktionswert genau dann Null ist, wenn ein $j \in \{1,\dots, d\}$ existiert, sodass $|x^{(j)} - x_{\bi}^{(j)}| \geq \frac{2a}{M}$ gilt. Andernfalls erhalten wir mit Bedingung~\textbf{ii)}, dass $B_{\bi}(x) > 0$ ist. 
\end{proof}

Hiermit erhalten wir, dass:
\begin{equation*}
\begin{split}
P_m(x) &= \sum_{\bi \in [M]^d} p_{\bi}^m(x) \prod_{j = 1}^{d} \bigg(1 - \frac{M}{2a} \cdot |x^{(j)} - x_{\mathbf{i}}^{(j)}|\bigg)_+ \\
& = \sum_{\bi \in [M]^d} p_{\bi}(x) \cdot B_{\bi}(x)
\end{split}
\end{equation*}
gilt und damit $P_m(x)$ eine Spline Interpolation von Taylorpolynomen von $m$ ist.
Die Wahl der Netzwerkarchitektur und der Werte aller Gewichte bis auf die aus der Ausgabeschicht ist durch folgendes Approximationsresultat durch eine lokale Spline Interpolation von Taylorpolynomen für $(p,C)$-glatte Funktionen auf dem Intervall $[-a, a]^d$ motiviert.
\begin{lem}
\label{lem:pcsmooth}
Sei $M \in \N$, $c > 0$, $a > 0$ und $f$ eine ($p, C$)-glatte Funktion, wobei $p = q + s$ mit $q \in \N_0, s \in (0,1]$ und $C > 0$ sind. Sei zudem $P_f(x)$ analog wie in (\ref{konvexkomb}) eine lokale Spline Interpolation von Taylorpolynomen von $f$. Dann gilt$\colon$
$$\sup_{x \in [-a, a]^d} |f(x) - P_f(x)|  \leq c \cdot \frac{1}{M^p}.$$
\end{lem}
\begin{proof}
Nach dem Satz über die Lagrange Form des Restglieds (REFERENZ) existiert ein $\xi \in [0, 1],$ so, dass 
\begin{equation}
\label{eq:lagrange}
\begin{split}
f(x) & = T_{x_{\mathbf{i}},q - 1}[f(x)] \\
& = \sum_{\substack{\bj \in [q]^d \\ |\bj|_1 \leq q - 1}}  \partial^{\bj}f(x_{\mathbf{i}}) \cdot \frac{ (x - x_{\mathbf{i}})^{\bj} }{\bj!} + \sum_{\substack{ \bj \in [q]^d \\|\bj|_1 = q}} \partial^{\bj}f(x_{\bi} + \xi(x - x_{\bi})) \frac{ (x - x_{\mathbf{i}})^{\bj} }{\bj!}.
\end{split}
\end{equation}
Nach der Basispline Eigenschaft aus Gleichung~(\ref{induktiongl1}) erhalten wir 
$$f(x) = \sum_{\bi \in [M]^d} f(x) \prod_{j = 1}^d \bigg(1 - \frac{M}{2a} \cdot |x^{(j)} - x_{\mathbf{i}}^{(j)}|\bigg)_+.$$ Zudem wissen wir, dass man immer den Abstand zu den nächsten $2^d$ Gitterpunkten betrachtet, da  $(1 - \frac{M}{2a} \cdot |x^{(j)} - x_{\mathbf{i}}^{(j)}|)_+ = 0$ immer dann gilt, wenn der Abstand zwischen $x^{(j)}$ und $x_{\mathbf{i}}^{(j)}$ größer als $\frac{2a}{M}$ ist. Daher ergibt sich:
\begin{equation}
\label{eq:absch}
\begin{split}
& \sum_{\substack{ \bj \in [q]^d \\ |\bj|_1 = q}} \partial^{\bj}f(x_{\bi} + \xi(x - x_{\bi})) \frac{ (x - x_{\mathbf{i}})^{\bj} }{\bj!} \leq \sum_{\substack{ \bj \in [q]^d \\ |\bj|_1 = q}} \partial^{\bj}f(x_{\bi} + \xi(x - x_{\bi})) \frac{1}{\bj!} \cdot \bigg(\frac{2a}{M}\bigg)^q.
\end{split}
\end{equation} 
Mithilfe der Dreiecksungleichung und der Definition von $P_f(x)$ erhalten wir:
\begin{equation}
\label{eq:drei}
\begin{split}
& |f(x) - P_f(x)| \\
& \leq \sum_{\bi \in [M]^d} \bigg|f(x) - \sum_{\substack{\mathbf{j} \in [q]^d \\|\mathbf{j}|_1 \leq q}} \partial^{\mathbf{j}}f(x_{\mathbf{i}}) \cdot \frac{(x - x_{\mathbf{i}})^{\mathbf{j}}}{\mathbf{j}!}\bigg| \prod_{j = 1}^d \bigg(1 - \frac{M}{2a} \cdot |x^{(j)} - x_{\mathbf{i}}^{(j)}|\bigg)_+.
\end{split}
\end{equation}
Nach Gleichung~(\ref{eq:lagrange}) erhalten wir:
\begin{equation}
\label{eq:sum}
\begin{split}
& \bigg|f(x) - \sum_{\substack{\mathbf{j} \in [q]^d \\|\mathbf{j}|_1 \leq q}} \partial^{\mathbf{j}}f(x_{\mathbf{i}}) \cdot \frac{(x - x_{\mathbf{i}})^{\mathbf{j}}}{\mathbf{j}!}\bigg| \\ 
& = \bigg| \sum_{\substack{\bj \in [q]^d \\ |\bj|_1 \leq q - 1}}  \partial^{\bj}f(x_{\mathbf{i}}) \cdot \frac{ (x - x_{\mathbf{i}})^{\bj} }{\bj!} + \sum_{\substack{ \bj \in [q]^d \\|\bj|_1 = q}} \partial^{\bj}f(x_{\bi} + \xi(x - x_{\bi})) \frac{ (x - x_{\mathbf{i}})^{\bj} }{\bj!} \\
& \qquad \qquad \qquad \qquad \qquad \qquad \qquad \qquad \qquad - \sum_{\substack{\mathbf{j} \in [q]^d \\|\mathbf{j}|_1 \leq q}} \partial^{\mathbf{j}}f(x_{\mathbf{i}}) \cdot \frac{(x - x_{\mathbf{i}})^{\mathbf{j}}}{\mathbf{j}!}\bigg| \\
& = \bigg| \sum_{\substack{ \bj \in [q]^d \\|\bj|_1 = q}} \partial^{\bj}f(x_{\bi} + \xi(x - x_{\bi})) \frac{ (x - x_{\mathbf{i}})^{\bj} }{\bj!} - \sum_{\substack{ \bj \in [q]^d \\|\bj|_1 = q}} \partial^{\bj}f(x_{\bi}) \frac{ (x - x_{\mathbf{i}})^{\bj} }{\bj!}\bigg|.
\end{split}
\end{equation}
Aus der $(p,C)$-Glattheit von $f$ nach Definition~\ref{def:pc} mit $x_{\bi} + \xi(x - x_{\mathbf{i}}), x_{\bi} \in \R^d$ folgt durch (\ref{eq:absch}) und (\ref{eq:sum}):
\begin{equation}
\label{eq:last}
\begin{split}
& \bigg| \sum_{\substack{ \bj \in [q]^d \\|\bj|_1 = q}} \partial^{\bj}f(x_{\bi} + \xi(x - x_{\bi})) \frac{ (x - x_{\mathbf{i}})^{\bj} }{\bj!} - \sum_{\substack{ \bj \in [q]^d \\|\bj|_1 = q}} \partial^{\bj}f(x_{\bi}) \frac{ (x - x_{\mathbf{i}})^{\bj} }{\bj!}\bigg| \\
& \leq \bigg(\frac{2a}{M}\bigg)^q \|x_{\bi} + \xi(x - x_{\mathbf{i}}) - x_{\bi} \|^s \cdot C.
\end{split}
\end{equation}
Fassen wir die Gleichungen~(\ref{eq:drei}), (\ref{eq:sum}) und (\ref{eq:last}) zusammen, erhalten wir:
\begin{equation*}
\begin{split}
|f(x) - P_f(x)| & \leq \sum_{\bi \in [M]^d} \bigg|f(x) - \sum_{\substack{\mathbf{j} \in [q]^d \\|\mathbf{j}|_1 \leq q}} \partial^{\mathbf{j}}f(x_{\mathbf{i}}) \cdot \frac{(x - x_{\mathbf{i}})^{\mathbf{j}}}{\mathbf{j}!}\bigg| \prod_{j = 1}^d \bigg(1 - \frac{M}{2a} \cdot |x^{(j)} - x_{\mathbf{i}}^{(j)}|\bigg)_+ \\
& \leq \bigg(\frac{2a}{M}\bigg)^q \|x_{\bi} + \xi(x - x_{\mathbf{i}}) - x_{\bi} \|^s \cdot C \cdot \sum_{\bi \in [M]^d} \prod_{j = 1}^d \bigg(1 - \frac{M}{2a} \cdot |x^{(j)} - x_{\mathbf{i}}^{(j)}|\bigg)_+ \\
& \leq \bigg(\frac{2a}{M}\bigg)^q \cdot \bigg(\frac{2a}{M}\bigg)^s \cdot d^{s/2} \cdot C\\
& = c \cdot \frac{1}{M^p},
\end{split}
\end{equation*}
wobei wir bei der letzten Gleichheit Bedingung~(\ref{induktiongl1}) und $q + s = p$ verwendet und 
$$c = (2a)^p \cdot d^{s/2} \cdot C,$$
gewählt haben.
\end{proof}

$P_m(x)$ lässt sich in die Form 
$$\sum_{\bi \in [M]^d} \sum_{\substack{ \bj \in [q]^d\\|\bj|_1 \leq q}} a_{\bi, \bj} \cdot (x - x_{\mathbf{i}})^{\bj} \prod_{j = 1}^d \bigg(1 - \frac{M}{2a} \cdot |x^{(j)} - x_{\mathbf{i}}^{(j)}|\bigg)_+$$
durch geeignet gewählte $a_{\bi, \bj} \in \R$ bringen, da sich jedes $p_i(x)$ der form so und so ... lässt sich $P(x)$ umschreiben ... (TBD)
Als nächstes wollen wir geeignete neuronale Netze $f_{\net, \bj, \bi}$ definieren, die die Funktionen
$$x \mapsto (x - x_{\mathbf{i}})^{\bj} \prod_{j = 1}^d \bigg(1 - \frac{M}{2a} \cdot |x^{(j)} - x_{\mathbf{i}}^{(j)}|\bigg)_+$$ approximieren. Zudem möchten wir die Netzwerkarchitektur so wählen, dass neuronale Netze der Form
$$\sum_{\bi \in [M]^d} \sum_{\substack{ \bj \in [q]^d\\|\bj|_1 \leq q}} a_{\mathbf{i}, \bj} \cdot f_{\net,\bj,\mathbf{i}}(x) \qquad (a_{\mathbf{i},\bj} \in \R)$$ in ihr enthalten sind.
Um dies zu erreichen, sei $$\sigma(x) = \frac{1}{(1 + \exp(-x))} \quad (x \in \R)$$ der logistische Squasher~(\ref{logsquasher}), wählen $R \geq 1$ und definieren die folgenden neuronale Netze:

Das neuronale Netz
\begin{equation}
\label{def:fid}
f_{\id}(x) = 4R \cdot \sigma\Big(\frac{x}{R}\Big) - 2R,
\end{equation}
welches, wie in Lemma~\ref{lem:1} gezeigt, die Funktion $f(x) = x$ approximiert und in Abbildung~\ref{fig:fid} veranschaulicht wird.
\begin{figure}[htp]
\centering
\begin{tikzpicture}[
plain/.style={
  draw=none,
  fill=none,
  },
net/.style={
  matrix of nodes,
  nodes={
    draw,
    circle,
    inner sep=6pt
    },
  nodes in empty cells,
  column sep=0.6cm,
  row sep=-5pt
  },
>=latex
]
\matrix[net] (mat)
{
|[plain]| \parbox{1cm}{\centering Input\\layer} & |[plain]| \parbox{1cm}{\centering Hidden\\layer} & |[plain]| \parbox{1cm}{\centering Output\\layer} \\
& & \\
};
\foreach \ai [count=\mi ]in {2}
  \draw[<-] (mat-\ai-1) -- node[above] {$x$} +(-1cm,0);
\foreach \ai in {2}
{\foreach \aii in {2}
  \draw[->] (mat-\ai-1) -- (mat-\aii-2);
}
\foreach \ai in {2}
  \draw[->] (mat-\ai-2) -- (mat-2-3);
\draw[->] (mat-2-3) -- node[above] {$f_{\id}(x)$} +(2.5cm,0);
\end{tikzpicture}

\caption{Neuronales Netz $f_{\id}(x) = 4R \cdot \sigma(\frac{x}{R}) - 2R$}
\label{fig:fid}
\end{figure}

Das neuronale Netz 
\begin{equation}
\label{def:fmult}
\begin{split}
f_{\mult}(x, y) = \frac{R^2}{4} \cdot \frac{(1 + \exp(- 1))^3}{\exp(-2) - \exp(-1)} \cdot & \bigg(\sigma\Big(\frac{2(x + y)}{R} + 1\Big) - 2 \cdot \sigma \Big(\frac{x + y}{R} + 1\Big) \\
& - \sigma\Big(\frac{2(x - y)}{R} + 1\Big) + 2 \cdot \sigma\Big(\frac{x - y}{R} + 1\Big)\bigg),
\end{split}
\end{equation}
welches, wie in Lemma~\ref{lem:2} gezeigt, die Funktion $f(x, y) = x \cdot y$ approximiert und in Abbildung~\ref{fig:fmult} veranschaulicht wird.
\begin{figure}[htp]
\centering
\begin{tikzpicture}[
plain/.style={
  draw=none,
  fill=none,
  },
net/.style={
  matrix of nodes,
  nodes={
    draw,
    circle,
    inner sep=6pt
    },
  nodes in empty cells,
  column sep=0.6cm,
  row sep=-5pt
  },
>=latex
]
\matrix[net] (mat)
{
|[plain]| \parbox{1cm}{\centering Input\\layer} & |[plain]| \parbox{1cm}{\centering Hidden\\layer} & |[plain]| \parbox{1cm}{\centering Output\\layer} \\
|[plain]| &  & |[plain]|\\
|[plain]| & |[plain]| & |[plain]| \\
& & |[plain]| \\
|[plain]| & |[plain]| & \\
& & |[plain]| \\
|[plain]| & |[plain]| & |[plain]| \\
|[plain]| &  & |[plain]|\\
};
  \draw[<-] (mat-4-1) -- node[above] {$x$} +(-1cm,0);
  \draw[<-] (mat-6-1) -- node[above] {$y$} +(-1cm,0);
\foreach \ai in {4,6}
{\foreach \aii in {2,4,6,8}
  \draw[->] (mat-\ai-1) -- (mat-\aii-2);
}
\foreach \ai in {2,4,6,8}
  \draw[->] (mat-\ai-2) -- (mat-5-3);
\draw[->] (mat-5-3) -- node[above] {$f_{\mult}(x, y)$} +(2.5cm,0);
\end{tikzpicture}

\caption{Neuronales Netz $f_{\mult}(x) = ...$}
\label{fig:fmult}
\end{figure}

Das neuronale Netz 
\begin{equation}
\label{def:frelu}
f_{ReLu}(x) = f_{\mult}(f_{\id}(x), \sigma(R \cdot x)),
\end{equation}
welches, wie in Lemma~\ref{lem:3} gezeigt, die Funktion $f(x) = x_+$ approximiert und in Abbildung~\ref{fig:frelu} veranschaulicht wird
\begin{figure}[htp]
\centering
\begin{tikzpicture}[
plain/.style={
  draw=none,
  fill=none,
  },
net/.style={
  matrix of nodes,
  nodes={
    draw,
    circle,
    inner sep=6pt
    },
  nodes in empty cells,
  column sep=0.6cm,
  row sep=-5pt
  },
>=latex
]
\matrix[net] (mat)
{
|[plain]| \parbox{1cm}{\centering Input\\layer} & |[plain]| \parbox{1.5cm}{\centering Hidden\\layer 1} & |[plain]| \parbox{1.5cm}{\centering Hidden\\layer 2} & |[plain]| \parbox{1cm}{\centering Output\\layer} \\
|[plain]| & |[plain]| &  & |[plain]|\\
|[plain]| & |[plain]| & |[plain]| & |[plain]| \\
|[plain]| & & & |[plain]| \\
& |[plain]| & |[plain]| & \\
|[plain]| & & & |[plain]| \\
|[plain]| & |[plain]| & |[plain]| & |[plain]| \\
|[plain]| & |[plain]| &  & |[plain]|\\
};
  \draw[<-] (mat-5-1) -- node[above] {$x$} +(-1cm,0);
\foreach \ai in {5}
{\foreach \aii in {4,6}
  \draw[->] (mat-\ai-1) -- (mat-\aii-2);
}  
\foreach \ai in {4,6}
{\foreach \aii in {2,4,6,8}
  \draw[->] (mat-\ai-2) -- (mat-\aii-3);
}
\foreach \ai in {2,4,6,8}
  \draw[->] (mat-\ai-3) -- (mat-5-4);
\draw[->] (mat-5-4) -- node[above] {$f_{\ReLU}(x, y)$} +(2.5cm,0);
\end{tikzpicture}

\caption{Neuronales Netz $f_{\ReLU}(x) = ...$}
\label{fig:frelu}
\end{figure}
und schließlich das neuronale Netz 
\begin{equation}
\label{def:fhat}
f_{\mathrm{hat},y}(x) = f_{ReLu}\bigg(\frac{M}{2a} \cdot (x - y) + 1\bigg) - 2 \cdot f_{ReLu}\bigg(\frac{M}{2a} \cdot (x - y)\bigg) +  f_{ReLu}\bigg(\frac{M}{2a} \cdot (x - y) - 1\bigg),
\end{equation}
welches, wie in Lemma~\ref{lem:4} gezeigt, für fixes $y \in \R$ die Funktion $$f(x) = \bigg(1 - \bigg(\frac{M}{2a}\bigg) \cdot |x - y|\bigg)_+$$ approximiert und in Abbildung~\ref{fig:fhat} veranschaulicht wird.
\begin{figure}[htp]
\centering
\begin{tikzpicture}[
plain/.style={
  draw=none,
  fill=none,
  },
net/.style={
  matrix of nodes,
  nodes={
    draw,
    circle,
    inner sep=6pt
    },
  nodes in empty cells,
  column sep=0.6cm,
  row sep=-5pt
  },
>=latex
]
\matrix[net] (mat)
{
|[plain]| \parbox{1cm}{\centering Input\\layer} & |[plain]| \parbox{1.5cm}{\centering Hidden\\layer 1} & |[plain]| \parbox{1.5cm}{\centering Hidden\\layer 2} & |[plain]| \parbox{1cm}{\centering Output\\layer} \\
|[plain]| & |[plain]| &  & |[plain]|\\
|[plain]| & |[plain]| & |[plain]| & |[plain]| \\
|[plain]| & & & |[plain]| \\
|[plain]| & |[plain]| & |[plain]| & |[plain]| \\
|[plain]| & & & |[plain]| \\
|[plain]| & |[plain]| & |[plain]| & |[plain]| \\
|[plain]| & |[plain]| &  & |[plain]|\\
|[plain]| & |[plain]| & |[plain]| & |[plain]| \\
|[plain]| & |[plain]| & |[plain]| & |[plain]| \\
|[plain]| & |[plain]| &  & |[plain]|\\
|[plain]| & |[plain]| & |[plain]| & |[plain]| \\
|[plain]| & & & |[plain]| \\
& |[plain]| & |[plain]| & \\
|[plain]| & & & |[plain]| \\
|[plain]| & |[plain]| & |[plain]| & |[plain]| \\
|[plain]| & |[plain]| &  & |[plain]|\\
|[plain]| & |[plain]| & |[plain]| & |[plain]| \\
|[plain]| & |[plain]| & |[plain]| & |[plain]| \\
|[plain]| & |[plain]| &  & |[plain]|\\
|[plain]| & |[plain]| & |[plain]| & |[plain]| \\
|[plain]| & & & |[plain]| \\
|[plain]| & |[plain]| & |[plain]| & |[plain]| \\
|[plain]| & & & |[plain]| \\
|[plain]| & |[plain]| & |[plain]| & |[plain]| \\
|[plain]| & |[plain]| &  & |[plain]|\\
};
  \draw[<-] (mat-14-1) -- node[above] {$x$} +(-1cm,0);
\foreach \ai in {14}
{\foreach \aii in {4,6,13,15,22,24}
  \draw[->] (mat-\ai-1) -- (mat-\aii-2);
}  
\foreach \ai in {4,6}
{\foreach \aii in {2,4,6,8}
  \draw[->] (mat-\ai-2) -- (mat-\aii-3);
}
\foreach \ai in {13,15}
{\foreach \aii in {11,13,15,17}
  \draw[->] (mat-\ai-2) -- (mat-\aii-3);
}
\foreach \ai in {22,24}
{\foreach \aii in {20,22,24,26}
  \draw[->] (mat-\ai-2) -- (mat-\aii-3);
}
\foreach \ai in {2,4,6,8,11,13,15,17,20,22,24,26}
  \draw[->] (mat-\ai-3) -- (mat-14-4);
\draw[->] (mat-14-4) -- node[above] {$f_{\mathrm{\mathrm{hat}},y}(x)$} +(2.5cm,0);
\end{tikzpicture}

\caption{Neuronales Netz $f_{\mathrm{\mathrm{hat}},y}(x) = ...$}
\label{fig:fhat}
\end{figure} 
Mit diesen neuronalen Netzen können wir nun $f_{\net,\bj,\bi}$ rekursiv definieren. Dafür wählen wir $N \geq q$, setzen $s = \lceil\log_2(N + d)\rceil$ und definieren für $\bj \in [N]^d$, $\bi \in [M]^d$ und $k \in \{1,\dots,(M + 1)^d\}$:  
\begin{align}
\label{fnet}
f_{\net,\bj,\mathbf{i}}(x) & = f_1^{(0)}(x). \nonumber
\intertext{wobei} \nonumber
f_k^{(l)}(x) & = f_{\mult}\Big(f_{2k - 1}^{(l + 1)}(x),f_{2k}^{(l + 1)}(x)\Big) \nonumber
\intertext{für $k \in \{1, 2, \dots, 2^l \}$ und $l \in\{0,\dots,s - 1\}, $ und}
f_k^{(s)}(x) & = f_{\id}(f_{\id}(x^{(l)} - x_{\mathbf{i}_k}^{(l)})) \nonumber
\intertext{für $j_1 + j_2 + \dots + j_{l-1} + 1 \leq k \leq j_1 + j_2 + \dots + j_l$ und $1 \leq l \leq d$ und}
f_{j_1 + j_2 + \dots + j_d + k}^{(s)}(x) & = f_{\mathrm{hat}, x_{\mathbf{i}_k}^{(k)}}(x^{(k)})
\intertext{für $1 \leq k \leq d$ und}
f_k^{(s)}(x) & = 1 \nonumber
\end{align} 
für $j_1 + j_2 + \dots + j_d + d + 1 \leq k \leq 2^s.$
 
Da das neuronale Netz $f_{\net,\bj,\bi}$ aus mehrere neuronalen Netzen zusammengebaut wurde, lässt sich dadurch auch die Anzahl an Schichten und Neuronen pro Schicht durch diese Struktur erklären. Aus (\ref{fnet}) entnimmt man, dass $f_{\mathrm{net}}$ $s + 2$ verborgene Schichten, durch $s$-maliges Anwenden von $f_{\mult}$ und einer Anwendung von $f_{\mathrm{hat}}$, hat. Da $f_{\mathrm{hat}}$ zwei verborgene Schichten besitzt ergibt sich daraus die Anzahl an verborgenen Schichten von $f_{\mathrm{net}}$. Für die Anzahl an Neuronen für die jeweiligen Schichten können wir nur eine oberen Schranke angeben, da ... (TBD)
Die Anzahl der Neuronen pro verborgener Schicht von $f_{\mathrm{net}}$ ergeben sich wie folgt:
\begin{itemize}
\item Die erste verborgene Schicht enthält maximal $3 \cdot 2 \cdot 2^s = 6 \cdot 2^s$ Neuronen, da dies die erste verborgene Schicht von $f_{\mathrm{hat}}$ ist und maximal $2^s$ mal aufgerufen wird. 
\item Die zweite verborgene Schicht maximal $3 \cdot 4 \cdot 2^s = 12\cdot 2^s$ Neuronen, da dies die zweite verborgene Schicht von $f_{\mathrm{hat}}$ ist und maximal $2^s$-mal aufgerufen wird.
\item Die verborgenen Schichten $3,\dots,s + 2$ enthalten maximal   $2 \cdot 2^s, 2^s, \dots, 8, 4$ Neuronen, da wir $s$-mal $f_{\mult}$ ineinander geschachtelt aufrufen. 
\end{itemize}  
Da man bei fully connected neuronalen Netzen die Gewichte der Verbindungen zwischen zwei Neuronen auf Null setzen kann, können auch nicht fully connected neuronale Netze in dieser Klasse enthalten sein. Daher liegt auch $f_{\net,\bj,\bi}$ in der der Klasse aller fully connected neuronaler Netze, mit $s + 2$ verbogenen Schichten mit jeweils $24 \cdot (N + d)$ Neuronen pro Schicht, da für die größte Anzahl an Neuronen in einer Schicht $$12 \cdot 2^s = 12 \cdot 2^{\lceil\log_2(N + d)\rceil} \leq 12 \cdot 2^{\log(N + d) + 1} = 24 \cdot (N + d)$$ gilt. Weiterhin erkennt man durch die Zusammensetzung der neuronalen Netze, dass alle Gewichte im Betrag durch $c \cdot \max\{\frac{M}{2a}, R^2\}$ beschränkt sind, wobei $c > 0$ ist. 

\section{Definition der Gewichte der Ausgabeschicht}

Wir definieren unseren Neuronale-Netze-Regressionsschätzer $\tilde{m}_n(x)$  durch:
\begin{equation}
\label{estimate}
\tilde{m}_n(x) = \sum_{\bi \in [M]^d} \sum_{\substack{\bj \in [N]^d\\|\bj|_1 \leq N}} a_{\mathbf{i},\bj} \cdot f_{\net,\bj,\mathbf{i}}(x),
\end{equation}
wobei $n$ die Größe unserer gegebenen Datenmenge~(\ref{dataset}) ist und wir die Koeffizienten $a_{\mathbf{i},\bj}$ durch Minimierung des Funktionals 
\begin{equation}
\label{min} 
\begin{split}
\varphi(a_{\bi,\bj}) & \coloneqq \sum_{i = 1}^n|Y_i - \tilde{m}_n(X_i)|^2 + \lambda\|(a_{\bi,\bj})\|_2^2 \\
& =\frac{1}{n} \sum_{i = 1}^n|Y_i - \tilde{m}_n(X_i)|^2 + \frac{c}{n} \cdot \sum_{\bi \in [M]^d} \sum_{\substack{\bj \in [N]^d\\|\bj|_1 \leq N}} a_{\mathbf{i},\bj}^2
\end{split}
\end{equation}
für eine Konstante $c > 0$ und Regularitätsterm $\lambda\|(a_{\bi,\bj})\|_2^2$ erhalten. Dieses regularisierte lineare Kleinste-Quadrate Schätzung erhalten wir durch die Lösung des, in (\ref{les}) folgenden, linearen Gleichungssystems. Dafür definieren wir uns die Menge 
\begin{align*} 
\{U_s \mid s = 1,\dots,S\} = \Bigl\{f_{\net,\bj,\bi}(x) \mid \bi \in [M]^d \text{ und } |\bj|_1 \leq N \, \text{ mit }\, \bj \in [N]^d \Bigr\}
\end{align*}
wobei
$$ S = \big|[M]^d\big| \cdot  \binom{N + d}{d} = (M + 1)^d \cdot \binom{N + d}{d}$$ die Kardinalität der Menge ist .
Diese Kardinalität erhalten wir mit einem Kombinatorik Argument.
Wir wissen, dass es insgesamt $(M + 1)^d$ Möglichkeiten gibt $d$ viele Zahlen aus einer Menge mit der Größe $(M + 1)$ zu ziehen mit Zurücklegen und da wir Vektoren betrachten und die Komponenten nicht vertauschbar sind ist auch die Reihenfolge der Ziehung zu beachten.
Für jede dieser $(M + 1)^d$ Möglichkeiten ist noch zu beachten, dass wir zusätzlich $d$ mal aus einer Menge mit $(N + 1)$ vielen Zahlen ziehen müssen und gleichzeitig die Bedingung dass die Summe der gezogen $d$ Elemente zwischen Null und $N$ liegt.
Gesucht ist also 
$$\big|\Bigl\{\bj \in [N]^d : |\bj|_1 \leq N \Bigr\}\big| \eqqcolon H.$$ 
Wir stellen fest, dass:
\begin{equation*}
\begin{split}
& \Bigl\{\bj \in [N]^d : |\bj|_1 \leq N \Bigr\} \\
& = \Bigl\{\bj \in [N]^d : |\bj|_1 = 0 \Bigr\}
 \cup \Bigl\{\bj \in [N]^d : |\bj|_1 = 1 \Bigr\}
 \cup \dots 
 \cup \Bigl\{\bj \in [N]^d : |\bj|_1 = N\Bigr\}
\end{split}
\end{equation*}
gilt. Mit Lemma~\ref{lem:kombi} wissen wir, dass für $d, N \in \N$ und $k \in \N_0$:
$$\big|\Bigl\{\bj \in [N]^d : |\bj|_1 = k \Bigr\}\big| = \binom{d + k - 1}{k}$$ gilt.
Damit erhalten wir:
$$|H| = \sum_{k = 0}^N \binom{N - 1 + k}{k} = \binom{N + d}{d},$$
mit dem \emph{Hockey Stick Lemma}
$$\binom{n}{k} = \sum_{i = 0}^k \binom{n - k - 1 + i}{i} \quad (k, n \in \N  \text{ mit  } k < n).$$
Wir setzen nun 
$$ \mathbf{U} = (U_s(X_i))_{1\leq i \leq n,1\leq s \leq S} \quad \text{und} \quad \mathbf{Y} = (Y_i)_{i = 1,\dots,n}.$$
Im folgenden Lemma bestimmen wir den Koeffizientenvektor unseres Schätzers~(\ref{estimate}).
\begin{lem}
Der Koeffizientenvektor unseres Schätzers (\ref{estimate}) ist die eindeutige Lösung des linearen Gleichungssystems 
\begin{equation}	
\label{les}
\bigg(\frac{1}{n}\mathbf{U}^T\mathbf{U} + \frac{c}{n} \cdot \mathbf{1} \bigg) \mathbf{a} = \frac{1}{n} \mathbf{U}^T\mathbf{Y}
\end{equation}
Hierbei ist $\mathbf{1}$ eine $S \times S$-Einheitsmatrix und $\ba \in \R^S$, wobei $\ba$ wie in Gleichung~(\ref{min}) den Ausdruck:
\begin{equation}
\label{eq:min}
\begin{split}
& \frac{1}{n} \sum_{i = 1}^n|Y_i - \tilde{m}_n(X_i)|^2 + \frac{c}{n} \cdot \sum_{\bi \in [M]^d} \sum_{\substack{\bj \in [N]^d\\|\bj|_1 \leq N}} a_{\mathbf{i},\bj}^2 \\
& = \frac{1}{n}(\mathbf{Y} - \mathbf{U}\mathbf{a})^T(\mathbf{Y} - \mathbf{U}\mathbf{a}) + \frac{c}{n} \cdot \mathbf{a}^T\mathbf{a}
\end{split}
\end{equation}
minimiert.
\end{lem}
\begin{proof}
Der Schätzer aus Gleichung~(\ref{estimate}) lässt sich umschreiben zu 
$$ \tilde{m}_n(x) = \sum_{s = 1}^S a_s \cdot U_s(x)$$
mit $(a_s)_{s = 1,\dots,S} = \mathbf{a}\in \R^S$. Da $\mathbf{Y}^T\mathbf{U}\mathbf{a} = \mathbf{a}^T\mathbf{U}^T\mathbf{Y}$ gilt, da dieser Ausdruck eine reelle Zahl und damit insbesondere symmetrisch ist erhalten wir mit Gleichung~(\ref{eq:min}):
\begin{equation}
\label{eq:matrix}
\begin{split}
& \frac{1}{n}(\mathbf{Y} - \mathbf{U}\mathbf{a})^T(\mathbf{Y} - \mathbf{U}\mathbf{a}) + \frac{c}{n} \cdot \mathbf{a}^T\mathbf{a} \\
& = \frac{1}{n}(\mathbf{Y}^T\mathbf{Y} - \mathbf{Y}^T\mathbf{U}\mathbf{a} - \mathbf{a}^T\mathbf{U}^T\mathbf{Y} + \mathbf{a}^T\mathbf{U}^T\mathbf{U}\mathbf{a}) + \frac{c}{n} \cdot \mathbf{a}^T\mathbf{a} \\
& = \frac{1}{n}(\mathbf{Y}^T\mathbf{Y} - 2\mathbf{Y}^T\mathbf{U}\mathbf{a}) + \mathbf{a}^T\bigg(\frac{1}{n} \mathbf{U}^T\mathbf{U} + \frac{c}{n} \cdot \mathbf{1}\bigg) \mathbf{a}.
\end{split}
\end{equation} 
Die Matrix $\mathbf{U}^T\mathbf{U} \in \R^{S \times S}$ ist positiv semidefinit, denn aufgrund der Verschiebungseigenschaft des Standardskalarprodukts gilt für alle $x \in \R$:
$$\langle x, \mathbf{U^T}\mathbf{U} x\rangle = \langle \mathbf{U} x, \mathbf{U} x\rangle \geq 0.$$
Zudem wissen wir dass $\frac{c}{n}\mathbf{1}$ durch die Wahl von $c$ nur positive Eigenwerte besitzt und damit positiv definit ist.  
Daher wissen wir, dass die Matrix
$$\mathbf{A} \coloneqq \frac{1}{n}\mathbf{U}^T\mathbf{U} + \frac{c}{n} \cdot \mathbf{1}$$ also Summe einer positiv semidefiniten und einer positiv definiten Matrix nur positive Eigenwerte besitzt (REFERENZ), damit also positiv definit und damit insbesondere invertierbar ist. Zudem ist die Matrix $\mathbf{A}$ symmetrisch. 
Mit $$\mathbf{b} = \frac{1}{n} \cdot \mathbf{A}^{-1}\mathbf{U}^T\mathbf{Y} \in \R^S$$ und $$\mathbf{b}^T\mathbf{A}\mathbf{a} = \mathbf{a}^T\mathbf{A}\mathbf{b} = \mathbf{Y}^T\mathbf{U}\mathbf{a},$$ was aus der Symmetrie von $\mathbf{A}$ folgt, erhalten wir in Gleichungs~(\ref{eq:matrix}):
\begin{equation*}
\begin{split}
& \frac{1}{n}(\mathbf{Y}^T\mathbf{Y} - 2\mathbf{Y}^T\mathbf{U}\mathbf{a}) + \mathbf{a}^T\bigg(\frac{1}{n} \mathbf{U}^T\mathbf{U} + \frac{c}{n} \cdot \mathbf{1}\bigg) \mathbf{a} \\
& = \frac{1}{n}(\mathbf{Y}^T\mathbf{Y} - 2\mathbf{Y}^T\mathbf{U}\mathbf{a}) + \mathbf{a}^T \bA \mathbf{a} \\
& = \frac{1}{n}\bY^T\bY - \frac{1}{n}\bb^T\bA\ba - \frac{1}{n}\ba^T\bA\bb + \ba^T\bA\ba \\
& = \frac{1}{n}\bY^T\bY - \frac{1}{n}\bb^T\bA\ba - \frac{1}{n}\ba^T\bA\bb + \frac{1}{n}\bb^T\bU^T\bY - \frac{1}{n}\bY^T\bU\bb + \ba^T\bA\ba \\
& = \frac{1}{n}\bY^T\bY - \frac{1}{n}\bb^T\bA\ba - \frac{1}{n}\ba^T\bA\bb + \frac{1}{n}\bb^T\bA\bA^{-1}\bU^T\bY - \frac{1}{n}\bY^T\bU\bb + \ba^T\bA\ba \\
& = \frac{1}{n}\bY^T\bY - \frac{1}{n}\bb^T\bA\ba - \frac{1}{n}\ba^T\bA\bb + \frac{1}{n}\bb^T\bA\bb - \frac{1}{n^2}\bY^T\bU\bA^{-1}\bU^T\bY + \ba^T\bA\ba \\
& = \mathbf{a}^T\mathbf{A}\mathbf{a} - \frac{1}{n}\mathbf{b}^T\mathbf{A}\mathbf{a} - \frac{1}{n}\mathbf{a}^T\mathbf{A}\mathbf{b} + \mathbf{b}^T\mathbf{A}\mathbf{b} + \frac{1}{n}\mathbf{Y}^T\mathbf{Y} - \frac{1}{n^2}\mathbf{Y}^T\mathbf{U}\mathbf{A}^{-1}\mathbf{U}^T\mathbf{Y} \\
& = (\mathbf{a} - \frac{1}{n} \cdot \mathbf{A}^{-1}\mathbf{U}^T\mathbf{Y})^T \mathbf{A} (\mathbf{a} - \frac{1}{n} \cdot \mathbf{A}^{-1} \mathbf{U}^T\mathbf{Y}) - \frac{1}{n}\mathbf{Y}^T\mathbf{Y} - \frac{1}{n^2}\mathbf{Y}^T\mathbf{U}\mathbf{A}^{-1}\mathbf{U}^T\mathbf{Y}.
\end{split} 
\end{equation*} 
Die letzte Gleichung wird für $\mathbf{a} = \frac{1}{n} \cdot \mathbf{A}^{-1}\mathbf{U}^T\mathbf{Y}$ minimal, 
da wir wissen dass $\mathbf{A}$ positiv definit ist und damit $x^T\mathbf{A}x > 0$ für alle $x \in \R^S$ mit $x \neq 0$ gilt und $(\mathbf{a} - \mathbf{b})^T\mathbf{A}(\mathbf{a} - \mathbf{b}) = 0$ genau dann, wenn $\mathbf{a} = \mathbf{b}$ gilt.
Dies zeigt also, dass der Koeffizientenvektor unseres Schätzers~(\ref{estimate}) die eindeutige Lösung des linearen Gleichungssystems~(\ref{les}) ist.
\end{proof}
%\begin{bemnumber}
%Da der Koeffizientenvektor $\mathbf{a}$ die Gleichung~(\ref{eq:min}) minimiert, erhalten wir, wenn wir den Koeffizientenvektor gleich Null setzen:
%$$\frac{1}{n}(\mathbf{Y} - \mathbf{U}\mathbf{a})^T(\mathbf{Y} - \mathbf{U}\mathbf{a}) + \frac{c}{n} \cdot \mathbf{a}^T\mathbf{a} \leq \frac{1}{n} \sum_{i = 1}^n Y_i^2,$$
%was uns erlaubt eine obere Schranke für den absoluten Wert unserer Koeffizienten abzuleiten. (TBD) WOFÜR BRAUCH ICH DAS DENN? WEGLASSEN!
%\end{bemnumber}
%\chapter{Resultat zur Konvergenzgeschwindigkeit}
\label{chap:3}

In dieser Arbeit behandeln wir Neuronale-Netze-Regressionsschätzer im Kontext der nichtparametrischen Regression mit zufälligem Design. 

Im Gegensatz zur parametrischen Regression ist bei der nichtparametrischen die Bauart der zu schätzenden Funktion komplett unbekannt, was den Vorteil besitzt, dass weniger Annahmen getroffen werden müssen, man aber dadurch noch mehr Daten benötigt, um eine Funktion zu schätzen.

Bei der nichtparametrischen Regressionsschätzung seien $(X, Y), (X_1, Y_1), (X_2, Y_2), \dots$ unabhängig identisch verteilte (u.\@i.\@v.\@) $\R^d \times \R$-wertige Zufallsvariablen mit $\E[Y^2] < \infty$ und $d \in \N$. Zudem sei $m\colon \R^d \to \R$ definiert durch $m(x) = \E[Y \mid X = x]$ die zugehörige Regressionsfunktion. Ausgehend von der Stichprobe
$$ (X_1, Y_1),\dots,(X_n, Y_n),$$ 
mit Stichprobengröße $n \in \N$, soll die Regressionsfunktion $m$ geschätzt werden.

Das Problem der Regressionsschätzung bei zufälligem Design lässt sich wie folgt erläutern: 

In Anwendungsfällen ist üblicherweise die Verteilung von $(X, Y)$ unbekannt, daher kann $m(x) = \E[Y \mid X = x]$ nicht berechnet werden. Oft ist es aber möglich, Werte von $(X, Y)$ zu beobachten. Ziel ist es dann, daraus die Regressionsfunktion $m$ zu schätzen. Im Hinblick auf die Minimierung des $L_2$-Risikos sollte dabei der $L_2$-Fehler der Schätzfunktion möglichst klein sein. 

Für das $L_2$-Risiko einer beliebigen messbaren Funktion $f\colon \R^d \to \R$  gilt:
$$\E[|f(X) - Y|^2] = \E[|m(X) - Y|^2] + \int_{\R^d}|f(x) - m(x)|^2 \mathds{P}_X (dx),$$
d.h. der mittlere quadratische Vorhersagefehler einer Funktion ist darstellbar als Summe des $L_2$-Risikos der Regressionsfunktion (unvermeidbarer Fehler) und des $L_2$-Fehlers der entsteht aufgrund der Verwendung von $f$ an Stelle von $m$ bei der Vorhersage bzw.\@ Approximation des Wertes von Y.

Formal führt das daher auf folgende Problemstellung:

$(X, Y), (X_1, Y_1), (X_2, Y_2), \dots$ seien u.\@i.\@v.\@ $\R^d \times \R$ wertige Zufallsvariablen mit $\E[Y^2] < \infty$ und $m\colon\R^d \to \R$ definiert durch $m(x) = \E[Y \mid X = x]$ sei die zugehörige Regressionsfunktion. Gegeben sei die Datenmenge 
$$ \mathcal{D}_n = \{(X_1, Y_1),\dots,(X_n, Y_n)\}.$$
Gesucht ist eine Schätzung 
$$m_n(\cdot) = m_n(\cdot, \mathcal{D}_n)\colon\R^d \to \R $$
von $m$, für die der $L_2$-Fehler 
$$\int |m_n(x) - m(x)|^2 \mathds{P}_X(dx)$$
möglichst \glqq klein\grqq\@ ist. \cite{györfi2002}

In diesem Kapitel stellen wir das Hauptresultat dieser Arbeit vor.
Ziel im Folgenden ist es, eine Abschätzung des erwarteten $L_2$-Fehlers 
$$\E \int |m_n(x) - m(x)|^2  \mathds{P}_X(dx)$$
im Falle unseres Neuronale-Netze-Regresssionsschätzers aus Kapitel~\ref{chap:2} mit einer $(p,C)$-glatten Regressionsfunktion $m$ herzuleiten.

\begin{thm}\label{optstop}
Angenommen die Verteilung von $Y$ erfüllt 
$$ \E\Big[\mathrm{e}^{c_1 \cdot |Y|^2}\Big] < \infty$$
für eine Konstante $c_1 > 0$ und die Verteilung von $X$ besitzt einen beschränkten Träger $\supp(\mathds{P}_X)$. Sei $m(x) = \E[Y \mid X = x]$ die zu dem Tupel $(X, Y)$ gehörige Regressionsfunktion.
Angenommen $m$ ist $(p,C)$-glatt, mit $q \in \N_0$, $s \in (0,1]$ und $C > 0$. 

Wir betrachten unseren Neuronale-Netze-Regressionsschätzer $\tilde{m}_n$ aus Kapitel~\ref{subsec:2.2} der aus mehreren neuronalen Netzen konstruiert wurde. Für diesen wählen wir unsere Aktivierungsfunktion $\sigma$ als den logistische squasher aus Gleichung~(\ref{logsquasher}). Aufgrund der Konstruktion unseres Schätzers wählen wir zudem $N \geq q, M = M_n = \lceil c_2 \cdot n^{1/(2p + d)}\rceil$ mit $c_2 >0$ und unabhängig von $n$, $R = R_n = n^{d + 4}$ und $a = a_n = (\log n)^{1/(6(N + d))}.$
Sei $\beta_n = c_3 \cdot \log(n)$ für eine hinreichend große und von $n$ unabhängige Konstante $c_3 > 0$ und sei $m_n$ gegeben durch
$$m_n(x) \coloneqq T_{\beta_n}\tilde{m}_n (x)$$
mit $T_{\beta}z = \max\{\min\{z, \beta\}, -\beta\}$ für $z \in \R$ und $\beta > 0.$ 

Dann erhalten wir für hinreichend großes $n$:
$$\E \int |m_n(x) - m(x)|^2  \mathds{P}_X(dx) \leq c_{\text{fin}} \cdot (\log n)^3 \cdot n^{- \frac{2p}{2p + d}},$$
wobei $c_{\text{fin}} > 0$ ist und nicht von $n$ abhängt.
\end{thm}
Die nächsten Lemmata benötigen wir für den Beweis unseres Hauptresultats, einer Aussage über die Konvergenzgeschwindigkeit unseres Neuronale-Netze-Schätzers. Diese Lemmata werden hier nur der Vollständigkeit halber und ohne Beweis aufgeführt.
\begin{lem}[\cite{kohler19}, Lemma 5]
\label{lem:5}
Sei $M \in \N$ und $\sigma\colon \R \to [0, 1]$ $2$-zulässig nach Definition~\ref{nzulässig}.
Sei $a \geq 1$ und $R \in \R$ mit
\begin{equation*}
\begin{split}
R & \geq \max\biggl\{\frac{\|\sigma''\|_{\infty} \cdot (M + 1)}{2 \cdot |\sigma'(t_{\sigma})|}, \frac{9 \cdot \|\sigma''\|_{\infty} \cdot a}{|\sigma'(t_{\sigma})|}, \\
& \quad \frac{20 \cdot \|\sigma'''\|_{\infty}}{3 \cdot |\sigma''(t_{\sigma})|} \cdot 3^{3 \cdot 3^s} \cdot a^{3 \cdot 2^s}, 1792 \cdot \frac{\max\{\|\sigma''\|_{\infty},\|\sigma'''\|_{\infty}, 1\}}{\min\{2 \cdot |\sigma'(t_{\sigma})|, |\sigma''(t_{\sigma})|, 1\}} \cdot M^3 \biggr\}
\end{split}
\end{equation*}
und sei $y \in [-a, a]^d.$ Sei $N \in \N$ und $\bj \in [\N_0]^d$ so, dass $|\bj|_1 \leq N$ gilt und wir setzen $s = \lceil\log_2(N + d)\rceil$. Sei $f_{\id}, f_{\mult}$ und $f_{\mathrm{hat}, z}$ (für $z \in \R$) die neuronalen Netze aus Lemma~\ref{lem:1}, Lemma~\ref{lem:2} und Lemma~\ref{lem:4}. Wir definieren das Netz $f_{net,\bj,y}$ durch:
\begin{align*}
\label{fnet}
& f_{net,\bj,y}(x) = f_1^{(0)}(x). \\
\intertext{wobei} 
& f_k^{(l)}(x) = f_{\mult}\Big(f_{2k - 1}^{(l + 1)}(x),f_{2k}^{(l + 1)}(x)\Big) \\
\intertext{für $1 \leq k \leq 2^l$ und $0 \leq l \leq s - 1$ und} 
& f_k^{(s)}(x) = f_{\id}(f_{\id}(x^{(l)} - y^{(l)}))  \\
\intertext{für $j_1 + j_2 + \dots + j_{l-1} + 1 \leq k \leq j_1 + j_2 + \dots + j_l$ und $1 \leq l \leq d$ und} 
& f_{|\bj|_1 + k}^{(s)}(x) = f_{\mathrm{hat},y^{(k)}}(x^{(k)}) \\
\intertext{für $1 \leq k \leq d$ und} 
& f_k^{(s)}(x) = 1 \\
\intertext{für $|\bj|_1+ d + 1 \leq k \leq 2^s.$}
\end{align*} 
Dann erhalten wir für $x \in [-a, a]^d$:
\begin{equation*}
\begin{split}
& \bigg|f_{net,\bj,y}(x) - (x^{(1)} - y^{(1)})^{j_1} \cdots (x^{(d)} - y^{(d)})^{j_d} \prod_{j = 1}^d (1 - \frac{M}{2a} \cdot |x^{(j)} - y^{(j)}|)_+\bigg| \\
& \leq c \cdot 3^{3 \cdot 3^s} \cdot a^{3 \cdot 2^s} \cdot M^3 \cdot \frac{1}{R},
\end{split}
\end{equation*}
für eine von $n$ unabhängige Konstante $c > 0$.
\end{lem}
  \begin{lem}[\cite{kohler19}, Lemma 8]
  \label{lem:8}
Sei $\beta_n = c_1 \cdot \log(n)$ für eine hinreichend große Konstante $c_1 > 0$. Angenommen die Verteilung von $Y$ erfüllt 
$$ \E\Big[\mathrm{e}^{c_2 \cdot |Y|^2}\Big] < \infty$$
für eine Konstante $c_2 > 0$ und dass der Betrag der Regressionsfunktion $m$ beschränkt ist. Sei $\mathcal{F}_n$ eine Menge von Funktionen $f\colon \R^d \to \R$ und wir nehmen an, dass der Schätzer $m_n$ 
$$m_n = T_{\beta_n}\tilde{m}_n$$ 
erfüllt, mit 
$$\tilde{m}_n(\cdot) = \tilde{m}_n(\cdot,(X_1, Y_1),\dots,(X_n, Y_n)) \in \mathcal{F}_n$$
und 
$$\frac{1}{n} \sum_{i = 1}^n |Y_i - \tilde{m}_n(X_i)|^2 \leq \min_{l \in \Theta_n}\bigg(\frac{1}{n}\sum_{i = 1}^n |Y_i - g_{n,l}(X_i)|^2 + \text{pen}_n(g_n,l)\bigg)$$
mit einer nichtleeren Parametermenge $\Theta_n$, zufällige Funktionen $g_{n,l}\colon \R^d \to \R$ und deterministischen penalty Terme $\text{pen}_n(g_{n,l}) \geq 0$, wobei die zufälligen Funktionen $g_{n,l}\colon \R^d \to \R$ nur von den Zufallsvariablen
$$\mathbf{b}_1^{(1)},\dots,\mathbf{b}_r^{(1)},\dots,\mathbf{b}_1^{(I_n)},\dots,\mathbf{b}_r^{(I_n)},$$
abhängen, die unabhängig von $(X_1, Y_1), (X_2, Y_2),\dots$ sind.
Dann erfüllt $m_n\colon$
\begin{equation*}
\begin{split}
& \E \int |m_n(x) - m(x)|^2 \mathds{P}_X(dx) \\
& \leq \frac{c \cdot \log(n)^2 \cdot \big(\log\big(\sup_{x_1^n \in (\supp(X))^n}\mathcal{N}_1\big(\frac{1}{n \cdot \beta_n},\mathcal{F}_n,x_1^n\big)\big) + 1\big)}{n} \\
& \quad + 2 \cdot \E\bigg(\min_{l \in \Theta_n} \int |g_{n,l}(x) - m(x)|^2 \mathds{P}_X(dx) + \text{pen}_n(g_{n,l})\bigg),
\end{split}
\end{equation*}
für $n > 1$ und eine von $n$ unabhängige Konstante $c > 0$.
  \end{lem}
Das nächste Lemma benötigen wir um eine Schranke für die Überdeckungszahl $\mathcal{N}_1\big(\frac{1}{n \cdot \beta_n},\mathcal{F}_n,x_1^n\big)$ zu finden.
\begin{lem}[\cite{kohler19}, Lemma 9]
\label{lem:9}
Sei $a > 0$ und $d, N, J_n \in \N$ so, dass $J_n \leq n^{c_{1}}$ und setze $\beta_n = c_2 \cdot \log(n).$
Sei $\sigma$ 2-zulässig nach Definition~\ref{nzulässig}. Sei $\mathcal{F}$ die Menge aller Funktionen die durch Definition~\ref{def:nn} definiert sind mit $k_1 = k_2 = \cdots = k_L = 24 \cdot (N + d)$ und einer Beschränkung der Betrag der Gewichte durch $c_{3} \cdot n^{c_{4}}$.
$$ \mathcal{F}^{(J_n)} = \biggl\{\sum_{j = 1}^{J_n} a_j \cdot f_j : f_j \in \mathcal{F} \quad \text{und} \quad \sum_{j = 1}^{J_n} a_j^2 \leq c_{5} \cdot n^{c_{6}}\biggr\}.$$
Dann gilt für $n > 1:$
$$\log\bigg(\sup_{x_1^n\in[-a,a]^{d \cdot n}} \mathcal{N}_1\bigg(\frac{1}{n \cdot \beta_n}, \mathcal{F}^{(J_n)},x_1^n\bigg)\bigg) \leq c \cdot \log(n) \cdot J_n,$$
für eine Konstante $c$ die nur von $L, N, a$ und $d$ abhängt.
\end{lem}
Mit diesen Hilfsresultaten können wir nun unser Hauptresultat beweisen.
\begin{proof}[Beweis von Satz~\ref{optstop}]
Da nach Voraussetzung $\supp(\mathds{P}_X)$ beschränkt ist und der Träger $\supp(\mathds{P}_X)$ immer abgeschlossen ist, wissen wir, dass $\supp(\mathds{P}_X)$ kompakt ist. Da $m$ als ($p,C$)-glatte Funktion insbesondere stetig ist wissen wir das sie auf einer kompakten Menge ihr Maximum und Minimum annimmt. Dadurch können wir $n$ so groß wählen, dass ohne Beschränkung der Allgemeinheit $\|m\|_{\infty} \leq \beta_n$ gilt.

Sei $\mathcal{F}$ die Menge aller Funktion aus Definition~\ref{def:nn} mit Aktivierungsfunktion $\sigma$ und $L = s + 2 = \lceil\log_2(N + d)\rceil + 2,$ mit $k_1 = k_2 = \cdots = k_L = 24 \cdot (N + d)$ und der Eigenschaft, dass der Betrag der Gewichte durch $n^{c_{4}}$, mit einer Konstante $c_4 > 0$ beschränkt ist. Sei  für eine Konstante $c_5 > 0$
$$ \mathcal{F}^{(J_n)} \coloneqq \biggl\{\sum_{j = 1}^{J_n} a_j \cdot f_j \mid f_j \in \mathcal{F} \text{ und } \sum_{j = 1}^{J_n} a_j^2 \leq c_5 \cdot n \biggr\}$$
wobei
\begin{equation}
\label{constantc21}
\begin{split}
c_5 \coloneqq \max\Biggl\{\frac{1 + \E[Y^2]}{c_6}, & c_2^d \cdot (N + 1)^d \cdot \max\biggl\{\Big| \frac{1}{\bj!} \cdot \partial^{\bj}m(x_{\mathbf{i}})\Big|^2 \bigg| \, \bj \in [q]^d,\, |\bj|_1 \leq q\biggr\}\Biggr\},
\end{split}
\end{equation}
mit $c_6$ als Konstante des Regularitätsterms aus Gleichung~(\ref{min}). Da wir uns in der nichtparametrischen Regressionsschätzung befinden gilt unter anderem die Bedingung $\E[Y^2] < \infty$, und daher ist $c_5$ auch wohldefiniert.
Weiterhin ist
$$J_n = (M_n + 1)^d \cdot \big|\{\bj \mid \bj \in [N]^d, |\bj|_1 \leq N\}\big|,$$ 
die Kardinalität der Menge $\mathcal{F}^{(J_n)}$. 
Ohne die Restriktion $|\bj|_1 \leq N$ lässt sich $$\big|\{\bj \mid \bj \in [N]^d\}\big|$$ durch eine Analogie zu  einem Urnenexperiment bestimmen. Wir betrachten die Anzahl an Möglichkeiten, wie man $d$-Mal mit Zurücklegen (da auch mehrere Komponenten den gleichen Wert haben können) und mit Beachtung der Reihenfolge (da wir einen Vektor betrachten und die Komponenten nicht vertauschen können) aus $(N + 1)$ Kugeln ziehen kann. Durch das Weglassen der Restriktion $|\bj|_1 \leq N$ erhalten wir:
\begin{equation}
\label{jn}
J_n \leq (M_n + 1)^d \cdot (N + 1)^d.
\end{equation}
Da $m$ nach Voraussetzung $(p,C)$-glatt ist, folgt:
\begin{equation}
\label{bound}
z \coloneqq \max_{\bi \in [M]^d,\, \bj \in [q]^d,\, |\bj|_1 \leq q} \big|\partial^{\bj}m(x_{\mathbf{i}})\big| < \infty,
\end{equation}
denn das größte Elemente muss auch beschränkt sein, wenn der Abstand zweier beliebiger Elemente beschränkt ist.
Sei 
$$g_n(x) = \sum_{\bi \in [M_n]^d} \sum_{\substack{ \bj \in [q]^d \\ |\bj|_1 \leq q}} \frac{1}{\bj!} \cdot \partial^{\bj}m(x_{\mathbf{i}}) \cdot f_{\net,\bj,\mathbf{i}}(x).$$
Da nach Konstruktion $f_{\net,\bj,\mathbf{i}} \in \mathcal{F}$ ist, folgt mit Gleichung~(\ref{constantc21}), dass für $n$ hinreichend groß
\begin{equation*}
\begin{split}
\sum_{\bi \in [M_n]^d} \sum_{\substack{ \bj \in [q]^d \\ |\bj|_1 \leq q}} \big| \frac{1}{\bj!} \cdot \partial^{\bj}m(x_{\mathbf{i}})\big|^2 & \leq (M_n + 1)^d (N + 1)^d \cdot \big|\frac{1}{\bj!} \cdot \partial^{\bj}m(x_{\mathbf{i}})\big|^2 \\
& \leq (2 c_2 \cdot n^{1/2p + d})^d \cdot (N + 1)^d \cdot \big|\frac{1}{\bj!} \cdot \partial^{\bj}m(x_{\mathbf{i}})\big|^2 \\
& \leq c_5 \cdot n,
\end{split}
\end{equation*}
gilt und damit $g_n$ in $\mathcal{F}^{(J_n)}$ liegt. 
Wir wählen das Ereignis 
\begin{equation}
\label{event}
A_n \coloneqq \bigg[\frac{1}{n} \sum_{i = 1}^n Y_i^2 \leq 1 + \E[Y^2]\bigg].
\end{equation}
Wir wissen, dass aufgrund der Unabhängigkeit und identischen Verteiltheit der $\R^d \times \R$-wertigen Zufallsvariablen  $(X, Y), (X_1, Y_1), (X_2, Y_2), \dots$ auch die Zufallsvariablen $Y_1,\dots,Y_n$ unabhängig und identisch verteilt sind. 
Daraus folgern wir $\E\big[\frac{1}{n} \sum_{i = 1}^n Y_i^2\big] = \E[Y^2]$ mit der Linearität des Erwartungswerts.
Mit Hilfe der Monotonie der Wahrscheinlichkeitsfunktion $\mathds{P}$ und der Chebyshev-Ungleichung für $\epsilon = 1$ (\cite{Klenke2013}, Satz 5.11) erhalten wir:
\begin{equation*}
\begin{split}
\mathds{P}(A_n^{\mathsf{c}}) & = \mathds{P}(\frac{1}{n}\sum_{i=1}^n Y_i^2 - \E[Y^2] \geq 1) \\
& \leq \mathds{P}(\Big|\frac{1}{n}\sum_{i=1}^n Y_i^2 - \E[Y^2]\Big| \geq 1) \\
& \leq \mathds{V}[\frac{1}{n}\sum_{i = 1}^nY_i^2].
\end{split}
\end{equation*}
Da die Zufallsvariablen $Y_1,\dots,Y_n$ u.\@i.\@v.\@ sind, folgt mit den Rechenregeln der Varianz:
\begin{equation}
\label{tscheby}
\begin{split}
\mathds{P}(A_n^{\mathsf{c}}) & \leq \frac{n \cdot  \mathds{V}[Y^2]}{n^2} \\
& = \frac{\mathds{V}[Y^2]}{n} \\
& = \frac{c_7}{n},
\end{split}
\end{equation}
wobei $c_7 \coloneqq \mathds{V}[Y^2]$ ist.

Sei $$\hat{m}_n \coloneqq \mathds{1}_{A_n}m_n + \mathds{1}_{A_n}T_{\beta_n}g_n$$
mit $m_n= T_{\beta_n}\tilde{m}_n$.

Für zwei beliebige reelle Zahlen $u, v \in \R$ gilt durch $0 \leq (u - v)^2 = u^2 + v^2 - 2uv$:
$$u^2 + v^2 \geq 2uv$$ und damit schließlich:
\begin{equation}
\label{ungl}
\begin{split}
|(u - v)^2| & = |u^2 - 2uv + v^2| \\
& \leq  u^2 + 2uv + v^2 \\
& \leq 2u^2 + 2v^2.
\end{split}
\end{equation}
Mit Ungleichung~(\ref{ungl}) und durch die Unabhängigkeit von $A_n$ von den Zufallsvariablen $X, X_1, \dots, X_n$ erhalten wir:
\begin{equation}
\label{anc}
\begin{split}
 \E \bigg[\int |m_n(x) - m(x)|^2 \mathds{P}_X(dx) \mathds{1}_{A_n^{\mathsf{c}}}\bigg] & =  \E \bigg[\int |m_n(x) - m(x)|^2 \mathds{P}_X(dx) \bigg] \cdot \mathds{P}(A_n^{\mathsf{c}}) \\
 & \leq \E \bigg[\int 2m_n(x)^2 + 2m(x)^2 \mathds{P}_X(dx)\bigg] \cdot \mathds{P}(A_n^{\mathsf{c}})\\
 & \leq \E \bigg[\int  2\beta_n^2 + 2\beta_n^2 \mathds{P}_X(dx)\bigg] \cdot \mathds{P}(A_n^{\mathsf{c}})\\
 & = 4\beta_n^2 \cdot \mathds{P}(A_n^{\mathsf{c}}) \\
 & \stackrel{(\ref{tscheby})}{\leq} \frac{4 \cdot c_7 \cdot \beta_n^2}{n}.
\end{split}
\end{equation}
Da wir nach Bemerkung~(\ref{mtildebeschraenkt}) wissen, dass $\tilde{m}_n$ beschränkt ist, ist $m_n$ nach Konstruktion ebenfalls beschränkt. Wir haben daher bei Ungleichung~(\ref{anc}) zudem verwendet, dass wir $n$ und $c_3$ so groß wählen, dass $\max\{\|m\|_{\infty}, \|m_n\|_{\infty}\} < \beta_n$ gilt. Bei der letzten Gleichung haben wir schließlich verwendet dass $\beta_n$ deterministisch und $\mathds{P}(X \in \supp(\mathds{P}_X)) = 1$ ist.
Durch unsere Definition von $\hat{m}_n$ erhalten wir durch die Monotonie des Erwartungswerts und einer Abschätzung über den ganzen Raum:
\begin{equation}
\label{an}
\begin{split}
\E \bigg[\int |m_n(x) - m(x)|^2 \mathds{P}_X(dx) \mathds{1}_{A_n}\bigg] & = \E \bigg[\int |\hat{m}_n(x) - m(x)|^2 \mathds{P}_X(dx) \mathds{1}_{A_n}\bigg] \\
& \leq \E \bigg[\int |\hat{m}_n(x) - m(x)|^2 \mathds{P}_X(dx)\bigg].
\end{split}
\end{equation}
Zusammen mit (\ref{tscheby}), (\ref{anc}), (\ref{an}) und der Linearität des Erwartungswerts erhalten wir dann:
\begin{equation}
\label{originaleq}
\begin{split}
\E \int |m_n(x) - m(x)|^2 \mathds{P}_X(dx) & = \E \bigg[\int |m_n(x) - m(x)|^2 \mathds{P}_X(dx) \cdot (\mathds{1}_{A_n^{\mathsf{c}}} + \mathds{1}_{A_n})\bigg] \\
& = \E \bigg[\int |m_n(x) - m(x)|^2 \mathds{P}_X(dx) \mathds{1}_{A_n^{\mathsf{c}}}\bigg] \\
& \quad + \E \bigg[\int |m_n(x) - m(x)|^2 \mathds{P}_X(dx) \mathds{1}_{A_n}\bigg] \\
& \leq \frac{4 \cdot c_7 \cdot \beta_n^2}{n} + \E \bigg[\int |\hat{m}_n(x) - m(x)|^2 \mathds{P}_X(dx)\bigg].
\end{split}
\end{equation}
Wir zeigen nun $\tilde{m}_n \in \mathcal{F}^{(J_n)}$.

Nach Gleichung~(\ref{umschreiben}) können wir unseren Schätzer $\tilde{m}_n$ darstellen durch:
$$\tilde{m}_n(x) = \sum_{j = 1}^{J_n}\hat{a}_j \cdot f_j$$
für geeignete $f_j \in \mathcal{F}$ und $\hat{a}_j$ welche 
\begin{equation*}
\begin{split}
\frac{c_6}{n}\sum_{j = 1}^{J_n} \hat{a}_j^2 & = \frac{c_6}{n} \sum_{\bi \in [M]^d} \sum_{\substack{\bj \in [q]^d \\ |\bj|_1 \leq q}} a_{\mathbf{i},\bj}^2 \\
& \leq \frac{1}{n} \sum_{i = 1}^n|Y_i - \tilde{m}_n(X_i)|^2 + \frac{c_6}{n} \sum_{\bi \in [M]^d} \sum_{\substack{\bj \in [q]^d \\ |\bj|_1 \leq q}} a_{\mathbf{i},\bj}^2 \\
& \leq \sum_{i = 1}^n Y_i^2,
\end{split}
\end{equation*}
erfüllen, wobei wir bei der letzten Ungleichung wie in Kapitel~\ref{chap:2} die minimierende Eigenschaft von $a_{\mathbf{i},\bj}$ verwendet haben und zum Schluss die Koeffizienten Null gesetzt haben. Da $c_6 > 0$ ist, erhalten wir dass die Koeffizienten $\hat{a}_j$ die Eigenschaft
$$\sum_{j = 1}^{J_n} \hat{a}_j^2  \leq \frac{1}{n}\sum_{i = 1}^n Y_i^2 \cdot \frac{n}{c_6}$$
erfüllen müssen.
Aus der Definition des Ereignissen $A_n$ in Gleichung~(\ref{event}) und der Definition der Konstante $c_5$ in Gleichung~(\ref{constantc21}) erhalten wir dann die Abschätzung
$$\sum_{j = 1}^{J_n}\hat{a}_j^2 \leq \frac{1 + \E[Y^2]}{c_6} \cdot n \leq c_5 \cdot n,$$
woraus durch $f_j \in \mathcal{F}$ schließlich $\tilde{m}_n \in \mathcal{F}^{(J_n)}$ folgt.

Daher setzen wir $\hat{m}_n \coloneqq T_{\beta_n}\bar{m}_n$ für $\bar{m}_n \in \{\tilde{m}_n, g_n\} \subseteq \mathcal{F}^{(J_n)}$.
Die Funktionen $\tilde{m}_n$ und $g_n$ unterscheiden sich in den Vorfaktoren von $f_j$. Die Koeffizienten $a_{\mathbf{i},\bj}$ von $\tilde{m}_n$ haben wir durch Minimierung des Funktionals in Gleichung~(\ref{min}) erhalten. Nach Voraussetzung ist $N \geq q$ und damit gilt dann insbesondere $\{0,\dots,q\} \subseteq \{0,\dots,N\}$. Bei der Minimierung des Funktionals hat man daher auch die Koeffizienten von $g_n$ betrachtet. Daher erhalten wir:
\begin{equation}
\begin{split}
\label{tilde}
& \frac{1}{n} \sum_{i = 1}^n|Y_i - \tilde{m}_n(X_i)|^2 \\
& \leq \frac{1}{n} \sum_{i = 1}^n|Y_i - \tilde{m}_n(X_i)|^2 + \frac{c_6}{n} \cdot \sum_{\bi \in [M_n]^d} \sum_{\substack{ \bj \in [N]^d \\ |\bj|_1 \leq N}} a_{\mathbf{i},\bj}^2 \\
& \leq \frac{1}{n} \sum_{i = 1}^n|Y_i - g_n(X_i)|^2 + \frac{c_6}{n} \cdot \sum_{\bi \in [M_n]^d} \sum_{\substack{ \bj \in [q]^d \\ |\bj|_1 \leq q}} \bigg|\frac{1}{\bj!} \cdot \partial^{\bj}m(x_{\mathbf{i}})\bigg|^2 \\
& = \frac{1}{n} \sum_{i = 1}^n|Y_i - g_n(X_i)|^2 + c_{8} \cdot \frac{(M_n + 1)^d}{n},
\end{split}
\end{equation}
mit $c_{8} = c_6 \cdot (q + 1)^d \cdot \big|\frac{1}{\bj!} \cdot \partial^{\bj}m(x_{\mathbf{i}})\big|^2$ als Konstante die unabhängig von $n$ ist.
Wir erhalten damit:
\begin{equation}
\label{bed}
\frac{1}{n} \sum_{i = 1}^n|Y_i - \bar{m}_n(X_i)|^2 \leq \frac{1}{n} \sum_{i = 1}^n|Y_i - g_n(X_i)|^2 + c_8 \cdot \frac{(M_n + 1)^d}{n},
\end{equation}
da für $\bar{m}_n = g_n$ die Ungleichung unmittelbar folgt.

Da $g_n$ nach Definition deterministisch, damit also unabhängig von $(X_1, Y_1), (X_2, Y_2),\dots$ ist, sind mit der einelementigen Parametermenge $\Theta_n$, der Funktion $g_{n,1} = g_n$, der Abschätzung~(\ref{bed}) für $\hat{m}_n = T_{\beta_n}\bar{m}_n$ mit $\hat{m}_n \in \mathcal{F}^{(J_n)}$ und dem \emph{penalty Term} $\text{pen}_n(g_{n,1}) = c_8 \cdot \frac{(M_n + 1)^d}{n} > 0$ die Voraussetzungen für Lemma~\ref{lem:8} erfüllt. Wir erhalten durch dessen Anwendung:
\begin{equation*}
\begin{split}
& \E \int |\hat{m}_n(x) - m(x)|^2 \mathds{P}_X(dx) \\
& \leq \frac{c_8 \cdot \log(n)^2 \cdot \big(\log\big(\sup_{x_1^n \in (\supp(X))^n}\mathcal{N}_1\big(\frac{1}{n \cdot \beta_n},\mathcal{F}^{(J_n)},x_1^n\big)\big) + 1\big)}{n} \\
& \quad + 2 \cdot \E\bigg(\min_{l \in \Theta_n} \int |g_{n,l}(x) - m(x)|^2 \mathds{P}_X(dx) + c_8 \cdot \frac{(M_n + 1)^d}{n}\bigg) \\
& = \frac{c_8 \cdot \log(n)^2 \cdot \big(\log\big(\sup_{x_1^n \in (\supp(X))^n}\mathcal{N}_1\big(\frac{1}{n \cdot \beta_n},\mathcal{F}^{(J_n)},x_1^n\big)\big) + 1\big)}{n} \\
& \quad + 2 \int |g_{n}(x) - m(x)|^2 \mathds{P}_X(dx) + 2 \cdot c_8 \cdot \frac{(M_n + 1)^d}{n},
\end{split}
\end{equation*}
wobei wir bei der letzten Gleichheit verwendet haben, dass der letzte Summand deterministisch ist. Zudem wissen wir, dass $c_8$ unabhängig von $n$ ist und $n > 1$, da wir $n$ hinreichend groß wählen.
Als nächstes überprüfen wir die Voraussetzungen von Lemma~\ref{lem:9}, um damit dann die letzte Gleichung mit der Überdeckungszahl $\mathcal{N}_1$weiter abzuschätzen.
Nach Voraussetzung ist $\beta_n = c_3 \cdot \log(n)$ und $a_n = (\log n)^{1/(6(N + d))} > 0$ für hinreichend großes $n$. Nach Voraussetzung sind zudem $d, N, J_n \in \N$ und es gilt nach Gleichung~(\ref{jn}): 
$$J_n \leq (M_n + 1)^d \cdot (N + 1)^d \leq n^{c},$$
für hinreichend großes $n$ und Konstante $c > 0$. Wir betrachten hier den logistischen squasher $\sigma$ welcher nach Lemma~\ref{lem:logsquasher} insbesondere 2-zulässig ist. Da die hier betrachtete Menge von Funktionen $\mathcal{F}^{(J_n)}$ identisch mit der aus Lemma~\ref{lem:9} ist, sind nun alle Voraussetzungen für Lemma~\ref{lem:9} erfüllt. Nach Voraussetzung wissen wir, dass $\supp(\mathds{P}_X)$ beschränkt ist, und wir können $n$ so groß wählen, dass wir ohne Beschränkung der Allgemeinheit annehmen können, dass $\supp(\mathds{P}_X) = \{x \in \R^d \mid \forall \epsilon > 0 : \mathds{P}_X(S_{\epsilon}(x)) > 0\} \subseteq [-a_n, a_n]^d$ ist, mit $S_{\epsilon}$ als $\epsilon$-Umgebung um $x \in \R^d$. In der nächsten Ungleichungskette fassen wir die Konstanten die sich von $c_8$ unterscheiden und die unabhängig von $n$ sind, zu einer Konstante $c$ zusammen. Wir erhalten damit durch Lemma~\ref{lem:9} für hinreichend großes $n$:
\begin{equation}
\label{lem9sol}
\begin{split}
& \frac{c_8 \cdot \log(n)^2 \cdot \big(\log\big(\sup_{x_1^n \in (\supp(X))^n}\mathcal{N}_1\big(\frac{1}{n \cdot \beta_n},\mathcal{F}^{(J_n)},x_1^n\big)\big) + 1\big)}{n} \\
& \leq \frac{c_8 \cdot \log(n)^2 \cdot \big(\big(c \cdot \log(n) \cdot (M_n + 1)^d \cdot (N + 1)^d \big) + 1\big)}{n} \\
& \leq \frac{c_8 \cdot \log(n)^2 \cdot \big(2 \cdot c \cdot \log(n) \cdot (M_n + 1)^d \cdot (N + 1)^d\big)}{n} \\
& \leq c_{9} \cdot \frac{\log(n)^3 \cdot (M_n + 1)^d \cdot (N + 1)^d}{n},
\end{split}
\end{equation}
wobei $c_9 \coloneqq c_8 \cdot c$ eine von $n$ unabhängige Konstante ist.
Sei 
\begin{equation*}
\begin{split}
P_{m,n}(x) \coloneqq \sum_{\bi \in [M_n]^d} \sum_{\substack{ \bj \in [q]^d \\ |\bj|_1 \leq q}} \frac{1}{\bj!} \cdot \partial^{\bj} m (x_{\mathbf{i}}) \cdot (x - x_{\mathbf{i}})^{\bj} \prod_{j = 1}^d\Big(1 - \frac{M_n}{2a_n} \cdot |x^{(j)} - x_{\mathbf{i}}^{(j)}|\Big)_+. 
\end{split}
\end{equation*}
Mit Ungleichung~(\ref{ungl}), zusammen mit einer Nulladdition und der Linearität des Integrals erhalten wir:
\begin{equation*}
\begin{split}
& \int |g_n(x) - m(x)|^2 \mathds{P}_X(dx) \\
& = \int |g_n(x) - P_{m,n}(x) + P_{m,n}(x) - m(x)|^2 \mathds{P}_X(dx) \\
& \leq \int 2 |g_n(x) - P_{m,n}(x)|^2 + 2 |P_{m,n}(x) - m(x)|^2 \mathds{P}_X(dx).
\end{split}
\end{equation*}
Aus der Supremumseigenschaft folgt
\begin{equation}
\label{gnmx}
\begin{split}
& \int |g_n(x) - m(x)|^2 \mathds{P}_X(dx) \\
& \leq 2 \int \sup_{x \in [-a_n, a_n]^d} |g_n(x) - P_{m,n}(x)|^2 \mathds{P}_X(dx) + 2 \int \sup_{x \in [-a_n, a_n]^d} |P_{m,n}(x) - m(x)|^2 \mathds{P}_X(dx) \\
& = 2 \sup_{x \in [-a_n, a_n]^d} |g_n(x) - P_{m,n}(x)|^2 + 2 \sup_{x \in [-a_n, a_n]^d} |P_{m,n}(x) - m(x)|^2, 
\end{split}
\end{equation}
wobei wir im letzten Schritt $\supp(\mathds{P}_X) \subseteq [-a_n, a_n]^d$ und $\mathds{P}(X \in \supp(\mathds{P}_X)) = 1$ verwendet haben. Um die letzten beiden Summanden der Ungleichung~(\ref{gnmx}) weiterhin abzuschätzen möchten wir Lemma~\ref{lem:5} anwenden. Dafür überprüfen wir, ob dafür alle Voraussetzungen erfüllt sind. Wir betrachten wieder den logistischen squasher $\sigma$ aus Gleichung~\ref{logsquasher}, welcher nach Lemma~\ref{lem:logsquasher} insbesondere 2-zulässig ist. Zudem ist für hinreichend großes $n$ die Bedingung 
\begin{equation*}
\begin{split}
R_n & \geq \max\biggl\{\frac{\|\sigma''\|_{\infty} \cdot (M_n + 1)}{2 \cdot |\sigma'(t_{\sigma})|}, \frac{9 \cdot \|\sigma''\|_{\infty} \cdot a_n}{|\sigma'(t_{\sigma})|}, \\
& \quad \frac{20 \cdot \|\sigma'''\|_{\infty}}{3 \cdot |\sigma''(t_{\sigma})|} \cdot 3^{3 \cdot 3^s} \cdot a^{3 \cdot 2^s}, 1792 \cdot \frac{\max\{\|\sigma''\|_{\infty},\|\sigma'''\|_{\infty}, 1\}}{\min\{2 \cdot |\sigma'(t_{\sigma})|, |\sigma''(t_{\sigma})|, 1\}} \cdot M_n^3 \biggr\}
\end{split}
\end{equation*}
erfüllt und da unser neuronales Netz (\ref{fnet}) mit $x_{\mathbf{i}} \in [-a_n, a_n]^d$ identisch mit der Definition aus Lemma~\ref{lem:5} ist, sind alle Voraussetzungen für Lemma~\ref{lem:logsquasher} erfüllt.
Wir erhalten damit für $x \in [-a_n ,a_n]^d$ und $n$ hinreichend groß:
\begin{equation}
\label{1stsum}
\begin{split}
& |g_n(x) - P_{m,n}(x)| \\
& = \bigg| \sum_{\bi \in [M_n]^d} \sum_{\substack{ \bj \in [q]^d \\ |\bj|_1 \leq q}} \frac{1}{\bj!} \cdot \partial^{\bj} m (x_{\mathbf{i}}) \bigg| \cdot \bigg| f_{\net,\bj,\mathbf{i}}(x) - (x - x_{\mathbf{i}})^{\bj} \cdot \prod_{j = 1}^d(1 - \frac{M_n}{2a_n} \cdot |x^{(j)} - x_{\mathbf{i}}^{(j)}|)_+ \bigg| \\
& \stackrel{(\ref{constantc21})}{\leq} (M_n + 1)^d \cdot (q + 1)^d \cdot z \cdot \bigg| f_{\net,\bj,\mathbf{i}}(x) - (x - x_{\mathbf{i}})^{\bj} \cdot \prod_{j = 1}^d(1 - \frac{M_n}{2a_n} \cdot |x^{(j)} - x_{\mathbf{i}}^{(j)}|)_+ \bigg| \\
& \leq (M_n + 1)^d \cdot (q + 1)^d \cdot z \cdot c \cdot 3^{3 \cdot 3^s} \cdot a_n^{3 \cdot 2^s} \cdot M_n^3 \cdot \frac{1}{R_n} \\
& \leq (M_n + 1)^d \cdot (q + 1)^d \cdot c\cdot a_n^{3 \cdot (N + d) \cdot 2} \cdot \frac{M_n^3}{R_n} \\
& =  (M_n + 1)^d \cdot (q + 1)^d \cdot c\cdot \log(n) \cdot \frac{M_n^3}{R_n},
\end{split}
\end{equation}
wobei wir in $c$ alle von $n$ unabhängige Konstante zusammenfassen, welche insbesondere die Konstante aus Lemma~\ref{lem:logsquasher} enthält.
Wie haben unter anderem verwendet, dass für hinreichend großes $n$:
$$a_n^{2^{\lceil\log_2(N + d)\rceil}} \leq a_n^{2^{\log_2(N + d) + 1}} = a_n^{(N + d) \cdot 2},$$ gilt. Im letzten Schritt haben wir in Ungleichung~(\ref{1stsum}) die Definition von $a_n$ eingesetzt. 
Da nach Konstruktion $a_n > 0$, $m$ $(p, C)$-glatt und $P_{m,n}(x)$ nach Lemma~\ref{lem:loccon} eine Spline Interpolation von Taylorpolynomen von $m$ ist, erhalten wir mit Lemma~\ref{lem:pcsmooth}:
\begin{equation}
\label{2ndsum}
\begin{split}
|P_{m,n}(x) - m(x)| \leq c_{11} \cdot \frac{a_n^p}{M_n^p} \leq c_{11} \cdot \log(n) \cdot \frac{1}{M_n^p}.
\end{split}
\end{equation}
In dieser Ungleichung ist $c_{11} > 0$ eine von $n$ unabhängige Konstante aus Lemma~\ref{lem:pcsmooth} und wir haben zudem verwendet, dass:
$$a_n^p = a_n^{q + s} \leq a_n^{N + d} \leq a_n^{6 \cdot (N + d)} = \log(n),$$
für $p = q + s$ für hinreichend großes $n$ gilt, da nach Voraussetzung $N \geq q$ und $d \geq s$ mit $s \in (0, 1]$ ist.
Durch Quadrieren bleiben die Ungleichungen~(\ref{1stsum}) und (\ref{2ndsum}) auch erhalten und da in beiden Ungleichungen die rechte Seite unabhängig von $x$ ist, gelten die Ungleichungen ebenfalls für das Supremum über $x \in [-a_n,a_n]^d$. Durch Einsetzen der Definitionen von $M_n$ und $R_n$ erhalten wir für $n$ hinreichend groß:
\begin{equation}
\label{3rdsum}
\begin{split}
\sup_{x \in [-a_n, a_n]^d} |g_n(x) - P_{m,n}(x)|^2 & \leq \bigg((M_n + 1)^d \cdot (q + 1)^d \cdot c \cdot \log(n) \cdot \frac{M_n^3}{R_n}\bigg)^2 \\
& \leq c^2 \cdot (M_n + 1)^{2d} \cdot \log(n)^2 \cdot \frac{M_n^6}{R_n^2} \\
& \leq  c_{10}^2 \cdot (M_n + 1)^{2d} \cdot \log(n)^2 \cdot \frac{(M_n + 1)^{6d}}{R_n^2} \\
& \leq  c^2 \cdot \log(n)^2 \cdot \frac{(M_n + 1)^{8d}}{R_n^2} \\
& \leq c_{10} \cdot \frac{n^{\frac{8d}{2p + d}}}{n^{2d + 8}} \cdot \log(n)^2 \\
& = c_{10} \cdot n^{\frac{8d}{2p + d} - 2d -8} \cdot \log(n)^2 \\
& \leq c_{10} \cdot n^{\frac{8d}{2p + d}  -8\frac{2p + d}{2p + d}} \cdot \log(n)^2 \\
& = c_{10} \cdot n^{-\frac{16p}{2p + d}} \cdot \log(n)^2 \\
& \leq c_{10} \cdot n^{-\frac{2p}{2p + d}} \cdot \log(n)^3,
\end{split}
\end{equation}
mit $c_{10} = c^2 \cdot (2c_2)^{8d} > 0$ und unabhängig von $n$ ist. Bei der letzten Ungleichung haben wir verwendet, dass $\frac{16p}{2p + d} > \frac{2p}{2p + d}$, da $p > 0$ ist und $\log(n)^2 < \log(n)^3$ für $n$ hinreichend groß gilt. Ebenfalls erhalten wir:
\begin{equation}
\label{4thsum}
\begin{split}
\sup_{x \in [-a_n, a_n]^d} |P_{m,n}(x) - m(x)|^2 & \leq \bigg(c_{11} \cdot \log(n) \cdot \frac{1}{M_n^p}\bigg)^2 \\
& \leq c_{11}^2 \cdot \log(n)^2 \cdot c_{2}^{-2p} \cdot n^{-\frac{2p}{2p + d}} \\
& \leq \bigg(\frac{c_{11}}{c_2^p}\bigg)^2 \cdot \log(n)^3 \cdot n^{-\frac{2p}{2p + d}}.
\end{split}
\end{equation}
Mit analogem Vorgehen erhalten wir für Ungleichung~(\ref{lem9sol}):
\begin{equation}
\begin{split}
& \frac{c_8 \cdot \log(n)^2 \cdot \big(\log\big(\sup_{x_1^n \in (\supp(X))^n}\mathcal{N}_1\big(\frac{1}{n \cdot \beta_n},\mathcal{F}^{(J_n)},x_1^n\big)\big) + 1\big)}{n} \\
& \leq c_9 \cdot \frac{\log(n)^3 \cdot (M_n + 1)^d \cdot (N + 1)^d}{n} \\
& \leq c_9 \cdot (N+1)^d \cdot \log(n)^3 \cdot \frac{(c_2 \cdot n^{\frac{1}{2p + d}} + 2)^d}{n} \\
& \leq c_9 \cdot (N+1)^d \cdot 2^d \cdot \log(n)^3 \cdot \frac{c_2^d\cdot n^{\frac{d}{2p + d}}}{n} \\
& = c_9 \cdot (N+1)^d \cdot 2^d \cdot c_2^d \cdot \log(n)^3 \cdot n^{\frac{d}{2p + d} - 1} \\
& = c_{12} \cdot \log(n)^3 \cdot n^{-\frac{2p}{2p + d}},
\end{split}
\end{equation}
mit einer von $n$ unabhängigen Konstante $c_{12} \coloneqq c_9 \cdot (N + 1)^d \cdot 2^d \cdot c_2^d > 0.$
Zudem erhalten wir, da für $n$ hinreichend groß $\log(n)^3 > 1$ gilt, mit analogem Vorgehen:
\begin{equation}
\label{5thsum}
\begin{split}
c_8 \cdot \frac{(M_n + 1)^d}{n} & \leq c_{13} \cdot \frac{n^{\frac{d}{2p + d}}}{n} \\
& \leq c_{13} \cdot \log(n)^3 \cdot n^{- \frac{2p}{2p + d}} 
\end{split}
\end{equation} 
und 
\begin{equation}
\label{6thsum}
\begin{split}
\frac{4 \cdot c_7 \cdot \beta_n^2}{n} & \leq c_{14} \cdot \log(n)^3 \cdot n^{-1} \\
& \leq c_{14} \cdot \log(n)^3 \cdot n^{- \frac{2p}{2p + d}},
\end{split}
\end{equation} 
mit von $n$ unabhängigen Konstanten $c_{13} \coloneqq c_8 \cdot (2c_2)^d > 0$ und $c_{14} \coloneqq 4 \cdot c_7 \cdot c_3^2 > 0$.
Nun haben wir alle Summanden von Ungleichung ($\ref{originaleq}$) abgeschätzt und erhalten schließlich mit Ungleichungen (\ref{1stsum}) - (\ref{6thsum}):
\begin{equation*}
\E \int |m_n(x) - m(x)|^2 \mathds{P}_X(dx) \leq c_{\text{fin}} \cdot \log(n)^3 \cdot n^{- \frac{2p}{2p + d}},
\end{equation*}
mit 
$$c_{\text{fin}} = c_{14} + 2 \cdot c_{13} + c_{12} + 2 \cdot \Bigg(2 \cdot \bigg(\frac{c_{11}}{c_2^p}\bigg)^2 + 2 \cdot  c_{10}\Bigg),$$
wobei $c_{\text{fin}}$ als Summe nichtnegativer und positiver Konstanten, die unabhängig von $n$ sind, nichtnegativ und unabhängig von $n$ ist.
Damit haben wir unser Hauptresultat bewiesen.
\end{proof}
%\chapter{Anwendungsbeispiel auf simulierte Daten}
\label{chap:4}

In diesem Kapitel betrachten wir die Leistung unseres Neuronale-Netze"=Regressionsschätzers aus Kapitel~\ref{chap:2}, bei endlicher Stichprobengröße auf simulierte Daten in $Python$~\cite[Version 3.7.3]{van1995python} mithilfe der $Keras$ Bibliothek \cite{chollet2015keras}.

Die simulierten Daten, welche wir verwenden werden, sehen wie folgt aus:
Wir wählen $X$ gleichverteilt auf $[-2, 2]^d$, wobei $d$ die Dimension des Inputs ist, zudem wählen wir $\epsilon$ standardnormalverteilt und unabhängig von $X$ und wir definieren $Y$ durch:
$$Y = m_j(X) + \sigma \cdot \lambda_j \cdot \epsilon,$$ 
wobei $m_j \colon [-2, 2]^d \to \R$ für  $j \in \{1, 2\}$ den Funktionen:
$$ m_1(x) =  \sin\big(0.2 \cdot x^2\big) + \exp(0.5 \cdot x) + x^3$$
und
$$ m_2(x_0, x_1) = \sin\big(\sqrt[2]{x_0^2 + x_1^2}\big)$$
entspricht.
Den Skalierungsfaktor $\lambda_j > 0$ wählen wir als Interquartilsabstand \emph{(IQR)} einer Stichprobe der Größe $N = 800$ von $m(X)$. Für den Rauschfaktor $\sigma$ gilt $\sigma \in \{0.05, 0.1\}.$ Mit diesen Daten lässt sich nun auch $Y$ darstellen.

Um die Leistung unseres Neuronale-Netze-Regressionsschätzers zu überprüfen, haben wir erstmals die Funktion $m_1$ und die jeweilige Schätzung durch unseren Neuronale-Netze-Regressionsschätzer für $\sigma = 0.05$ zeichnen lassen. 
Visuell kann man mit Abbildung~\ref{fig:subfig1.a.4} erkennen, dass der Schätzer eine sehr gute Approximation der Funktion liefert, aber um genauer beurteilen zu können wie gut die Schätzung wirklich ist und wie gut unser Schätzer im Vergleich zu anderen Schätzern abschneidet, betrachten wir in Tabelle~\ref{tab:truthTablesm1} und Tabelle~\ref{tab:truthTablesm2} den Interquartilsabstand und den Median des skalierten $L_2$-Fehlers der einzelnen Schätzer. 

Unser Vorgehen zum Vergleich der drei hier betrachteten Regressionsschätzer gestaltet sich wie folgt:
Wir bestimmen als Erstes den $L_2$-Fehler der einzelnen Schätzer $m_{n,i}$ für $i = 1,2,3$ approximativ durch den empirisch arithmetischen $L_2$-Fehler $\epsilon_{L_2,N}(m_{n,i})$ auf einer unabhängigen Stichprobe von $X$ der Größe $N = 1000$. 

Da wir unsere Regressionsfunktionen kennen und der Fehler stark vom Verhalten der Funktion $m_j$ abhängt, betrachten wir den empirischen $L_2$-Fehler im Verhältnis zum simpelsten Schätzer von $m_j$, einer konstanten Funktion. Der Wert dieses \textit{konstanten Schätzers} bestimmen wir, in dem wir das empirische Mittel der beobachteten Daten $Y$ nehmen. Wir erhalten damit ein skaliertes Fehlermaß $\epsilon_{L_2,N}(m_{n,i})/\bar{\epsilon}_{L_2,N}(avg)$ mit $\bar{\epsilon}_{L_2,N}(avg)$ als Median von $50$ unabhängigen Realisierungen von $\epsilon_{L_2,N}$. 

Wir bestimmten $\epsilon_{L_2,N}(\cdot)$ als empirisches Mittel von $25$ quadratischen Fehlern der Schätzung des konstanten Schätzers. Dieses skalierte Fehlermaß ist so zu deuten, dass ein großer Fehler durch einen der drei Regressionsschätzer im Falle, dass der Fehler des konstanten Schätzers klein ist, auf eine noch schlechtere Leistung hindeutet. Der Fehler $\epsilon_{L_2,N}(m_{n,i})$ wird also durch $\bar{\epsilon}_{L_2,N}(avg)$ gewichtet.

Die resultierenden skalierten Fehler hängen noch von der Stichprobe von $(X, Y)$ ab und um diese Werte besser vergleichen zu können, führen wir die Fehlerberechnung jeweils $50$ mal durch und geben dann den Median und den Interquartilsabstand für die Schätzung der betrachteten Regressionsschätzer aus.
Wir teilen für jeden Schätzer die Stichprobe in ein \textit{learning sample} der Größe $n_l = 0.8 \cdot N$ und in ein \textit{testing sample} der Größe $n_t = 0.2 \cdot N$ auf. Wir bestimmen die Parameterwerte für die Schätzer auf dem learning sample und bestimmen das korrespondiere $L_2$-Risiko auf dem testing sample.

Unser erster Schätzer \textit{fc\_neural\_1\_estimate} ist ein neuronales Netz mit einer verborgenen Schicht. Dieser Schätzer hat eine feste Anzahl an Neuronen die wir aus der Menge $\{5, 10, 25, 50, 75\}$ auswählen die bei der Simulation zu einem minimalen empirischen $L_2$-Risiko führt.

Unser zweiter Schätzer \textit{nearest\_neighbor\_estimate} ist ein Nächste-Nachbar-Schätzer bei der die Anzahl an nächsten Nachbarn so aus der Menge $\{1, 2, 3\} \cup \{4, 8, 12, 16, \dots, 4 \cdot \lfloor\frac{n_l}{4}\rfloor\}$ ausgewählt wird, dass dieser zu minimalen empirischen $L_2$-Risiko führt.

Unser letzter Schätzer ist der in Kapitel~\ref{chap:2} vorgestellte Neuronale-Netze-Regressionsschätzer den wir mit \textit{new\_neural\_network\_estimate} bezeichnen. Hier haben wir die Parameter, je nachdem, welche Regressionsfunktion wir betrachtet haben, entsprechend angepasst, da wir zum Beispiel den Grad der zu schätzenden Funktion kennen und dies bei unserem Schätzer berücksichtigen können. Für die Schätzung von $m_1$ setzten wir: $d = 1$, $N = 3$, $q = 2$, $R = 10^6$, $a = 2$, $M = 2$.
Für die Schätzung von $m_2$ wählten wir: $d = 2$, $N = 2$, $q = 2$, $R = 10^6$, $a = 2$, $M = 2$.

\begin{figure}
    \begin{subfigure}[b]{0.5\textwidth}
        \centering
        \scalebox{0.9}{
          \begin{tikzpicture}
\begin{axis}[%
scatter/classes={%
a={mark=o,draw=black}},
legend style={fill opacity=0.8, draw opacity=1, text opacity=1, at={(0.03,0.97)}, anchor=north west, draw=white!80!black},
tick align=outside,
tick pos=left,
title={\quad},
x grid style={white!69.0196078431373!black},
xlabel={ },
xmin=-2.20852591900454, xmax=2.21306049239568,
xtick style={color=black},
y grid style={white!69.0196078431373!black},
ylabel={ },
ymin=-7.86405361849716, ymax=12.4391808546764,
ytick style={color=black}
]
\addplot[blue, mark=*,only marks,mark size=1]%
table{
x y 
1.98342671961567 11.2067083316341
-0.164883410431491 0.921819789666956
1.40142299960581 5.15033590554207
1.47929179834368 5.7561559124411
1.19715514411127 3.81799531027161
-1.05745151470862 -0.371309570367478
-0.696648319326161 0.464684877420012
1.82247761300943 9.15712037755566
1.33480759235371 4.67626070341997
0.862438587068488 2.32882713322422
-1.16503690445439 -0.754686954628516
0.417168238066338 1.33933090794867
1.15596892097842 3.59120483869643
0.850205757253619 2.28838577184768
-1.30994868597998 -1.39187573215822
-1.55346692657437 -2.82488317686155
1.61669575447566 6.96902026146192
0.500046589667835 1.45907874969249
1.28496886028727 4.34712794233107
-1.20227096727067 -0.904559178268851
0.366681884355046 1.27741420051878
0.818260589820523 2.18688479377424
-1.23845976488916 -1.05920098696907
0.697189507824902 1.8530215419852
1.93926361424016 10.6132554825254
1.04107639834298 3.02636787410901
-0.22353824693392 0.893074337960662
0.668246244317835 1.78431442012492
0.151928097424038 1.08704702436641
-1.07641959705448 -0.433765220730816
1.44298781246897 5.46662361852968
-0.284110387060426 0.860783387392703
-0.44017790192429 0.755901668314636
0.352464522063814 1.26134584362801
1.18574235128237 3.75381711403657
1.66572862979816 7.4485981217319
-1.77196434008901 -4.56390224034042
-1.42077197081339 -1.98366518029347
1.90156609340709 10.1254900733363
-0.399226079743066 0.787289288568677
-0.502642406493864 0.701288589370527
0.570484726981828 1.58079483029362
-0.626439688866295 0.563662303293445
0.646913187691053 1.73622883021686
-1.70011590874738 -3.94019890361793
0.188560902818183 1.11268318790154
0.613581309391728 1.66528403436379
1.53866334933713 6.25708686606659
-1.90745710621872 -5.88963818524853
0.41517370699518 1.33673477925145
0.327838052279217 1.23484812561957
1.73772873954887 8.19951335349887
-1.78984900196206 -4.72748294080677
1.52854808178337 6.16929527049623
0.00832534691972997 1.00418578874098
1.53059984449369 6.18702069900853
-1.45719188829622 -2.19959638868634
0.291257080120115 1.19843927737529
1.19523677300415 3.8071354073554
1.13842913470803 3.49861808490207
1.2366636859145 4.04822570623427
1.89871430519567 10.0892868982552
0.019953194670169 1.01011409945552
1.82122210851877 9.14233718077853
-0.298230517123074 0.852732115538147
0.382945539678767 1.29651503281945
1.08047569362621 3.20916222764976
-0.100962219519459 0.951781416117594
-0.409006788513409 0.780081708579802
1.41020530945337 5.21583204082174
-0.862424320245129 0.156477928786038
-1.98328806650409 -6.72214450231026
-0.696177128652387 0.465406125702621
-0.0933688730806197 0.955318102465428
1.94870249960763 10.73808188061
-0.936123458958722 -0.01976601438959
-1.58312442994531 -3.03409351460663
1.4664267507308 5.65211917053288
1.57141891136512 6.54839610341355
-1.47051135974113 -2.28133185267719
0.874729152014026 2.37035416122943
1.03751084984896 3.01036865668898
1.48030997994332 5.76445750267437
-1.54541524782003 -2.76946312915615
0.552369097636586 1.54760978330716
-1.68042557562017 -3.77839948923686
-1.51061355679375 -2.53657210298722
-0.167528933067055 0.920559140081827
-0.787421015424107 0.310011987252715
1.6501030154187 7.29301335433562
-0.0814367755752747 0.960885772415491
0.0795951024473589 1.04237142562678
0.979738939834169 2.76334381388308
-0.380817770934829 0.800394418604914
0.400591024706831 1.31813698650762
-0.4141113404742 0.776250097560955
1.91403104307527 10.2848760382045
0.920870673865245 2.53445330225573
-1.30315603463426 -1.35866671389968
0.0275601242433647 1.01404829046996
-0.927503900072607 0.00222689952751143
0.0368974953304853 1.01894249538715
1.15977317392 3.61159817433105
0.514097432066963 1.4818167241679
1.17984224265819 3.72104216577677
1.58572120040197 6.6789936011996
-1.37713701802779 -1.73918304115634
-0.472661706237887 0.728589446354698
-1.95506901984573 -6.40446611244112
-1.1779484446276 -0.805616762687941
1.23713048801279 4.05102163809569
0.455615485804151 1.39192818075103
1.59249283776102 6.74155542309861
-0.503448403039978 0.700525301295737
-0.446140509435029 0.75105598593073
0.0961164638478107 1.05196738610277
1.85889095909979 9.59383788518873
0.961953261926548 2.69181699566452
-0.383322867187997 0.79864530138838
0.24213138507616 1.15461989372253
0.788212033700322 2.09669287956382
0.84070370731225 2.25757827157573
1.40004883629895 5.14015179344157
0.107393086597664 1.05870960384272
0.225480591396121 1.14097309250214
-1.67936142766505 -3.7697647231671
-0.0816940927046601 0.960765514872752
1.47773128825186 5.74345186910455
0.401713010033748 1.31954418061981
0.422412491054748 1.34621800253768
0.776700220186323 2.06346076723564
0.317989295994632 1.22470791713495
0.816536031039081 2.18157160974445
1.53962127691442 6.26545353631168
-1.29064269548569 -1.29838208894096
-1.47198666263898 -2.29047742640869
-1.18769903957274 -0.844807189958559
1.5617847179266 6.46159449041297
-0.137017676939132 0.934967618307204
0.10162263173636 1.05523925887493
1.42737547993898 5.34593888574134
-0.755076052164864 0.368829213132414
1.71133977541391 7.91773667028815
-0.4721794317973 0.729011711839228
1.89344301858784 10.0226243465763
0.867958702524959 2.34736672383903
-0.411724457536147 0.778047811177945
0.824378736637617 2.20587035898719
0.170515291427536 1.09977047759294
0.979189264724758 2.76110191698977
-0.897121719560843 0.0767896021307181
1.55034434863224 6.35973527019784
-1.67062801778883 -3.69932121498425
-1.24614106788035 -1.09319208537185
1.54412256063696 6.30489082553266
1.43769151607769 5.42542699235533
-1.6456012323085 -3.50158560570617
-0.635373759957425 0.551982650209956
0.0798070563154254 1.04249250741436
-2.21592582434482e-05 0.999988920530453
-1.56549856175875 -2.90878956620822
-1.85995182081513 -5.40180355705662
-0.744024022342477 0.387963590115266
1.95878262880175 10.8725914329787
1.52876440755115 6.17116216956417
0.252423090403078 1.16334903121843
-0.393738335259513 0.791257880640355
0.449465165244776 1.38318113689044
-1.68386371444268 -3.80637391755432
1.00910500968525 2.88606233001523
-1.12050763647973 -0.587298253264637
0.230997555210016 1.14543091365568
0.0349769356316969 1.01792975394193
-1.56076894216133 -2.87564939931168
-1.45121512537244 -2.16340493618994
-1.16545297771441 -0.756311132410629
-0.0191205019479455 0.990551431326679
-1.86008081192998 -5.4030939017934
-0.436428121653843 0.758911768315441
1.49957184981221 5.9233919160222
-1.44219942190571 -2.10937609184844
1.01529613831035 2.91268092874956
-1.50140183330494 -2.47671254155102
-0.732092139869074 0.408086177914741
-1.31471126281736 -1.41536454673876
-1.16144104543941 -0.740697216143942
-0.689196716081129 0.475997455288745
-0.217022486630988 0.896366930296396
1.62508170186513 7.04925029897457
-0.0406198149252917 0.980157921905294
1.63119848795736 7.1082339788836
0.261617455458398 1.1713441338436
-0.0226785485704886 0.988815972333353
1.93031313461916 10.4958902043135
0.10579027870623 1.05774135817833
-1.34359720927249 -1.56148709026429
1.08272594869317 3.21993898042956
-1.13811010512611 -0.651956813151084
0.157670660808171 1.09091781606985
0.693694852953544 1.84450954897994
-0.413991721747403 0.776340444640307
-0.3496582994086 0.821300655555392
-1.494906796828 -2.43494997466264
0.0812226478153182 1.04330249998073
-1.87946646412891 -5.59912494096078
0.135232078225477 1.07608502949923
-0.0314802740817721 0.984550095861464
-0.536661666154805 0.667662425434589
0.939570817014385 2.6047405093355
1.9606433016915 10.8975565658301
-1.31276353932712 -1.4057380074589
-1.30676601901061 -1.3762731425726
-0.303871457144716 0.849451172880855
-0.643511730589387 0.541119017871324
-0.527466345137777 0.677041842362992
-0.522841202805422 0.681676958703923
0.4944545179478 1.45023449606521
-1.70114092695177 -3.94872752787689
-1.17864666831437 -0.808402179553694
1.70768647325131 7.87932950534975
1.57303031948494 6.56300623145007
0.553581378198257 1.54978869652802
-1.49462820991928 -2.43316685250796
0.0532637575658805 1.02770819278823
-1.03292019870819 -0.293649505337755
-0.0539046615190384 0.973832153044759
-1.4019635634533 -1.87640980895707
0.867566356157913 2.34604301414628
-1.05398034586825 -0.360109127135227
-1.13058478157364 -0.624075386023385
0.588011197017629 1.61419683576319
0.883913094472204 2.40197685981143
1.7032464005838 7.83284702573416
0.723806315564655 1.91984629454512
1.35906175870283 4.84427052411395
-1.23544243096618 -1.04596172691295
-1.54329111912356 -2.75493982502582
1.16098779424438 3.61813290630781
-1.87095954053839 -5.51258428050454
-1.2782629197376 -1.23986916511837
-1.05654172495337 -0.368367143763051
1.24353360832467 4.08955390230973
1.19028282951638 3.7792262679224
-0.775725735545598 0.33177186288789
1.02057273878203 2.93557518071213
0.821921250495339 2.19821881120506
0.918229962933942 2.52470603700035
-0.984864184033527 -0.151361080185913
1.93401749117918 10.5443463915569
1.83315978839409 9.28363700124014
0.533232250852013 1.51399333097939
-0.722014352421549 0.424656658899597
0.481969685265481 1.43090303239757
1.17847167126153 3.71346785213066
-1.79048490297368 -4.73336148187232
-0.945697114342179 -0.0446368331238406
1.60064007186754 6.81744991856981
1.93687669965452 10.5818612825148
1.35197568049207 4.79464642514883
1.91147516287336 10.2520425549434
-0.00460781786329045 0.997702891593491
-1.39948950363334 -1.86251399232994
1.04873511695863 3.06103594471185
0.000913924775346242 1.00045723462595
-0.50513266090352 0.69892519214253
-1.26704247379017 -1.18779508918314
1.76648926011692 8.515401523615
0.157856915922508 1.09104424813471
1.40787527714754 5.19838618820376
-0.240710854768439 0.88424614078834
-1.50153204953699 -2.47755356296565
-1.98239180697612 -6.71190905798603
-0.274480732689117 0.866148734835799
-1.63517535024967 -3.42099958712273
0.960111664862905 2.68452965066336
-0.166975897002565 0.920822887277161
1.54394939093015 6.30336988746059
-1.48090289946801 -2.34614449639195
1.23336476648575 4.02851752917348
-0.400362192403418 0.786460973883875
-1.87413140169175 -5.54475667677901
0.336611711534724 1.24409863604828
1.18872743268058 3.77050370068962
-1.7294971642558 -4.18886360380201
1.45013892943371 5.52266642043676
0.0983398967377997 1.05328401410495
1.82859103308664 9.22936432902555
-0.251038744131117 0.878821858955889
0.0752656353562249 1.03990925556896
0.0812009707648484 1.04329007921359
1.94778947733205 10.7259599229546
1.44701878178583 5.49815477228648
-0.327412551641523 0.835331216623138
0.243300205866964 1.15559974805297
0.092634846159414 1.04991799090477
1.99004222629651 11.2976793178811
0.370987182526909 1.28239482533905
1.12746366204799 3.44192570040845
-0.88103525140069 0.11444499583188
-1.87311208155932 -5.53440536407914
1.53244608154239 6.2030062191922
-0.460133374851041 0.739391819239952
0.307337448869911 1.21402461928519
0.369331774016094 1.28047335921953
-0.234023743918722 0.887711030862566
1.04808634766716 3.05808318471249
-1.71866655486462 -4.09618344887114
-0.541554732849246 0.662580975842837
-1.48954836729529 -2.40077032291578
-0.0346183802921285 0.983037951949156
1.66262513718082 7.41748978740735
-0.0632924638352459 0.969396910959773
0.0976192149694759 1.05285658014762
1.64277341714372 7.22092429870893
-1.23064117143324 -1.02502582306785
-1.2119829712324 -0.945178959503772
1.82543703976219 9.19203877696444
-1.96283859036236 -6.49099090907211
-1.07753916455694 -0.437518252974556
0.401431131407247 1.3191902799471
1.11313051112005 3.36918797402718
-1.82353738645624 -5.04487061859584
-1.02282270803728 -0.262684033496767
1.34249571990793 4.72895418337375
-1.7423388101225 -4.30030535815532
0.574932133692535 1.58914867255579
0.356056594784311 1.26535121241934
1.85225865167069 9.5131424338517
-0.398757442351753 0.787630287557302
0.689987650483266 1.83554534383615
-0.293455390498176 0.855480226380854
-1.36443451041482 -1.67085494806716
1.35065286985791 4.78543213551069
1.01572198192603 2.91452148801636
-1.34644750916868 -1.57624963302143
0.934098979805452 2.58394372921836
-0.337375961698936 0.82913400048838
0.64371991587656 1.72921343065672
-1.12505418398717 -0.603812272410027
-0.0424605535851552 0.97927752676906
-1.03874715299427 -0.311782480194157
-0.934567390971795 -0.0157678258945499
1.06185524551419 3.1213904800115
-1.68614727733944 -3.82501880552636
-0.0939923135416958 0.955027594191503
0.186353265908155 1.1110726428412
-0.702833029741399 0.455143382722217
0.669467597030042 1.78713226709335
-0.989846282352082 -0.16552193374252
0.105206333004412 1.05738943070462
-0.327016013168481 0.835575044028545
0.687399875760051 1.82932776989665
1.84677034716651 9.4467555282193
1.08694083448046 3.24022372148457
1.91270842735118 10.2678755681039
-1.84876330693187 -5.29058235734343
0.140188037147681 1.07929462015681
0.340988226479485 1.24879088314535
1.63241684639954 7.12002930432158
-0.134435062793881 0.936176722358918
1.89935988024498 10.0974738907713
1.10456045506315 3.32642942521779
-0.669863931935491 0.504430082672366
1.20253788657643 3.848623510478
-0.476896000953684 0.72485957082047
-0.864160922436813 0.152624313681391
0.779743629721808 2.07217627840376
0.44847566998307 1.38178594372062
1.88317123466855 9.8936753259772
0.966481764673039 2.70983107632884
0.977908789987251 2.7558871775324
-1.52375847203518 -2.62327699740979
1.13406912530685 3.4759673999773
0.522353840992912 1.49552708713273
-1.06822671124227 -0.406527917387101
-0.114355250680581 0.945546284394021
1.89308302690468 10.0180838654405
0.845248577878253 2.27224791568968
0.122425695926454 1.06795770242727
0.127977871489443 1.071452329096
0.01625186865687 1.00821615648482
-1.08796984143618 -0.472845104422563
-0.320666677251147 0.839450527601591
0.334320179485394 1.24166256257872
-1.05886861846269 -0.375902385164731
-0.564026324025872 0.638414711051878
-1.59115775979802 -3.0921520343033
-0.624686875303529 0.565923843870914
0.0413789150922659 1.02131826057146
1.31657491972719 4.55336406862912
-0.726992150093075 0.416520105094092
0.0524606842524733 1.02727218818635
-0.492141039700758 0.71109076518107
0.978275827209485 2.75738082493985
1.14130644500639 3.51364516032462
0.0392796360480538 1.02020313016701
1.30721738204172 4.49141064784605
0.490133920204322 1.44348012355915
-0.706442978901193 0.4495096644889
-0.993964710262591 -0.177329563835979
0.943620152780988 2.62025423675144
-0.0932557921766692 0.955370798670078
1.6365425027311 7.16008680515812
0.973324880119798 2.73730832363678
-0.725054451856926 0.419698555277889
-1.96417313354495 -6.50592462133396
1.31324497867966 4.53123083543522
1.87324432504556 9.77024625410757
0.303189087378309 1.20994224416354
1.9389374422235 10.6089613728723
0.501357612206841 1.46116901188386
-1.720056540905 -4.1080112039577
1.31781128519504 4.56160630738453
0.0951516110926338 1.05139794858651
0.408192815810979 1.32774813968817
0.513811453765586 1.48134650425401
-1.23574332705072 -1.04727912774199
-1.45264147142963 -2.17201479615298
-0.227918340918484 0.890843932569515
0.768794189473504 2.04105392984846
-0.383261040749235 0.79868859764658
-1.52535647365474 -2.63392131109938
0.978901900798691 2.75993067808586
0.989417826238432 2.80315006002761
-1.36387886985462 -1.66789491236366
1.30076519383918 4.44913205476502
-0.61920844932297 0.572929632714456
1.33769433552972 4.69598515376487
-0.74829966725375 0.380618179060386
1.65991412952015 7.3903995937643
-1.54882463937474 -2.79285887150011
-1.89409203935268 -5.74981453952136
0.116957207902192 1.06455795146663
1.17623321690576 3.70112896015788
-1.35372154548933 -1.61420715774001
-1.88141472578473 -5.61905888338989
1.32924807506156 4.63847962469226
-1.77134487512033 -4.55829675676622
-0.128741369799476 0.938838418393211
0.51095333379939 1.4766641046124
-1.67747670266001 -3.75449901365524
1.29669044236686 4.42261577546358
1.34952912501012 4.77761666284377
1.33266062176777 4.66163835456556
1.12465453846378 3.42754849841394
0.4823513930431 1.43148567258717
1.56244015327284 6.46747005636869
0.11081670699023 1.06078907799291
1.47420921577403 5.71486476538048
-1.40118206014657 -1.87201510170116
0.634664924286693 1.70957467919594
-0.684074619393981 0.483658496059001
1.60007325210646 6.81214767306192
0.871416435428499 2.35907231047904
-0.983390175523971 -0.147197135538043
-0.681323886845001 0.487734387895415
-1.26150283611157 -1.16241909746674
1.67379753998392 7.52995968901337
-0.987930545268803 -0.160060852020812
0.454921140058064 1.39093417506421
1.51394338581359 6.04432436735546
-0.804954691722807 0.276317438492785
-0.830724433208845 0.224398470191638
0.335858029439155 1.24329584466934
-0.159847298344431 0.924212755588489
0.0812119985668653 1.04329639798377
-0.117494185355901 0.944084196411911
0.497555841122187 1.45512515695177
0.640924836877307 1.72311134065059
-0.123471924446762 0.941297774124052
1.15598732778889 3.59130324256049
-1.27272624906405 -1.21406008321854
-0.76949998857843 0.343125991938706
1.4162089827075 5.26101451251004
-1.91041766129389 -5.92088694414514
1.99504962001563 11.3668989718006
-1.00490416987749 -0.209143338340367
-1.14258809682261 -0.668718416842979
1.18209823855401 3.73354181514181
0.128884868780194 1.07202738126087
-1.94570137893154 -6.30108914105676
0.27471991687839 1.18306806252273
-1.72671969673579 -4.16498229116897
0.187211122317864 1.11169743641647
0.48198115722144 1.43092053557093
-1.98132421803257 -6.69972954213555
0.758502825129475 2.01238871783111
1.34931125357486 4.77610269536153
0.245180659219023 1.15718232866859
-0.276183458320058 0.865207129827379
1.39456510983045 5.09968311504284
-1.6804327463332 -3.77845771215393
-0.665754239595305 0.510306575023814
0.76684140612743 2.03557122891776
1.09948744175177 3.30137507524853
-1.7534846008282 -4.39840634527731
-1.83044998652393 -5.11151071125484
1.37370201959309 4.94821688942207
-0.382340179934817 0.799332698920512
1.84016263074689 9.36729469457565
0.288068510119739 1.19542452478037
-0.114543656857639 0.945458542257138
0.206795309294432 1.1263285058384
-1.57769003775003 -2.99515593056534
-1.03763955741927 -0.308320801104982
-1.52268760032851 -2.61615653065817
1.11941760456987 3.4009040701208
-1.33311090638393 -1.50771004489711
0.128759478350397 1.07194781107472
-0.202242477740643 0.903731634603166
0.878360806223666 2.3827980126127
1.19855610565421 3.82594465441631
-0.95885680090605 -0.0795950750614648
1.9731762658412 11.0668248504742
-1.30507394687268 -1.36800862623331
-0.0157669575178914 0.992193314134026
1.88887607888191 9.96513724318019
-1.01174714032224 -0.229379091520865
-1.56039182316205 -2.87301570032441
-0.926950247777248 0.00362671646421875
1.64418618486032 7.23477519532048
-1.20043385996112 -0.896947485134618
-1.59792808416967 -3.14154749994605
1.69730156795686 7.77094735900795
0.575625029518244 1.59045760970327
1.4325210476095 5.38546272461785
-0.821655978972185 0.242999595633413
-0.810842030128184 0.264710330683507
-1.76436741471195 -4.49543613243629
-0.101373628416222 0.951589874479197
0.645687235329547 1.73352993640201
0.10576329327991 1.05772508508171
0.497152732945724 1.45448746150772
-1.20605108820682 -0.920293170938932
1.65818341412337 7.37314604751803
-1.87962426988887 -5.60073796632451
1.33566435853625 4.68210715495086
0.282054103858054 1.18980470698956
-1.26250466931175 -1.16699204691647
-1.67476380995526 -3.73258686461563
1.01967995195807 2.93168806191011
-1.67840964447672 -3.76205118380057
-1.50130451410046 -2.47608408679851
0.141109062399635 1.07989515924757
1.20995178365977 3.89118827835135
-0.531511081284835 0.6729432404821
-1.54006889061477 -2.73298567027789
1.66007705123738 7.39202540581725
1.38566787929341 5.03460666896519
1.42172807729341 5.30284416594051
1.19501735684592 3.80589515622314
-0.0613276058947694 0.970323119939249
0.588196574174466 1.61455706145536
-1.86528429810381 -5.45530051391384
0.327998076046251 1.23501500011517
1.94392556510098 10.6747725129444
0.391687289083808 1.30710792726173
0.953213581470196 2.6574299461252
-0.727575792369922 0.415559950077793
1.21439774087505 3.91692509422767
1.65999716563432 7.3912281841227
0.382087929756178 1.29548822962649
-0.16287160981814 0.92277676837371
-0.0542029915768536 0.973690800842405
-1.63260596193285 -3.40129984579709
1.12690113690572 3.43904191500925
-0.602201080415101 0.594083125950106
1.26792747852857 4.23945539583106
1.30847145779741 4.49966965275129
-0.718738337998323 0.429960531379342
-1.86119536711375 -5.4142508988663
-0.356727609725149 0.81669094026702
1.90459290272936 10.1640216503881
-1.44227588332194 -2.10983146250815
-0.391557970599605 0.792819909722158
0.631524958158446 1.70285186854789
-1.8673507757779 -5.47611714719273
1.42661454192866 5.34011495936151
0.170205151974584 1.09955348313825
-1.83568026427185 -5.16227865205378
1.2073940863486 3.87645422732492
0.706156293104354 1.87513804754791
-1.98203177551578 -6.70780014014512
1.36502793573962 4.88639889620432
0.399224089177922 1.31642786191096
-1.0253447363566 -0.270364174212614
1.50929202137678 6.00496342833847
-1.63337973901535 -3.40722582261244
0.16217533939344 1.09399145885371
-0.98542617423463 -0.152951737923894
1.5360467607192 6.23427976294006
-0.483081925171916 0.719337382260632
-1.58770633174206 -3.06713467478903
0.754998431678347 2.00275622643676
-0.315694310614774 0.842448548232049
0.80731770988626 2.15345551594596
-0.520999610898754 0.683507318681412
1.49195497855709 5.86011399750875
-1.10514956412394 -0.532465223812011
1.97857357935404 11.1403174392499
0.201904204556214 1.12260731318804
-1.48063273802369 -2.34444782161577
-1.74765842069912 -4.3469663733447
0.97582411111886 2.74742059343922
-1.94025775141107 -6.24148801641387
0.378690645694897 1.29144231745017
-0.675642050925532 0.496069283877749
-0.85135959252331 0.180704325857191
1.29651347872175 4.42146741003185
1.38442162021466 5.0255486562329
0.0993659004083431 1.05389365645253
0.322056803667919 1.22886488173402
-0.319240158445763 0.840313911505586
-0.584549292170264 0.615110609184494
-1.12598759652605 -0.60721864335168
1.41262589486618 5.2340088706993
0.646121632266878 1.73448545340135
1.78052185613418 8.67289265981815
0.0456857711294889 1.02361857502653
0.650561123062864 1.74430075452423
0.444925785300273 1.37680795769407
0.360910773533672 1.27082206027476
-0.00804773722748919 0.995996648292271
0.377237311240174 1.28972199915939
1.52091849859923 6.10374749026291
-0.0373496792893628 0.981725351704962
0.862206665085895 2.32805217881156
-0.507190355652468 0.696960910121683
-1.86205179057238 -5.42283330831689
-1.70459448861946 -3.97754027439972
-0.151804594967045 0.928017392461753
-1.68436080041345 -3.81042813744005
0.187659122172576 1.1120242490766
1.3554496588658 4.8189190431175
-0.487869327729311 0.715003204906045
0.842675192251204 2.26392694634303
-0.912200762028982 0.0403535280591013
-0.946369416514677 -0.0464010410197027
0.330462441073498 1.23759349000862
-0.0195847733926913 0.990324603273036
1.26237990986932 4.20493049782254
-0.675334511443002 0.496517206261075
-0.229399784212416 0.890086353839939
1.9604831935885 10.8954066810646
0.653098747548911 1.74995226855662
-1.99271462098377 -6.83038081816128
1.52604705488264 6.14774493496474
0.265785713873526 1.17503120638017
-0.162967637142707 0.922731122134607
-0.283790790788197 0.860963045149854
-1.36121907666447 -1.65375883389091
-1.31329929847602 -1.40838314891627
1.81527813187679 9.07259527626234
-0.442821583254181 0.753762240778201
-1.35568360263731 -1.62451551362428
-1.51220574381476 -2.54699354010741
-1.88865198249977 -5.69348252498278
-0.0584710604719203 0.971671561732019
0.381577876562059 1.29487859540097
-0.606990213147065 0.58821653307224
-0.749626560934504 0.378324034734459
0.00894689778701219 1.00450019527111
-0.134416985656438 0.936185181356841
0.947482538246278 2.63514997624858
0.652193660151404 1.74793312291132
-1.4971053276791 -2.44904543319222
1.73588142591092 8.17953815015466
1.82488208996298 9.18548319269702
-0.433621813289647 0.76114591416293
-1.85764198404764 -5.37872869313915
0.189957707840087 1.11370674618603
-0.0839730204312175 0.959700875691003
1.24728831958757 4.1123056418777
1.65468022478909 7.33832011275106
-0.586495676362482 0.612837325483901
1.79589799743536 8.84802129209566
0.428791423729258 1.35471550491822
-1.21645298698983 -0.964090333743602
1.5431165224241 6.29605904121035
0.178665194044994 1.10553169571189
-1.65418653612307 -3.5687317524992
1.59677966307949 6.78140367909103
0.933307944837329 2.58095305728585
-1.77917297787843 -4.62942925053352
-1.02449591391644 -0.267775315313738
0.586894557125167 1.61203007076959
-0.522527848756777 0.681989010527291
-0.839039065755164 0.207022944691082
-0.985881123675554 -0.154240675866536
-0.464749211709719 0.735452035855211
-1.64426334947883 -3.491186075905
0.467548491374151 1.40927248703334
1.92460116275829 10.4214987297265
-1.6443087962847 -3.49153905651549
-0.491293941305216 0.711870026008429
1.91222188641349 10.2616270266605
-0.62453546156841 0.566118745447308
0.300401507594917 1.20722329431868
-0.408259673873164 0.780638446049794
-0.386727514305038 0.796251031923926
-1.2807629327012 -1.25159562141562
-0.35007346853807 0.821032011779048
1.81278321825848 9.04344311739489
1.04336477412012 3.03668320841017
-1.32183939999296 -1.45083724437438
0.292077450222015 1.19921892546656
0.994071082116497 2.82251080980646
0.177466744750081 1.10467731384257
-0.505717794748531 0.698367674676974
-1.43761031855531 -2.08213415727262
-1.89149919300563 -5.7229237053933
1.28583777853235 4.35268374821217
-1.68115385495103 -3.7843153794068
-1.4142254050391 -1.94600789683803
1.52526984373002 6.14106061529525
0.612917731722653 1.6639222092179
-0.607365102169935 0.587754345462488
0.0446778208985652 1.02307869555183
-0.96209360813814 -0.0883318304875763
1.30375431185376 4.46867408429554
1.44134726467672 5.4538346503775
1.00844224545156 2.88322832443489
-1.59164327974008 -3.09568014470601
-1.93898777050765 -6.22763304253179
1.75462960692508 8.38402365592018
0.409361680989284 1.32924207079661
1.96570363363382 10.9656676715992
-1.89042953048247 -5.71185225186763
0.634605413505302 1.70944684791535
1.51117328592788 6.02085759885574
-0.336130600952831 0.829916161502131
-1.70081948401835 -3.94605184186047
-1.17900918117107 -0.809849619160207
-1.88055737118247 -5.61028146583971
1.7226892869659 8.03798589758147
-0.28444331034676 0.860596119807494
0.256235525537556 1.1666413705637
1.70534277107545 7.85476685433527
1.12385801800205 3.42348268523638
-1.47077804764093 -2.28298371381439
-1.61806822565318 -3.29102278955293
1.90294759424218 10.1430630966633
1.44451129494468 5.47852266160623
-1.01616478290891 -0.242580965266955
1.15493317240572 3.5856718757953
-1.40498242439259 -1.89343221858113
-0.452277827687501 0.745991134070497
-1.02618668656894 -0.272936092067389
0.16135892421457 1.0934319415492
-1.38881021937692 -1.80309489411771
0.83749527307802 2.24729452697683
-1.50103564033491 -2.4743482179168
-0.486808076009357 0.71596857791788
-0.538027218351843 0.666250669752277
1.9974166483278 11.3997283821791
0.0678536146367503 1.03574211366182
-0.0899177932002528 0.956926814390687
-0.427960738075485 0.765604944431546
1.60332836723081 6.84264236473997
-1.91821035570889 -6.00361962545922
-0.484358263630535 0.718187066115336
-1.14455694548876 -0.676128497006041
0.0304102520839233 1.01553439168116
-1.35890685005729 -1.64151467560062
-0.79930465921805 0.28731719612153
1.73862147320117 8.2091801257019
1.45803599732521 5.58511646755393
1.89547191274214 10.0482431762309
0.690931727750082 1.83782178852227
-0.621921289478656 0.569472296962388
-1.85804910216955 -5.38279143389878
-1.90519163173854 -5.86579373854702
-1.16833564944958 -0.767594943367109
-1.97462737655031 -6.62363958121218
-0.318014178865605 0.84105380251823
-1.6721081152242 -3.71120685604103
0.12506651417845 1.06961442749644
-0.598865775068001 0.598127669210699
1.77408588270415 8.60038455461797
-1.28722274061527 -1.28210603288769
1.97264479348539 11.0596075319345
-1.38334936445139 -1.77306192632327
-1.4234852858919 -1.99937517996845
0.68138823441224 1.81500941674427
1.04276100159047 3.03395800254456
0.363443058450483 1.2737027863533
-0.484681708491861 0.717894959660252
-0.42417904656689 0.768548456579527
-0.463869516349239 0.736206471705505
1.51603821310448 6.0621205930325
0.758851818256217 2.01335155326038
-0.361724413726992 0.813386601010214
1.9077891816332 10.2048298360685
-0.325593495457903 0.836447991161261
-1.4194697586773 -1.97614676213254
-1.98060600755018 -6.69154355462185
-0.74386957637504 0.388227584497721
0.0550760056004314 1.02869441717498
-1.07639607264541 -0.433686442070373
-1.98146995411248 -6.70139135895703
-0.614558293171367 0.578802301776621
0.580519533847571 1.59976104188703
0.85530197722539 2.30512646438327
1.16025561474269 3.61419236913617
-0.856557722320832 0.169393066404631
-1.15325753072591 -0.70917200147902
-0.239167870624584 0.885048828266123
-0.763721926225796 0.353522802407235
-0.977358490724609 -0.13027982803673
0.778036169125887 2.06728038238241
0.884340845547061 2.40346215552389
0.328138275778158 1.2351612579192
-1.55099225540116 -2.80778775223326
-1.03940188301237 -0.313832078552333
-1.72324163505083 -4.13518796186514
-1.84654800214979 -5.26872481567198
-1.43760044903241 -2.08207575751688
1.33241634702983 4.65997723453523
0.00470785624302161 1.00236123791202
-0.756956890576916 0.36552483747915
1.47817407575479 5.74705420485713
1.7488753934927 8.32084679314365
0.24029622852773 1.15308731260398
-0.229712543158756 0.889926204763276
0.936394382750414 2.5926445628988
-0.588247452495922 0.610781821734327
1.12089743121385 3.4084122555828
-1.68973056641353 -3.85438000159527
0.601490598378976 1.64077414366491
0.485060318584069 1.43563571921682
-0.743511201112161 0.388839796458524
-0.237997845978876 0.885656203601219
0.465159036576552 1.40575968640309
0.221364441760389 1.13768751303374
1.84453390617883 9.41980429808891
1.00647655101298 2.87484055951728
1.01000925130016 2.88993373969193
-1.87754777380837 -5.57953526649142
0.768807495091131 2.04109135728871
-0.126510324543307 0.939880107966774
1.59224876657191 6.73929233360148
-1.96901747277081 -6.56031066024765
-0.151585860600471 0.928120599126281
0.27514404486757 1.18345419050871
0.928969564002895 2.56462168594107
-1.13912071298838 -0.65572847467495
-1.3946007413604 -1.83520004087116
-1.53647125544187 -2.70858343004002
-0.853899259225281 0.175193415056966
-0.168066577292175 0.920302623347891
0.80680243856385 2.15189801728073
-0.568715904490896 0.633195480283646
1.90853719763859 10.214397764992
-0.788490064341025 0.30799463717303
0.442582735811799 1.37354566431466
0.459054233603626 1.39687552247798
0.39261749399079 1.30824877454604
-0.334717090759574 0.830801289010195
1.94086803536182 10.6343968181273
1.72074136428306 8.01724701749925
-0.292636510585159 0.855948854822363
-1.8825587450258 -5.63078399959954
-1.44751382991273 -2.14114206316334
-1.66380916147449 -3.64484099058998
0.401861302912721 1.31973046275312
0.623271629065938 1.68539340121019
-1.73144970591821 -4.20569915378075
-1.77371702050308 -4.57978392318897
0.955956320015098 2.66816793758185
-0.18897993598519 0.91023036546916
-1.6234885758441 -3.33190439423898
1.09818757474448 3.29498592215778
0.309124261139459 1.21579642671328
-0.0852514379276199 0.959103955205731
-1.20296402390791 -0.907436646342982
-0.270626161478439 0.868269331607198
0.580911905551383 1.60051121866171
1.89911009192989 10.0943055577833
-1.064461094195 -0.394142474879605
1.01285946066994 2.90217317390131
0.0402853617909917 1.0207368752534
-0.721932674649874 0.424789386080467
-1.66297073792276 -3.63817368047019
-1.81194227253934 -4.93425282463376
-0.65222601397195 0.529244626005021
-0.471450844379962 0.729648654986886
0.0370646534067927 1.01903079271421
-1.40749154055525 -1.90763621540311
0.695183252394338 1.84812752899872
-1.3887485131105 -1.80275420632957
1.85291059689277 9.52105182117024
-0.357923961782797 0.815903315065873
0.974884463317879 2.7436137850711
1.09839082028482 3.29598410091773
0.830736918330413 2.2258271038816
0.470309061762373 1.41335589628585
1.72836875306984 8.09869154159071
1.46857383792811 5.66937208921194
0.456457372499913 1.3931356304433
1.66877894800176 7.47927364775083
1.39834539959636 5.12755138331202
1.56275324449412 6.4702782522399
-0.325545359330059 0.836477482958752
1.85992830537649 9.60650597509558
-0.723639418786876 0.422010679340869
1.94494266655738 10.6882289495346
-0.603203955844345 0.592860405091233
1.7499954109815 8.33311478767895
-1.54049166796916 -2.73586090180703
0.259609595491829 1.16958206817038
-0.731632112601908 0.408851007811368
-0.0519325411599803 0.97476729017708
-1.72438866505661 -4.14500022640276
1.5064929250508 5.98137889887683
0.887629961378167 2.41492021599477
-1.11507085610526 -0.5677200891142
-0.917751339073211 0.0266596159564805
1.59838189136926 6.79634572283596
0.0866958965644469 1.04645605141118
0.511077028010305 1.47686610904852
1.18267057389016 3.73671928077183
-1.50837712909309 -2.52197132499597
-1.11106653271358 -0.553417700758948
1.81148448128094 9.02829616101279
1.50251810418365 5.94801949158777
0.814115306622483 2.17414199110334
0.790250422365434 2.10265274415143
1.63443981901696 7.13964875211235
-0.233454491645942 0.888004345247047
0.586663233387293 1.61158186572004
-0.0497866633920321 0.975786287788734
-1.61536089253957 -3.27070757076436
0.848178812923249 2.28176978444151
0.245000662573484 1.15703051541659
1.53488813686981 6.22420253630332
1.36834242044437 4.90994069397324
-0.704148960054754 0.453095256084725
0.00630860309351045 1.00316749235806
-1.40491432443655 -1.8930474153221
0.829387278753886 2.22157156070566
0.317830262660781 1.22454626287515
-1.90365540241397 -5.84965805885858
-1.62753556811879 -3.36260966263985
-1.07974178139327 -0.44492376214492
-0.0401652628917479 0.980375534538642
0.860426038380392 2.32211294404568
-1.97271544571729 -6.60201419432597
0.190591970662848 1.1141726925123
0.737104817700763 1.95457775769123
1.79080319347672 8.7896957865979
0.180385991615295 1.10676284928859
0.269402073604749 1.17826216971268
-1.21119745824843 -0.941869770922104
0.452501685883913 1.38748350398462
-0.672242667744425 0.501002119019167
-1.87624541070442 -5.56626179799371
0.697491518131789 1.85375997540669
-1.20317623202471 -0.908318352432709
1.03305973959424 2.99052061568452
1.93968190507009 10.6187642477155
-0.821861577213577 0.242581884621014
1.63669141219185 7.16153596429626
1.12811869078702 3.44528671234127
1.18993959765328 3.77729979941787
-0.25536082126297 0.876524214114186
-1.53401697010337 -2.69200261869658
-0.179743031420164 0.914703026826147
-1.02306687395788 -0.263426005993494
-1.74238098504039 -4.30067414845448
-0.468060562638022 0.732596940922769
-1.52858752321493 -2.65551241532027
-1.66643863626171 -3.66579564315443
-0.641421382308617 0.543930131873093
0.538245752971836 1.52265960139125
-0.955234916781065 -0.0698838993363239
1.80447172410411 8.94684189758768
-0.263185865738656 0.872320671481494
1.20310768144324 3.85187921672936
-0.759262351682819 0.361455230000449
-0.100655606763484 0.951924170797343
0.395716018306588 1.31206807179805
0.87211821829841 2.36145681851321
-1.15580354177307 -0.718933512086722
-1.46536166469255 -2.24955290054648
1.48295138677626 5.78604032278067
1.61459884644588 6.94907314634342
1.12568219517709 3.43280119124754
0.866338688902824 2.3419070079037
-1.88149196370742 -5.61985003530082
0.796995975427856 2.12253861289157
1.18993904227527 3.77729668298541
0.00326378068726907 1.00163538782081
0.99147879094986 2.81170713649919
-1.16854127436226 -0.768401909148824
-1.41411347488116 -1.94536708041133
-1.18129361556432 -0.818990900220248
1.75555598117783 8.39422907605019
1.4068679209575 5.19085919088058
-1.11794299148175 -0.578039886374277
1.35976299670166 4.84920562717139
-1.62864042117664 -3.37101946941616
0.837123791554514 2.24610769688773
1.90712812877881 10.1963798514224
-1.04264993613093 -0.3240363185946
-0.899533795903042 0.0710358364792234
0.131443054211328 1.07365571281068
1.58641001943151 6.68533601707391
0.133184681131325 1.07476973412852
-1.66808451108402 -3.67894626701431
0.958140299087838 2.67675334517056
-1.20071553721398 -0.898113080369225
-0.29853548062793 0.852555704806701
-1.59020194748787 -3.08521285924393
-1.14679993698604 -0.684600564244047
-1.80395518812676 -4.85889899518877
-0.137678202125405 0.9346581819364
-1.47745500338493 -2.32453774310355
0.939519736387501 2.60454548209869
-1.35213428522258 -1.60588969662356
1.55564777901188 6.40679070736085
1.3123367672864 4.52521084209601
-1.00625054930137 -0.213104313097389
0.22162258394559 1.13789256241124
0.480912021948272 1.42929135626533
-0.400866712337033 0.78609239569891
1.49035524462727 5.84689554278936
0.160235802126708 1.09266401209588
-0.0716438000134736 0.965470942110219
-0.888570383890701 0.0969612511986682
0.982840776990081 2.77603261006682
-1.69495795594223 -3.89744184996925
-1.4617497943475 -2.22739769747757
-0.42314483375042 0.769348589512897
1.59181588354367 6.73528003652307
0.591841901406039 1.62167046929035
-1.93229760790966 -6.1549561757401
-0.306711411046629 0.847784880556625
-0.138148433919949 0.934437841218547
-1.26896132597228 -1.19663623420301
1.68454711978169 7.63943621503765
1.1392782281719 3.50304601511557
-0.480386448395481 0.721754395289265
-0.490322690457073 0.712761417952586
-0.275439661364084 0.865618817515639
-0.54345663381348 0.660588689748085
0.613200605364189 1.66450249809784
0.688492804968752 1.83194973760757
-1.21377118055488 -0.952727981570386
-0.290464321290868 0.857188260010645
-0.477369279494698 0.724440156259689
0.527679401870274 1.5045097951213
1.52837326464824 6.16778693077292
1.31999322543593 4.57618457231067
1.49828207824019 5.91263737722987
0.141132714179094 1.07991059789785
1.27727139459846 4.29818992337106
-1.31957472473445 -1.43952592683124
1.07403596830212 3.17852442267588
-0.964152382379572 -0.0939174689590342
0.372903814950484 1.28462951769261
0.338964532666234 1.24661469892759
-1.90283120211612 -5.84101222823902
-0.538926179282687 0.665318605805233
0.640452057593696 1.72208273447832
-1.79462195146616 -4.77171112771805
0.845175689742974 2.27201170294667
-1.12367227739812 -0.598779210451571
1.38580255251099 5.03558633887015
0.781039940610238 2.07590383792586
0.742219428079976 1.96817826018767
-1.19406467829478 -0.870734118277661
-0.0449801113196626 0.978074598174748
0.632269387643565 1.70444168365252
0.95299971641987 2.65659470437528
-1.67434856904174 -3.72923934056497
0.295928098180434 1.2029004453019
-0.626471885314217 0.563620670988828
-0.214640428951248 0.897563375662694
1.14999846380373 3.55942407415894
-1.69462783017563 -3.89471430687958
1.00236726771143 2.85739090847459
1.63395337111733 7.13492709980748
-0.823298339067957 0.239657661120865
1.32830160638183 4.63207459559869
1.04968038564894 3.065343494382
-0.872053718616644 0.134932772431856
-1.35649547964011 -1.62878970674586
0.437074209182266 1.36594863663831
-1.95504569295756 -6.40420740716794
-1.08877311971084 -0.475592819257362
-0.490078402446558 0.712985271015062
-0.0849870563214972 0.959227383047594
-0.835507401722747 0.214440962342617
1.15150322718474 3.56740809740603
-0.348840406187581 0.821829140042039
0.32956554445898 1.23665318101434
-0.634505395875966 0.553129141425036
-0.712102981416007 0.440580147511722
-0.141181527238492 0.933015530624069
0.480824629828942 1.42915836771712
-0.993803048145222 -0.176864334335354
-0.350082585525117 0.821026109561412
0.516649780001659 1.48602722259306
-0.5906743819141 0.607919149502502
-1.38542315227203 -1.78443923472196
0.0552461398158584 1.02878717007279
-0.160301599759029 0.923997242450794
1.61316880277441 6.93549586466522
-1.01829581720307 -0.248988405125319
-0.605550343710735 0.589987699174208
-1.13095129953082 -0.625425045182286
-0.430802229699371 0.763374751301627
0.481892362246438 1.43078507028784
-1.3585046072138 -1.63938887888464
0.738234268448013 1.95756946471052
1.01394446821935 2.90684705473204
-1.41398252492884 -1.94461750198599
-1.61322129466368 -3.25470147546836
-0.652286033515564 0.529161966391323
0.933830736067949 2.58292913135864
1.43087486679923 5.37279130342767
-1.00525884447349 -0.210185800461815
-1.50577099100418 -2.50501195390687
1.69439119087545 7.74078341437282
0.10925386070688 1.05983790235327
-0.188831907928934 0.910302370929738
1.57725839578374 6.60146623821642
-1.65640169917705 -3.58617255705754
-0.0101443338755236 0.994960212353824
-1.08708920938024 -0.469837268188092
0.441869766673601 1.37255664793776
-0.564949950927475 0.637391656904085
-0.823574016021408 0.239095543015953
0.174998644970009 1.1029256423352
0.35957055618162 1.26930486866708
1.01517511591745 2.91215807920767
-1.01068379458088 -0.226217569950526
-1.01107275652972 -0.227373296108133
1.59442120696011 6.75945724907352
-0.345592888348411 0.823917773481681
-1.29280590229753 -1.30872129958885
0.908709804141768 2.48992821737328
0.251582943895579 1.16262780597153
-1.67068398963368 -3.69977029502073
-0.358884765262246 0.815269170072609
0.543566188111855 1.53196512814761
0.39052183014819 1.30568228138866
-1.78482711120297 -4.68120920110513
0.392427898563614 1.30801603114462
-1.75534148417263 -4.41487490434709
1.8503004408122 9.48941536715725
-0.919870775031479 0.0213903298721088
1.47400951506076 5.71324744708682
1.75263123441605 8.36204118570609
-0.198501908355358 0.905574384373035
1.30090062834669 4.45001583074007
-1.66407712761412 -3.64697335799929
-1.84703443163292 -5.27351959062834
0.849949429186349 2.28754778244004
0.212031550491898 1.13036315308368
0.328070932321069 1.23509099799962
-0.76219608701357 0.356245856880229
-1.05034528540007 -0.348455087434208
0.130191110209795 1.07285764383031
0.238265181418163 1.15139947479866
-0.1604445162542 0.923929431107584
-1.88544272017285 -5.66040746381983
-0.355263481040195 0.817651875296108
-1.62197081642769 -3.32042900348838
-0.27260628005381 0.867181855396237
-0.0393374727000664 0.98077204637976
-1.67029587641437 -3.69665697078896
-1.4710559598251 -2.28470574025771
0.3082855060622 1.21496370491836
-0.172468446827101 0.918198235346273
-1.34416057620319 -1.56439999322047
-1.33564904373112 -1.52064952502959
0.802023719528068 2.13752419665797
0.59189749511968 1.62177939336193
-0.325807182637529 0.836317032850372
0.929845640955968 2.56790994971066
-0.220918870588491 0.894401548446393
-1.71225608825882 -4.0418882790471
1.8403591192216 9.36965021921371
1.12938560758863 3.4517965635835
-1.5047675750776 -2.49849804363137
-0.861877904777332 0.157687562511912
0.613015809930273 1.66412336752378
1.99783392103624 11.4055229812418
-0.600432925376096 0.596231459695451
-1.45190828144078 -2.16758691243222
-1.48904013730822 -2.39754134298098
-0.285818161709389 0.859821491004982
-0.440240268030195 0.755851362680721
1.94114199811838 10.6380099391004
1.5203777216456 6.09912332861589
1.42674268688303 5.34109535400439
1.53771130486377 6.24878063047949
1.08882233533862 3.24932053509569
0.737204249571775 1.95484086941443
-0.658852549666758 0.520045515942364
-1.56974238807888 -2.93869925768835
0.573850104841899 1.58710864281915
-0.150629813670606 0.928571540966119
-1.59348005323711 -3.10904716797659
-1.78069008218441 -4.64328957281517
-1.23127998733112 -1.027802140198
-1.22026092161847 -0.980308555640824
0.882795537168005 2.39810156317915
-0.930186148595834 -0.00457659538527411
1.79227171174474 8.80647713121272
0.766752410763582 2.03532184929355
0.773207551344699 2.05352048702182
0.260074786136564 1.16998950568125
0.252140468136878 1.16310624064939
-1.14290600744467 -0.669913248840446
-0.374213118765592 0.804955680857456
-0.191181658633414 0.909158058002845
0.810589974522598 2.163381259695
-1.41514599508773 -1.9512822678602
0.395078496897364 1.31127984273892
1.16524505127672 3.64112744273495
-1.03632649241488 -0.304226042565756
1.18948005796077 3.77472198587331
-0.270297527364671 0.868449433874931
-0.198422000604987 0.90561366269916
-0.857131906047643 0.16813602458199
-1.47168682695009 -2.28861721392019
1.30970381039638 4.50779879595879
1.08896952136146 3.25003325106136
-0.253744816342775 0.877385261341134
1.60809238143707 6.8874700276336
-1.46164224840128 -2.22673969804884
-0.523674976294005 0.680845431379947
-0.436105436776697 0.75916946927637
1.95670181580846 10.8447229790432
-0.29578081090816 0.854145237454634
-0.146498271875552 0.930517489092381
-1.48114813660415 -2.34768518142267
-0.360618287314252 0.814121407141682
-0.137969553008917 0.934521666109314
0.692041230923785 1.84050269395564
0.848541326972606 2.28295127115617
0.306602934648166 1.21329862270282
-1.99440478525603 -6.84990030388248
-0.163046902412025 0.922693441264199
0.807028705077833 2.15258176473391
-1.38021215283323 -1.75591511891502
-0.391259388345681 0.793033169097501
-1.28762850152753 -1.28403263820684
-0.174382470327815 0.917280783766521
1.13131840297683 3.46175128155318
-1.37045581518671 -1.70308653950514
1.88195365105748 9.87847330015004
0.42451589737284 1.34900533570564
-0.250619191395958 0.879044006890007
0.203502318902969 1.12381816913033
-1.58694498485549 -3.06163097800262
0.506640548475837 1.46965705070567
0.74023425649113 1.96288327802345
1.21036770304808 3.8935892176777
-0.783689800415043 0.317015675037404
0.36004969780675 1.26984669049933
-0.00841423401782704 0.995815284654669
-1.36225289641879 -1.65924676348695
-0.169846151402795 0.919452780658226
-1.80167404054627 -4.83750329606028
-1.71064649794186 -4.02832064119088
-0.0144324412344345 0.992848406647661
-1.71446370448892 -4.06053936154386
0.588542852677619 1.61523034349259
1.69359471543941 7.7325445269266
0.960172150386684 2.68476864465872
-0.501004418258892 0.702834934789813
-1.73543242629262 -4.24016056347248
-1.18738325377207 -0.843528056777747
-0.784046358824987 0.316348896817618
-1.41258173898454 -1.9366078844798
1.71385367062331 7.94424996569319
-1.95982876404981 -6.45738784953313
-1.74772858892011 -4.34758380647548
-1.8435030868594 -5.23877001754084
0.662633333223705 1.77145552470409
-1.97857398165206 -6.66841640626697
1.46737028666602 5.65969553365888
1.40820450129305 5.20084818512847
-0.26203923453244 0.872940237525545
0.446314292887601 1.37874996770754
0.423303278566297 1.34739667713506
0.654264776554472 1.75255916982409
0.195648317957575 1.11791351216775
0.869789244443124 2.3535548140646
-0.972812372591263 -0.117657788190493
0.728326427169269 1.93154966766546
-0.615050256398702 0.578184193752807
-1.68378682359237 -3.80574701635623
1.32645410070877 4.61959452246244
-0.00306303876824643 0.998471500496445
-1.92613251195332 -6.08844338856703
-1.7678595137034 -4.5268330652448
-0.294710498461862 0.854760461822923
-1.4156965362227 -1.95443979345889
1.73060191716862 8.12265833310048
-1.19280754552839 -0.865592394287253
1.31589340472998 4.54882640739453
-1.20936667062426 -0.934173394631358
0.824988959339324 2.2077756649389
1.44994772354708 5.52116167829113
1.66091128157801 7.40035470103932
-0.822159008525746 0.241977269724585
-0.783898941357204 0.316624637232964
0.535165215213903 1.51732290557903
1.37416286857846 4.9515201528893
-1.26707642429524 -1.18795128714585
-0.594686170782149 0.603148756327027
1.50327056291737 5.95432281026179
1.92323778364012 10.4038007624767
-0.0661064161920821 0.968072202724847
1.7973699615765 8.86492744905595
1.78822759931664 8.76032265014644
-0.0355696491027953 0.982580429304874
0.710778270422538 1.88669431924725
1.79510199198771 8.83888910864229
-1.6013204280294 -3.16645825549723
0.827550060828212 2.21579542506338
0.710290387126795 1.88546942765681
0.75328106345806 1.99805941204832
-0.602679413986287 0.59350031537745
-0.736116787059058 0.401360263426038
-1.75442056855673 -4.40670289919965
-1.43469212012865 -2.06490171684732
-1.42457901983685 -2.00572488115273
-1.85067490959863 -5.30948688938268
-0.356802885485278 0.816641446786753
-0.483117331157287 0.719305522849883
-0.737512856820079 0.399012610082135
1.93790301004031 10.595351464905
-1.21266198885831 -0.948042898510603
1.53227608060943 6.20153286471949
0.94099363195935 2.61017957606926
-0.00887345133348605 0.995588151009568
-1.15415909725177 -0.712623858434803
0.591980211072964 1.62194148219783
1.36550586974274 4.88978744815124
1.70357448545846 7.83627435027178
-1.70289120626857 -3.96331499568938
-1.40104805470438 -1.8712620281133
-1.61851806470054 -3.29440499476517
-0.63112002838442 0.557575573319032
-0.565084838129248 0.637242049113841
0.851954413494163 2.29411277330813
0.48982569862781 1.4430008972904
1.45141946350045 5.53275280275556
1.29955424466575 4.44123700840332
-0.84757109003723 0.188870428958125
-0.0599663634505343 0.970965410863567
1.24627519307932 4.10615513282809
0.778027223959338 2.06725477506621
1.11303740065966 3.36872048346061
-1.07369452877235 -0.424661437034536
-1.06772890975252 -0.404885801541526
-1.61984117288767 -3.3043641509849
-1.56478897897108 -2.9038045844714
0.935177122726119 2.58802629462615
-0.511888877383642 0.692436685658468
0.453053721775183 1.38826905851069
-1.10552539167778 -0.533789646913369
-1.82968366491081 -5.1040974467213
-1.66590564833054 -3.6615427251593
-1.92322212449706 -6.05719723651022
-1.64755790625506 -3.51682616610233
0.536198102612343 1.51910810086017
-0.107249325761568 0.948854634668092
1.96901847512502 11.0104557393676
0.223121280965427 1.13908566639894
-1.70559709939951 -3.98592738382609
-1.25120198878827 -1.11581428511409
-1.86325846684435 -5.43493952243064
-0.117696588412903 0.943989897472763
1.82415943911388 9.17695190231721
0.562110201979332 1.56528687742743
-0.560464756988403 0.642337358042499
1.36413404164186 4.88006667367168
1.85764949536136 9.57869371898181
-0.572732680065321 0.628675610461671
1.31909501625392 4.57017833613826
1.31768933845374 4.56079275955768
0.953529119384244 2.65866280935027
-1.62428029635491 -3.33789907091573
1.23452868466645 4.03546076426489
-1.22147326342045 -0.985492891674448
1.32006706182028 4.57667862051968
-1.75140090095866 -4.37996875561887
-0.879055494390023 0.118993780269021
0.54827145310572 1.54028869924352
-1.03646261797876 -0.304650088505529
-0.29642091653693 0.853776665310185
1.69965456488001 7.79540165753882
-1.92626251476466 -6.08984139395354
1.25754503364581 4.17505045600754
-1.39510351177812 -1.83800024057679
0.768149445576419 2.03924146447786
0.629209199577733 1.69792232313776
-0.398388833061027 0.787898226117644
-1.85205832248426 -5.32319312482858
-0.788306477157378 0.30834141495755
1.92336307531936 10.405426230607
-0.70940929986311 0.444844585697191
-1.62290632240129 -3.3274995398133
-1.90084325800046 -5.82019063857491
1.89309179361197 10.0181944192673
0.470274000325968 1.41330386505985
1.62265830770733 7.02598987553108
0.893485980468911 2.43548364637751
-0.733495310758882 0.405748290948773
0.148485290971542 1.0847514676653
-0.150844352111013 0.928470370725297
-1.91783695783904 -5.99963945317911
0.190400907466076 1.11403225548728
0.399679743270707 1.3169969369697
1.096754005881 3.28795402126858
-1.49633334111525 -2.44409122497096
-0.958128006818784 -0.0776354898908592
1.98316190841294 11.2030781834086
0.235138767152799 1.14881836851237
0.159138615959582 1.09191580441678
-0.50351942722194 0.700457964971038
-0.747017559231104 0.382828317131004
0.954233168495334 2.66141598552957
-0.984114947122056 -0.149243088442066
0.566261707929518 1.57293847659294
0.0268518303870944 1.01368001238554
-1.5080912666153 -2.52010818864839
-0.698445924122061 0.461925923461316
1.55686017968182 6.41758757741351
1.93469424921284 10.553216972581
1.53776104519807 6.24921437317124
-0.862739132584211 0.15578038107934
0.79029759576884 2.10279094050976
0.0471274247975035 1.02439239875053
1.78987721790029 8.77912690052206
-0.0431857369949249 0.978931048381096
-0.929862394078688 -0.00375346834793933
1.92044127119492 10.36756965867
1.65432556888321 7.33480166918088
0.33085216554685 1.23800274922351
0.70593915826521 1.87459777061375
0.529588679409994 1.50775701177613
1.15961162919433 3.61072991364129
0.927783599136133 2.56017804399492
1.51679446472447 6.06855578095598
0.505843770609108 1.46837017949939
-1.86608403340888 -5.46335098835564
-0.13183306760037 0.937393734241051
-1.5948633922361 -3.1191350481614
-0.402491431752823 0.784902332976656
-1.2133053277603 -0.950759255996364
-0.399992927369038 0.786730451090819
-1.12793902639114 -0.614357814757456
-1.81371503806416 -4.95107099648
0.999099306135273 2.84359567367731
0.252238154701149 1.16319013979868
0.724810795768328 1.92243806366124
1.35877053448433 4.84222226357485
1.44383629208362 5.47324790942361
-0.52201647213084 0.682497725271347
0.792357051294308 2.10883613132269
1.13142344393241 3.46229309621455
-0.909827386308057 0.0461623830673111
0.627588348397961 1.69448642159623
-1.66249879225256 -3.63442369352429
1.0244559071169 2.95254617827994
-1.06115849580953 -0.38334870945599
-1.08590118719238 -0.465786940173083
0.547207888067801 1.53839951004046
1.52819489594 6.16624825942774
0.638538035700954 1.7179289325443
-1.4204616749999 -1.98187240908629
0.966847435128181 2.71129155995665
1.75335340411573 8.36998004208666
-0.502963356595071 0.700984834758228
1.93133594447325 10.5092527985002
0.538941134559069 1.52386945789964
0.188749750665735 1.11282136647531
-0.292593598496842 0.855973391432946
-1.37039122070766 -1.70273926532393
0.102400835197629 1.0557047712163
-1.53332126060227 -2.68731226056737
0.891235759838326 2.42755720186516
1.52170743807571 6.11049884802857
-0.844376562263295 0.195705543756215
-0.375488262401924 0.80408067223377
1.4696309220224 5.67788246054487
1.13816047065756 3.49721817323933
-1.84272956088872 -5.2311765844534
-1.74749904667497 -4.34556417819544
0.644838063213352 1.73166457646492
0.482666975213408 1.43196777559443
-1.89561475321412 -5.76564239081856
0.0640935205283264 1.03365067786788
-0.993936303270147 -0.17724780418363
-0.182062122929765 0.913583883048714
-1.76248041563369 -4.47852345492373
1.6180302702965 6.98173882657229
-1.72862303592487 -4.18133915761499
1.26946417620312 4.24906447899531
0.0232857615563313 1.01183199442652
0.376062841662597 1.2883363499593
-1.92558451797562 -6.08255259910049
-0.301932119684734 0.85058343404128
-0.615330242664868 0.577832079638759
-0.574937648601093 0.626174929312303
1.24780110159707 4.11542187846027
1.62037688557016 7.0041481320128
0.776891640072805 2.06400746289965
0.805897067176439 2.14916497188128
-1.04699449756928 -0.337780348737123
0.2413457555599 1.15396291866958
0.779938874450223 2.07273712296313
1.05148057441272 3.07356436028159
-0.60208950984249 0.594218965773506
-0.662847385953129 0.514428181055222
1.47783716222536 5.74431304459151
0.677094877773097 1.80489057535815
-1.36138557463422 -1.65464211025583
0.826952051915089 2.21391948052634
-0.616460223808683 0.576408519698541
1.90049331762246 10.1118598699591
1.8857421665975 9.92583229362722
-1.99232133463139 -6.82584373830562
0.347339572283597 1.25569397803319
-1.94164832298215 -6.25668016754697
-0.87580042214737 0.126432453372567
0.618131148387185 1.67467389041367
1.64951506747112 7.28720973548482
0.856870484129805 2.31030965843759
0.75816365172245 2.01145359287336
-0.630893720992384 0.557871496703937
-1.68440412323232 -3.81078159312989
-0.801909424927161 0.28226301307832
1.67684929108731 7.56091329136224
-1.77823002811071 -4.62082666321174
-1.50476161515709 -2.49845937961117
-1.07092277187912 -0.415446994640916
0.554453463411948 1.55135983180469
0.343066363886576 1.25103729993777
1.78407551388253 8.71312912961346
-1.04397837982574 -0.328227301390198
0.562222038737691 1.56549207455857
-0.866328340604409 0.147795038857713
1.25613287223244 4.16635985946709
-1.90746279943368 -5.88969818106689
-1.287035879447 -1.28121919610003
-0.539949798500574 0.664254689315366
-0.862636863904048 0.15600703402093
1.01478429583278 2.91047031699819
1.23975984222251 4.06680364423564
-1.19235655588129 -0.863750409325774
-0.831963437570925 0.221828798485852
1.40265973241274 5.15951629228891
0.0194303791313675 1.0098453788844
1.40372315003464 5.1674213568491
-0.528505640610852 0.675992745163049
-1.6564367176555 -3.58644865323731
1.69314550365526 7.7279008288673
1.18427001205538 3.74561264251057
-1.49048114937126 -2.40670243203198
-1.5021541457011 -2.48157349939488
-0.973580992960761 -0.119784081701067
0.780153688626938 2.0733544205084
-1.97175917019071 -6.59121430968625
1.70812572752589 7.88393967866811
1.71164628801791 7.92096566685086
-1.46901377285881 -2.2720670262631
-1.23705343296793 -1.05302246946098
-1.31701534670063 -1.42678900543212
0.973916433523538 2.73969810519761
0.0138870558554998 1.00700893829624
1.24736888638784 4.1127951120542
1.85797901316208 9.58271162347587
-1.55414649990604 -2.82958745285144
-1.26664102180807 -1.1859487335488
-0.491671869592991 0.711522572045034
1.56608847925319 6.50025427042704
0.106986326153375 1.05846356533392
-0.524782636566406 0.67973800984612
-1.28628653522886 -1.27766539355675
1.96238202836105 10.9209239716033
-1.51631884062425 -2.57401806958537
0.649325034600645 1.74155870983148
0.724807417907873 1.92242933945989
0.0888190227673209 1.04768881497555
-0.727130728089694 0.416292246427935
1.61375447966397 6.94105389685173
1.85902495714243 9.59547356068222
-0.26205808636835 0.872930061358466
-1.91095938475336 -5.92661571731319
-1.84063806586138 -5.21067802010906
-0.83203390035015 0.221682455008366
1.77395758284144 8.59894387399551
1.39522424578815 5.10453286882605
0.991061263606985 2.80997126875214
1.75867999543365 8.42871557387581
-0.145078200731571 0.931185345721736
0.424811330532875 1.34939798162108
-1.97676835775563 -6.64790737169224
-0.0224534136102359 0.988925588632502
-1.79084339270371 -4.73667740838314
0.623736702573969 1.68636904571803
1.74134746699276 8.23875263443377
-1.22329593745987 -0.993306226618697
-0.68863543648969 0.476841503255705
-0.841565476998836 0.201682049751289
-0.114665665274081 0.945401721077095
-1.27052530195143 -1.203861842856
1.78013467986155 8.66851747260619
1.86228752760561 9.63536391582172
1.54718747645037 6.33185935875192
-0.0767255669261084 0.96308944111748
-1.35816488303435 -1.63759446223621
-0.494991410247095 0.708456202517538
-1.45696110771646 -2.19819336995244
1.06594787433349 3.14046823450914
1.93612990897787 10.5720532579514
1.90943433588774 10.2258819915518
-1.45262100952012 -2.1718911615037
0.785661940921628 2.08926897040826
1.48129203807223 5.77247401344108
-1.78019018746414 -4.63871983208713
1.95664099860872 10.8439092498593
0.477914979210265 1.4247462952397
-0.515843724751158 0.688586128698253
1.50482410118571 5.96735421244212
-0.434201139132909 0.760685999633963
1.09245878750959 3.26697566808596
-1.35868571455372 -1.64034585027071
-0.758148874059312 0.363423392727458
-0.786945539293618 0.310907702980332
0.52878601489933 1.50639014501486
-1.63848464065965 -3.44646537614018
1.36992278408109 4.92120022086491
1.61674374079115 6.96947727260168
1.0830007265137 3.22125744692358
-0.492413148033551 0.710840088759813
-1.06776075633125 -0.404990811139046
0.349569286229697 1.25814391282824
-1.76802641988979 -4.5283368888717
0.454907214792726 1.39091425702831
-1.90638792522486 -5.87837759132325
-1.54731497238487 -2.78248642605251
-1.43140345060785 -2.04556579019779
1.81777930286695 9.10189243174927
0.0749928195697884 1.0397548078862
1.89915503061025 10.0948755081904
-0.456716963407882 0.742278124899361
-0.831663785685692 0.2224508961276
-1.21310084781679 -0.949895577925508
1.56068431452565 6.45173980936193
0.918894025714693 2.52715306974637
-1.51902640644904 -2.59188884866268
-1.85715800237426 -5.37390130183896
-0.19956111548409 0.905053393591488
0.127893873663956 1.07139913277355
0.78402794741665 2.0845307395796
0.808996478858507 2.1585402118705
1.22984588684676 4.00759330424166
0.234445791603679 1.14824903644456
-1.2782780393036 -1.2399399481654
-0.609232095571956 0.585446173369471
-1.03990220203699 -0.315399958059745
-1.63056543630134 -3.38569990115128
1.67575051841743 7.54975703877128
0.790850052170221 2.10441029419865
1.43535004703875 5.40729800660216
1.28305857717661 4.3349362174954
0.602211871837049 1.64221837714499
1.17002221696631 3.66709775231917
-1.96969016459752 -6.56788478965851
-0.420503314401308 0.771382845430355
0.00566025559792926 1.00284072543611
1.68187917657893 7.61214922970133
-1.32969440699063 -1.49036988534212
-1.21559370058737 -0.960444327997432
-1.51750003742957 -2.58180640932343
1.78995853088708 8.78005459910523
-0.533564886410108 0.670845822719905
-0.863322864649673 0.154485756564612
-0.485369859142183 0.7172726778564
-1.90882702469775 -5.90408524932761
1.41520006110501 5.25339829886051
-1.35924644196202 -1.64331034841882
-0.861356600799521 0.158840316600257
-1.31755553564965 -1.4294732040879
-1.11214461415033 -0.557258553461082
0.383861508736479 1.29761414507274
1.41257472447696 5.23362405936077
1.66732817080819 7.46467150976683
0.315562558070873 1.22224832688126
1.15482684194415 3.5851043270083
0.883803781254541 2.40159746566727
1.34039284239157 4.71448965561803
-1.96618885461039 -6.52852077074172
0.520093164339951 1.49174708496151
-0.059431820558177 0.971217772926455
1.20056066320666 3.8373461174504
-1.26527759983197 -1.17968671330802
-1.19603151765713 -0.878799840768772
0.908964190174309 2.49085015515867
1.64444117719047 7.23727740584927
1.83226136004263 9.27294537353598
-0.618081054980039 0.574359623139687
-1.79611021553569 -4.78555162805232
-1.45813243987457 -2.20531905080718
0.0797175808956805 1.04244138701734
-1.68602571107216 -3.82402493587793
-0.83913379187387 0.206823209072277
0.369839788247009 1.28106217047193
-1.96293142430682 -6.49202905123444
-0.977438612004603 -0.130503269498393
-0.764717547176466 0.351740938790488
0.665388103370584 1.77774792956026
-0.95405496559106 -0.0667349281815306
1.74365110010926 8.26380751553717
0.88937319263139 2.42101968929262
1.44053858124633 5.44753982488186
-1.06214485291108 -0.386565669563593
1.61677967387832 6.96981950711904
-0.969711271564239 -0.109110859264739
-1.10111582664108 -0.518304852801946
0.479358068620247 1.4269307348515
1.00147119269776 2.85360105200195
-1.50614978008554 -2.50747323363271
0.135087463666287 1.07599192289944
-0.10437397163449 0.950193113552419
-0.865024945595861 0.150701782569106
-0.227882217765365 0.890862384696654
1.74774498479461 8.30847914790494
-1.06749011983779 -0.404098616506272
0.447194619714659 1.3799846034777
0.191706917504417 1.11499353633612
-1.53467718360765 -2.69645765329396
-0.0303931324022266 0.985074991992057
1.97064170860828 11.0324372794247
-0.338622466202244 0.828348925746348
-0.763123127650191 0.354592550778447
-1.78187740239317 -4.6541539279079
-1.76095886117017 -4.46491322418032
1.05509605695447 3.09014445357528
0.471792586686954 1.41556143616822
-1.63713799322961 -3.43608993064268
1.21639985543668 3.92856695231781
-0.209052526897512 0.900355406759493
1.17052390765357 3.66983542516921
-1.68780168124837 -3.83855902194721
-1.3546328996739 -1.61899157436293
1.48663823198821 5.8162777913109
-1.65635437320755 -3.58579944386302
-1.67608702535521 -3.74326547549861
1.21350843499642 3.91176433969345
-1.87888741434495 -5.59320853108333
0.395408335680854 1.31168749807443
-1.79419211586412 -4.76771815281607
-1.51250302176874 -2.54894179033519
-1.61684686950874 -3.28184943102874
1.58631980276513 6.68450505598533
-0.802594996628898 0.280927956280322
-0.0794859431731667 0.96179783037884
-0.732966009715986 0.406631070533624
-1.01893260734773 -0.250907994009945
-1.63790006061609 -3.44195925555662
-0.615478446559309 0.577645597833662
0.343695035731867 1.25171922039455
0.57779178361288 1.59456375401596
-0.699937440091468 0.459627865535633
0.707100514414448 1.87749020069058
-0.886972889362756 0.100690569025272
0.815133489524867 2.17726293510365
0.764683353719457 2.02953595477432
-0.34376078044843 0.825089264565482
-0.335400168623273 0.830373901387892
-1.60973235621782 -3.22869365261734
-0.0625530951816025 0.969745318511178
-1.56187486241968 -2.8833802995486
0.235851754428173 1.14940519024992
-0.604510140635392 0.591263300128981
-1.52595209584025 -2.63789453670762
-0.0167539530500709 0.991709448812335
-0.48133754314143 0.720903460043628
-0.331574810527902 0.832758980354371
1.11301569467954 3.3686115109984
0.26433173410415 1.17374059837321
1.86404601882561 9.65691630769074
-0.351661789064046 0.82000187970048
1.30528355617809 4.47870152488599
1.13849983059223 3.49898654693974
-0.700254175629474 0.459138814559712
-1.16323425742782 -0.747663206520722
-1.09317503071268 -0.490719432982972
-1.50947056014758 -2.52910445464645
0.369611341071494 1.28079729679421
-1.93280424225842 -6.16044165455544
-1.58252649628825 -3.02979597900591
-0.763232969381421 0.354396427529145
0.697362563986819 1.85344461974421
-0.78885638453812 0.307302274213614
-1.97001006573632 -6.57148858056276
0.41282324688506 1.33369168895236
0.0915867832461297 1.04930396860131
1.38481369648395 5.02839681198353
-1.37682768681298 -1.73750406988822
-1.81277009698659 -4.94210215315662
-1.5436799434384 -2.7575953075135
0.31214230664009 1.21880760764692
0.745415545892666 1.9767462646854
0.908678008631643 2.48981301339038
1.58766883354634 6.69693929600202
-0.115972246424081 0.944793175590471
-1.77203737735337 -4.56456341357926
1.99893935864718 11.4208845267762
-1.08983879770738 -0.479244124305624
0.606564105436654 1.65098108377588
-0.59347945779573 0.604588689282053
-1.75972116339784 -4.45385988010386
1.41433405500788 5.24686843839027
-0.565038180623414 0.637293804248749
0.629341526646942 1.69820335474672
0.359092699084816 1.26876515243835
1.31209455722298 4.52360657942307
0.0231655473192554 1.01176987386168
-0.91855554467885 0.0246628562726551
1.76618437749433 8.51200436784866
-1.88883365455313 -5.69535834021965
-0.35877884193425 0.815339150943887
-1.28223361298719 -1.258515120528
-1.74717176179298 -4.34268550251242
-0.0381931754896185 0.981320627546728
0.28756260071101 1.19494845761518
0.149413701940852 1.08536867189767
0.408579070397023 1.32824133588349
-1.49914779368727 -2.46217781434149
0.0744843979020424 1.0394671982087
0.838133071460148 2.24933406340543
-1.25106277067889 -1.11518956498603
-1.94527337085676 -6.29639042233159
1.63989670287326 7.19278579532602
0.947947137930402 2.63694823975164
-1.95022300713495 -6.35085915138298
-1.14036041717668 -0.660364012151103
1.59891211959265 6.8012963388368
-0.557883632982971 0.645158125220656
-1.90969707059121 -5.91327184497705
1.08206173751063 3.21675416302705
1.34230080876573 4.72761185919997
-1.36304571536481 -1.66346099869758
0.198751080498258 1.1202323785615
0.0171134301767624 1.00865701431544
1.09076257328937 3.25872843042728
-1.22084344547998 -0.98279834030487
1.95397912147747 10.8083379921804
1.84483292780524 9.42340441154085
1.90754866825673 10.2017548388507
0.988413745574071 2.79899146699252
-1.16113479370591 -0.739509628040418
-0.898757806765679 0.0728899639194056
-1.72453369303265 -4.14624182096285
-0.218035835765646 0.895856789997443
0.269600038029095 1.17843990359055
-1.78591422979284 -4.69120367724842
0.658024340641895 1.76100770281471
1.80974732277935 9.00806624621475
-1.40394985040404 -1.88760160308266
-1.50873158780901 -2.52428253507125
0.307451961694163 1.21413792736538
-0.528884650568429 0.675609467041245
1.4386350581012 5.43274695287675
1.09033381004139 3.25664705752175
0.522694119614134 1.49609776021739
1.54715171273416 6.33154412831436
-0.669302419702281 0.505236420592782
-0.59255965317619 0.605683348425251
1.52492262642214 6.13807633920665
0.309944458185623 1.21661245618531
0.494040954130513 1.44958500461109
-0.721085446062989 0.426164658692831
-0.695056255808788 0.467118618177397
-1.07749284575177 -0.437362833929012
-0.791133218524787 0.302986357719246
1.53796642661311 6.25100558859299
-1.03394103954002 -0.296812271055293
-0.559812087764417 0.643052368011132
-0.633968339936727 0.553836990474224
-1.08897230869268 -0.476274773578637
-1.83406499079466 -5.14656803373782
-0.815162051023353 0.256098168103319
0.143125881657654 1.08121468711607
1.22397845948015 3.97292829861374
0.725829322232748 1.92507132628676
1.3715998567677 4.9331732876944
-1.52342357214668 -2.62104908437208
-1.93042075917092 -6.13466094661306
-0.616595189844368 0.576238222099196
0.213993461034763 1.13188857365299
1.39803812909352 5.12528131282752
0.100701243153916 1.05468909117862
1.35925049059371 4.84559833145911
-0.554927920621903 0.648365664774076
1.97667711505462 11.1144531365842
1.57800292613626 6.6082575605031
0.645395105646288 1.73288784568775
0.974282898146801 2.74117972496166
-1.03708996103011 -0.306605699392774
1.69454522744885 7.74237759001367
0.375377173129433 1.28752927895184
-0.655548801477311 0.524650300775287
-0.953868870865906 -0.0662389530943418
0.282829356831491 1.19052423629332
1.34300982734016 4.73249637214456
-1.17640153921017 -0.799457174811928
1.18335643107566 3.74053038779944
-1.58585802253435 -3.05378273525992
-0.406269387633428 0.782116550369699
1.64426823954281 7.23558031363243
1.8751683113906 9.79407752265142
1.6271576806091 7.06922491586328
-1.41956760269503 -1.97671119129604
0.91561499477998 2.51509692453709
1.42657150446897 5.33978572828725
-0.226384349697887 0.891626661442817
-1.37240178410673 -1.71356383293874
0.329503820648372 1.2365885486861
1.76835687279365 8.53623420599155
-0.0643571965988285 0.968895435777672
-1.04472205333709 -0.330577886582466
-1.54851585294466 -2.7907356294578
1.74531312125836 8.28192067639441
-1.29753201911022 -1.33142964625767
1.03058135477804 2.97952910077753
-0.739307202417413 0.395984147242487
1.07074212816424 3.16296938985897
-0.276871417668281 0.864825836028826
1.59428070989477 6.75815166865391
1.9775394033217 11.1262076505164
-0.322668723171736 0.838234307567092
-0.755764004994456 0.367622206163626
1.65383353457481 7.32992254260252
-0.144282128392698 0.931559525911121
0.402805096754445 1.32091764161098
1.24507461810011 4.09887761868257
0.735468533287763 1.95025524288497
-1.64111697855334 -3.46679688829703
0.0579891524950362 1.03028656200424
1.67760140437525 7.56855725419944
-1.69368273555164 -3.88691181889737
-1.24377755022998 -1.08268901549638
-0.432681352834765 0.761891095706417
-1.35297342057784 -1.61028445695437
1.92355852750648 10.4079623031266
-1.57448990466284 -2.9723539336749
-1.34222151991915 -1.55438427031929
0.289194162205326 1.19648599819658
-0.286787433879757 0.859274129623031
-0.656062103608636 0.523937269374979
0.536164444597334 1.51904986165716
-1.22738074864422 -1.010899981434
-1.29505040068755 -1.31948533683892
1.26892424565769 4.24568599818437
-1.92474243500583 -6.07350719887432
1.55085128086431 6.36422090095643
0.680577102095658 1.81309088302144
-0.298382311684437 0.852644321350996
-1.65984865582475 -3.61340658154561
-1.68973754295019 -3.85443729120498
-0.899730465091301 0.0705654553515266
1.79002475993214 8.78081026076583
1.24563099265532 4.10224871154197
1.57497112340899 6.58063788084533
1.37769729650331 4.97691775303615
-0.281410350029088 0.862297708674295
1.93491822661699 10.5561539738156
0.0363111485001473 1.01863296478277
-0.332987432240669 0.831880589168509
-1.80974006411356 -4.91340769524829
0.694896077357561 1.84742862080097
1.95457147794537 10.8162462688562
1.27308420348949 4.27177908733688
-0.874512743956169 0.129361319704098
0.914187743108128 2.50987033058807
-0.244938422010026 0.882036783630122
1.9907673585679 11.3076838361221
-0.13346042304654 0.936632714681418
-0.68233701334333 0.486236293949049
0.548963450448359 1.54152031847289
-0.808553351871949 0.269240248340317
1.04654601741391 3.05108454552819
-0.119795110934064 0.943012023037743
0.701049984621372 1.86249454770044
-0.226734136952503 0.891448333304594
-1.43358458674436 -2.05837995413147
0.374885551613246 1.28695146707067
1.68338246621681 7.62751502901291
    };
    \addlegendentry{m\_1(x)}
    \end{axis}
    \end{tikzpicture}}
        \caption{Zu schätzende Funktion $m_1$. \\\qquad  }
        \label{fig:subfig8}
    \end{subfigure}
    \begin{subfigure}[b]{0.5\textwidth}
    \centering
    \scalebox{0.9}{
           \begin{tikzpicture}
\begin{axis}[%
scatter/classes={%
a={mark=o,draw=black}},
legend style={fill opacity=0.8, draw opacity=1, text opacity=1, at={(0.03,0.97)}, anchor=north west, draw=white!80!black},
tick align=outside,
tick pos=left,
title={\quad},
x grid style={white!69.0196078431373!black},
xlabel={ },
xmin=-2.20852591900454, xmax=2.21306049239568,
xtick style={color=black},
y grid style={white!69.0196078431373!black},
ylabel={ },
ymin=-7.86405361849716, ymax=12.4391808546764,
ytick style={color=black}
]
\addplot[red, mark=*,only marks,mark size=1]%
table{
x                      y
-1.52762947004398 -2.63665679852607
1.74584146706459 8.30313418504723
-0.523478826928232 0.684981455327291
-1.73063347132838 -4.16754055824511
0.199724526357971 1.12413986692501
1.56583537106678 6.50212130573281
-0.463057519114948 0.747665654157118
1.27292037122199 4.25911843213007
-0.561300740541583 0.653996446870297
0.640763194967071 1.72674751050313
-1.73791342939199 -4.24096264611894
0.131982542924304 1.07163562445192
-1.83168405496555 -5.10376196338133
1.09632486787013 3.28860190499629
1.14869409291258 3.54804856022076
1.71062342784982 7.90152466385944
1.54063884052671 6.26233498072305
0.96036949117845 2.67188471681226
-0.36582225744249 0.829228691385662
1.03292095360504 2.98417928381415
1.57196496539069 6.54609386202508
0.316383914894403 1.23104885644691
-1.25339289598648 -1.13554910601212
-0.46394682067964 0.740299643711383
-0.998800252147865 -0.168876133428104
1.03613462412239 3.00699952570108
-1.64532444822729 -3.49293200198847
-1.97015264204063 -6.5769121469698
1.10611827964268 3.3325927944455
-1.40144531975216 -1.84518158676689
-1.71485775358785 -4.06883587132901
-1.55651601432609 -2.85466248166607
0.704149753887608 1.87137436411416
-1.35391349938216 -1.62114783985896
-0.19713428931697 0.91320639815536
0.890553629834266 2.41963800191651
-0.748635890255846 0.367757908994129
-0.0825495453396563 0.955542429515598
-1.52863234199303 -2.64905992413016
-1.06308264251695 -0.412431701954534
1.34829445246638 4.75822063973227
1.21677724541683 3.93402096649771
0.321064511493674 1.22796917544757
-0.79448963572049 0.286639405408638
1.29306522904297 4.38129536324555
0.991145476362639 2.81469205516578
0.408046183683614 1.33070158460102
1.50470887336654 5.9575393202239
0.241965856580386 1.14976159314043
-0.698625089638838 0.462631462074469
0.69277772759134 1.84475566662164
1.54249469168855 6.29223000018984
-1.37940963447366 -1.77397785165444
0.328552328283534 1.2358697609747
-1.76899063762277 -4.55181464197594
-0.1615118341785 0.917358084831994
0.979551973019646 2.76474222344531
-1.41400818915192 -1.93743608707598
0.920422122884637 2.5371272327229
-1.58813906344299 -3.08034023934407
0.567973302168688 1.574798756124
-1.7027550272339 -3.92661639975712
-1.66143052087379 -3.6310451778419
0.509387499079396 1.46884150329936
-0.649372469635067 0.531357119197573
1.27714939756953 4.29644865621161
1.99062920502101 11.3301714102496
-1.77767653093197 -4.60832173114591
0.95231808494159 2.65334945903267
-0.729339734657573 0.394415405360015
-0.69929992108586 0.469586842821553
1.62359205741297 7.03622721245334
-1.26350376928392 -1.15873251814128
-1.80495773362803 -4.87445067342779
1.4755075644755 5.72000202565798
-0.582140577898955 0.62882666918393
0.391595936124218 1.30969317341825
1.68359095958506 7.62901889841519
1.55014341230999 6.35669199872224
-0.461713591692503 0.760474080303689
0.837995109576847 2.25388599012062
0.0647111588263334 1.03231598849297
0.449501515148342 1.39080169134864
0.842450703096236 2.25731626328023
-0.662607691037783 0.514319433635468
-1.73347656811602 -4.21427270283707
1.97125407913951 11.0374500899417
1.06434018406928 3.1398766859045
1.24774144636567 4.10652488237163
0.301552235654667 1.22610695051265
-1.419054264094 -1.96851225748391
0.770735380506118 2.05529699022521
0.673663615667701 1.79826450863936
-1.20717318186629 -0.921117900268635
-0.410538504386376 0.765068471502614
-0.997783575545104 -0.190561556633821
-1.9247004506919 -6.06488872678494
1.67307707998314 7.51858011026587
1.13094409339522 3.46594756951591
1.05978614977254 3.1119281304365
0.368497683846726 1.27437039746974
1.19041117219778 3.77115160864485
1.53532876144759 6.21911480195912
0.977879415140715 2.75127965038619
-0.661130676558663 0.515515881624548
0.372560417149814 1.28014251982014
-1.41713935044321 -1.96759316741115
-0.721990537872462 0.427308795775999
-0.593304843043549 0.610478409782568
1.39782862521232 5.11259491405215
1.31945403787629 4.55774479220443
1.29797173081997 4.42915493648867
0.500221801881652 1.46197989702367
-1.22858048710803 -1.00927000019683
1.84901594258915 9.47648534922581
1.5150131259747 6.05407785409292
-1.22889583462019 -1.02559434722004
-1.4115062222231 -1.9348610640109
-1.33406864789987 -1.49095685339282
0.212784793779644 1.13119633500023
-1.35572868353146 -1.63622217917435
1.34923257581794 4.77328323295919
-0.117934710702453 0.944959877079549
1.74259492028609 8.25448605318237
1.79584748839092 8.8397486108736
-1.56683538172848 -2.9161474951396
-0.100670899179379 0.954539281248498
1.3130637409383 4.52476010808466
0.413802060319139 1.33744361353339
1.03952918317515 3.01003625023142
0.355758361430128 1.26871554742047
1.06587471434403 3.13952285378779
1.47827077764151 5.73417393797497
-1.423248862284 -2.00477990059306
-1.64138318843414 -3.44920568469485
0.858985773610331 2.31055719751639
0.30940079900652 1.21678960560482
-1.59532819132151 -3.11544889268605
0.674828392358402 1.79411719788582
-1.99791465316728 -6.92318765043389
-0.983915817371269 -0.150434192007041
1.05210422602657 3.08040038211607
-0.460520520996104 0.753445163781151
0.793577616390077 2.11102106409086
-0.916701943346689 0.0293686277642946
1.80568436851851 8.9740041455143
-0.41905914986584 0.775491286151664
1.4905136786334 5.85347534864318
0.659003889297849 1.76823560366373
1.33264972737708 4.64719307545267
0.367019067455185 1.27683076268824
0.284838881991474 1.19735122698525
-1.56586227893797 -2.91574117100978
0.865300493711527 2.33095887814049
-0.939391307968867 -0.027203254178731
-1.88319560164792 -5.6496600168737
-0.525375767528593 0.684583616019016
0.496152544580833 1.45177679093498
0.386055497083448 1.29164860153159
0.351299629919691 1.26288874499817
-1.65598010248621 -3.56432636874693
0.751560061269861 1.99687186057058
0.284661315211381 1.19235145843037
-1.09215628143104 -0.491573300112702
-1.10348826232006 -0.524640770030218
-1.00871955963086 -0.225586793602744
-1.54583226518551 -2.74904177692399
-1.55838780436906 -2.85682837025492
-1.03843932105837 -0.325694840270278
-0.946724153501291 -0.06049006215608
-0.654071049279125 0.525404355982477
-1.90500466748741 -5.86320237196585
-1.27255094354131 -1.20661238758745
-1.98026347546296 -6.67989024206411
1.63566363171233 7.14486566814899
0.0339713196359335 1.0093287980423
-0.50459247118296 0.703641359276434
1.81559893738786 9.07136373457222
-1.63792963751805 -3.43811065019935
1.81879951740473 9.11907419867196
1.78184093727109 8.68248370406747
-1.7721536426822 -4.55105725684285
-1.18740832757247 -0.855914984597267
1.38038656570182 4.99110878004799
-1.38783512640697 -1.78378026045502
0.222412885401944 1.12854784914632
0.756559645781282 2.00559445533501
-1.7219989420926 -4.1241360273851
-0.471086587199301 0.734105533462656
-0.227738431702199 0.89333543925946
-1.92935794945476 -6.12418892694745
1.1844290411186 3.73565569764499
-1.71501087455885 -4.0376542850963
-0.385903061117554 0.806071659467239
-0.890520314626398 0.0740593743595064
1.67287119794209 7.51115796744165
0.83918431164092 2.25650874404071
1.90289877922406 10.132856650788
1.03833740641711 3.01299290037405
0.789709020700827 2.10355162038978
0.683611022039177 1.82285368297838
-0.179507048827252 0.925081437742517
-0.000584604517955611 0.984340468812395
-0.520405730169149 0.693339417279343
1.92715956982417 10.4668434718492
0.400860795137212 1.32987293072187
1.19890765610203 3.81451374961079
-1.54460728333877 -2.74787026367663
-1.06248141493513 -0.406999040383797
-1.28672179531481 -1.26478539543189
1.25307005082172 4.1424154599753
1.42377930122459 5.31239619932744
1.70117716962273 7.82129096593986
0.179317666451924 1.12103935376769
-0.684304672411856 0.492541754858381
-0.868768205503995 0.139903632995063
-1.51891316821996 -2.57567499366925
1.36361253017276 4.86933885074933
0.0562416104519134 1.02593905383052
0.0758099949952324 1.03392765224946
-1.58113161808202 -3.01793810810587
-0.538980074292445 0.677534726364365
0.99270343344348 2.8231022276446
-0.409786318725182 0.789398428413281
-0.281142307585388 0.868639961816714
-1.69745666265164 -3.92386559000922
-0.370783589091509 0.804295606864759
-0.230187126417397 0.901809719586254
0.454070535352471 1.39744018922443
-1.49888350765559 -2.45181224972991
0.494126089847875 1.44129943661414
0.609683100231467 1.65241912894958
-0.416834592236302 0.790211066977293
-1.32562012581223 -1.46790146868228
-0.535019952703819 0.663209586913815
1.46284070470964 5.61791567029054
1.79853558911958 8.87579227219735
0.543964275193791 1.54269241151363
-1.12101167563638 -0.587903117238959
0.634585312093048 1.71296042187226
-1.42338564384549 -1.99995903396004
1.34323872565678 4.74249413175406
0.740479740191689 1.96380089060307
0.530420671662911 1.51309076141407
-0.652511983931383 0.521107180103227
-1.48154716489781 -2.34954351397374
-1.51975770883372 -2.59116390660472
-0.357623802961932 0.83381141128213
-0.873939423227781 0.125932421557286
0.819777580708701 2.1998943411007
-0.771908094224033 0.330724932627523
-1.25533466724185 -1.14123378490591
1.38048627129223 4.99804329532372
-1.41426900723107 -1.92000749326875
-0.931515071393668 -0.0230607378401753
-1.83612664651561 -5.19848319937411
1.91787399816632 10.3457425819134
0.822087947034545 2.20295568925398
1.58408998110077 6.67853165613535
-1.47666690758891 -2.32386558025554
1.59307040077941 6.74324021433242
1.82604877768948 9.19739450525777
-0.0643848869724568 0.961139915466772
-0.170247791580812 0.930901724014497
0.576780822247429 1.59740172217391
1.23866843250204 4.06837261252609
-1.59773061856609 -3.13496249784628
0.779140264391865 2.07272054088885
0.122981305271473 1.07007161647795
1.16018101227257 3.62053567845117
1.41737534405674 5.2489009736737
-0.125273475886996 0.945827754840978
-0.0752832772171206 0.969968969047981
1.60062278839547 6.80197787067499
0.672159440578955 1.78812444699988
1.74816251900159 8.29682205654595
1.52974515968788 6.17933789924982
0.818853077332124 2.19652244374451
1.17010851194556 3.67232407707197
0.669500173210326 1.77942747846152
-1.3832662134864 -1.7597563579927
-1.5056123093124 -2.4761572617853
0.676128404865428 1.80489610197411
0.354286326377876 1.26228019772993
0.110812810663127 1.05287390989139
1.85823905265138 9.59712772884433
-0.296690209866944 0.872379305015782
-1.87852985180451 -5.58878670482915
-0.561567284182094 0.647239410899878
-1.34797827041401 -1.58171359060356
-1.75450255430753 -4.40024872280045
-1.65803450911699 -3.59862420326409
0.899307716725425 2.45310985507215
-1.46062397932672 -2.20590832094584
-1.35397611455221 -1.61410551820944
-1.55462920986326 -2.83631479759145
-0.325694520641469 0.844135792066377
0.160385508824723 1.09385491556522
-0.0461815331773878 0.972733177334569
-1.99131143951082 -6.82898935026358
-0.974657684229868 -0.126816225683027
-1.1961708928454 -0.871963915385551
-1.08638984261912 -0.46817607178617
1.87310279781129 9.77910038124017
1.91185810187506 10.2346177433489
-1.14125893032288 -0.651143612967894
0.528322311654728 1.50440518127949
1.72555442756985 8.07019342203641
1.49516179008768 5.88943360017781
1.76908731382206 8.54319474425647
-0.98536492396836 -0.15700092510946
0.822905514907454 2.20331143461938
-1.9175384908525 -5.98579056838029
-0.634376195176915 0.557193522846342
0.803957485520665 2.1390454266119
1.8795561718596 9.83452954613513
0.247638382220518 1.15926441782563
-1.71883857673055 -4.08546459916181
0.391903377619543 1.31389415981742
1.03697635833877 2.99252053203022
0.46546627101158 1.40460445522065
-0.872255355483843 0.12452919017408
1.59963770381937 6.79624745217679
1.0905089794862 3.25222839325567
1.48364736751732 5.77669503094874
0.552152149540717 1.54712047919933
1.16305110024444 3.62821178894314
-0.764162395575607 0.353582332679796
-1.7996558729966 -4.79060711409738
0.257994234634654 1.159234494009
0.975819193435273 2.74312471916832
-0.374068096434548 0.811941703820846
0.732361402874449 1.94289288266549
0.841017152420385 2.25753945477373
-0.317325089675505 0.852558721839276
-1.14729985819558 -0.679405812472515
-0.84115142375628 0.201590836168429
-0.195889906726145 0.914464415551292
1.2198790586966 3.94647009708091
0.153095278361347 1.075407139745
-0.721646086143089 0.418113274550875
-0.276444233603721 0.874891330907061
0.997336602107659 2.83914267877571
-0.619971894263125 0.570644133479019
-0.987367251982856 -0.162659397984739
-0.854665346688174 0.163321947347233
0.191712135688579 1.11008010260121
-0.574488728271056 0.632500385245556
0.00863800784050106 1.00372572314337
0.307398788784521 1.21625487746565
0.342728984756977 1.25802380214618
-1.68981021594347 -3.85635974658295
0.784208901049624 2.08772138448435
-0.208369439125148 0.913988845682418
0.526648200820822 1.49948454846173
1.24202756976366 4.07702962445207
-0.727519971518694 0.416619346588576
-0.222858774302506 0.899970690899377
1.42817273628317 5.34391409346648
1.38769399086101 5.04512706115714
-0.41333614984213 0.781506781386756
1.30000904085596 4.43810676743973
-0.0370187504556303 0.973900549312141
-0.381829049097851 0.8006227077616
1.29945975418766 4.43145515107368
1.69798018089684 7.76034437975662
-0.978298825314186 -0.12745101956886
-1.66927279889651 -3.67091481981135
0.606203457135684 1.65452552319128
0.211301964437266 1.12170481513714
-1.00052625391862 -0.198810328948178
1.19345373102677 3.78511354553779
1.57544991206892 6.57360192278671
-0.77491023651087 0.332263920464899
0.159718334398355 1.08833688203287
1.45665398404387 5.56727126547789
-0.499484712478448 0.701589152523506
-1.16220758413343 -0.752574336865251
-1.02136658157845 -0.258844629313983
-1.06674862945718 -0.392686634605278
-0.744757010347744 0.379338217369195
0.623435205569478 1.68225418863393
-0.0599727657251976 0.964843293712371
-1.6688985298492 -3.67191162329495
0.00795195473762167 1.00095435444599
-1.58371970380259 -3.04223877021898
-0.814488977715953 0.25400468163354
1.62784684311749 7.06018606462141
1.10659216693293 3.33647695876888
0.746100055298261 1.97888290220787
0.004435567456055 0.993685051475752
0.273048244481866 1.17750929462584
-0.674494714349047 0.491965498520511
-0.630136777628819 0.559016438722326
0.404649191564105 1.32802772062289
0.867192246641215 2.34518405432025
-1.53462890319051 -2.69450269874947
-1.65797847179726 -3.61466929551865
0.936206661941811 2.59443165886383
-1.78249170379428 -4.64964642499287
-0.989577984804854 -0.170243909850942
1.74689924571509 8.30477247920927
1.36878773996584 4.91833067708624
1.85181571111122 9.50239345194755
-1.74027129194973 -4.26582547182024
1.32195369830832 4.57287695579038
0.779295097985163 2.08547213848138
0.527134401728108 1.49717159191067
-0.483385457457882 0.721013992566675
-1.98054469090637 -6.67787951840296
1.19919838050934 3.83794979312015
0.247147941406696 1.15388799111552
0.0973894309066896 1.05036888897258
-1.41470298624829 -1.91867719201216
1.22947358771133 4.00233387501312
-1.94803996825382 -6.34115402871578
-1.37522052503804 -1.70312718709698
-1.69508729758743 -3.90422830135512
-1.26400612715092 -1.17882495195713
0.177514997367698 1.10070307970624
-0.669640851903173 0.506276290801674
0.67112924906392 1.79045933660594
1.50698404505808 5.98164127150539
0.0503843989658583 1.00362100010182
-0.479836536391115 0.717791570047364
1.90580750567755 10.180201922895
-1.95045988746964 -6.37125640924513
-1.23772859446057 -1.04768655971777
-1.2491992534694 -1.09917023861416
1.43020405218296 5.35155575358987
-1.7958406700768 -4.77516846470591
0.6793613225212 1.8200682595746
-1.46040150163991 -2.20800476249903
1.73771962842721 8.21560467911202
-0.84730566648012 0.181473749770297
-1.00894166771783 -0.226811317886355
-1.00859420626261 -0.222584462396821
0.958638844017508 2.67815991960788
1.95044573476592 10.766490129105
0.0346133021748543 1.00994873094911
-1.05062638029796 -0.363128077667757
-0.0934869135279128 0.949219762303069
0.210153076911402 1.12020283425862
1.06304622638943 3.11894277339172
-1.4913720899944 -2.39900388339497
1.60616168357588 6.85479758268493
0.330881009604624 1.23690985922366
-1.41434542654993 -1.96193001839419
0.54827261775172 1.53292166739015
-1.05020516352889 -0.355253516364615
1.99084203167489 11.2963039571817
0.284498672856374 1.18803669814785
1.20635571406931 3.86687589245562
1.92884887664005 10.4605744830663
-0.787836212255184 0.299734103526356
-0.0665370859942191 0.958109369188055
0.266365751071064 1.17365638382964
0.284821733614269 1.19864673269428
1.79658899705357 8.84717765083633
1.36800009261001 4.90984160895382
1.71913709666286 7.99583556443286
-1.64691824710695 -3.47709331465477
-0.634577139004624 0.55915870436024
0.907611460044221 2.48490211512442
-1.31752746635936 -1.43220825506977
-1.25799147179641 -1.15868706829474
0.164578863753904 1.09458335531589
1.5724728529926 6.56004765078573
-1.61334505165664 -3.25225593055653
-0.884101245130284 0.0929428769915395
1.0785857201919 3.20491586289296
-1.93700500597912 -6.19805198060372
1.57399805208916 6.56459550935528
-1.71744427363904 -4.0708018208458
-1.05808128539432 -0.384282057908827
0.335865050135636 1.24121539994481
1.67249976327608 7.53111403985285
1.60127860974401 6.80875163915128
-0.517147816897207 0.69333488423439
0.641324841192098 1.72460209308304
0.676786295025203 1.80733330661154
-1.82794401593092 -5.09710000332624
-0.283331862331161 0.87196176011591
-0.0639641528388717 0.966729126839163
-0.506504957395847 0.707835320922324
-0.916554913203721 0.0288747985266691
-1.74864875263597 -4.32050250007612
-1.22417132700887 -0.990462261382918
-0.480427977583389 0.737086525022445
0.522720217749657 1.49225304695111
1.79977796292236 8.87678850979209
0.821654982549636 2.2002667735956
0.705451304778161 1.87003509143419
1.91752469627165 10.3449889069322
0.874637643884115 2.38294790188029
1.51463407476122 6.03663515683795
0.151964456223034 1.08602328208579
1.04920814119334 3.068717656264
-1.40228679868784 -1.8703882159323
-1.77433062051123 -4.56898824875179
-0.0807199561361247 0.950590991918523
-0.554066225390007 0.676013409967546
0.206793034906914 1.12162283249017
1.75495320295264 8.37161770096163
-0.58976643547149 0.6122332590153
-0.570265564796054 0.632088878798852
-1.95652629667675 -6.43236712937771
-1.4809655079249 -2.34962300810275
1.24656860874358 4.09436510331362
1.04805738906209 3.06644566689696
0.245998012934822 1.14616974570374
0.712250733621294 1.8985145336873
-0.178999679781882 0.937689573199967
-0.476724521976923 0.732967579431614
1.57873780306787 6.60316766960317
0.492905548214301 1.44673427787968
1.96692074715421 10.9935413794022
-0.866109375642834 0.137806550330372
-0.902579722541617 0.0533495254779157
-0.4836989537125 0.721480405653948
-1.13860622829876 -0.65812494496624
-1.68674389282483 -3.8269478035141
-0.696380684896323 0.462339397493708
-0.474718614621026 0.746846022399951
-0.55910870278528 0.649211265691872
1.82989710207585 9.26165707381093
-0.747639397853016 0.384910459455232
-1.21735597090758 -0.980617456389869
0.370162748415718 1.29014014231753
1.52374645718827 6.10901973485953
-0.973366382246718 -0.131009116090045
1.16171016526292 3.61697168798142
-0.862987828535092 0.161078907132451
0.621063996734559 1.67901339083638
1.65446335556188 7.34017160890469
1.77094426801913 8.54839534495773
-1.94388635474945 -6.28311770959015
0.761200781796524 2.0338355887882
0.0561172373005312 1.0213583364538
-1.06957567386867 -0.411001296441611
-0.0905481850787533 0.959183499127584
0.927040219844024 2.55424325573772
1.8673231985492 9.70422640604761
0.192108321913312 1.10176166062749
0.513399284106047 1.48857011176581
-1.94821049397743 -6.31319466428297
-1.1214590808955 -0.592631271826917
1.87606393659345 9.80087088545134
1.84450991390793 9.4353368625487
-0.049463327121563 0.977076978786467
-1.18160442071139 -0.836295878091666
0.304089057238058 1.20940477365912
-1.59280051013905 -3.09419187436563
-0.949890380048008 -0.0644734829889833
-1.55324124736082 -2.83858147926448
0.447202411694953 1.38807365079743
-1.34200036036863 -1.5262171878444
0.973498301084059 2.75044470769863
0.177052608250697 1.10709754591426
0.0405965700044342 1.01925157212804
-0.752339810673941 0.37289273972447
0.168360846647789 1.0924382059155
-0.147385650667953 0.935925901041053
-0.0765706138124309 0.966310494501428
-0.834073223482545 0.189431722979178
1.77400258345304 8.58815547506922
-1.10923369352111 -0.537029360030512
1.47731730787747 5.73058332341255
-1.97426591551265 -6.65678088227076
-0.46465919100746 0.743422947401185
0.255453290103878 1.18141978600141
0.401946609783347 1.33135091302226
0.7064603645273 1.87840239998449
0.195040572032665 1.11690200145903
0.492725116630803 1.45754533122342
-0.703358333592255 0.454883515873727
-1.77071940640349 -4.54100131963547
0.0263778409957616 1.0146304100638
1.5817323511687 6.64653301342941
0.188973280062564 1.12135526154422
1.00105500765951 2.85977988013338
0.665601496451699 1.78226977757198
-1.0758927035242 -0.439445577528765
-0.628467864929659 0.552521016834506
-0.462845639931512 0.743649977634857
0.875025914541497 2.37374222271564
-1.41383683667054 -1.94667718809135
-1.6697111938579 -3.68486173375494
1.09603549914051 3.28567858678585
1.62489834081238 7.03690114243045
0.416808458329184 1.33498770793675
-1.91346899374008 -5.96256237772706
0.48890415351061 1.44286865771006
-1.82702980036481 -5.05551592861472
-0.271371741436216 0.870328942223097
0.0838712538685789 1.04189708615533
1.24655440116632 4.10024175055629
-1.46183141783321 -2.22374166989591
1.28908473130799 4.36630170768102
-1.10855990085608 -0.562063462639219
0.933803589297539 2.58521487037548
0.744090662949441 1.9707947494304
-1.18542851197 -0.840491037489215
1.09662847195115 3.28482021885626
-0.891904969622721 0.0831156169344429
-0.753767622372932 0.369836018957886
-0.419699453159715 0.788081674304745
1.92783611898138 10.4584854988399
1.18434629483526 3.74481770501953
-0.850048612766949 0.176560415789933
0.308458399857949 1.21462782639147
-0.104193997970978 0.946005898372418
-0.799972736430903 0.285927701009198
0.452767708690214 1.39766719368028
0.597412080517727 1.64373993259383
0.593059513923374 1.62427200463893
-1.40250057451328 -1.8789264057244
0.798607175412876 2.13331105315066
1.93257294274525 10.5355389630815
1.99854849151513 11.4244399794226
-0.641668732087585 0.56475366966752
1.30016400132094 4.43814008524292
-1.5326104928455 -2.67493541839968
1.37502262423487 4.94986951331791
0.693800378736968 1.8515716644279
1.84517395726995 9.43599452837678
1.92798071190122 10.4753358247517
-1.40772643694339 -1.88214066370521
-1.28601120370353 -1.26647654672643
-1.26627685651145 -1.18784224817778
1.80731301646996 8.98912832188046
0.90850383276711 2.4863705914669
0.0979589637191518 1.05077299541241
-1.49663558808682 -2.44076085587208
-0.912437419464073 0.0290584337091611
0.877357979495537 2.38865486323248
0.166820627653863 1.09218585818911
-1.77159986131449 -4.5633205102713
1.30437703030859 4.47265186631175
-1.3499326931183 -1.60225917189635
-0.596182841232093 0.608163392297305
-1.87620318613111 -5.58351911620609
1.39855464279318 5.11359242740474
1.75283927641585 8.36942259144701
1.73162706612773 8.13826599634615
-0.766719221853512 0.340864384170174
-1.92279238876334 -6.02321863912057
-0.69318058382554 0.479668680603467
0.00831948249165038 0.999751590490404
1.48305548603697 5.76142487230455
0.153118246334721 1.08972946434602
0.58443191408571 1.61378691004878
0.236693925468985 1.14339170496782
0.646310871944384 1.74030084845718
-0.398390991704618 0.808366816187522
1.62011872039472 6.98953178832817
0.163170880205424 1.09101052685306
-1.05611175254839 -0.360223134163143
-0.450061164749038 0.760914089150395
1.08455697511961 3.23431779839021
1.07031030216604 3.15892859213176
-1.13369677432869 -0.640947173373567
0.954869425914125 2.6738604387378
0.190796234602346 1.09960866645781
0.34599075164786 1.25596865017729
-0.0417235161195215 0.971710379783804
1.6996522021834 7.79550916409845
-1.98068498864511 -6.68203605801189
-1.9115004631109 -5.93364414254743
1.53382768152136 6.22353879720709
1.31412729272728 4.51961993490372
-0.685326460329759 0.47418665950724
1.97256061584938 11.0547994580786
0.520690061019052 1.49111868744377
0.431498523136836 1.35789622587648
-1.54362657653036 -2.75460998020516
1.30076943936791 4.43975314877167
0.110705357748875 1.0486479677216
0.0688186092254197 1.02119784489148
1.67413450951731 7.53341721455009
-0.985207968505172 -0.160684540430415
-0.806089795979227 0.284444176437579
-1.85237667046108 -5.33233210062879
-0.920371820742431 0.0164668433527722
1.38141443382448 4.98863015282167
1.4044209338851 5.16669611818724
-0.542691600248609 0.675735850532089
1.53677101777 6.23346451860097
-1.85850522725861 -5.37882491327555
-0.444161392146742 0.763771158560042
-0.688501376551688 0.478927137327461
-0.593839572757993 0.606594772417755
-0.47441079581978 0.741345605082063
1.8831706950332 9.91285698162258
1.09169014263716 3.2517433215441
1.25389927608542 4.15534414662229
-1.59454200048405 -3.10600040525396
-0.442711305472325 0.762550665143076
1.95806353082891 10.8724704838653
-1.78460742639196 -4.68605248225971
0.40343560525519 1.3213194545941
0.652220564362389 1.7449551608856
-0.380836178645997 0.813053290144024
0.301719509164299 1.20512017543195
-1.21097150035761 -0.945192574788016
-0.029223722261007 0.986101783902391
-0.0358890414286468 0.967455226587691
-1.73775276817218 -4.2353480015459
-0.974754376510504 -0.1267811188567
1.44203622854204 5.45062670730321
-1.59890642565992 -3.14658288961961
0.795623049507605 2.12756096210532
-0.150922502332556 0.928688080958424
-1.53940041912788 -2.71723181405971
0.0487454727564445 1.01873554029746
1.96324448991474 10.9345743710465
1.66781793345191 7.45481886645438
1.73884391761524 8.21231629267577
-1.92817622064621 -6.11153730590321
-0.45376515808234 0.739633548853883
1.56230883155345 6.45409691855904
0.54828212047107 1.54415654792256
0.442437873944731 1.37678215206925
0.192524261196695 1.11685238085851
-1.18547716382506 -0.833172875923745
0.131479945441129 1.07644065677529
-1.09917176544636 -0.529014837658806
-0.470282670431948 0.730487027283281
-0.398534332506407 0.809759655887622
-0.425917514759681 0.774729820255762
-1.20315602945709 -0.901621628900615
-0.658854777674013 0.520887961468839
1.90279869634687 10.1535144403424
0.692627430585937 1.85850380394309
1.75702419283417 8.40248244830576
-0.218196717797614 0.908040475854448
-1.78916450307182 -4.7072198549223
0.129070596408469 1.07891023355875
0.605038800031316 1.65138168952477
-0.296316256804238 0.867483154147729
0.0164951209081132 0.988124447555131
1.21346412797034 3.90729131164088
0.574011131001371 1.59115392326239
0.462290671416839 1.40614477192918
0.40410819275721 1.32633180954348
-0.481595493587628 0.733282230587212
-1.31782552436269 -1.43093933687504
-1.5194837321693 -2.58315848179824
-1.62116265544403 -3.30344496903047
1.26310081801528 4.20151526945381
0.127202279000513 1.07081464024437
-1.68520250790848 -3.77793523091417
-0.603284843250533 0.581234017543621
-1.71673296848097 -4.07683475528603
-0.532651132542052 0.672261847207747
0.261358540370296 1.18512451859773
0.558581812877678 1.55736841573768
0.666199452291657 1.77515665700216
-1.00171427776995 -0.204103549526015
0.216887403452102 1.12488879547923
-1.66491700448442 -3.63645953599789
1.62371233364995 7.03647864252876
1.89298814563213 10.0075273939971
1.99696540172668 11.3985956317353
1.49857458851087 5.91810351458466
0.0259271242716275 0.99490040667827
1.25066053839307 4.12725819342969
0.741554342666891 1.97024401050884
-0.211134376905973 0.904265588112284
-0.21293296225532 0.902208372935693
1.17443099534439 3.68928772174936
-1.81670082980822 -4.96647251739284
0.236732509242013 1.13315323922089
0.643870709298875 1.73772272122517
-1.72227657824281 -4.13521216446846
1.63583921447564 7.14238562238487
-1.68186374417258 -3.78531224191223
-1.48826071885157 -2.37873968411393
-1.49310781433945 -2.39855010937928
-1.10696078276606 -0.543277142424544
1.12452649312155 3.42889377843165
0.783650456241737 2.09083865222502
-0.161911748499459 0.931137194158423
0.628007684776644 1.68442530419979
-0.243924773625153 0.893116388521999
0.834772814765674 2.2351337827488
-0.223841541824117 0.900516913726931
1.07132477350497 3.16306477307216
-1.18266111651542 -0.832822823434673
-0.428500086152374 0.778049775291733
0.90169194696583 2.46648254129569
1.09745444799678 3.28796605696869
0.523909850278403 1.50151748649013
1.73717403201051 8.19682065367904
-1.22113069512167 -0.995120760996621
-1.20955849396782 -0.953255169415617
-0.399486529556002 0.808781879310853
1.45060068843825 5.53087848795961
0.495836028886981 1.44879256164344
-1.54563926520025 -2.76119064927738
0.664636176471893 1.78163979951814
-0.84536772623423 0.195978403180779
1.84784612583578 9.444972752919
-0.102234273221019 0.960155648523784
-0.92791030503047 -0.00201279096713414
-1.58862031670824 -3.09677838620582
-1.40597534148641 -1.90921472245161
-0.0567426886504618 0.967309510794382
1.95112557211344 10.7836771757371
-1.16202891862751 -0.743642158251616
1.21626034128942 3.92421811806867
-0.80452779989253 0.26609097967275
1.92498870321335 10.4303878655003
-1.92515688031637 -6.10821307405095
-0.852504867179035 0.170600669614789
1.55705189462523 6.42445316743795
-1.51788427122346 -2.55981081658546
1.3689591098897 4.90303972175912
1.75972680655051 8.43497694100864
-1.65957374583011 -3.62665673598408
0.270920890606869 1.16949282304471
-1.16885208042951 -0.763987666325113
-1.38447964389345 -1.78579113660629
0.177339566960526 1.10697129299923
0.543295442544644 1.53305438086109
1.81238014720433 9.03883120895235
1.71910415517618 7.98951188473036
-1.38574884153756 -1.80453591712449
-0.866168705486111 0.13155224998541
0.56542842101778 1.57148386121625
-0.999358479402314 -0.195291929222586
0.517148701909557 1.48294034813878
0.385827101406312 1.29470435488792
1.5126698410979 6.00707691823652
-0.847485819932245 0.181307650065231
1.37674588490367 4.96367311113197
-0.836170270341625 0.192930597010434
0.772951559632365 2.04951697969711
1.59583203319957 6.773635455285
-1.82180810974138 -5.05045903144996
-1.69300170047109 -3.85896485817725
1.44794960970744 5.4999476493892
-1.66109627169685 -3.59464056731823
-1.7573382699851 -4.43118186833144
-1.10719019711833 -0.532950161422905
-1.51256371702266 -2.55738152032505
-1.2165312811155 -0.966154532794749
0.879251925733673 2.38729259303622
-1.87054596694108 -5.51283413585285
-0.441182691897304 0.771939195941864
-0.19493987501971 0.911954170378767
-1.37728571964468 -1.72844840209457
-1.82483177421839 -5.06024739494962
-0.702010881428292 0.451843329310979
0.326122906767451 1.23172693663433
1.16851737260205 3.66289673262974
-0.214871572242574 0.902305420744069
1.31468918333912 4.5327764118568
0.62424958546635 1.68687541631648
-0.41248770407115 0.78026988517054
0.462518280013717 1.40117426976985
-0.876558973447731 0.121068561770403
-0.445309663339691 0.765122542337629
0.286003384268398 1.19420332378833
-1.38403511544209 -1.77359364826383
1.88068779520718 9.85376193718621
-1.6185680719818 -3.27523195342785
0.152180559735233 1.07754121780664
-0.109158171926152 0.94797019633471
0.24651978919582 1.14965891306799
0.628965416458113 1.69882274744483
1.91935588472357 10.3540013605604
-0.0373830886566422 0.976154099548856
-1.68054947050608 -3.76688627554973
-1.87451027850194 -5.53804599169035
-0.555079240761019 0.66428892526408
1.43463336473141 5.38729732564786
0.296471880579543 1.21214083308364
0.596403994269193 1.63471382573155
-1.9599885782736 -6.46862056778319
-1.18377169591759 -0.835764211058809
-1.06598170603066 -0.406347917054451
-0.0744541532140537 0.970544054476839
-1.76337010414204 -4.47980815903837
-1.09667441590765 -0.497958138448634
1.62826999259492 7.09346422247264
-1.66496683308716 -3.65859501527713
1.1752654100337 3.6867880267327
-0.711313673950508 0.441159206437018
0.949181662762235 2.64913485004861
0.23509675462388 1.14631030410878
-0.132282055480576 0.951380021157754
1.15162918388688 3.56590005483656
-0.825309704656912 0.235428161007987
0.0416208378466676 1.01916426816218
1.28237397513989 4.32803419372785
1.67804663468333 7.57560936511758
1.12847037380373 3.44583730827938
-0.666677668801422 0.500535466902455
0.0771935567965247 1.0251579835824
0.382430291118693 1.30125913420781
-1.75396408351992 -4.38684640189132
-1.27713484050294 -1.23728764579277
1.26371408079828 4.20747471043581
0.631521703007308 1.71097872801138
-1.9927489392525 -6.84210703404694
0.429833334690461 1.36756426971277
0.885000450618656 2.40899164257351
-0.928739515085039 -0.00112897243548949
-1.79090100863839 -4.70858495483779
-0.190402521964136 0.920291720165465
-1.02277075142579 -0.254672518218796
-1.98986414127296 -6.81967763997728
-1.56202189656636 -2.8700002969686
-0.823305383673487 0.240385903454283
-0.16535725880631 0.924942878158168
1.60375707727229 6.83113643466909
0.473228880387381 1.41554736679962
0.265718505149918 1.17273273686391
-0.0597395478923923 0.971368513214256
0.300775466668908 1.20781495684076
-1.05886021200778 -0.380369566542633
1.57753799391147 6.60566982934968
1.2678036676601 4.24054628185889
0.408969622909501 1.33434109200796
-0.802400565314632 0.289453107342089
0.797390436275073 2.11921434311487
0.269292520108907 1.17748066233802
1.09863981296441 3.30238844951017
1.66270555828155 7.39872202244838
-0.0678474745177926 0.961163352850507
1.15694791286823 3.59120166028488
-0.811165066606464 0.267983807143246
-0.948634762579282 -0.0490322431644212
0.98120715897871 2.7593947116178
-0.228724559837756 0.900312851955522
-1.37228948746877 -1.72229450435978
1.34818409339734 4.7540661109
-0.63180800074183 0.559632310925091
0.969694764923288 2.72146353859967
-0.15563669746887 0.925302312133414
0.732961867207918 1.94489665160677
0.985039768614193 2.77722119348587
0.320869401360116 1.22146660615769
1.73657869247921 8.1709769068018
-1.48569607335105 -2.35751632668528
-1.67052307021378 -3.67684260041769
1.80986213936798 9.00751847410175
0.269901033496012 1.17497245115711
1.86964035709248 9.72053894079343
-0.618116427779094 0.57948319057898
-0.559013229623111 0.640029635332599
1.86053966331359 9.62113266205902
1.83639477816247 9.32089258884026
-0.144643275194269 0.934702694488004
-0.661368634771165 0.518005380333269
-1.49996211674519 -2.47060303978315
1.61565620264765 6.9556738208655
-0.525992789448132 0.690153670706809
0.616639082134211 1.67233716321322
1.11048851688358 3.34846455523854
0.490327890360109 1.44696062751948
-1.70630821016035 -3.97942357433033
-1.08432082336125 -0.447866032025121
-0.470247615489478 0.737040127491929
-1.12330616920672 -0.592688203625215
-1.01151194857804 -0.22531529845062
0.93537143033542 2.58899557709597
-0.780955892956971 0.32494943131485
-1.64128116553575 -3.45549750121784
1.19263378023814 3.78334139440739
1.10397943127227 3.32477236366575
0.564707687955063 1.57096864135258
-1.07973399874306 -0.462843967286802
1.20149098397145 3.84041420340939
0.695933299560299 1.84328315610021
0.495176883716653 1.45111124871371
-1.67550375125215 -3.72282087079692
-0.529803028699704 0.67391141147433
-0.445856021780258 0.746872603821456
0.433812532010605 1.36821956949014
1.54202085374851 6.27014082733516
-1.32375381765984 -1.47198532477877
-1.62134922807424 -3.32157292176848
1.51933493963951 6.08477696358128
0.471822244683532 1.42329463587884
-1.7933624880284 -4.7584821763407
0.596742359496119 1.63294109206395
0.552130519421354 1.55492174251572
1.47710227591609 5.73121027778143
0.00302994761583175 0.998517527547431
-0.236203654112023 0.890720008570176
0.970165770754015 2.72487587611247
-1.34277493490018 -1.560845574621
0.206076871985257 1.1257152390431
1.46412316541014 5.63221243541157
-1.80601682524143 -4.87953733752454
-1.391266729644 -1.80715428918603
1.88631551186505 9.94610960451085
-1.10746678059245 -0.531319186183653
1.05027379222569 3.07342814811358
-0.542183740128499 0.661768181124191
0.0271656137687568 1.00594390041955
0.298087731551128 1.20094009138946
0.434215781068596 1.37296325698884
0.853518173197716 2.29721032854808
-1.02478655719921 -0.278179128954753
0.930280869478582 2.56653063252946
1.282077171204 4.33841398892212
-0.867967962223037 0.137217640904829
-1.77505278071821 -4.57906190273618
-0.485529113216811 0.70779431482014
-0.935709399072525 -0.0297072239274747
-1.7124804972844 -4.03787929313712
1.7597432806956 8.43146031674183
0.756160878053701 2.00445575272812
-0.00706127845737115 0.981964054736906
1.8189676207847 9.12539100938857
-1.43708373555605 -2.07284379552645
-0.391867214393602 0.81005629802083
-0.125622315968765 0.948908915942806
-0.227608808523419 0.901222918334211
-0.771183404944305 0.339516244786058
1.46047134279569 5.5976889022238
0.467766545173386 1.40672102796338
1.95848779282826 10.8896618409186
0.864335094856651 2.33375476079497
-1.63506328198021 -3.41060197363597
1.71568012752807 7.94562910214681
1.39272912489522 5.08961515331874
-1.22465164856736 -0.999370845561272
-1.02616643763754 -0.281109670883781
0.442734584377311 1.3636453659406
1.22412292980746 3.97718119164815
1.97421205804188 11.0926394512552
1.48145057673646 5.75859204672228
0.819496054168538 2.19628102305374
1.4980055219723 5.90529065694513
-1.18573895693203 -0.831044404668815
-1.99278274877207 -6.8494856711783
-1.91941536557922 -6.03689343835302
0.131241931655927 1.0600975339464
-0.580657766359237 0.621323641971421
-0.888088359400202 0.0879748598093739
-0.499052602078298 0.710622739621337
1.71820673011863 7.97388064367148
-0.316414951511492 0.850172933138876
0.393727886022705 1.30739732012249
1.00589706885763 2.87467650861009
0.57688395039735 1.58492314006907
-0.426697312104449 0.778529184768187
0.375189921780094 1.28553206713643
-0.813630157701015 0.272796922860551
1.2372573491537 4.04711291749065
-1.53832901940263 -2.71516049093013
1.35062857979094 4.78353830796967
-1.33694973116921 -1.5097033149936
0.564734755108878 1.58133376269604
-1.03761236511533 -0.304364625652098
-1.7347823407846 -4.19876453668963
0.945711070922103 2.63280043844309
-1.95258942984414 -6.39542582524028
0.0826143577681773 1.04656434275939
-0.511421420547976 0.700390007367014
-0.689362310384055 0.484189815360375
0.954490563853327 2.66189127695141
1.66434830933929 7.42551155767517
1.83138414083698 9.2535906483007
-1.02932672246069 -0.302454377177136
-1.04661156110905 -0.349194696608168
0.805645153296944 2.14727751708212
1.77401916557127 8.58561787796415
-1.40235998658197 -1.85626609020103
-1.0095984473455 -0.219611478121682
-0.270919141607282 0.871927628762499
1.2570941367174 4.15531025075
-1.86122822980676 -5.41679154352957
1.84661759143133 9.47138804968631
-0.90332609244693 0.0560399363281534
1.39563280875743 5.09338234894223
-1.98651168890195 -6.75881861620785
-1.9996278349795 -6.9117747602424
1.35079226518494 4.77394550653772
-1.79705482162434 -4.77476543435062
0.733839749036802 1.95167204064777
0.825728470288538 2.20865121947308
-1.999966065817 -6.9613984030445
-0.843299266844224 0.208944452179643
1.24300068592906 4.09229023697789
-0.568510673671339 0.628967296580392
1.51908286676269 6.08838913446026
-1.87867682212743 -5.56365333032191
0.51612906772595 1.48055749728735
-0.859068978962361 0.154159105057622
1.94861230872647 10.7340006651497
0.0534360058709065 1.02654421919225
1.15519390843326 3.59722956870279
-1.74527861230335 -4.29824285096523
1.66524052084331 7.43899384484359
-0.325981587735932 0.836218281066738
-0.752689738584645 0.372297506566086
0.128978317546487 1.07630978592339
-0.5621304215494 0.641110563792218
0.955494604929388 2.66243994646497
-0.555682648203415 0.657877220211347
-1.94548297516413 -6.30660150158887
-1.43026669699117 -2.03133524617525
0.356931141155441 1.26497851786683
-1.77861129362684 -4.64027561767086
0.986070692175463 2.79238643308185
-1.12135417504645 -0.583873973582881
-0.874192151028235 0.121619859150448
1.64662777642106 7.23880275344668
-1.36031739836864 -1.63023469181487
1.43962972373616 5.43679065708496
0.452585949580314 1.38020612816226
-1.76731956908791 -4.51087184778122
0.407565146203995 1.32590473134467
1.98898051125618 11.2966858096881
-0.598156498020338 0.603843415879494
-0.331884353093339 0.843716252343694
1.9650023216185 10.9636901570038
1.47465463159076 5.70579701648684
1.5204818278323 6.09842556468083
0.382297270832394 1.29907946604983
-0.308472069208821 0.86452974033168
0.464213345558603 1.40962225619966
-0.313121928977522 0.851139529144977
0.404197074317428 1.32297254888198
0.106771342877163 1.05328070645161
1.23841881398682 4.05574132595654
-0.282973759325347 0.86927530549534
-1.73525840893654 -4.24571728575114
1.20188368354561 3.84484408833895
-0.429258807295439 0.781389744582732
-0.748427956519946 0.38035750152032
0.0804517475120305 1.0282753335618
1.97951469313928 11.1614158365608
0.124479142005626 1.07152735282127
-0.68196947072659 0.478051156977546
0.223057422568356 1.12770066107423
0.469529140163063 1.42042477119935
1.81420517068728 9.06347397745196
1.09682487909601 3.29017544589656
0.914181350597536 2.51081156794732
-1.44096664107871 -2.09895744289672
0.659760682063931 1.77014541399254
1.04063674372985 3.0245021484701
-0.410649290239818 0.772472104365948
1.18214697910478 3.73390336265526
0.672531383775982 1.7818910668283
0.833061843126512 2.23739284468989
-0.617849314672589 0.576766640827125
-1.07259635030233 -0.439464162323758
0.391562842556896 1.29947440262709
-1.32972617708255 -1.48797880023631
1.70646260742689 7.86914077390867
-1.30706265571546 -1.35497932712865
0.187434893599617 1.11386759115505
-1.73820046861836 -4.24171387296479
-0.953602384085788 -0.0634957917389822
1.29875352212481 4.43489366176754
-0.103768824953191 0.943596096010008
1.78570975428274 8.72670519279494
-1.53634762343696 -2.68365510961318
-0.845841128729985 0.179914162737134
1.96130458726239 10.9217955746843
-1.38123143847418 -1.74150095436957
-1.82357045233629 -5.05015042661105
1.17694940504624 3.70540656238348
-0.810272952936271 0.261757370921051
-0.991398036985326 -0.172896205243066
-0.961061734282657 -0.115018311465007
-0.214642191506338 0.912795507245668
-1.80569826891484 -4.85454330736145
-1.26660129022973 -1.19561802966693
1.71984917692486 7.99792161372587
0.855714417828803 2.30130420773794
-1.9857480554809 -6.77474122630028
-1.72758247087356 -4.12988884755024
0.400499487781533 1.31846776525232
0.114574489690743 1.064962908841
-0.779632733148505 0.32180448146621
1.41151444990222 5.21844113975241
-0.899575225078841 0.0588786822150742
1.24093702114785 4.07803491765101
-1.67889529411105 -3.76011061763772
0.261032694419074 1.16507301696946
0.145927319726146 1.08550642394716
1.72277198899299 8.03579300253467
-0.881853643689245 0.101790168217609
0.799560413958917 2.14568536438307
-0.0893640599908605 0.968936056859639
-1.62447012296027 -3.31089685950611
-0.818707933572382 0.235670890205997
1.41743799518074 5.25552410123793
-1.2788388253118 -1.22818174476524
1.02741670432943 2.96620221224974
-1.75287345650198 -4.40168189949109
0.638185008569665 1.72221985914351
-0.566120518358522 0.650562467976798
-0.400109521832959 0.81283133976865
1.28489563036768 4.35346886871966
0.458876720195964 1.39828053433356
1.66049162432052 7.38465069349612
1.42024196877238 5.28362108847663
1.48326877753877 5.78490256315112
-0.588650003265751 0.610370712375689
-0.927367872736143 -0.0114696038192499
-1.02199237234107 -0.286769881922343
-1.81491938247115 -4.96943751358194
0.198105630543795 1.10133534625769
1.05566735595332 3.0952690904825
0.216283757072662 1.13320891762436
0.0952747408355625 1.04483947637352
1.88356421222412 9.88305662489422
-1.50949082880058 -2.53185428408598
-1.61612601851821 -3.25947685794374
-0.355798702545588 0.822434457477091
0.389974202292196 1.3071554623711
-0.346350159364043 0.837496587477261
-0.731906551736923 0.4231017017737
0.406415533261588 1.32501739654885
0.718342248995067 1.90723924704651
0.0356803131817549 0.995954228879776
0.32991610507093 1.23539495995339
1.42183905565164 5.29867881100847
1.11622051534856 3.38301214580797
1.92780392950029 10.4727861730491
0.674874856790324 1.79951472002745
-1.23425008498926 -1.05251765071292
-1.3974131057156 -1.82541063839319
-1.43099673513162 -2.02815987997362
1.82551334714146 9.17010534096805
1.79738352841659 8.8557328134358
-1.88660253039814 -5.67393300807767
0.26405973820585 1.17900471321331
-1.65724856526241 -3.59570924152451
-0.656847438749077 0.546819432088693
-1.85985870708892 -5.41022761275473
-1.92669888790178 -6.10120496681913
1.29510024203283 4.41730776155684
1.21689615980552 3.92920219065836
-1.48920215574009 -2.38781776744757
0.169083726631659 1.10055457991979
0.0312554397253666 1.0056786571507
0.117235240035499 1.05365418816086
1.61530394923782 6.94640681714602
-1.08506533136718 -0.471932555403605
-1.50655348044929 -2.51834677153378
0.542047461016566 1.52317408442638
-1.49196218005367 -2.40584651618691
-1.62953219152099 -3.38030741319258
0.653952115193265 1.75577793997946
1.16589918784997 3.64381270220462
-1.77321322818014 -4.56469465484214
-1.44063688957111 -2.09918107331092
-0.786776926887563 0.303985331212182
1.56757347418732 6.51428252945648
-1.46068373870209 -2.20864555530647
1.51617594077454 6.04677205451449
-0.576313733093636 0.634815342485373
1.34090731590764 4.71533528035314
0.521243193049925 1.50107230710231
-1.23207534268808 -1.02999184790934
-1.19751185489956 -0.896932767351459
1.42298322693977 5.31141137619722
-1.84564276286538 -5.24061024182928
-1.31186592728896 -1.38985918746882
-0.922858524377307 0.0126962294143067
1.60549991204032 6.85970352455356
1.1792532466363 3.70185784389861
-1.45845606087949 -2.18931104790498
1.97842696473836 11.1559755107588
-0.889592956516454 0.0881258849970086
-0.159391035282863 0.93137945195953
1.41663934420473 5.27030191748982
1.72248227506626 8.03246200431787
-1.76016432377812 -4.42610150387204
-1.82777538997854 -5.0940822768977
-1.78206005009434 -4.67457738753414
1.46006475348857 5.58108255384519
-0.938027545334918 -0.0241690864101531
-1.89713331915919 -5.79558084756152
-0.569747605280394 0.63397940911139
0.124837037063116 1.06821918165989
-0.684232287818634 0.491088808639237
-1.83385619125808 -5.15769800925519
-1.26330160100964 -1.17849027537144
1.37217992254188 4.93689519001865
-0.849073677961068 0.167156637959228
-0.0294605789009212 0.979995443089457
-1.75569733601588 -4.39528696677677
-1.8642781527774 -5.42629141379155
1.00355790811262 2.85991001296316
0.215009660003008 1.12504404521865
-1.32660668313371 -1.47256222141998
1.38069310867288 5.01125720875715
1.86924894822072 9.71332902292094
-0.190122966059632 0.921551851526446
-0.152887041387415 0.925827416569911
-0.0428840413177665 0.979160071851424
1.65420267720902 7.33863979494432
-0.900901126415482 0.0560309237293428
-1.90168442933663 -5.83625996430053
-0.779415341955448 0.322351078873821
-1.45067678889729 -2.15145455279476
-0.998250686285214 -0.171300019597086
1.66790312588376 7.46235873854451
0.28681978794021 1.19691243674556
-0.484523640375409 0.722736324751854
-1.6555459953985 -3.57807298722176
-1.5455072230979 -2.77803191414267
-0.688771224907164 0.483977747667041
-1.13893090991256 -0.648782171993345
1.74139892441216 8.24162070521052
1.49893617150363 5.90031251014836
-1.08765843446177 -0.487572985944805
1.7835459703997 8.69772460480078
-1.10577444808059 -0.535344053874771
-0.905368974653909 0.0580592417348249
-1.75191942613126 -4.38392739531254
1.93356751038377 10.5522770070376
0.9575879972616 2.67385742229081
-0.796119159539778 0.275863167409546
-1.17064460231169 -0.78528781265218
0.526940860518066 1.50994238785575
-1.18559700128084 -0.831772801428173
-0.0335359284307351 0.983857790419807
-0.818512156031854 0.245098220106386
1.04792454807403 3.05831772259283
0.127512311118942 1.07848873391252
1.15092609088398 3.55360510416694
1.37906280384393 4.97418955877048
-0.0473423907324682 0.96888765579321
-1.81903180933046 -4.97334268297504
-0.67902321115456 0.493726309635186
-0.798945544129008 0.283558204400877
0.598411739509126 1.63323878745168
0.301391070356309 1.20040807529484
1.47141508416494 5.68106703834634
0.883240667186142 2.40434783390999
0.156345836899162 1.0835720610823
0.749642134354526 1.98493438217823
1.49307307902606 5.86243439468387
1.0051133217067 2.87566950658169
1.22060969375902 3.94798487387617
-0.360249201898012 0.818908380812852
0.925829811942897 2.56283549067487
-1.3265763354304 -1.47290186967351
-1.79276493501864 -4.74355242164767
0.666266542292513 1.77844030769164
-0.95305747324246 -0.088946051202931
1.92203859375713 10.3913976509573
1.60206068652581 6.82197328090426
-0.424429206266169 0.787192414010418
-1.80954533229095 -4.92551946312734
1.91009723362869 10.2448499958386
1.68732742763117 7.67365132459873
-1.63516046158355 -3.40894235039875
-0.459631207281084 0.753705430502516
-0.522873646404107 0.68248796383917
1.59857006948995 6.79465092547699
1.03205890678154 2.98085490625002
1.59776386078881 6.78937914715549
1.15190579811508 3.57596356932545
-1.84684367541041 -5.26939484200659
1.26353314362285 4.2062199699076
1.35352808434662 4.81378397228656
1.40621035681587 5.18601440856323
-0.0272575569321387 0.971616125992548
-1.94627739820394 -6.30517712192766
-0.320949431762209 0.853498206307889
-1.59390277865162 -3.10199050406274
1.03232614482696 2.98732112354215
-0.288476953259873 0.869825858368562
-0.111058309972998 0.954146397701777
1.5736115516678 6.55485856481457
-0.860588998866321 0.15642658047556
1.39521178802958 5.09339707406236
0.0423370688837927 1.01865161261728
-1.91965920542692 -6.00410035946931
1.61115595534765 6.91347479932424
1.4840555846706 5.77788686376411
1.34085225582999 4.72451646530128
0.595661324885158 1.63354063840449
-0.138836632916785 0.931557022836676
0.156291343415457 1.07807508598658
-0.199381414145481 0.918076848212177
-0.559238193346088 0.649944599417361
0.0290059939260847 1.00886399855458
-0.698820135672535 0.46917212264515
-0.87237557527917 0.123745015732272
-1.58195014208742 -3.01173280754049
-1.9689231345982 -6.53962061378453
-0.410543849357809 0.774965062938156
1.3643124949043 4.88762618880428
0.823704960749702 2.20013420747358
0.674083822044872 1.78666363937089
-0.0325492464892863 0.963388674789698
1.94758130859223 10.721755324163
0.612287849542449 1.66267016705614
-0.690234865295915 0.477630742418045
-1.48019927598632 -2.33678801276316
-1.13292069294096 -0.626449522089767
-1.16648749781285 -0.747041932511638
1.3021658674702 4.43795673604559
-1.1482368174498 -0.678728219680581
-0.328263345948502 0.836237638536323
1.07923511541064 3.20608357727173
-1.02190446706387 -0.265216820748055
1.24442771631325 4.07920038204131
-1.65796184062742 -3.57106322333758
-0.846881701827113 0.178251765787283
-0.195611341148256 0.91061814545446
-1.25153965425888 -1.11718835843621
1.09206189549521 3.25655635179998
0.787016617069217 2.09829918227663
0.638861992142554 1.71771868617151
0.399893941960438 1.31866926108832
-0.761502475471771 0.351803496175628
1.0527372198015 3.08520865018572
-0.0214878931824463 0.975799084075638
1.0568998378917 3.10330396030909
-1.46873068032241 -2.27852019339702
1.54491279479151 6.31880603915764
-1.30793049787522 -1.41517139604437
0.864324612137999 2.33487841538761
0.208219021114417 1.11889714003132
1.70114170828014 7.81190050955589
1.06753378547388 3.1404300027043
-1.11001488165708 -0.566475339272693
0.900206776703087 2.45633244911888
-0.971142814603263 -0.114440229600136
-0.588370612047532 0.622154241959811
-1.36107915505973 -1.66834163500672
-0.580659726949028 0.62587338231478
-1.71050106264944 -3.98719767008242
-0.838727982174718 0.200416642319025
-0.248301588894354 0.89533448102124
-0.989899385886345 -0.170950406653217
1.45858345806112 5.58376216154754
-0.827374754255352 0.227448429646508
-1.28680681315339 -1.26625444942153
-0.560470344749795 0.65092888444907
1.34103519777274 4.72632065303982
0.854124982911279 2.30360432367002
0.727499969443455 1.93095709065338
-0.831947031058761 0.233215952035243
1.45247471287064 5.52279789355315
-1.98284281776198 -6.72926578696236
-1.03189856907482 -0.304522152519759
-0.340538292167053 0.843966930170733
1.13223942151434 3.4621697814609
-0.25423408601707 0.874988364124284
1.15671650638182 3.59695541308704
-1.67804201036224 -3.76369997152317
-0.55737849175 0.641312139375349
1.22284522480552 3.96610427024584
0.303291470771417 1.21458446684847
-1.74160573469583 -4.2921974139226
1.51465156566754 6.06014896988245
1.85431344649103 9.53042038623335
-0.958264059974216 -0.0779086575288621
-0.181367974652247 0.921003247623564
0.354703920530524 1.2681057921221
1.5957858792884 6.77614092967193
1.88579874047252 9.93153684685712
-0.708036601941509 0.459243219841501
1.23034559038518 4.00717188621409
-0.842252021191427 0.192630036329356
-1.71655053080056 -4.07657785364739
-0.881557445068461 0.0991392922692949
0.26514318781149 1.16717793776953
0.232989434111001 1.1482522835511
-1.82715605861806 -5.06253576136783
1.1497195166075 3.54793790777054
0.00757987893164103 0.997700018868247
-0.61739697037267 0.5735170991148
1.50647848392169 5.99546984833221
1.39166169440554 5.09328381826458
-1.96514402110851 -6.52215451332572
-0.838737524950258 0.194299200340564
0.138339494100933 1.06186840384019
-0.213443547284491 0.904404608021305
-1.46046481821283 -2.2080135401958
0.609831401488206 1.66300067914825
-1.26495257795061 -1.1561069134446
-1.16687188815107 -0.778955859164267
1.27510839832142 4.28933685985629
1.14154528439393 3.50920383108221
-0.778864741792832 0.327268260447748
1.63754071403956 7.15964537774718
-0.643346611279531 0.549356031200366
1.67127472441855 7.50541451384146
0.111885802770817 1.05179206964457
0.882107696941235 2.39247794919952
0.91120952438826 2.50577477205716
0.0119238114203624 1.00140676882301
-1.38083405007942 -1.77009842808448
-1.90610125215633 -5.86952783934642
0.175143291573622 1.09829986463524
-1.61207890964624 -3.24037009230353
1.5512325813554 6.35356088789581
1.65048286681893 7.28520831121552
0.188159596204922 1.11595697837751
-0.0382019512333382 0.975155447665413
1.4540122369886 5.55071321768358
0.432059746865559 1.35824805871841
1.67199448097685 7.50303519222687
0.257309056406947 1.1677196249847
1.06166139347091 3.10876862857168
1.17483889410732 3.68968549532277
-1.26118992652891 -1.16115793532479
0.382797281524637 1.30668725064353
0.889016236667162 2.42748823070005
1.8514711371672 9.50847084255472
-0.0193655465799174 0.975163024843158
0.496029176730886 1.45262652122023
0.870030643099832 2.34669385276523
-1.76505887531598 -4.51997178880659
-1.49081964038111 -2.41529704137638
1.56004730908803 6.42839628279718
-0.566720010329685 0.649859838998548
-0.000648148169025031 0.994814329710489
-1.36290341870177 -1.65567881249887
-1.93639048032525 -6.21241344520818
-0.981602974126864 -0.149832691429882
0.853315405048908 2.30571530395305
1.82818383675168 9.24487137613128
1.97604295543575 11.1219852538549
-0.545612270990488 0.667582385127273
-0.155994484131787 0.926496782955658
-0.674605850503229 0.496796399880046
0.0708290577944006 1.03235432288971
-0.71072226618099 0.445681361775277
1.99239397624334 11.3604064288272
-1.72307522658745 -4.13217941900431
-1.3528138003879 -1.62026796776499
1.34167153626317 4.71668204199445
-1.46411382512332 -2.22783314293238
0.39152830503833 1.30147245218844
-0.908375413808081 0.0288199848965107
-1.94765578353395 -6.31316002076307
0.85951902865061 2.31840828190778
0.44754279041818 1.38058500588807
1.0067674631368 2.86274815808275
0.326349266204949 1.23939686078795
-0.540902655479463 0.660366168807756
1.60714789447584 6.88435406585097
1.28467265919759 4.33791827404069
0.418555441252085 1.34600156677161
-1.19450012731273 -0.887995026790252
-1.65377266173249 -3.54407195701156
1.58540353990366 6.6678132010446
-0.944656579769672 -0.0624543370852131
-1.07655252991842 -0.441720469500251
-1.59444925894351 -3.11427511740346
1.71387417043291 7.95638094156658
-1.26811232935142 -1.19141718266164
0.729303047022139 1.93695961928968
-0.691313419934584 0.474337501633529
-0.333711135403621 0.846686242271427
0.743107835090247 1.97670809350009
1.92591919534008 10.4377933417182
0.473148974962963 1.41366645468216
0.828791786541936 2.21895247598261
0.3345891138474 1.23955509502903
1.39616687331661 5.10666305825946
-0.107275505718615 0.950290375938965
-0.666267250289023 0.505426329319135
1.09289645024831 3.26411664524251
0.765859836476295 2.03488783891601
-0.076768188464758 0.964850464205859
1.03012343761375 2.97897332767509
0.759985684719819 2.02728103831829
0.843133608500746 2.26551528180674
0.0185066691335232 0.997453045371327
0.81233786159402 2.17018711706274
-1.65239328606655 -3.5672486883531
1.13197953472704 3.45242391677667
1.4057918676896 5.17409631571188
-1.54510507422204 -2.76209138479126
-1.39250318746843 -1.81557285837887
0.951075027305254 2.65346207742809
0.552736525770104 1.55705750348945
0.75878836535122 2.01236637732636
-1.16362024732662 -0.756914679674317
-1.11217610609222 -0.564261933287231
0.29255119146397 1.20298932043435
0.479236055210636 1.42928091158792
0.496025141932312 1.44951455884393
1.36609464527703 4.88068154557778
0.648538620336751 1.74753746436698
-0.71533225144047 0.433695281884308
-0.286770268938291 0.871823398856179
-1.95954722302836 -6.44929625957437
0.0739036573349594 1.03753704458365
1.72379301526689 8.04513583149042
0.783420714537717 2.08166388865928
0.0340859451793007 1.0132308773025
-1.74428102155042 -4.32525431791812
-0.281295155032678 0.863976874564564
0.0382462130772492 1.0180676212334
0.0983137141314621 1.05255513583602
1.20290104395388 3.85108434055182
-1.29481912896917 -1.30459771930081
-0.472474670183368 0.74509676010242
0.995401823040681 2.83115558456327
1.75319062696321 8.36567906912779
1.93603985687797 10.5723349974804
-1.26379812183442 -1.18303648649257
-1.18369692778102 -0.834816459090062
-0.152466784570433 0.920038060380442
0.720518130186124 1.9218428064194
0.894330754604709 2.43343420666996
-0.366421366664984 0.819732553669731
-1.42045248163259 -1.98412367895874
-1.91774835420265 -5.9759897017323
-1.20154198012988 -0.904325167068791
-1.33402732468644 -1.49111658514942
-1.50006564953072 -2.45603628183577
-1.66573140633805 -3.65543141562523
-1.5756079616456 -2.95917698659341
-0.392833384024106 0.800043458612075
0.240577748960603 1.16030800232116
-0.331661484133019 0.843324377323934
1.53620483944148 6.22849660991294
1.65714912226404 7.36720614017756
-0.0807157883464811 0.950704651969986
-1.77898436190461 -4.62025313227013
1.67585648664319 7.54823942403045
1.52477266006196 6.1187316608363
1.91797135475423 10.3417542070811
1.04023328929125 3.02492428773039
0.152241395502356 1.07670621900785
0.267046725387243 1.18055550785525
-1.56787123692317 -2.90639158108232
1.79178200361487 8.79783347146289
-1.23832791286711 -1.05843066198729
-0.635983682522904 0.554167833607614
1.60367031537368 6.83510040594383
-1.43481455291798 -2.05228495787823
-1.34739155945832 -1.5820111799188
1.04189068051751 3.02340727837807
0.115200091567095 1.05584975421508
-1.56120307441528 -2.85405514439073
1.48105121191228 5.76149935691055
0.603419123923464 1.64265840860011
1.48273036672582 5.78463909421347
-1.85380337466648 -5.33897504761528
-1.57355092300985 -2.95677414336917
1.47731221536187 5.7269004017349
1.07094623968228 3.1601858383006
1.31681938520243 4.55115143856544
-0.187214975522651 0.926397790816692
-1.76804267009889 -4.52164572165682
-1.77431958323804 -4.56931709534463
-0.871989619000863 0.124054020531704
-1.96752937866711 -6.54008401687717
-1.38883862653797 -1.80297219382232
-0.327448169569295 0.834872236128242
-0.690183808076499 0.480177524505144
1.34261236925349 4.72281373472197
-0.234061054505653 0.909474538769884
-1.85406584301217 -5.34226013439623
-1.11703539684548 -0.580381249126704
-1.67582335398275 -3.71999950522554
0.620588046902504 1.67816478821467
0.0395400978842821 1.01551799149946
-0.381418696904148 0.817207506981972
0.323309013266111 1.23600467523781
1.87142246633171 9.74870087342748
0.914983635785953 2.51724871692523
-1.98702218833896 -6.7279904564761
1.23668201082246 4.04673891861823
-1.61541830725903 -3.26578016768269
-0.536681404637868 0.6646827534073
1.03184549074563 2.97977288547089
-1.43017253565283 -2.04495330253704
1.17441696521193 3.68646265301762
-1.80765575659796 -4.87649258477049
0.474712397966165 1.42047152054773
-1.34420669032224 -1.56172252528927
-1.49750533480565 -2.44561101717458
0.513620072846923 1.48015025720336
-0.181096461189659 0.916883890875462
-1.40583515211248 -1.87647974602402
-1.06265554006933 -0.390387309735581
1.65282905824244 7.29493166451245
1.66613735323944 7.43914211227543
0.382870698033663 1.29193673507371
-0.590190346697981 0.616899363724993
0.674235447856505 1.80008301025492
1.10470625697213 3.32133411715915
0.298221482999279 1.20649053030239
-1.77844578663561 -4.61137131191489
-0.251845902583502 0.893052687740206
-0.526362363048593 0.677052334110852
-1.9285328629071 -6.13142398806669
0.900312461603144 2.45765483573309
-1.16068571417867 -0.725991110077104
0.788977412787463 2.10122441798541
-0.618023494494006 0.590162006628996
0.654998547681636 1.75288201335859
1.18336330269809 3.73071029040782
1.31524089609012 4.53769166491941
0.536173360459293 1.51739394892642
-1.73823963116706 -4.2535981564127
-0.937413290639119 -0.0188526248895549
0.90753282930239 2.47685833734623
0.931842474996466 2.57309961283514
0.83486413335356 2.22652023320826
-1.94351384966993 -6.29627520712077
0.490866577973689 1.4484370249403
0.0556246733436576 1.0154465907918
0.430478704730366 1.36158879573413
0.328921568783732 1.23664315263756
-0.0505001537468308 0.964305599299581
-0.701819536199354 0.449999116700887
-0.788744350627737 0.300664730132906
-0.883644406730378 0.0877547887264485
0.777204138795909 2.08098907305371
0.404023384471348 1.33036729637145
1.76799442449208 8.52326451192898
-0.5371473705932 0.668579714078399
1.43450835051504 5.39236754399515
1.59992625107588 6.80108342442818
-0.0145317781918641 0.973809921922492
-0.858960912599482 0.161889779176557
1.64230744253453 7.22339541450917
1.66998788212528 7.4942110818753
1.68929584078918 7.69716483625987
-0.543860238380735 0.666641322435849
0.256868355219774 1.16414014155318
0.376161311061836 1.27664840274201
0.921408393990831 2.53631401686555
0.325325670266654 1.21472848356356
-0.973116046489098 -0.126730500007163
-0.132667533932598 0.931236035795649
-0.810451709626184 0.267886963544335
-1.06430194672865 -0.405383782335702
-1.23502125439532 -1.0397945704866
0.616029416380815 1.67516031035656
1.78745526782503 8.76227787418742
1.27931288916847 4.29798436573562
-1.96119267589845 -6.46255441288198
0.918509441967202 2.5320154212315
-0.306940973882254 0.858326142702147
0.248852555225668 1.16555921034989
1.35316821236359 4.79508743215688
0.371451472248911 1.28008006402277
-0.0794440959535696 0.956005125018601
-0.341680082781666 0.844119094518881
-0.955971468985779 -0.0852165990259612
0.230866145284785 1.14945751241327
-0.0534421397766298 0.969914319455691
-1.43957202939354 -2.09791697219914
0.882489151516282 2.40664510197794
-0.285313963934104 0.880573605380485
0.505935499872656 1.46343379814224
0.128261451177393 1.07816194692088
-0.586429144990519 0.629530107770523
1.23866600370325 4.06559185397949
0.608518453790227 1.65703005032368
-1.35766524022517 -1.63538355867661
-0.37282628974832 0.799384448782269
-0.419963720245379 0.791005201947998
1.8799898503159 9.87557110974731
1.61158609663821 6.91480338354218
0.664953488897894 1.78236988607563
0.0381909020082576 1.0174443363817
-0.761353550774154 0.350851006317331
-0.938026164897382 -0.024747382537045
-1.92236700880279 -6.08004135770381
-1.07957429545738 -0.462580093288629
-1.34622173676857 -1.57747081911211
-1.85871116212759 -5.3772350007477
-1.22803534582444 -1.02576141340772
0.953203011928426 2.66258882259319
-0.0857827511584714 0.960057081465361
0.650724796933962 1.75055291964157
-0.689073718481559 0.479183477617647
-0.584716312051022 0.620072841706773
0.894474614097228 2.44429276115889
0.79865611645924 2.13584353133673
1.77078759706236 8.5596311596453
-1.03304944262268 -0.296886437274474
-1.08241491519637 -0.46778174459126
0.133838570248027 1.06426121203456
0.704280169796623 1.86488251242162
-1.45523923732272 -2.1786120769966
1.86704022144864 9.70734519076174
0.576394656998143 1.59523376434528
0.33815543603102 1.24583696334757
-1.70620039544918 -3.97625834036472
1.75644320812458 8.40681772519677
-1.12422479455467 -0.607728519914606
1.88394526347206 9.89404180276951
1.34119140011109 4.71394366242212
-0.866659018725533 0.13342644160029
1.42751332053673 5.33351738479513
-0.911531703377452 0.0316959680587874
0.48508057268549 1.44718587288087
-1.64027078153502 -3.46066925693651
-1.25139406637923 -1.11242585131484
-0.248185975672745 0.897321916366227
-1.28391033792466 -1.28402663811459
1.02742077214036 2.95866689031474
-1.00469821661397 -0.200025010181867
-1.4839169486916 -2.36476061584865
-0.388063615699482 0.791749832117417
-0.351330238210304 0.829615571305214
-0.919162594924984 0.0179998135616073
-1.8300296068598 -5.09523489832429
0.261126692492549 1.16381348789023
-0.434250311374255 0.760268745883126
1.71127701069478 7.91589387950685
-1.74396833588139 -4.29086220370818
1.99084425319697 11.3249840251345
1.96740657851932 11.0128349752937
0.00519723053157728 1.00314128653432
-1.0300674545336 -0.292760531150931
-0.679871421735784 0.497273222659833
0.346398392238981 1.26222464776238
1.86986837689732 9.72720250684752
0.681415810398582 1.82471633810466
-1.77451887462146 -4.5905765817216
0.407108934350533 1.33052102469306
0.234319557119511 1.1441133039476
0.680255483899439 1.82443646987879
-1.46806636854791 -2.24633672359298
1.37432472355754 4.94560301954173
-0.0444665624422655 0.967684314859427
0.946510929382118 2.62683889570353
-0.348613222191393 0.840643230744185
-1.16085167224828 -0.752985307128517
1.52131471522301 6.10602184516401
-0.651129722285528 0.522096171976065
-0.974782964327846 -0.130492050354493
1.38933505926152 5.06701358419365
-0.405141209582377 0.788636045619015
-1.07427216939099 -0.439348576157792
0.287925799214639 1.20249444515467
1.99219825484057 11.3300781732418
-1.97984787642545 -6.67440166993524
-1.16710498101462 -0.766106919849088
0.0880439799314834 1.04798139144606
1.65176684501771 7.31505205040777
-0.417609509438676 0.781733792741341
-0.383443127401129 0.804453481067922
1.01000388646387 2.88434280725849
-1.32293438896487 -1.47265021133377
-0.99370936187379 -0.173512541750945
-0.404147661596262 0.789443384056748
-0.105702670056522 0.956508392161341
0.219292884334177 1.12972481089977
0.693092568668425 1.84314500579793
-1.49873438446745 -2.42808076757983
-1.71917638967302 -4.11445980893961
1.92275290133258 10.4060905577529
0.557888885180773 1.55207039096411
0.52577897427859 1.50941467333881
1.19379941638111 3.80067611582445
-1.77051760303498 -4.55469830292259
-1.45032825014804 -2.14115707845768
-1.23795852109155 -1.04877741172978
-1.27410510105383 -1.22052728026993
-0.468704783116342 0.748660167008506
-0.886149958616648 0.108477253199736
-0.545832084491524 0.677538442559796
-0.313696873254803 0.868101603557104
1.87479353239457 9.8019035285421
-1.96168063254752 -6.46823505732823
1.96440026679769 10.9681385483752
-0.662848989822634 0.526116383312213
1.20039942205076 3.83742650014747
-1.54721858629723 -2.79417593391984
-1.63364823827773 -3.39429408125394
1.57659720521503 6.58152772647782
1.21200313281051 3.90426995352027
-1.83574674575522 -5.16861220219619
0.566320040889504 1.57474620086365
0.972707935299989 2.74195158574218
-0.134613050994941 0.938027907445401
-1.49592388715788 -2.40495280111003
0.51808299716567 1.48451456005604
0.327121432995443 1.24307092490041
-0.0117717724824216 0.98816938612114
0.746398682964023 1.9860250421071
-0.160139522955972 0.926953715518133
-0.426027531146679 0.77686273662751
-1.69477632375206 -3.89258825943529
-1.05936959994091 -0.374798803510702
-1.1240287678796 -0.591747658886704
-0.512397011144832 0.697801414844525
0.1151272306282 1.06359410266455
-0.757649192903382 0.365441869230087
1.62827887860161 7.072337004527
-0.956679625097431 -0.0770897233418417
1.48797995432694 5.81793273760318
-1.21173687511944 -0.942758423942144
0.760812354488629 2.0248690601893
1.98588523589092 11.2553963644691
-1.23865491219261 -1.06419656255286
-1.16524180671887 -0.754310282516707
-1.74637746992985 -4.33715527923672
-1.77321833494424 -4.56462033849922
-1.33398224041622 -1.52005657626804
-0.326070805576438 0.851593511677802
0.256559935709201 1.16176848434099
-0.542645927277806 0.672885495947295
-0.971484127948466 -0.110229010415843
-0.230176814593698 0.901620005359611
-0.646934580672173 0.533070758814424
0.611919937847676 1.66399803907441
1.39915511741084 5.1443280764869
-0.313547573274856 0.864464716742631
-0.524558158878761 0.689965272099956
-1.76473060004123 -4.48645796483076
1.58150824165447 6.61680769943754
1.88474461581527 9.94195904014881
-0.403718263151636 0.805022294489218
0.549520352891367 1.52742518740342
-0.535840524475981 0.676708758476658
0.292676339923188 1.2041758733986
-0.066385069686731 0.949511916891023
1.36977344495431 4.908901766673
0.80996594457818 2.15548891267912
1.99370354145932 11.3561875597128
1.18161107639942 3.71275108063636
0.754047140010336 2.00558932006914
-1.88210308314506 -5.62783733524903
0.732503302683575 1.94171425844159
-0.105199099402469 0.946278688183064
0.208559701528764 1.12431237278009
-1.02994191709767 -0.28548779700985
-0.478799425230845 0.729109789391446
0.364156171385083 1.27495206463488
0.237431538421525 1.14882985176487
1.56652522547804 6.4928423735079
1.11753751981401 3.38651113266133
0.559766732161253 1.56331223476094
-1.05253303063569 -0.3624041255737
-1.13665085908765 -0.623529463788013
1.33676157043896 4.68030204116338
-1.25855088245022 -1.15405200948732
-0.728781804869114 0.418785600098235
0.332252831390097 1.24183419559873
1.75266336276354 8.34030576346291
-0.40671405183982 0.79036816665767
0.662488101353502 1.77281326360528
-0.527760172621683 0.689230694951209
-1.98480003264102 -6.71804058404472
-1.22220688253495 -0.988741786050339
-1.32471942362105 -1.46071557583938
-1.17864005024045 -0.811013937624279
1.52887589294808 6.17035021378948
0.414821241537061 1.32832430744308
1.35356563000651 4.80017499293894
-0.449204147693734 0.750592488597213
-0.00290057132472299 0.984750224695719
0.324368361284517 1.22909634265887
-0.521848177920158 0.692713691182687
1.1533626147704 3.57618636017796
-1.10677675050792 -0.56455294394963
-0.318060382757006 0.846300804142516
1.87730983895634 9.82142181400893
-1.96746098497336 -6.54993728917854
-1.93965862086028 -6.25476790416755
-0.459787076254924 0.752114786759366
-1.31348172702031 -1.43242130060961
-1.14514328231759 -0.673155888416164
0.682021135275135 1.81691810531443
0.281605217003446 1.18687597092314
-0.918989576387054 0.0173575925587022
1.86009333852669 9.5959629724972
0.422851653026241 1.35683263155697
-1.18753310918885 -0.857527370097498
-1.05991208496293 -0.386661829729422
0.0788137458913938 1.02237543048186
1.81703979429071 9.09427555649752
-1.77523326667127 -4.60471049943597
-0.134965858344948 0.939760171760071
-0.774427290685961 0.330664922257729
1.84912287266476 9.48631256447653
-0.0473840523157882 0.975596591749443
-0.310445770973367 0.855826804192119
-1.85040726666538 -5.28123090480442
0.0195484778286565 1.00949379122544
-0.32148485732909 0.85171566029906
1.00341624502118 2.85243197422087
-0.0413769862398357 0.968545930426121
-1.50428660641862 -2.4734882991516
-0.315956120587869 0.850716428317477
0.379662199511765 1.28561243609028
1.26056841548983 4.17614570857556
0.117109705720536 1.0719757158902
-0.72026478263644 0.426526965944616
1.76009071973641 8.44462688866445
};
\addlegendentry{new\_neural\_network\_estimate(x)}
\end{axis}
\end{tikzpicture}}
        \caption{Schätzung von $m_1$ durch unseren Neuronale-Netze-Regressionsschätzer.}   
        \label{fig:subfig9}
    \end{subfigure}
       \hspace{0.4cm}
    \begin{subfigure}[b]{1\textwidth}
    \centering
    \scalebox{0.9}{
	\begin{tikzpicture}
\begin{axis}[%
scatter/classes={%
a={mark=o,draw=black}},
legend style={fill opacity=0.8, draw opacity=1, text opacity=1, at={(0.03,0.97)}, anchor=north west, draw=white!80!black},
tick align=outside,
tick pos=left,
title={\quad},
x grid style={white!69.0196078431373!black},
xlabel={ },
xmin=-2.20852591900454, xmax=2.21306049239568,
xtick style={color=black},
y grid style={white!69.0196078431373!black},
ylabel={ },
ymin=-7.86405361849716, ymax=12.4391808546764,
ytick style={color=black}
]
\addplot[blue, mark=*,only marks,mark size=1]%
table{
x y 
1.98342671961567 11.2067083316341
-0.164883410431491 0.921819789666956
1.40142299960581 5.15033590554207
1.47929179834368 5.7561559124411
1.19715514411127 3.81799531027161
-1.05745151470862 -0.371309570367478
-0.696648319326161 0.464684877420012
1.82247761300943 9.15712037755566
1.33480759235371 4.67626070341997
0.862438587068488 2.32882713322422
-1.16503690445439 -0.754686954628516
0.417168238066338 1.33933090794867
1.15596892097842 3.59120483869643
0.850205757253619 2.28838577184768
-1.30994868597998 -1.39187573215822
-1.55346692657437 -2.82488317686155
1.61669575447566 6.96902026146192
0.500046589667835 1.45907874969249
1.28496886028727 4.34712794233107
-1.20227096727067 -0.904559178268851
0.366681884355046 1.27741420051878
0.818260589820523 2.18688479377424
-1.23845976488916 -1.05920098696907
0.697189507824902 1.8530215419852
1.93926361424016 10.6132554825254
1.04107639834298 3.02636787410901
-0.22353824693392 0.893074337960662
0.668246244317835 1.78431442012492
0.151928097424038 1.08704702436641
-1.07641959705448 -0.433765220730816
1.44298781246897 5.46662361852968
-0.284110387060426 0.860783387392703
-0.44017790192429 0.755901668314636
0.352464522063814 1.26134584362801
1.18574235128237 3.75381711403657
1.66572862979816 7.4485981217319
-1.77196434008901 -4.56390224034042
-1.42077197081339 -1.98366518029347
1.90156609340709 10.1254900733363
-0.399226079743066 0.787289288568677
-0.502642406493864 0.701288589370527
0.570484726981828 1.58079483029362
-0.626439688866295 0.563662303293445
0.646913187691053 1.73622883021686
-1.70011590874738 -3.94019890361793
0.188560902818183 1.11268318790154
0.613581309391728 1.66528403436379
1.53866334933713 6.25708686606659
-1.90745710621872 -5.88963818524853
0.41517370699518 1.33673477925145
0.327838052279217 1.23484812561957
1.73772873954887 8.19951335349887
-1.78984900196206 -4.72748294080677
1.52854808178337 6.16929527049623
0.00832534691972997 1.00418578874098
1.53059984449369 6.18702069900853
-1.45719188829622 -2.19959638868634
0.291257080120115 1.19843927737529
1.19523677300415 3.8071354073554
1.13842913470803 3.49861808490207
1.2366636859145 4.04822570623427
1.89871430519567 10.0892868982552
0.019953194670169 1.01011409945552
1.82122210851877 9.14233718077853
-0.298230517123074 0.852732115538147
0.382945539678767 1.29651503281945
1.08047569362621 3.20916222764976
-0.100962219519459 0.951781416117594
-0.409006788513409 0.780081708579802
1.41020530945337 5.21583204082174
-0.862424320245129 0.156477928786038
-1.98328806650409 -6.72214450231026
-0.696177128652387 0.465406125702621
-0.0933688730806197 0.955318102465428
1.94870249960763 10.73808188061
-0.936123458958722 -0.01976601438959
-1.58312442994531 -3.03409351460663
1.4664267507308 5.65211917053288
1.57141891136512 6.54839610341355
-1.47051135974113 -2.28133185267719
0.874729152014026 2.37035416122943
1.03751084984896 3.01036865668898
1.48030997994332 5.76445750267437
-1.54541524782003 -2.76946312915615
0.552369097636586 1.54760978330716
-1.68042557562017 -3.77839948923686
-1.51061355679375 -2.53657210298722
-0.167528933067055 0.920559140081827
-0.787421015424107 0.310011987252715
1.6501030154187 7.29301335433562
-0.0814367755752747 0.960885772415491
0.0795951024473589 1.04237142562678
0.979738939834169 2.76334381388308
-0.380817770934829 0.800394418604914
0.400591024706831 1.31813698650762
-0.4141113404742 0.776250097560955
1.91403104307527 10.2848760382045
0.920870673865245 2.53445330225573
-1.30315603463426 -1.35866671389968
0.0275601242433647 1.01404829046996
-0.927503900072607 0.00222689952751143
0.0368974953304853 1.01894249538715
1.15977317392 3.61159817433105
0.514097432066963 1.4818167241679
1.17984224265819 3.72104216577677
1.58572120040197 6.6789936011996
-1.37713701802779 -1.73918304115634
-0.472661706237887 0.728589446354698
-1.95506901984573 -6.40446611244112
-1.1779484446276 -0.805616762687941
1.23713048801279 4.05102163809569
0.455615485804151 1.39192818075103
1.59249283776102 6.74155542309861
-0.503448403039978 0.700525301295737
-0.446140509435029 0.75105598593073
0.0961164638478107 1.05196738610277
1.85889095909979 9.59383788518873
0.961953261926548 2.69181699566452
-0.383322867187997 0.79864530138838
0.24213138507616 1.15461989372253
0.788212033700322 2.09669287956382
0.84070370731225 2.25757827157573
1.40004883629895 5.14015179344157
0.107393086597664 1.05870960384272
0.225480591396121 1.14097309250214
-1.67936142766505 -3.7697647231671
-0.0816940927046601 0.960765514872752
1.47773128825186 5.74345186910455
0.401713010033748 1.31954418061981
0.422412491054748 1.34621800253768
0.776700220186323 2.06346076723564
0.317989295994632 1.22470791713495
0.816536031039081 2.18157160974445
1.53962127691442 6.26545353631168
-1.29064269548569 -1.29838208894096
-1.47198666263898 -2.29047742640869
-1.18769903957274 -0.844807189958559
1.5617847179266 6.46159449041297
-0.137017676939132 0.934967618307204
0.10162263173636 1.05523925887493
1.42737547993898 5.34593888574134
-0.755076052164864 0.368829213132414
1.71133977541391 7.91773667028815
-0.4721794317973 0.729011711839228
1.89344301858784 10.0226243465763
0.867958702524959 2.34736672383903
-0.411724457536147 0.778047811177945
0.824378736637617 2.20587035898719
0.170515291427536 1.09977047759294
0.979189264724758 2.76110191698977
-0.897121719560843 0.0767896021307181
1.55034434863224 6.35973527019784
-1.67062801778883 -3.69932121498425
-1.24614106788035 -1.09319208537185
1.54412256063696 6.30489082553266
1.43769151607769 5.42542699235533
-1.6456012323085 -3.50158560570617
-0.635373759957425 0.551982650209956
0.0798070563154254 1.04249250741436
-2.21592582434482e-05 0.999988920530453
-1.56549856175875 -2.90878956620822
-1.85995182081513 -5.40180355705662
-0.744024022342477 0.387963590115266
1.95878262880175 10.8725914329787
1.52876440755115 6.17116216956417
0.252423090403078 1.16334903121843
-0.393738335259513 0.791257880640355
0.449465165244776 1.38318113689044
-1.68386371444268 -3.80637391755432
1.00910500968525 2.88606233001523
-1.12050763647973 -0.587298253264637
0.230997555210016 1.14543091365568
0.0349769356316969 1.01792975394193
-1.56076894216133 -2.87564939931168
-1.45121512537244 -2.16340493618994
-1.16545297771441 -0.756311132410629
-0.0191205019479455 0.990551431326679
-1.86008081192998 -5.4030939017934
-0.436428121653843 0.758911768315441
1.49957184981221 5.9233919160222
-1.44219942190571 -2.10937609184844
1.01529613831035 2.91268092874956
-1.50140183330494 -2.47671254155102
-0.732092139869074 0.408086177914741
-1.31471126281736 -1.41536454673876
-1.16144104543941 -0.740697216143942
-0.689196716081129 0.475997455288745
-0.217022486630988 0.896366930296396
1.62508170186513 7.04925029897457
-0.0406198149252917 0.980157921905294
1.63119848795736 7.1082339788836
0.261617455458398 1.1713441338436
-0.0226785485704886 0.988815972333353
1.93031313461916 10.4958902043135
0.10579027870623 1.05774135817833
-1.34359720927249 -1.56148709026429
1.08272594869317 3.21993898042956
-1.13811010512611 -0.651956813151084
0.157670660808171 1.09091781606985
0.693694852953544 1.84450954897994
-0.413991721747403 0.776340444640307
-0.3496582994086 0.821300655555392
-1.494906796828 -2.43494997466264
0.0812226478153182 1.04330249998073
-1.87946646412891 -5.59912494096078
0.135232078225477 1.07608502949923
-0.0314802740817721 0.984550095861464
-0.536661666154805 0.667662425434589
0.939570817014385 2.6047405093355
1.9606433016915 10.8975565658301
-1.31276353932712 -1.4057380074589
-1.30676601901061 -1.3762731425726
-0.303871457144716 0.849451172880855
-0.643511730589387 0.541119017871324
-0.527466345137777 0.677041842362992
-0.522841202805422 0.681676958703923
0.4944545179478 1.45023449606521
-1.70114092695177 -3.94872752787689
-1.17864666831437 -0.808402179553694
1.70768647325131 7.87932950534975
1.57303031948494 6.56300623145007
0.553581378198257 1.54978869652802
-1.49462820991928 -2.43316685250796
0.0532637575658805 1.02770819278823
-1.03292019870819 -0.293649505337755
-0.0539046615190384 0.973832153044759
-1.4019635634533 -1.87640980895707
0.867566356157913 2.34604301414628
-1.05398034586825 -0.360109127135227
-1.13058478157364 -0.624075386023385
0.588011197017629 1.61419683576319
0.883913094472204 2.40197685981143
1.7032464005838 7.83284702573416
0.723806315564655 1.91984629454512
1.35906175870283 4.84427052411395
-1.23544243096618 -1.04596172691295
-1.54329111912356 -2.75493982502582
1.16098779424438 3.61813290630781
-1.87095954053839 -5.51258428050454
-1.2782629197376 -1.23986916511837
-1.05654172495337 -0.368367143763051
1.24353360832467 4.08955390230973
1.19028282951638 3.7792262679224
-0.775725735545598 0.33177186288789
1.02057273878203 2.93557518071213
0.821921250495339 2.19821881120506
0.918229962933942 2.52470603700035
-0.984864184033527 -0.151361080185913
1.93401749117918 10.5443463915569
1.83315978839409 9.28363700124014
0.533232250852013 1.51399333097939
-0.722014352421549 0.424656658899597
0.481969685265481 1.43090303239757
1.17847167126153 3.71346785213066
-1.79048490297368 -4.73336148187232
-0.945697114342179 -0.0446368331238406
1.60064007186754 6.81744991856981
1.93687669965452 10.5818612825148
1.35197568049207 4.79464642514883
1.91147516287336 10.2520425549434
-0.00460781786329045 0.997702891593491
-1.39948950363334 -1.86251399232994
1.04873511695863 3.06103594471185
0.000913924775346242 1.00045723462595
-0.50513266090352 0.69892519214253
-1.26704247379017 -1.18779508918314
1.76648926011692 8.515401523615
0.157856915922508 1.09104424813471
1.40787527714754 5.19838618820376
-0.240710854768439 0.88424614078834
-1.50153204953699 -2.47755356296565
-1.98239180697612 -6.71190905798603
-0.274480732689117 0.866148734835799
-1.63517535024967 -3.42099958712273
0.960111664862905 2.68452965066336
-0.166975897002565 0.920822887277161
1.54394939093015 6.30336988746059
-1.48090289946801 -2.34614449639195
1.23336476648575 4.02851752917348
-0.400362192403418 0.786460973883875
-1.87413140169175 -5.54475667677901
0.336611711534724 1.24409863604828
1.18872743268058 3.77050370068962
-1.7294971642558 -4.18886360380201
1.45013892943371 5.52266642043676
0.0983398967377997 1.05328401410495
1.82859103308664 9.22936432902555
-0.251038744131117 0.878821858955889
0.0752656353562249 1.03990925556896
0.0812009707648484 1.04329007921359
1.94778947733205 10.7259599229546
1.44701878178583 5.49815477228648
-0.327412551641523 0.835331216623138
0.243300205866964 1.15559974805297
0.092634846159414 1.04991799090477
1.99004222629651 11.2976793178811
0.370987182526909 1.28239482533905
1.12746366204799 3.44192570040845
-0.88103525140069 0.11444499583188
-1.87311208155932 -5.53440536407914
1.53244608154239 6.2030062191922
-0.460133374851041 0.739391819239952
0.307337448869911 1.21402461928519
0.369331774016094 1.28047335921953
-0.234023743918722 0.887711030862566
1.04808634766716 3.05808318471249
-1.71866655486462 -4.09618344887114
-0.541554732849246 0.662580975842837
-1.48954836729529 -2.40077032291578
-0.0346183802921285 0.983037951949156
1.66262513718082 7.41748978740735
-0.0632924638352459 0.969396910959773
0.0976192149694759 1.05285658014762
1.64277341714372 7.22092429870893
-1.23064117143324 -1.02502582306785
-1.2119829712324 -0.945178959503772
1.82543703976219 9.19203877696444
-1.96283859036236 -6.49099090907211
-1.07753916455694 -0.437518252974556
0.401431131407247 1.3191902799471
1.11313051112005 3.36918797402718
-1.82353738645624 -5.04487061859584
-1.02282270803728 -0.262684033496767
1.34249571990793 4.72895418337375
-1.7423388101225 -4.30030535815532
0.574932133692535 1.58914867255579
0.356056594784311 1.26535121241934
1.85225865167069 9.5131424338517
-0.398757442351753 0.787630287557302
0.689987650483266 1.83554534383615
-0.293455390498176 0.855480226380854
-1.36443451041482 -1.67085494806716
1.35065286985791 4.78543213551069
1.01572198192603 2.91452148801636
-1.34644750916868 -1.57624963302143
0.934098979805452 2.58394372921836
-0.337375961698936 0.82913400048838
0.64371991587656 1.72921343065672
-1.12505418398717 -0.603812272410027
-0.0424605535851552 0.97927752676906
-1.03874715299427 -0.311782480194157
-0.934567390971795 -0.0157678258945499
1.06185524551419 3.1213904800115
-1.68614727733944 -3.82501880552636
-0.0939923135416958 0.955027594191503
0.186353265908155 1.1110726428412
-0.702833029741399 0.455143382722217
0.669467597030042 1.78713226709335
-0.989846282352082 -0.16552193374252
0.105206333004412 1.05738943070462
-0.327016013168481 0.835575044028545
0.687399875760051 1.82932776989665
1.84677034716651 9.4467555282193
1.08694083448046 3.24022372148457
1.91270842735118 10.2678755681039
-1.84876330693187 -5.29058235734343
0.140188037147681 1.07929462015681
0.340988226479485 1.24879088314535
1.63241684639954 7.12002930432158
-0.134435062793881 0.936176722358918
1.89935988024498 10.0974738907713
1.10456045506315 3.32642942521779
-0.669863931935491 0.504430082672366
1.20253788657643 3.848623510478
-0.476896000953684 0.72485957082047
-0.864160922436813 0.152624313681391
0.779743629721808 2.07217627840376
0.44847566998307 1.38178594372062
1.88317123466855 9.8936753259772
0.966481764673039 2.70983107632884
0.977908789987251 2.7558871775324
-1.52375847203518 -2.62327699740979
1.13406912530685 3.4759673999773
0.522353840992912 1.49552708713273
-1.06822671124227 -0.406527917387101
-0.114355250680581 0.945546284394021
1.89308302690468 10.0180838654405
0.845248577878253 2.27224791568968
0.122425695926454 1.06795770242727
0.127977871489443 1.071452329096
0.01625186865687 1.00821615648482
-1.08796984143618 -0.472845104422563
-0.320666677251147 0.839450527601591
0.334320179485394 1.24166256257872
-1.05886861846269 -0.375902385164731
-0.564026324025872 0.638414711051878
-1.59115775979802 -3.0921520343033
-0.624686875303529 0.565923843870914
0.0413789150922659 1.02131826057146
1.31657491972719 4.55336406862912
-0.726992150093075 0.416520105094092
0.0524606842524733 1.02727218818635
-0.492141039700758 0.71109076518107
0.978275827209485 2.75738082493985
1.14130644500639 3.51364516032462
0.0392796360480538 1.02020313016701
1.30721738204172 4.49141064784605
0.490133920204322 1.44348012355915
-0.706442978901193 0.4495096644889
-0.993964710262591 -0.177329563835979
0.943620152780988 2.62025423675144
-0.0932557921766692 0.955370798670078
1.6365425027311 7.16008680515812
0.973324880119798 2.73730832363678
-0.725054451856926 0.419698555277889
-1.96417313354495 -6.50592462133396
1.31324497867966 4.53123083543522
1.87324432504556 9.77024625410757
0.303189087378309 1.20994224416354
1.9389374422235 10.6089613728723
0.501357612206841 1.46116901188386
-1.720056540905 -4.1080112039577
1.31781128519504 4.56160630738453
0.0951516110926338 1.05139794858651
0.408192815810979 1.32774813968817
0.513811453765586 1.48134650425401
-1.23574332705072 -1.04727912774199
-1.45264147142963 -2.17201479615298
-0.227918340918484 0.890843932569515
0.768794189473504 2.04105392984846
-0.383261040749235 0.79868859764658
-1.52535647365474 -2.63392131109938
0.978901900798691 2.75993067808586
0.989417826238432 2.80315006002761
-1.36387886985462 -1.66789491236366
1.30076519383918 4.44913205476502
-0.61920844932297 0.572929632714456
1.33769433552972 4.69598515376487
-0.74829966725375 0.380618179060386
1.65991412952015 7.3903995937643
-1.54882463937474 -2.79285887150011
-1.89409203935268 -5.74981453952136
0.116957207902192 1.06455795146663
1.17623321690576 3.70112896015788
-1.35372154548933 -1.61420715774001
-1.88141472578473 -5.61905888338989
1.32924807506156 4.63847962469226
-1.77134487512033 -4.55829675676622
-0.128741369799476 0.938838418393211
0.51095333379939 1.4766641046124
-1.67747670266001 -3.75449901365524
1.29669044236686 4.42261577546358
1.34952912501012 4.77761666284377
1.33266062176777 4.66163835456556
1.12465453846378 3.42754849841394
0.4823513930431 1.43148567258717
1.56244015327284 6.46747005636869
0.11081670699023 1.06078907799291
1.47420921577403 5.71486476538048
-1.40118206014657 -1.87201510170116
0.634664924286693 1.70957467919594
-0.684074619393981 0.483658496059001
1.60007325210646 6.81214767306192
0.871416435428499 2.35907231047904
-0.983390175523971 -0.147197135538043
-0.681323886845001 0.487734387895415
-1.26150283611157 -1.16241909746674
1.67379753998392 7.52995968901337
-0.987930545268803 -0.160060852020812
0.454921140058064 1.39093417506421
1.51394338581359 6.04432436735546
-0.804954691722807 0.276317438492785
-0.830724433208845 0.224398470191638
0.335858029439155 1.24329584466934
-0.159847298344431 0.924212755588489
0.0812119985668653 1.04329639798377
-0.117494185355901 0.944084196411911
0.497555841122187 1.45512515695177
0.640924836877307 1.72311134065059
-0.123471924446762 0.941297774124052
1.15598732778889 3.59130324256049
-1.27272624906405 -1.21406008321854
-0.76949998857843 0.343125991938706
1.4162089827075 5.26101451251004
-1.91041766129389 -5.92088694414514
1.99504962001563 11.3668989718006
-1.00490416987749 -0.209143338340367
-1.14258809682261 -0.668718416842979
1.18209823855401 3.73354181514181
0.128884868780194 1.07202738126087
-1.94570137893154 -6.30108914105676
0.27471991687839 1.18306806252273
-1.72671969673579 -4.16498229116897
0.187211122317864 1.11169743641647
0.48198115722144 1.43092053557093
-1.98132421803257 -6.69972954213555
0.758502825129475 2.01238871783111
1.34931125357486 4.77610269536153
0.245180659219023 1.15718232866859
-0.276183458320058 0.865207129827379
1.39456510983045 5.09968311504284
-1.6804327463332 -3.77845771215393
-0.665754239595305 0.510306575023814
0.76684140612743 2.03557122891776
1.09948744175177 3.30137507524853
-1.7534846008282 -4.39840634527731
-1.83044998652393 -5.11151071125484
1.37370201959309 4.94821688942207
-0.382340179934817 0.799332698920512
1.84016263074689 9.36729469457565
0.288068510119739 1.19542452478037
-0.114543656857639 0.945458542257138
0.206795309294432 1.1263285058384
-1.57769003775003 -2.99515593056534
-1.03763955741927 -0.308320801104982
-1.52268760032851 -2.61615653065817
1.11941760456987 3.4009040701208
-1.33311090638393 -1.50771004489711
0.128759478350397 1.07194781107472
-0.202242477740643 0.903731634603166
0.878360806223666 2.3827980126127
1.19855610565421 3.82594465441631
-0.95885680090605 -0.0795950750614648
1.9731762658412 11.0668248504742
-1.30507394687268 -1.36800862623331
-0.0157669575178914 0.992193314134026
1.88887607888191 9.96513724318019
-1.01174714032224 -0.229379091520865
-1.56039182316205 -2.87301570032441
-0.926950247777248 0.00362671646421875
1.64418618486032 7.23477519532048
-1.20043385996112 -0.896947485134618
-1.59792808416967 -3.14154749994605
1.69730156795686 7.77094735900795
0.575625029518244 1.59045760970327
1.4325210476095 5.38546272461785
-0.821655978972185 0.242999595633413
-0.810842030128184 0.264710330683507
-1.76436741471195 -4.49543613243629
-0.101373628416222 0.951589874479197
0.645687235329547 1.73352993640201
0.10576329327991 1.05772508508171
0.497152732945724 1.45448746150772
-1.20605108820682 -0.920293170938932
1.65818341412337 7.37314604751803
-1.87962426988887 -5.60073796632451
1.33566435853625 4.68210715495086
0.282054103858054 1.18980470698956
-1.26250466931175 -1.16699204691647
-1.67476380995526 -3.73258686461563
1.01967995195807 2.93168806191011
-1.67840964447672 -3.76205118380057
-1.50130451410046 -2.47608408679851
0.141109062399635 1.07989515924757
1.20995178365977 3.89118827835135
-0.531511081284835 0.6729432404821
-1.54006889061477 -2.73298567027789
1.66007705123738 7.39202540581725
1.38566787929341 5.03460666896519
1.42172807729341 5.30284416594051
1.19501735684592 3.80589515622314
-0.0613276058947694 0.970323119939249
0.588196574174466 1.61455706145536
-1.86528429810381 -5.45530051391384
0.327998076046251 1.23501500011517
1.94392556510098 10.6747725129444
0.391687289083808 1.30710792726173
0.953213581470196 2.6574299461252
-0.727575792369922 0.415559950077793
1.21439774087505 3.91692509422767
1.65999716563432 7.3912281841227
0.382087929756178 1.29548822962649
-0.16287160981814 0.92277676837371
-0.0542029915768536 0.973690800842405
-1.63260596193285 -3.40129984579709
1.12690113690572 3.43904191500925
-0.602201080415101 0.594083125950106
1.26792747852857 4.23945539583106
1.30847145779741 4.49966965275129
-0.718738337998323 0.429960531379342
-1.86119536711375 -5.4142508988663
-0.356727609725149 0.81669094026702
1.90459290272936 10.1640216503881
-1.44227588332194 -2.10983146250815
-0.391557970599605 0.792819909722158
0.631524958158446 1.70285186854789
-1.8673507757779 -5.47611714719273
1.42661454192866 5.34011495936151
0.170205151974584 1.09955348313825
-1.83568026427185 -5.16227865205378
1.2073940863486 3.87645422732492
0.706156293104354 1.87513804754791
-1.98203177551578 -6.70780014014512
1.36502793573962 4.88639889620432
0.399224089177922 1.31642786191096
-1.0253447363566 -0.270364174212614
1.50929202137678 6.00496342833847
-1.63337973901535 -3.40722582261244
0.16217533939344 1.09399145885371
-0.98542617423463 -0.152951737923894
1.5360467607192 6.23427976294006
-0.483081925171916 0.719337382260632
-1.58770633174206 -3.06713467478903
0.754998431678347 2.00275622643676
-0.315694310614774 0.842448548232049
0.80731770988626 2.15345551594596
-0.520999610898754 0.683507318681412
1.49195497855709 5.86011399750875
-1.10514956412394 -0.532465223812011
1.97857357935404 11.1403174392499
0.201904204556214 1.12260731318804
-1.48063273802369 -2.34444782161577
-1.74765842069912 -4.3469663733447
0.97582411111886 2.74742059343922
-1.94025775141107 -6.24148801641387
0.378690645694897 1.29144231745017
-0.675642050925532 0.496069283877749
-0.85135959252331 0.180704325857191
1.29651347872175 4.42146741003185
1.38442162021466 5.0255486562329
0.0993659004083431 1.05389365645253
0.322056803667919 1.22886488173402
-0.319240158445763 0.840313911505586
-0.584549292170264 0.615110609184494
-1.12598759652605 -0.60721864335168
1.41262589486618 5.2340088706993
0.646121632266878 1.73448545340135
1.78052185613418 8.67289265981815
0.0456857711294889 1.02361857502653
0.650561123062864 1.74430075452423
0.444925785300273 1.37680795769407
0.360910773533672 1.27082206027476
-0.00804773722748919 0.995996648292271
0.377237311240174 1.28972199915939
1.52091849859923 6.10374749026291
-0.0373496792893628 0.981725351704962
0.862206665085895 2.32805217881156
-0.507190355652468 0.696960910121683
-1.86205179057238 -5.42283330831689
-1.70459448861946 -3.97754027439972
-0.151804594967045 0.928017392461753
-1.68436080041345 -3.81042813744005
0.187659122172576 1.1120242490766
1.3554496588658 4.8189190431175
-0.487869327729311 0.715003204906045
0.842675192251204 2.26392694634303
-0.912200762028982 0.0403535280591013
-0.946369416514677 -0.0464010410197027
0.330462441073498 1.23759349000862
-0.0195847733926913 0.990324603273036
1.26237990986932 4.20493049782254
-0.675334511443002 0.496517206261075
-0.229399784212416 0.890086353839939
1.9604831935885 10.8954066810646
0.653098747548911 1.74995226855662
-1.99271462098377 -6.83038081816128
1.52604705488264 6.14774493496474
0.265785713873526 1.17503120638017
-0.162967637142707 0.922731122134607
-0.283790790788197 0.860963045149854
-1.36121907666447 -1.65375883389091
-1.31329929847602 -1.40838314891627
1.81527813187679 9.07259527626234
-0.442821583254181 0.753762240778201
-1.35568360263731 -1.62451551362428
-1.51220574381476 -2.54699354010741
-1.88865198249977 -5.69348252498278
-0.0584710604719203 0.971671561732019
0.381577876562059 1.29487859540097
-0.606990213147065 0.58821653307224
-0.749626560934504 0.378324034734459
0.00894689778701219 1.00450019527111
-0.134416985656438 0.936185181356841
0.947482538246278 2.63514997624858
0.652193660151404 1.74793312291132
-1.4971053276791 -2.44904543319222
1.73588142591092 8.17953815015466
1.82488208996298 9.18548319269702
-0.433621813289647 0.76114591416293
-1.85764198404764 -5.37872869313915
0.189957707840087 1.11370674618603
-0.0839730204312175 0.959700875691003
1.24728831958757 4.1123056418777
1.65468022478909 7.33832011275106
-0.586495676362482 0.612837325483901
1.79589799743536 8.84802129209566
0.428791423729258 1.35471550491822
-1.21645298698983 -0.964090333743602
1.5431165224241 6.29605904121035
0.178665194044994 1.10553169571189
-1.65418653612307 -3.5687317524992
1.59677966307949 6.78140367909103
0.933307944837329 2.58095305728585
-1.77917297787843 -4.62942925053352
-1.02449591391644 -0.267775315313738
0.586894557125167 1.61203007076959
-0.522527848756777 0.681989010527291
-0.839039065755164 0.207022944691082
-0.985881123675554 -0.154240675866536
-0.464749211709719 0.735452035855211
-1.64426334947883 -3.491186075905
0.467548491374151 1.40927248703334
1.92460116275829 10.4214987297265
-1.6443087962847 -3.49153905651549
-0.491293941305216 0.711870026008429
1.91222188641349 10.2616270266605
-0.62453546156841 0.566118745447308
0.300401507594917 1.20722329431868
-0.408259673873164 0.780638446049794
-0.386727514305038 0.796251031923926
-1.2807629327012 -1.25159562141562
-0.35007346853807 0.821032011779048
1.81278321825848 9.04344311739489
1.04336477412012 3.03668320841017
-1.32183939999296 -1.45083724437438
0.292077450222015 1.19921892546656
0.994071082116497 2.82251080980646
0.177466744750081 1.10467731384257
-0.505717794748531 0.698367674676974
-1.43761031855531 -2.08213415727262
-1.89149919300563 -5.7229237053933
1.28583777853235 4.35268374821217
-1.68115385495103 -3.7843153794068
-1.4142254050391 -1.94600789683803
1.52526984373002 6.14106061529525
0.612917731722653 1.6639222092179
-0.607365102169935 0.587754345462488
0.0446778208985652 1.02307869555183
-0.96209360813814 -0.0883318304875763
1.30375431185376 4.46867408429554
1.44134726467672 5.4538346503775
1.00844224545156 2.88322832443489
-1.59164327974008 -3.09568014470601
-1.93898777050765 -6.22763304253179
1.75462960692508 8.38402365592018
0.409361680989284 1.32924207079661
1.96570363363382 10.9656676715992
-1.89042953048247 -5.71185225186763
0.634605413505302 1.70944684791535
1.51117328592788 6.02085759885574
-0.336130600952831 0.829916161502131
-1.70081948401835 -3.94605184186047
-1.17900918117107 -0.809849619160207
-1.88055737118247 -5.61028146583971
1.7226892869659 8.03798589758147
-0.28444331034676 0.860596119807494
0.256235525537556 1.1666413705637
1.70534277107545 7.85476685433527
1.12385801800205 3.42348268523638
-1.47077804764093 -2.28298371381439
-1.61806822565318 -3.29102278955293
1.90294759424218 10.1430630966633
1.44451129494468 5.47852266160623
-1.01616478290891 -0.242580965266955
1.15493317240572 3.5856718757953
-1.40498242439259 -1.89343221858113
-0.452277827687501 0.745991134070497
-1.02618668656894 -0.272936092067389
0.16135892421457 1.0934319415492
-1.38881021937692 -1.80309489411771
0.83749527307802 2.24729452697683
-1.50103564033491 -2.4743482179168
-0.486808076009357 0.71596857791788
-0.538027218351843 0.666250669752277
1.9974166483278 11.3997283821791
0.0678536146367503 1.03574211366182
-0.0899177932002528 0.956926814390687
-0.427960738075485 0.765604944431546
1.60332836723081 6.84264236473997
-1.91821035570889 -6.00361962545922
-0.484358263630535 0.718187066115336
-1.14455694548876 -0.676128497006041
0.0304102520839233 1.01553439168116
-1.35890685005729 -1.64151467560062
-0.79930465921805 0.28731719612153
1.73862147320117 8.2091801257019
1.45803599732521 5.58511646755393
1.89547191274214 10.0482431762309
0.690931727750082 1.83782178852227
-0.621921289478656 0.569472296962388
-1.85804910216955 -5.38279143389878
-1.90519163173854 -5.86579373854702
-1.16833564944958 -0.767594943367109
-1.97462737655031 -6.62363958121218
-0.318014178865605 0.84105380251823
-1.6721081152242 -3.71120685604103
0.12506651417845 1.06961442749644
-0.598865775068001 0.598127669210699
1.77408588270415 8.60038455461797
-1.28722274061527 -1.28210603288769
1.97264479348539 11.0596075319345
-1.38334936445139 -1.77306192632327
-1.4234852858919 -1.99937517996845
0.68138823441224 1.81500941674427
1.04276100159047 3.03395800254456
0.363443058450483 1.2737027863533
-0.484681708491861 0.717894959660252
-0.42417904656689 0.768548456579527
-0.463869516349239 0.736206471705505
1.51603821310448 6.0621205930325
0.758851818256217 2.01335155326038
-0.361724413726992 0.813386601010214
1.9077891816332 10.2048298360685
-0.325593495457903 0.836447991161261
-1.4194697586773 -1.97614676213254
-1.98060600755018 -6.69154355462185
-0.74386957637504 0.388227584497721
0.0550760056004314 1.02869441717498
-1.07639607264541 -0.433686442070373
-1.98146995411248 -6.70139135895703
-0.614558293171367 0.578802301776621
0.580519533847571 1.59976104188703
0.85530197722539 2.30512646438327
1.16025561474269 3.61419236913617
-0.856557722320832 0.169393066404631
-1.15325753072591 -0.70917200147902
-0.239167870624584 0.885048828266123
-0.763721926225796 0.353522802407235
-0.977358490724609 -0.13027982803673
0.778036169125887 2.06728038238241
0.884340845547061 2.40346215552389
0.328138275778158 1.2351612579192
-1.55099225540116 -2.80778775223326
-1.03940188301237 -0.313832078552333
-1.72324163505083 -4.13518796186514
-1.84654800214979 -5.26872481567198
-1.43760044903241 -2.08207575751688
1.33241634702983 4.65997723453523
0.00470785624302161 1.00236123791202
-0.756956890576916 0.36552483747915
1.47817407575479 5.74705420485713
1.7488753934927 8.32084679314365
0.24029622852773 1.15308731260398
-0.229712543158756 0.889926204763276
0.936394382750414 2.5926445628988
-0.588247452495922 0.610781821734327
1.12089743121385 3.4084122555828
-1.68973056641353 -3.85438000159527
0.601490598378976 1.64077414366491
0.485060318584069 1.43563571921682
-0.743511201112161 0.388839796458524
-0.237997845978876 0.885656203601219
0.465159036576552 1.40575968640309
0.221364441760389 1.13768751303374
1.84453390617883 9.41980429808891
1.00647655101298 2.87484055951728
1.01000925130016 2.88993373969193
-1.87754777380837 -5.57953526649142
0.768807495091131 2.04109135728871
-0.126510324543307 0.939880107966774
1.59224876657191 6.73929233360148
-1.96901747277081 -6.56031066024765
-0.151585860600471 0.928120599126281
0.27514404486757 1.18345419050871
0.928969564002895 2.56462168594107
-1.13912071298838 -0.65572847467495
-1.3946007413604 -1.83520004087116
-1.53647125544187 -2.70858343004002
-0.853899259225281 0.175193415056966
-0.168066577292175 0.920302623347891
0.80680243856385 2.15189801728073
-0.568715904490896 0.633195480283646
1.90853719763859 10.214397764992
-0.788490064341025 0.30799463717303
0.442582735811799 1.37354566431466
0.459054233603626 1.39687552247798
0.39261749399079 1.30824877454604
-0.334717090759574 0.830801289010195
1.94086803536182 10.6343968181273
1.72074136428306 8.01724701749925
-0.292636510585159 0.855948854822363
-1.8825587450258 -5.63078399959954
-1.44751382991273 -2.14114206316334
-1.66380916147449 -3.64484099058998
0.401861302912721 1.31973046275312
0.623271629065938 1.68539340121019
-1.73144970591821 -4.20569915378075
-1.77371702050308 -4.57978392318897
0.955956320015098 2.66816793758185
-0.18897993598519 0.91023036546916
-1.6234885758441 -3.33190439423898
1.09818757474448 3.29498592215778
0.309124261139459 1.21579642671328
-0.0852514379276199 0.959103955205731
-1.20296402390791 -0.907436646342982
-0.270626161478439 0.868269331607198
0.580911905551383 1.60051121866171
1.89911009192989 10.0943055577833
-1.064461094195 -0.394142474879605
1.01285946066994 2.90217317390131
0.0402853617909917 1.0207368752534
-0.721932674649874 0.424789386080467
-1.66297073792276 -3.63817368047019
-1.81194227253934 -4.93425282463376
-0.65222601397195 0.529244626005021
-0.471450844379962 0.729648654986886
0.0370646534067927 1.01903079271421
-1.40749154055525 -1.90763621540311
0.695183252394338 1.84812752899872
-1.3887485131105 -1.80275420632957
1.85291059689277 9.52105182117024
-0.357923961782797 0.815903315065873
0.974884463317879 2.7436137850711
1.09839082028482 3.29598410091773
0.830736918330413 2.2258271038816
0.470309061762373 1.41335589628585
1.72836875306984 8.09869154159071
1.46857383792811 5.66937208921194
0.456457372499913 1.3931356304433
1.66877894800176 7.47927364775083
1.39834539959636 5.12755138331202
1.56275324449412 6.4702782522399
-0.325545359330059 0.836477482958752
1.85992830537649 9.60650597509558
-0.723639418786876 0.422010679340869
1.94494266655738 10.6882289495346
-0.603203955844345 0.592860405091233
1.7499954109815 8.33311478767895
-1.54049166796916 -2.73586090180703
0.259609595491829 1.16958206817038
-0.731632112601908 0.408851007811368
-0.0519325411599803 0.97476729017708
-1.72438866505661 -4.14500022640276
1.5064929250508 5.98137889887683
0.887629961378167 2.41492021599477
-1.11507085610526 -0.5677200891142
-0.917751339073211 0.0266596159564805
1.59838189136926 6.79634572283596
0.0866958965644469 1.04645605141118
0.511077028010305 1.47686610904852
1.18267057389016 3.73671928077183
-1.50837712909309 -2.52197132499597
-1.11106653271358 -0.553417700758948
1.81148448128094 9.02829616101279
1.50251810418365 5.94801949158777
0.814115306622483 2.17414199110334
0.790250422365434 2.10265274415143
1.63443981901696 7.13964875211235
-0.233454491645942 0.888004345247047
0.586663233387293 1.61158186572004
-0.0497866633920321 0.975786287788734
-1.61536089253957 -3.27070757076436
0.848178812923249 2.28176978444151
0.245000662573484 1.15703051541659
1.53488813686981 6.22420253630332
1.36834242044437 4.90994069397324
-0.704148960054754 0.453095256084725
0.00630860309351045 1.00316749235806
-1.40491432443655 -1.8930474153221
0.829387278753886 2.22157156070566
0.317830262660781 1.22454626287515
-1.90365540241397 -5.84965805885858
-1.62753556811879 -3.36260966263985
-1.07974178139327 -0.44492376214492
-0.0401652628917479 0.980375534538642
0.860426038380392 2.32211294404568
-1.97271544571729 -6.60201419432597
0.190591970662848 1.1141726925123
0.737104817700763 1.95457775769123
1.79080319347672 8.7896957865979
0.180385991615295 1.10676284928859
0.269402073604749 1.17826216971268
-1.21119745824843 -0.941869770922104
0.452501685883913 1.38748350398462
-0.672242667744425 0.501002119019167
-1.87624541070442 -5.56626179799371
0.697491518131789 1.85375997540669
-1.20317623202471 -0.908318352432709
1.03305973959424 2.99052061568452
1.93968190507009 10.6187642477155
-0.821861577213577 0.242581884621014
1.63669141219185 7.16153596429626
1.12811869078702 3.44528671234127
1.18993959765328 3.77729979941787
-0.25536082126297 0.876524214114186
-1.53401697010337 -2.69200261869658
-0.179743031420164 0.914703026826147
-1.02306687395788 -0.263426005993494
-1.74238098504039 -4.30067414845448
-0.468060562638022 0.732596940922769
-1.52858752321493 -2.65551241532027
-1.66643863626171 -3.66579564315443
-0.641421382308617 0.543930131873093
0.538245752971836 1.52265960139125
-0.955234916781065 -0.0698838993363239
1.80447172410411 8.94684189758768
-0.263185865738656 0.872320671481494
1.20310768144324 3.85187921672936
-0.759262351682819 0.361455230000449
-0.100655606763484 0.951924170797343
0.395716018306588 1.31206807179805
0.87211821829841 2.36145681851321
-1.15580354177307 -0.718933512086722
-1.46536166469255 -2.24955290054648
1.48295138677626 5.78604032278067
1.61459884644588 6.94907314634342
1.12568219517709 3.43280119124754
0.866338688902824 2.3419070079037
-1.88149196370742 -5.61985003530082
0.796995975427856 2.12253861289157
1.18993904227527 3.77729668298541
0.00326378068726907 1.00163538782081
0.99147879094986 2.81170713649919
-1.16854127436226 -0.768401909148824
-1.41411347488116 -1.94536708041133
-1.18129361556432 -0.818990900220248
1.75555598117783 8.39422907605019
1.4068679209575 5.19085919088058
-1.11794299148175 -0.578039886374277
1.35976299670166 4.84920562717139
-1.62864042117664 -3.37101946941616
0.837123791554514 2.24610769688773
1.90712812877881 10.1963798514224
-1.04264993613093 -0.3240363185946
-0.899533795903042 0.0710358364792234
0.131443054211328 1.07365571281068
1.58641001943151 6.68533601707391
0.133184681131325 1.07476973412852
-1.66808451108402 -3.67894626701431
0.958140299087838 2.67675334517056
-1.20071553721398 -0.898113080369225
-0.29853548062793 0.852555704806701
-1.59020194748787 -3.08521285924393
-1.14679993698604 -0.684600564244047
-1.80395518812676 -4.85889899518877
-0.137678202125405 0.9346581819364
-1.47745500338493 -2.32453774310355
0.939519736387501 2.60454548209869
-1.35213428522258 -1.60588969662356
1.55564777901188 6.40679070736085
1.3123367672864 4.52521084209601
-1.00625054930137 -0.213104313097389
0.22162258394559 1.13789256241124
0.480912021948272 1.42929135626533
-0.400866712337033 0.78609239569891
1.49035524462727 5.84689554278936
0.160235802126708 1.09266401209588
-0.0716438000134736 0.965470942110219
-0.888570383890701 0.0969612511986682
0.982840776990081 2.77603261006682
-1.69495795594223 -3.89744184996925
-1.4617497943475 -2.22739769747757
-0.42314483375042 0.769348589512897
1.59181588354367 6.73528003652307
0.591841901406039 1.62167046929035
-1.93229760790966 -6.1549561757401
-0.306711411046629 0.847784880556625
-0.138148433919949 0.934437841218547
-1.26896132597228 -1.19663623420301
1.68454711978169 7.63943621503765
1.1392782281719 3.50304601511557
-0.480386448395481 0.721754395289265
-0.490322690457073 0.712761417952586
-0.275439661364084 0.865618817515639
-0.54345663381348 0.660588689748085
0.613200605364189 1.66450249809784
0.688492804968752 1.83194973760757
-1.21377118055488 -0.952727981570386
-0.290464321290868 0.857188260010645
-0.477369279494698 0.724440156259689
0.527679401870274 1.5045097951213
1.52837326464824 6.16778693077292
1.31999322543593 4.57618457231067
1.49828207824019 5.91263737722987
0.141132714179094 1.07991059789785
1.27727139459846 4.29818992337106
-1.31957472473445 -1.43952592683124
1.07403596830212 3.17852442267588
-0.964152382379572 -0.0939174689590342
0.372903814950484 1.28462951769261
0.338964532666234 1.24661469892759
-1.90283120211612 -5.84101222823902
-0.538926179282687 0.665318605805233
0.640452057593696 1.72208273447832
-1.79462195146616 -4.77171112771805
0.845175689742974 2.27201170294667
-1.12367227739812 -0.598779210451571
1.38580255251099 5.03558633887015
0.781039940610238 2.07590383792586
0.742219428079976 1.96817826018767
-1.19406467829478 -0.870734118277661
-0.0449801113196626 0.978074598174748
0.632269387643565 1.70444168365252
0.95299971641987 2.65659470437528
-1.67434856904174 -3.72923934056497
0.295928098180434 1.2029004453019
-0.626471885314217 0.563620670988828
-0.214640428951248 0.897563375662694
1.14999846380373 3.55942407415894
-1.69462783017563 -3.89471430687958
1.00236726771143 2.85739090847459
1.63395337111733 7.13492709980748
-0.823298339067957 0.239657661120865
1.32830160638183 4.63207459559869
1.04968038564894 3.065343494382
-0.872053718616644 0.134932772431856
-1.35649547964011 -1.62878970674586
0.437074209182266 1.36594863663831
-1.95504569295756 -6.40420740716794
-1.08877311971084 -0.475592819257362
-0.490078402446558 0.712985271015062
-0.0849870563214972 0.959227383047594
-0.835507401722747 0.214440962342617
1.15150322718474 3.56740809740603
-0.348840406187581 0.821829140042039
0.32956554445898 1.23665318101434
-0.634505395875966 0.553129141425036
-0.712102981416007 0.440580147511722
-0.141181527238492 0.933015530624069
0.480824629828942 1.42915836771712
-0.993803048145222 -0.176864334335354
-0.350082585525117 0.821026109561412
0.516649780001659 1.48602722259306
-0.5906743819141 0.607919149502502
-1.38542315227203 -1.78443923472196
0.0552461398158584 1.02878717007279
-0.160301599759029 0.923997242450794
1.61316880277441 6.93549586466522
-1.01829581720307 -0.248988405125319
-0.605550343710735 0.589987699174208
-1.13095129953082 -0.625425045182286
-0.430802229699371 0.763374751301627
0.481892362246438 1.43078507028784
-1.3585046072138 -1.63938887888464
0.738234268448013 1.95756946471052
1.01394446821935 2.90684705473204
-1.41398252492884 -1.94461750198599
-1.61322129466368 -3.25470147546836
-0.652286033515564 0.529161966391323
0.933830736067949 2.58292913135864
1.43087486679923 5.37279130342767
-1.00525884447349 -0.210185800461815
-1.50577099100418 -2.50501195390687
1.69439119087545 7.74078341437282
0.10925386070688 1.05983790235327
-0.188831907928934 0.910302370929738
1.57725839578374 6.60146623821642
-1.65640169917705 -3.58617255705754
-0.0101443338755236 0.994960212353824
-1.08708920938024 -0.469837268188092
0.441869766673601 1.37255664793776
-0.564949950927475 0.637391656904085
-0.823574016021408 0.239095543015953
0.174998644970009 1.1029256423352
0.35957055618162 1.26930486866708
1.01517511591745 2.91215807920767
-1.01068379458088 -0.226217569950526
-1.01107275652972 -0.227373296108133
1.59442120696011 6.75945724907352
-0.345592888348411 0.823917773481681
-1.29280590229753 -1.30872129958885
0.908709804141768 2.48992821737328
0.251582943895579 1.16262780597153
-1.67068398963368 -3.69977029502073
-0.358884765262246 0.815269170072609
0.543566188111855 1.53196512814761
0.39052183014819 1.30568228138866
-1.78482711120297 -4.68120920110513
0.392427898563614 1.30801603114462
-1.75534148417263 -4.41487490434709
1.8503004408122 9.48941536715725
-0.919870775031479 0.0213903298721088
1.47400951506076 5.71324744708682
1.75263123441605 8.36204118570609
-0.198501908355358 0.905574384373035
1.30090062834669 4.45001583074007
-1.66407712761412 -3.64697335799929
-1.84703443163292 -5.27351959062834
0.849949429186349 2.28754778244004
0.212031550491898 1.13036315308368
0.328070932321069 1.23509099799962
-0.76219608701357 0.356245856880229
-1.05034528540007 -0.348455087434208
0.130191110209795 1.07285764383031
0.238265181418163 1.15139947479866
-0.1604445162542 0.923929431107584
-1.88544272017285 -5.66040746381983
-0.355263481040195 0.817651875296108
-1.62197081642769 -3.32042900348838
-0.27260628005381 0.867181855396237
-0.0393374727000664 0.98077204637976
-1.67029587641437 -3.69665697078896
-1.4710559598251 -2.28470574025771
0.3082855060622 1.21496370491836
-0.172468446827101 0.918198235346273
-1.34416057620319 -1.56439999322047
-1.33564904373112 -1.52064952502959
0.802023719528068 2.13752419665797
0.59189749511968 1.62177939336193
-0.325807182637529 0.836317032850372
0.929845640955968 2.56790994971066
-0.220918870588491 0.894401548446393
-1.71225608825882 -4.0418882790471
1.8403591192216 9.36965021921371
1.12938560758863 3.4517965635835
-1.5047675750776 -2.49849804363137
-0.861877904777332 0.157687562511912
0.613015809930273 1.66412336752378
1.99783392103624 11.4055229812418
-0.600432925376096 0.596231459695451
-1.45190828144078 -2.16758691243222
-1.48904013730822 -2.39754134298098
-0.285818161709389 0.859821491004982
-0.440240268030195 0.755851362680721
1.94114199811838 10.6380099391004
1.5203777216456 6.09912332861589
1.42674268688303 5.34109535400439
1.53771130486377 6.24878063047949
1.08882233533862 3.24932053509569
0.737204249571775 1.95484086941443
-0.658852549666758 0.520045515942364
-1.56974238807888 -2.93869925768835
0.573850104841899 1.58710864281915
-0.150629813670606 0.928571540966119
-1.59348005323711 -3.10904716797659
-1.78069008218441 -4.64328957281517
-1.23127998733112 -1.027802140198
-1.22026092161847 -0.980308555640824
0.882795537168005 2.39810156317915
-0.930186148595834 -0.00457659538527411
1.79227171174474 8.80647713121272
0.766752410763582 2.03532184929355
0.773207551344699 2.05352048702182
0.260074786136564 1.16998950568125
0.252140468136878 1.16310624064939
-1.14290600744467 -0.669913248840446
-0.374213118765592 0.804955680857456
-0.191181658633414 0.909158058002845
0.810589974522598 2.163381259695
-1.41514599508773 -1.9512822678602
0.395078496897364 1.31127984273892
1.16524505127672 3.64112744273495
-1.03632649241488 -0.304226042565756
1.18948005796077 3.77472198587331
-0.270297527364671 0.868449433874931
-0.198422000604987 0.90561366269916
-0.857131906047643 0.16813602458199
-1.47168682695009 -2.28861721392019
1.30970381039638 4.50779879595879
1.08896952136146 3.25003325106136
-0.253744816342775 0.877385261341134
1.60809238143707 6.8874700276336
-1.46164224840128 -2.22673969804884
-0.523674976294005 0.680845431379947
-0.436105436776697 0.75916946927637
1.95670181580846 10.8447229790432
-0.29578081090816 0.854145237454634
-0.146498271875552 0.930517489092381
-1.48114813660415 -2.34768518142267
-0.360618287314252 0.814121407141682
-0.137969553008917 0.934521666109314
0.692041230923785 1.84050269395564
0.848541326972606 2.28295127115617
0.306602934648166 1.21329862270282
-1.99440478525603 -6.84990030388248
-0.163046902412025 0.922693441264199
0.807028705077833 2.15258176473391
-1.38021215283323 -1.75591511891502
-0.391259388345681 0.793033169097501
-1.28762850152753 -1.28403263820684
-0.174382470327815 0.917280783766521
1.13131840297683 3.46175128155318
-1.37045581518671 -1.70308653950514
1.88195365105748 9.87847330015004
0.42451589737284 1.34900533570564
-0.250619191395958 0.879044006890007
0.203502318902969 1.12381816913033
-1.58694498485549 -3.06163097800262
0.506640548475837 1.46965705070567
0.74023425649113 1.96288327802345
1.21036770304808 3.8935892176777
-0.783689800415043 0.317015675037404
0.36004969780675 1.26984669049933
-0.00841423401782704 0.995815284654669
-1.36225289641879 -1.65924676348695
-0.169846151402795 0.919452780658226
-1.80167404054627 -4.83750329606028
-1.71064649794186 -4.02832064119088
-0.0144324412344345 0.992848406647661
-1.71446370448892 -4.06053936154386
0.588542852677619 1.61523034349259
1.69359471543941 7.7325445269266
0.960172150386684 2.68476864465872
-0.501004418258892 0.702834934789813
-1.73543242629262 -4.24016056347248
-1.18738325377207 -0.843528056777747
-0.784046358824987 0.316348896817618
-1.41258173898454 -1.9366078844798
1.71385367062331 7.94424996569319
-1.95982876404981 -6.45738784953313
-1.74772858892011 -4.34758380647548
-1.8435030868594 -5.23877001754084
0.662633333223705 1.77145552470409
-1.97857398165206 -6.66841640626697
1.46737028666602 5.65969553365888
1.40820450129305 5.20084818512847
-0.26203923453244 0.872940237525545
0.446314292887601 1.37874996770754
0.423303278566297 1.34739667713506
0.654264776554472 1.75255916982409
0.195648317957575 1.11791351216775
0.869789244443124 2.3535548140646
-0.972812372591263 -0.117657788190493
0.728326427169269 1.93154966766546
-0.615050256398702 0.578184193752807
-1.68378682359237 -3.80574701635623
1.32645410070877 4.61959452246244
-0.00306303876824643 0.998471500496445
-1.92613251195332 -6.08844338856703
-1.7678595137034 -4.5268330652448
-0.294710498461862 0.854760461822923
-1.4156965362227 -1.95443979345889
1.73060191716862 8.12265833310048
-1.19280754552839 -0.865592394287253
1.31589340472998 4.54882640739453
-1.20936667062426 -0.934173394631358
0.824988959339324 2.2077756649389
1.44994772354708 5.52116167829113
1.66091128157801 7.40035470103932
-0.822159008525746 0.241977269724585
-0.783898941357204 0.316624637232964
0.535165215213903 1.51732290557903
1.37416286857846 4.9515201528893
-1.26707642429524 -1.18795128714585
-0.594686170782149 0.603148756327027
1.50327056291737 5.95432281026179
1.92323778364012 10.4038007624767
-0.0661064161920821 0.968072202724847
1.7973699615765 8.86492744905595
1.78822759931664 8.76032265014644
-0.0355696491027953 0.982580429304874
0.710778270422538 1.88669431924725
1.79510199198771 8.83888910864229
-1.6013204280294 -3.16645825549723
0.827550060828212 2.21579542506338
0.710290387126795 1.88546942765681
0.75328106345806 1.99805941204832
-0.602679413986287 0.59350031537745
-0.736116787059058 0.401360263426038
-1.75442056855673 -4.40670289919965
-1.43469212012865 -2.06490171684732
-1.42457901983685 -2.00572488115273
-1.85067490959863 -5.30948688938268
-0.356802885485278 0.816641446786753
-0.483117331157287 0.719305522849883
-0.737512856820079 0.399012610082135
1.93790301004031 10.595351464905
-1.21266198885831 -0.948042898510603
1.53227608060943 6.20153286471949
0.94099363195935 2.61017957606926
-0.00887345133348605 0.995588151009568
-1.15415909725177 -0.712623858434803
0.591980211072964 1.62194148219783
1.36550586974274 4.88978744815124
1.70357448545846 7.83627435027178
-1.70289120626857 -3.96331499568938
-1.40104805470438 -1.8712620281133
-1.61851806470054 -3.29440499476517
-0.63112002838442 0.557575573319032
-0.565084838129248 0.637242049113841
0.851954413494163 2.29411277330813
0.48982569862781 1.4430008972904
1.45141946350045 5.53275280275556
1.29955424466575 4.44123700840332
-0.84757109003723 0.188870428958125
-0.0599663634505343 0.970965410863567
1.24627519307932 4.10615513282809
0.778027223959338 2.06725477506621
1.11303740065966 3.36872048346061
-1.07369452877235 -0.424661437034536
-1.06772890975252 -0.404885801541526
-1.61984117288767 -3.3043641509849
-1.56478897897108 -2.9038045844714
0.935177122726119 2.58802629462615
-0.511888877383642 0.692436685658468
0.453053721775183 1.38826905851069
-1.10552539167778 -0.533789646913369
-1.82968366491081 -5.1040974467213
-1.66590564833054 -3.6615427251593
-1.92322212449706 -6.05719723651022
-1.64755790625506 -3.51682616610233
0.536198102612343 1.51910810086017
-0.107249325761568 0.948854634668092
1.96901847512502 11.0104557393676
0.223121280965427 1.13908566639894
-1.70559709939951 -3.98592738382609
-1.25120198878827 -1.11581428511409
-1.86325846684435 -5.43493952243064
-0.117696588412903 0.943989897472763
1.82415943911388 9.17695190231721
0.562110201979332 1.56528687742743
-0.560464756988403 0.642337358042499
1.36413404164186 4.88006667367168
1.85764949536136 9.57869371898181
-0.572732680065321 0.628675610461671
1.31909501625392 4.57017833613826
1.31768933845374 4.56079275955768
0.953529119384244 2.65866280935027
-1.62428029635491 -3.33789907091573
1.23452868466645 4.03546076426489
-1.22147326342045 -0.985492891674448
1.32006706182028 4.57667862051968
-1.75140090095866 -4.37996875561887
-0.879055494390023 0.118993780269021
0.54827145310572 1.54028869924352
-1.03646261797876 -0.304650088505529
-0.29642091653693 0.853776665310185
1.69965456488001 7.79540165753882
-1.92626251476466 -6.08984139395354
1.25754503364581 4.17505045600754
-1.39510351177812 -1.83800024057679
0.768149445576419 2.03924146447786
0.629209199577733 1.69792232313776
-0.398388833061027 0.787898226117644
-1.85205832248426 -5.32319312482858
-0.788306477157378 0.30834141495755
1.92336307531936 10.405426230607
-0.70940929986311 0.444844585697191
-1.62290632240129 -3.3274995398133
-1.90084325800046 -5.82019063857491
1.89309179361197 10.0181944192673
0.470274000325968 1.41330386505985
1.62265830770733 7.02598987553108
0.893485980468911 2.43548364637751
-0.733495310758882 0.405748290948773
0.148485290971542 1.0847514676653
-0.150844352111013 0.928470370725297
-1.91783695783904 -5.99963945317911
0.190400907466076 1.11403225548728
0.399679743270707 1.3169969369697
1.096754005881 3.28795402126858
-1.49633334111525 -2.44409122497096
-0.958128006818784 -0.0776354898908592
1.98316190841294 11.2030781834086
0.235138767152799 1.14881836851237
0.159138615959582 1.09191580441678
-0.50351942722194 0.700457964971038
-0.747017559231104 0.382828317131004
0.954233168495334 2.66141598552957
-0.984114947122056 -0.149243088442066
0.566261707929518 1.57293847659294
0.0268518303870944 1.01368001238554
-1.5080912666153 -2.52010818864839
-0.698445924122061 0.461925923461316
1.55686017968182 6.41758757741351
1.93469424921284 10.553216972581
1.53776104519807 6.24921437317124
-0.862739132584211 0.15578038107934
0.79029759576884 2.10279094050976
0.0471274247975035 1.02439239875053
1.78987721790029 8.77912690052206
-0.0431857369949249 0.978931048381096
-0.929862394078688 -0.00375346834793933
1.92044127119492 10.36756965867
1.65432556888321 7.33480166918088
0.33085216554685 1.23800274922351
0.70593915826521 1.87459777061375
0.529588679409994 1.50775701177613
1.15961162919433 3.61072991364129
0.927783599136133 2.56017804399492
1.51679446472447 6.06855578095598
0.505843770609108 1.46837017949939
-1.86608403340888 -5.46335098835564
-0.13183306760037 0.937393734241051
-1.5948633922361 -3.1191350481614
-0.402491431752823 0.784902332976656
-1.2133053277603 -0.950759255996364
-0.399992927369038 0.786730451090819
-1.12793902639114 -0.614357814757456
-1.81371503806416 -4.95107099648
0.999099306135273 2.84359567367731
0.252238154701149 1.16319013979868
0.724810795768328 1.92243806366124
1.35877053448433 4.84222226357485
1.44383629208362 5.47324790942361
-0.52201647213084 0.682497725271347
0.792357051294308 2.10883613132269
1.13142344393241 3.46229309621455
-0.909827386308057 0.0461623830673111
0.627588348397961 1.69448642159623
-1.66249879225256 -3.63442369352429
1.0244559071169 2.95254617827994
-1.06115849580953 -0.38334870945599
-1.08590118719238 -0.465786940173083
0.547207888067801 1.53839951004046
1.52819489594 6.16624825942774
0.638538035700954 1.7179289325443
-1.4204616749999 -1.98187240908629
0.966847435128181 2.71129155995665
1.75335340411573 8.36998004208666
-0.502963356595071 0.700984834758228
1.93133594447325 10.5092527985002
0.538941134559069 1.52386945789964
0.188749750665735 1.11282136647531
-0.292593598496842 0.855973391432946
-1.37039122070766 -1.70273926532393
0.102400835197629 1.0557047712163
-1.53332126060227 -2.68731226056737
0.891235759838326 2.42755720186516
1.52170743807571 6.11049884802857
-0.844376562263295 0.195705543756215
-0.375488262401924 0.80408067223377
1.4696309220224 5.67788246054487
1.13816047065756 3.49721817323933
-1.84272956088872 -5.2311765844534
-1.74749904667497 -4.34556417819544
0.644838063213352 1.73166457646492
0.482666975213408 1.43196777559443
-1.89561475321412 -5.76564239081856
0.0640935205283264 1.03365067786788
-0.993936303270147 -0.17724780418363
-0.182062122929765 0.913583883048714
-1.76248041563369 -4.47852345492373
1.6180302702965 6.98173882657229
-1.72862303592487 -4.18133915761499
1.26946417620312 4.24906447899531
0.0232857615563313 1.01183199442652
0.376062841662597 1.2883363499593
-1.92558451797562 -6.08255259910049
-0.301932119684734 0.85058343404128
-0.615330242664868 0.577832079638759
-0.574937648601093 0.626174929312303
1.24780110159707 4.11542187846027
1.62037688557016 7.0041481320128
0.776891640072805 2.06400746289965
0.805897067176439 2.14916497188128
-1.04699449756928 -0.337780348737123
0.2413457555599 1.15396291866958
0.779938874450223 2.07273712296313
1.05148057441272 3.07356436028159
-0.60208950984249 0.594218965773506
-0.662847385953129 0.514428181055222
1.47783716222536 5.74431304459151
0.677094877773097 1.80489057535815
-1.36138557463422 -1.65464211025583
0.826952051915089 2.21391948052634
-0.616460223808683 0.576408519698541
1.90049331762246 10.1118598699591
1.8857421665975 9.92583229362722
-1.99232133463139 -6.82584373830562
0.347339572283597 1.25569397803319
-1.94164832298215 -6.25668016754697
-0.87580042214737 0.126432453372567
0.618131148387185 1.67467389041367
1.64951506747112 7.28720973548482
0.856870484129805 2.31030965843759
0.75816365172245 2.01145359287336
-0.630893720992384 0.557871496703937
-1.68440412323232 -3.81078159312989
-0.801909424927161 0.28226301307832
1.67684929108731 7.56091329136224
-1.77823002811071 -4.62082666321174
-1.50476161515709 -2.49845937961117
-1.07092277187912 -0.415446994640916
0.554453463411948 1.55135983180469
0.343066363886576 1.25103729993777
1.78407551388253 8.71312912961346
-1.04397837982574 -0.328227301390198
0.562222038737691 1.56549207455857
-0.866328340604409 0.147795038857713
1.25613287223244 4.16635985946709
-1.90746279943368 -5.88969818106689
-1.287035879447 -1.28121919610003
-0.539949798500574 0.664254689315366
-0.862636863904048 0.15600703402093
1.01478429583278 2.91047031699819
1.23975984222251 4.06680364423564
-1.19235655588129 -0.863750409325774
-0.831963437570925 0.221828798485852
1.40265973241274 5.15951629228891
0.0194303791313675 1.0098453788844
1.40372315003464 5.1674213568491
-0.528505640610852 0.675992745163049
-1.6564367176555 -3.58644865323731
1.69314550365526 7.7279008288673
1.18427001205538 3.74561264251057
-1.49048114937126 -2.40670243203198
-1.5021541457011 -2.48157349939488
-0.973580992960761 -0.119784081701067
0.780153688626938 2.0733544205084
-1.97175917019071 -6.59121430968625
1.70812572752589 7.88393967866811
1.71164628801791 7.92096566685086
-1.46901377285881 -2.2720670262631
-1.23705343296793 -1.05302246946098
-1.31701534670063 -1.42678900543212
0.973916433523538 2.73969810519761
0.0138870558554998 1.00700893829624
1.24736888638784 4.1127951120542
1.85797901316208 9.58271162347587
-1.55414649990604 -2.82958745285144
-1.26664102180807 -1.1859487335488
-0.491671869592991 0.711522572045034
1.56608847925319 6.50025427042704
0.106986326153375 1.05846356533392
-0.524782636566406 0.67973800984612
-1.28628653522886 -1.27766539355675
1.96238202836105 10.9209239716033
-1.51631884062425 -2.57401806958537
0.649325034600645 1.74155870983148
0.724807417907873 1.92242933945989
0.0888190227673209 1.04768881497555
-0.727130728089694 0.416292246427935
1.61375447966397 6.94105389685173
1.85902495714243 9.59547356068222
-0.26205808636835 0.872930061358466
-1.91095938475336 -5.92661571731319
-1.84063806586138 -5.21067802010906
-0.83203390035015 0.221682455008366
1.77395758284144 8.59894387399551
1.39522424578815 5.10453286882605
0.991061263606985 2.80997126875214
1.75867999543365 8.42871557387581
-0.145078200731571 0.931185345721736
0.424811330532875 1.34939798162108
-1.97676835775563 -6.64790737169224
-0.0224534136102359 0.988925588632502
-1.79084339270371 -4.73667740838314
0.623736702573969 1.68636904571803
1.74134746699276 8.23875263443377
-1.22329593745987 -0.993306226618697
-0.68863543648969 0.476841503255705
-0.841565476998836 0.201682049751289
-0.114665665274081 0.945401721077095
-1.27052530195143 -1.203861842856
1.78013467986155 8.66851747260619
1.86228752760561 9.63536391582172
1.54718747645037 6.33185935875192
-0.0767255669261084 0.96308944111748
-1.35816488303435 -1.63759446223621
-0.494991410247095 0.708456202517538
-1.45696110771646 -2.19819336995244
1.06594787433349 3.14046823450914
1.93612990897787 10.5720532579514
1.90943433588774 10.2258819915518
-1.45262100952012 -2.1718911615037
0.785661940921628 2.08926897040826
1.48129203807223 5.77247401344108
-1.78019018746414 -4.63871983208713
1.95664099860872 10.8439092498593
0.477914979210265 1.4247462952397
-0.515843724751158 0.688586128698253
1.50482410118571 5.96735421244212
-0.434201139132909 0.760685999633963
1.09245878750959 3.26697566808596
-1.35868571455372 -1.64034585027071
-0.758148874059312 0.363423392727458
-0.786945539293618 0.310907702980332
0.52878601489933 1.50639014501486
-1.63848464065965 -3.44646537614018
1.36992278408109 4.92120022086491
1.61674374079115 6.96947727260168
1.0830007265137 3.22125744692358
-0.492413148033551 0.710840088759813
-1.06776075633125 -0.404990811139046
0.349569286229697 1.25814391282824
-1.76802641988979 -4.5283368888717
0.454907214792726 1.39091425702831
-1.90638792522486 -5.87837759132325
-1.54731497238487 -2.78248642605251
-1.43140345060785 -2.04556579019779
1.81777930286695 9.10189243174927
0.0749928195697884 1.0397548078862
1.89915503061025 10.0948755081904
-0.456716963407882 0.742278124899361
-0.831663785685692 0.2224508961276
-1.21310084781679 -0.949895577925508
1.56068431452565 6.45173980936193
0.918894025714693 2.52715306974637
-1.51902640644904 -2.59188884866268
-1.85715800237426 -5.37390130183896
-0.19956111548409 0.905053393591488
0.127893873663956 1.07139913277355
0.78402794741665 2.0845307395796
0.808996478858507 2.1585402118705
1.22984588684676 4.00759330424166
0.234445791603679 1.14824903644456
-1.2782780393036 -1.2399399481654
-0.609232095571956 0.585446173369471
-1.03990220203699 -0.315399958059745
-1.63056543630134 -3.38569990115128
1.67575051841743 7.54975703877128
0.790850052170221 2.10441029419865
1.43535004703875 5.40729800660216
1.28305857717661 4.3349362174954
0.602211871837049 1.64221837714499
1.17002221696631 3.66709775231917
-1.96969016459752 -6.56788478965851
-0.420503314401308 0.771382845430355
0.00566025559792926 1.00284072543611
1.68187917657893 7.61214922970133
-1.32969440699063 -1.49036988534212
-1.21559370058737 -0.960444327997432
-1.51750003742957 -2.58180640932343
1.78995853088708 8.78005459910523
-0.533564886410108 0.670845822719905
-0.863322864649673 0.154485756564612
-0.485369859142183 0.7172726778564
-1.90882702469775 -5.90408524932761
1.41520006110501 5.25339829886051
-1.35924644196202 -1.64331034841882
-0.861356600799521 0.158840316600257
-1.31755553564965 -1.4294732040879
-1.11214461415033 -0.557258553461082
0.383861508736479 1.29761414507274
1.41257472447696 5.23362405936077
1.66732817080819 7.46467150976683
0.315562558070873 1.22224832688126
1.15482684194415 3.5851043270083
0.883803781254541 2.40159746566727
1.34039284239157 4.71448965561803
-1.96618885461039 -6.52852077074172
0.520093164339951 1.49174708496151
-0.059431820558177 0.971217772926455
1.20056066320666 3.8373461174504
-1.26527759983197 -1.17968671330802
-1.19603151765713 -0.878799840768772
0.908964190174309 2.49085015515867
1.64444117719047 7.23727740584927
1.83226136004263 9.27294537353598
-0.618081054980039 0.574359623139687
-1.79611021553569 -4.78555162805232
-1.45813243987457 -2.20531905080718
0.0797175808956805 1.04244138701734
-1.68602571107216 -3.82402493587793
-0.83913379187387 0.206823209072277
0.369839788247009 1.28106217047193
-1.96293142430682 -6.49202905123444
-0.977438612004603 -0.130503269498393
-0.764717547176466 0.351740938790488
0.665388103370584 1.77774792956026
-0.95405496559106 -0.0667349281815306
1.74365110010926 8.26380751553717
0.88937319263139 2.42101968929262
1.44053858124633 5.44753982488186
-1.06214485291108 -0.386565669563593
1.61677967387832 6.96981950711904
-0.969711271564239 -0.109110859264739
-1.10111582664108 -0.518304852801946
0.479358068620247 1.4269307348515
1.00147119269776 2.85360105200195
-1.50614978008554 -2.50747323363271
0.135087463666287 1.07599192289944
-0.10437397163449 0.950193113552419
-0.865024945595861 0.150701782569106
-0.227882217765365 0.890862384696654
1.74774498479461 8.30847914790494
-1.06749011983779 -0.404098616506272
0.447194619714659 1.3799846034777
0.191706917504417 1.11499353633612
-1.53467718360765 -2.69645765329396
-0.0303931324022266 0.985074991992057
1.97064170860828 11.0324372794247
-0.338622466202244 0.828348925746348
-0.763123127650191 0.354592550778447
-1.78187740239317 -4.6541539279079
-1.76095886117017 -4.46491322418032
1.05509605695447 3.09014445357528
0.471792586686954 1.41556143616822
-1.63713799322961 -3.43608993064268
1.21639985543668 3.92856695231781
-0.209052526897512 0.900355406759493
1.17052390765357 3.66983542516921
-1.68780168124837 -3.83855902194721
-1.3546328996739 -1.61899157436293
1.48663823198821 5.8162777913109
-1.65635437320755 -3.58579944386302
-1.67608702535521 -3.74326547549861
1.21350843499642 3.91176433969345
-1.87888741434495 -5.59320853108333
0.395408335680854 1.31168749807443
-1.79419211586412 -4.76771815281607
-1.51250302176874 -2.54894179033519
-1.61684686950874 -3.28184943102874
1.58631980276513 6.68450505598533
-0.802594996628898 0.280927956280322
-0.0794859431731667 0.96179783037884
-0.732966009715986 0.406631070533624
-1.01893260734773 -0.250907994009945
-1.63790006061609 -3.44195925555662
-0.615478446559309 0.577645597833662
0.343695035731867 1.25171922039455
0.57779178361288 1.59456375401596
-0.699937440091468 0.459627865535633
0.707100514414448 1.87749020069058
-0.886972889362756 0.100690569025272
0.815133489524867 2.17726293510365
0.764683353719457 2.02953595477432
-0.34376078044843 0.825089264565482
-0.335400168623273 0.830373901387892
-1.60973235621782 -3.22869365261734
-0.0625530951816025 0.969745318511178
-1.56187486241968 -2.8833802995486
0.235851754428173 1.14940519024992
-0.604510140635392 0.591263300128981
-1.52595209584025 -2.63789453670762
-0.0167539530500709 0.991709448812335
-0.48133754314143 0.720903460043628
-0.331574810527902 0.832758980354371
1.11301569467954 3.3686115109984
0.26433173410415 1.17374059837321
1.86404601882561 9.65691630769074
-0.351661789064046 0.82000187970048
1.30528355617809 4.47870152488599
1.13849983059223 3.49898654693974
-0.700254175629474 0.459138814559712
-1.16323425742782 -0.747663206520722
-1.09317503071268 -0.490719432982972
-1.50947056014758 -2.52910445464645
0.369611341071494 1.28079729679421
-1.93280424225842 -6.16044165455544
-1.58252649628825 -3.02979597900591
-0.763232969381421 0.354396427529145
0.697362563986819 1.85344461974421
-0.78885638453812 0.307302274213614
-1.97001006573632 -6.57148858056276
0.41282324688506 1.33369168895236
0.0915867832461297 1.04930396860131
1.38481369648395 5.02839681198353
-1.37682768681298 -1.73750406988822
-1.81277009698659 -4.94210215315662
-1.5436799434384 -2.7575953075135
0.31214230664009 1.21880760764692
0.745415545892666 1.9767462646854
0.908678008631643 2.48981301339038
1.58766883354634 6.69693929600202
-0.115972246424081 0.944793175590471
-1.77203737735337 -4.56456341357926
1.99893935864718 11.4208845267762
-1.08983879770738 -0.479244124305624
0.606564105436654 1.65098108377588
-0.59347945779573 0.604588689282053
-1.75972116339784 -4.45385988010386
1.41433405500788 5.24686843839027
-0.565038180623414 0.637293804248749
0.629341526646942 1.69820335474672
0.359092699084816 1.26876515243835
1.31209455722298 4.52360657942307
0.0231655473192554 1.01176987386168
-0.91855554467885 0.0246628562726551
1.76618437749433 8.51200436784866
-1.88883365455313 -5.69535834021965
-0.35877884193425 0.815339150943887
-1.28223361298719 -1.258515120528
-1.74717176179298 -4.34268550251242
-0.0381931754896185 0.981320627546728
0.28756260071101 1.19494845761518
0.149413701940852 1.08536867189767
0.408579070397023 1.32824133588349
-1.49914779368727 -2.46217781434149
0.0744843979020424 1.0394671982087
0.838133071460148 2.24933406340543
-1.25106277067889 -1.11518956498603
-1.94527337085676 -6.29639042233159
1.63989670287326 7.19278579532602
0.947947137930402 2.63694823975164
-1.95022300713495 -6.35085915138298
-1.14036041717668 -0.660364012151103
1.59891211959265 6.8012963388368
-0.557883632982971 0.645158125220656
-1.90969707059121 -5.91327184497705
1.08206173751063 3.21675416302705
1.34230080876573 4.72761185919997
-1.36304571536481 -1.66346099869758
0.198751080498258 1.1202323785615
0.0171134301767624 1.00865701431544
1.09076257328937 3.25872843042728
-1.22084344547998 -0.98279834030487
1.95397912147747 10.8083379921804
1.84483292780524 9.42340441154085
1.90754866825673 10.2017548388507
0.988413745574071 2.79899146699252
-1.16113479370591 -0.739509628040418
-0.898757806765679 0.0728899639194056
-1.72453369303265 -4.14624182096285
-0.218035835765646 0.895856789997443
0.269600038029095 1.17843990359055
-1.78591422979284 -4.69120367724842
0.658024340641895 1.76100770281471
1.80974732277935 9.00806624621475
-1.40394985040404 -1.88760160308266
-1.50873158780901 -2.52428253507125
0.307451961694163 1.21413792736538
-0.528884650568429 0.675609467041245
1.4386350581012 5.43274695287675
1.09033381004139 3.25664705752175
0.522694119614134 1.49609776021739
1.54715171273416 6.33154412831436
-0.669302419702281 0.505236420592782
-0.59255965317619 0.605683348425251
1.52492262642214 6.13807633920665
0.309944458185623 1.21661245618531
0.494040954130513 1.44958500461109
-0.721085446062989 0.426164658692831
-0.695056255808788 0.467118618177397
-1.07749284575177 -0.437362833929012
-0.791133218524787 0.302986357719246
1.53796642661311 6.25100558859299
-1.03394103954002 -0.296812271055293
-0.559812087764417 0.643052368011132
-0.633968339936727 0.553836990474224
-1.08897230869268 -0.476274773578637
-1.83406499079466 -5.14656803373782
-0.815162051023353 0.256098168103319
0.143125881657654 1.08121468711607
1.22397845948015 3.97292829861374
0.725829322232748 1.92507132628676
1.3715998567677 4.9331732876944
-1.52342357214668 -2.62104908437208
-1.93042075917092 -6.13466094661306
-0.616595189844368 0.576238222099196
0.213993461034763 1.13188857365299
1.39803812909352 5.12528131282752
0.100701243153916 1.05468909117862
1.35925049059371 4.84559833145911
-0.554927920621903 0.648365664774076
1.97667711505462 11.1144531365842
1.57800292613626 6.6082575605031
0.645395105646288 1.73288784568775
0.974282898146801 2.74117972496166
-1.03708996103011 -0.306605699392774
1.69454522744885 7.74237759001367
0.375377173129433 1.28752927895184
-0.655548801477311 0.524650300775287
-0.953868870865906 -0.0662389530943418
0.282829356831491 1.19052423629332
1.34300982734016 4.73249637214456
-1.17640153921017 -0.799457174811928
1.18335643107566 3.74053038779944
-1.58585802253435 -3.05378273525992
-0.406269387633428 0.782116550369699
1.64426823954281 7.23558031363243
1.8751683113906 9.79407752265142
1.6271576806091 7.06922491586328
-1.41956760269503 -1.97671119129604
0.91561499477998 2.51509692453709
1.42657150446897 5.33978572828725
-0.226384349697887 0.891626661442817
-1.37240178410673 -1.71356383293874
0.329503820648372 1.2365885486861
1.76835687279365 8.53623420599155
-0.0643571965988285 0.968895435777672
-1.04472205333709 -0.330577886582466
-1.54851585294466 -2.7907356294578
1.74531312125836 8.28192067639441
-1.29753201911022 -1.33142964625767
1.03058135477804 2.97952910077753
-0.739307202417413 0.395984147242487
1.07074212816424 3.16296938985897
-0.276871417668281 0.864825836028826
1.59428070989477 6.75815166865391
1.9775394033217 11.1262076505164
-0.322668723171736 0.838234307567092
-0.755764004994456 0.367622206163626
1.65383353457481 7.32992254260252
-0.144282128392698 0.931559525911121
0.402805096754445 1.32091764161098
1.24507461810011 4.09887761868257
0.735468533287763 1.95025524288497
-1.64111697855334 -3.46679688829703
0.0579891524950362 1.03028656200424
1.67760140437525 7.56855725419944
-1.69368273555164 -3.88691181889737
-1.24377755022998 -1.08268901549638
-0.432681352834765 0.761891095706417
-1.35297342057784 -1.61028445695437
1.92355852750648 10.4079623031266
-1.57448990466284 -2.9723539336749
-1.34222151991915 -1.55438427031929
0.289194162205326 1.19648599819658
-0.286787433879757 0.859274129623031
-0.656062103608636 0.523937269374979
0.536164444597334 1.51904986165716
-1.22738074864422 -1.010899981434
-1.29505040068755 -1.31948533683892
1.26892424565769 4.24568599818437
-1.92474243500583 -6.07350719887432
1.55085128086431 6.36422090095643
0.680577102095658 1.81309088302144
-0.298382311684437 0.852644321350996
-1.65984865582475 -3.61340658154561
-1.68973754295019 -3.85443729120498
-0.899730465091301 0.0705654553515266
1.79002475993214 8.78081026076583
1.24563099265532 4.10224871154197
1.57497112340899 6.58063788084533
1.37769729650331 4.97691775303615
-0.281410350029088 0.862297708674295
1.93491822661699 10.5561539738156
0.0363111485001473 1.01863296478277
-0.332987432240669 0.831880589168509
-1.80974006411356 -4.91340769524829
0.694896077357561 1.84742862080097
1.95457147794537 10.8162462688562
1.27308420348949 4.27177908733688
-0.874512743956169 0.129361319704098
0.914187743108128 2.50987033058807
-0.244938422010026 0.882036783630122
1.9907673585679 11.3076838361221
-0.13346042304654 0.936632714681418
-0.68233701334333 0.486236293949049
0.548963450448359 1.54152031847289
-0.808553351871949 0.269240248340317
1.04654601741391 3.05108454552819
-0.119795110934064 0.943012023037743
0.701049984621372 1.86249454770044
-0.226734136952503 0.891448333304594
-1.43358458674436 -2.05837995413147
0.374885551613246 1.28695146707067
1.68338246621681 7.62751502901291
    };
    \addlegendentry{m\_1(x)}
\addplot[red, mark=*,only marks,mark size=1]%
table{
x                      y
-1.52762947004398 -2.63665679852607
1.74584146706459 8.30313418504723
-0.523478826928232 0.684981455327291
-1.73063347132838 -4.16754055824511
0.199724526357971 1.12413986692501
1.56583537106678 6.50212130573281
-0.463057519114948 0.747665654157118
1.27292037122199 4.25911843213007
-0.561300740541583 0.653996446870297
0.640763194967071 1.72674751050313
-1.73791342939199 -4.24096264611894
0.131982542924304 1.07163562445192
-1.83168405496555 -5.10376196338133
1.09632486787013 3.28860190499629
1.14869409291258 3.54804856022076
1.71062342784982 7.90152466385944
1.54063884052671 6.26233498072305
0.96036949117845 2.67188471681226
-0.36582225744249 0.829228691385662
1.03292095360504 2.98417928381415
1.57196496539069 6.54609386202508
0.316383914894403 1.23104885644691
-1.25339289598648 -1.13554910601212
-0.46394682067964 0.740299643711383
-0.998800252147865 -0.168876133428104
1.03613462412239 3.00699952570108
-1.64532444822729 -3.49293200198847
-1.97015264204063 -6.5769121469698
1.10611827964268 3.3325927944455
-1.40144531975216 -1.84518158676689
-1.71485775358785 -4.06883587132901
-1.55651601432609 -2.85466248166607
0.704149753887608 1.87137436411416
-1.35391349938216 -1.62114783985896
-0.19713428931697 0.91320639815536
0.890553629834266 2.41963800191651
-0.748635890255846 0.367757908994129
-0.0825495453396563 0.955542429515598
-1.52863234199303 -2.64905992413016
-1.06308264251695 -0.412431701954534
1.34829445246638 4.75822063973227
1.21677724541683 3.93402096649771
0.321064511493674 1.22796917544757
-0.79448963572049 0.286639405408638
1.29306522904297 4.38129536324555
0.991145476362639 2.81469205516578
0.408046183683614 1.33070158460102
1.50470887336654 5.9575393202239
0.241965856580386 1.14976159314043
-0.698625089638838 0.462631462074469
0.69277772759134 1.84475566662164
1.54249469168855 6.29223000018984
-1.37940963447366 -1.77397785165444
0.328552328283534 1.2358697609747
-1.76899063762277 -4.55181464197594
-0.1615118341785 0.917358084831994
0.979551973019646 2.76474222344531
-1.41400818915192 -1.93743608707598
0.920422122884637 2.5371272327229
-1.58813906344299 -3.08034023934407
0.567973302168688 1.574798756124
-1.7027550272339 -3.92661639975712
-1.66143052087379 -3.6310451778419
0.509387499079396 1.46884150329936
-0.649372469635067 0.531357119197573
1.27714939756953 4.29644865621161
1.99062920502101 11.3301714102496
-1.77767653093197 -4.60832173114591
0.95231808494159 2.65334945903267
-0.729339734657573 0.394415405360015
-0.69929992108586 0.469586842821553
1.62359205741297 7.03622721245334
-1.26350376928392 -1.15873251814128
-1.80495773362803 -4.87445067342779
1.4755075644755 5.72000202565798
-0.582140577898955 0.62882666918393
0.391595936124218 1.30969317341825
1.68359095958506 7.62901889841519
1.55014341230999 6.35669199872224
-0.461713591692503 0.760474080303689
0.837995109576847 2.25388599012062
0.0647111588263334 1.03231598849297
0.449501515148342 1.39080169134864
0.842450703096236 2.25731626328023
-0.662607691037783 0.514319433635468
-1.73347656811602 -4.21427270283707
1.97125407913951 11.0374500899417
1.06434018406928 3.1398766859045
1.24774144636567 4.10652488237163
0.301552235654667 1.22610695051265
-1.419054264094 -1.96851225748391
0.770735380506118 2.05529699022521
0.673663615667701 1.79826450863936
-1.20717318186629 -0.921117900268635
-0.410538504386376 0.765068471502614
-0.997783575545104 -0.190561556633821
-1.9247004506919 -6.06488872678494
1.67307707998314 7.51858011026587
1.13094409339522 3.46594756951591
1.05978614977254 3.1119281304365
0.368497683846726 1.27437039746974
1.19041117219778 3.77115160864485
1.53532876144759 6.21911480195912
0.977879415140715 2.75127965038619
-0.661130676558663 0.515515881624548
0.372560417149814 1.28014251982014
-1.41713935044321 -1.96759316741115
-0.721990537872462 0.427308795775999
-0.593304843043549 0.610478409782568
1.39782862521232 5.11259491405215
1.31945403787629 4.55774479220443
1.29797173081997 4.42915493648867
0.500221801881652 1.46197989702367
-1.22858048710803 -1.00927000019683
1.84901594258915 9.47648534922581
1.5150131259747 6.05407785409292
-1.22889583462019 -1.02559434722004
-1.4115062222231 -1.9348610640109
-1.33406864789987 -1.49095685339282
0.212784793779644 1.13119633500023
-1.35572868353146 -1.63622217917435
1.34923257581794 4.77328323295919
-0.117934710702453 0.944959877079549
1.74259492028609 8.25448605318237
1.79584748839092 8.8397486108736
-1.56683538172848 -2.9161474951396
-0.100670899179379 0.954539281248498
1.3130637409383 4.52476010808466
0.413802060319139 1.33744361353339
1.03952918317515 3.01003625023142
0.355758361430128 1.26871554742047
1.06587471434403 3.13952285378779
1.47827077764151 5.73417393797497
-1.423248862284 -2.00477990059306
-1.64138318843414 -3.44920568469485
0.858985773610331 2.31055719751639
0.30940079900652 1.21678960560482
-1.59532819132151 -3.11544889268605
0.674828392358402 1.79411719788582
-1.99791465316728 -6.92318765043389
-0.983915817371269 -0.150434192007041
1.05210422602657 3.08040038211607
-0.460520520996104 0.753445163781151
0.793577616390077 2.11102106409086
-0.916701943346689 0.0293686277642946
1.80568436851851 8.9740041455143
-0.41905914986584 0.775491286151664
1.4905136786334 5.85347534864318
0.659003889297849 1.76823560366373
1.33264972737708 4.64719307545267
0.367019067455185 1.27683076268824
0.284838881991474 1.19735122698525
-1.56586227893797 -2.91574117100978
0.865300493711527 2.33095887814049
-0.939391307968867 -0.027203254178731
-1.88319560164792 -5.6496600168737
-0.525375767528593 0.684583616019016
0.496152544580833 1.45177679093498
0.386055497083448 1.29164860153159
0.351299629919691 1.26288874499817
-1.65598010248621 -3.56432636874693
0.751560061269861 1.99687186057058
0.284661315211381 1.19235145843037
-1.09215628143104 -0.491573300112702
-1.10348826232006 -0.524640770030218
-1.00871955963086 -0.225586793602744
-1.54583226518551 -2.74904177692399
-1.55838780436906 -2.85682837025492
-1.03843932105837 -0.325694840270278
-0.946724153501291 -0.06049006215608
-0.654071049279125 0.525404355982477
-1.90500466748741 -5.86320237196585
-1.27255094354131 -1.20661238758745
-1.98026347546296 -6.67989024206411
1.63566363171233 7.14486566814899
0.0339713196359335 1.0093287980423
-0.50459247118296 0.703641359276434
1.81559893738786 9.07136373457222
-1.63792963751805 -3.43811065019935
1.81879951740473 9.11907419867196
1.78184093727109 8.68248370406747
-1.7721536426822 -4.55105725684285
-1.18740832757247 -0.855914984597267
1.38038656570182 4.99110878004799
-1.38783512640697 -1.78378026045502
0.222412885401944 1.12854784914632
0.756559645781282 2.00559445533501
-1.7219989420926 -4.1241360273851
-0.471086587199301 0.734105533462656
-0.227738431702199 0.89333543925946
-1.92935794945476 -6.12418892694745
1.1844290411186 3.73565569764499
-1.71501087455885 -4.0376542850963
-0.385903061117554 0.806071659467239
-0.890520314626398 0.0740593743595064
1.67287119794209 7.51115796744165
0.83918431164092 2.25650874404071
1.90289877922406 10.132856650788
1.03833740641711 3.01299290037405
0.789709020700827 2.10355162038978
0.683611022039177 1.82285368297838
-0.179507048827252 0.925081437742517
-0.000584604517955611 0.984340468812395
-0.520405730169149 0.693339417279343
1.92715956982417 10.4668434718492
0.400860795137212 1.32987293072187
1.19890765610203 3.81451374961079
-1.54460728333877 -2.74787026367663
-1.06248141493513 -0.406999040383797
-1.28672179531481 -1.26478539543189
1.25307005082172 4.1424154599753
1.42377930122459 5.31239619932744
1.70117716962273 7.82129096593986
0.179317666451924 1.12103935376769
-0.684304672411856 0.492541754858381
-0.868768205503995 0.139903632995063
-1.51891316821996 -2.57567499366925
1.36361253017276 4.86933885074933
0.0562416104519134 1.02593905383052
0.0758099949952324 1.03392765224946
-1.58113161808202 -3.01793810810587
-0.538980074292445 0.677534726364365
0.99270343344348 2.8231022276446
-0.409786318725182 0.789398428413281
-0.281142307585388 0.868639961816714
-1.69745666265164 -3.92386559000922
-0.370783589091509 0.804295606864759
-0.230187126417397 0.901809719586254
0.454070535352471 1.39744018922443
-1.49888350765559 -2.45181224972991
0.494126089847875 1.44129943661414
0.609683100231467 1.65241912894958
-0.416834592236302 0.790211066977293
-1.32562012581223 -1.46790146868228
-0.535019952703819 0.663209586913815
1.46284070470964 5.61791567029054
1.79853558911958 8.87579227219735
0.543964275193791 1.54269241151363
-1.12101167563638 -0.587903117238959
0.634585312093048 1.71296042187226
-1.42338564384549 -1.99995903396004
1.34323872565678 4.74249413175406
0.740479740191689 1.96380089060307
0.530420671662911 1.51309076141407
-0.652511983931383 0.521107180103227
-1.48154716489781 -2.34954351397374
-1.51975770883372 -2.59116390660472
-0.357623802961932 0.83381141128213
-0.873939423227781 0.125932421557286
0.819777580708701 2.1998943411007
-0.771908094224033 0.330724932627523
-1.25533466724185 -1.14123378490591
1.38048627129223 4.99804329532372
-1.41426900723107 -1.92000749326875
-0.931515071393668 -0.0230607378401753
-1.83612664651561 -5.19848319937411
1.91787399816632 10.3457425819134
0.822087947034545 2.20295568925398
1.58408998110077 6.67853165613535
-1.47666690758891 -2.32386558025554
1.59307040077941 6.74324021433242
1.82604877768948 9.19739450525777
-0.0643848869724568 0.961139915466772
-0.170247791580812 0.930901724014497
0.576780822247429 1.59740172217391
1.23866843250204 4.06837261252609
-1.59773061856609 -3.13496249784628
0.779140264391865 2.07272054088885
0.122981305271473 1.07007161647795
1.16018101227257 3.62053567845117
1.41737534405674 5.2489009736737
-0.125273475886996 0.945827754840978
-0.0752832772171206 0.969968969047981
1.60062278839547 6.80197787067499
0.672159440578955 1.78812444699988
1.74816251900159 8.29682205654595
1.52974515968788 6.17933789924982
0.818853077332124 2.19652244374451
1.17010851194556 3.67232407707197
0.669500173210326 1.77942747846152
-1.3832662134864 -1.7597563579927
-1.5056123093124 -2.4761572617853
0.676128404865428 1.80489610197411
0.354286326377876 1.26228019772993
0.110812810663127 1.05287390989139
1.85823905265138 9.59712772884433
-0.296690209866944 0.872379305015782
-1.87852985180451 -5.58878670482915
-0.561567284182094 0.647239410899878
-1.34797827041401 -1.58171359060356
-1.75450255430753 -4.40024872280045
-1.65803450911699 -3.59862420326409
0.899307716725425 2.45310985507215
-1.46062397932672 -2.20590832094584
-1.35397611455221 -1.61410551820944
-1.55462920986326 -2.83631479759145
-0.325694520641469 0.844135792066377
0.160385508824723 1.09385491556522
-0.0461815331773878 0.972733177334569
-1.99131143951082 -6.82898935026358
-0.974657684229868 -0.126816225683027
-1.1961708928454 -0.871963915385551
-1.08638984261912 -0.46817607178617
1.87310279781129 9.77910038124017
1.91185810187506 10.2346177433489
-1.14125893032288 -0.651143612967894
0.528322311654728 1.50440518127949
1.72555442756985 8.07019342203641
1.49516179008768 5.88943360017781
1.76908731382206 8.54319474425647
-0.98536492396836 -0.15700092510946
0.822905514907454 2.20331143461938
-1.9175384908525 -5.98579056838029
-0.634376195176915 0.557193522846342
0.803957485520665 2.1390454266119
1.8795561718596 9.83452954613513
0.247638382220518 1.15926441782563
-1.71883857673055 -4.08546459916181
0.391903377619543 1.31389415981742
1.03697635833877 2.99252053203022
0.46546627101158 1.40460445522065
-0.872255355483843 0.12452919017408
1.59963770381937 6.79624745217679
1.0905089794862 3.25222839325567
1.48364736751732 5.77669503094874
0.552152149540717 1.54712047919933
1.16305110024444 3.62821178894314
-0.764162395575607 0.353582332679796
-1.7996558729966 -4.79060711409738
0.257994234634654 1.159234494009
0.975819193435273 2.74312471916832
-0.374068096434548 0.811941703820846
0.732361402874449 1.94289288266549
0.841017152420385 2.25753945477373
-0.317325089675505 0.852558721839276
-1.14729985819558 -0.679405812472515
-0.84115142375628 0.201590836168429
-0.195889906726145 0.914464415551292
1.2198790586966 3.94647009708091
0.153095278361347 1.075407139745
-0.721646086143089 0.418113274550875
-0.276444233603721 0.874891330907061
0.997336602107659 2.83914267877571
-0.619971894263125 0.570644133479019
-0.987367251982856 -0.162659397984739
-0.854665346688174 0.163321947347233
0.191712135688579 1.11008010260121
-0.574488728271056 0.632500385245556
0.00863800784050106 1.00372572314337
0.307398788784521 1.21625487746565
0.342728984756977 1.25802380214618
-1.68981021594347 -3.85635974658295
0.784208901049624 2.08772138448435
-0.208369439125148 0.913988845682418
0.526648200820822 1.49948454846173
1.24202756976366 4.07702962445207
-0.727519971518694 0.416619346588576
-0.222858774302506 0.899970690899377
1.42817273628317 5.34391409346648
1.38769399086101 5.04512706115714
-0.41333614984213 0.781506781386756
1.30000904085596 4.43810676743973
-0.0370187504556303 0.973900549312141
-0.381829049097851 0.8006227077616
1.29945975418766 4.43145515107368
1.69798018089684 7.76034437975662
-0.978298825314186 -0.12745101956886
-1.66927279889651 -3.67091481981135
0.606203457135684 1.65452552319128
0.211301964437266 1.12170481513714
-1.00052625391862 -0.198810328948178
1.19345373102677 3.78511354553779
1.57544991206892 6.57360192278671
-0.77491023651087 0.332263920464899
0.159718334398355 1.08833688203287
1.45665398404387 5.56727126547789
-0.499484712478448 0.701589152523506
-1.16220758413343 -0.752574336865251
-1.02136658157845 -0.258844629313983
-1.06674862945718 -0.392686634605278
-0.744757010347744 0.379338217369195
0.623435205569478 1.68225418863393
-0.0599727657251976 0.964843293712371
-1.6688985298492 -3.67191162329495
0.00795195473762167 1.00095435444599
-1.58371970380259 -3.04223877021898
-0.814488977715953 0.25400468163354
1.62784684311749 7.06018606462141
1.10659216693293 3.33647695876888
0.746100055298261 1.97888290220787
0.004435567456055 0.993685051475752
0.273048244481866 1.17750929462584
-0.674494714349047 0.491965498520511
-0.630136777628819 0.559016438722326
0.404649191564105 1.32802772062289
0.867192246641215 2.34518405432025
-1.53462890319051 -2.69450269874947
-1.65797847179726 -3.61466929551865
0.936206661941811 2.59443165886383
-1.78249170379428 -4.64964642499287
-0.989577984804854 -0.170243909850942
1.74689924571509 8.30477247920927
1.36878773996584 4.91833067708624
1.85181571111122 9.50239345194755
-1.74027129194973 -4.26582547182024
1.32195369830832 4.57287695579038
0.779295097985163 2.08547213848138
0.527134401728108 1.49717159191067
-0.483385457457882 0.721013992566675
-1.98054469090637 -6.67787951840296
1.19919838050934 3.83794979312015
0.247147941406696 1.15388799111552
0.0973894309066896 1.05036888897258
-1.41470298624829 -1.91867719201216
1.22947358771133 4.00233387501312
-1.94803996825382 -6.34115402871578
-1.37522052503804 -1.70312718709698
-1.69508729758743 -3.90422830135512
-1.26400612715092 -1.17882495195713
0.177514997367698 1.10070307970624
-0.669640851903173 0.506276290801674
0.67112924906392 1.79045933660594
1.50698404505808 5.98164127150539
0.0503843989658583 1.00362100010182
-0.479836536391115 0.717791570047364
1.90580750567755 10.180201922895
-1.95045988746964 -6.37125640924513
-1.23772859446057 -1.04768655971777
-1.2491992534694 -1.09917023861416
1.43020405218296 5.35155575358987
-1.7958406700768 -4.77516846470591
0.6793613225212 1.8200682595746
-1.46040150163991 -2.20800476249903
1.73771962842721 8.21560467911202
-0.84730566648012 0.181473749770297
-1.00894166771783 -0.226811317886355
-1.00859420626261 -0.222584462396821
0.958638844017508 2.67815991960788
1.95044573476592 10.766490129105
0.0346133021748543 1.00994873094911
-1.05062638029796 -0.363128077667757
-0.0934869135279128 0.949219762303069
0.210153076911402 1.12020283425862
1.06304622638943 3.11894277339172
-1.4913720899944 -2.39900388339497
1.60616168357588 6.85479758268493
0.330881009604624 1.23690985922366
-1.41434542654993 -1.96193001839419
0.54827261775172 1.53292166739015
-1.05020516352889 -0.355253516364615
1.99084203167489 11.2963039571817
0.284498672856374 1.18803669814785
1.20635571406931 3.86687589245562
1.92884887664005 10.4605744830663
-0.787836212255184 0.299734103526356
-0.0665370859942191 0.958109369188055
0.266365751071064 1.17365638382964
0.284821733614269 1.19864673269428
1.79658899705357 8.84717765083633
1.36800009261001 4.90984160895382
1.71913709666286 7.99583556443286
-1.64691824710695 -3.47709331465477
-0.634577139004624 0.55915870436024
0.907611460044221 2.48490211512442
-1.31752746635936 -1.43220825506977
-1.25799147179641 -1.15868706829474
0.164578863753904 1.09458335531589
1.5724728529926 6.56004765078573
-1.61334505165664 -3.25225593055653
-0.884101245130284 0.0929428769915395
1.0785857201919 3.20491586289296
-1.93700500597912 -6.19805198060372
1.57399805208916 6.56459550935528
-1.71744427363904 -4.0708018208458
-1.05808128539432 -0.384282057908827
0.335865050135636 1.24121539994481
1.67249976327608 7.53111403985285
1.60127860974401 6.80875163915128
-0.517147816897207 0.69333488423439
0.641324841192098 1.72460209308304
0.676786295025203 1.80733330661154
-1.82794401593092 -5.09710000332624
-0.283331862331161 0.87196176011591
-0.0639641528388717 0.966729126839163
-0.506504957395847 0.707835320922324
-0.916554913203721 0.0288747985266691
-1.74864875263597 -4.32050250007612
-1.22417132700887 -0.990462261382918
-0.480427977583389 0.737086525022445
0.522720217749657 1.49225304695111
1.79977796292236 8.87678850979209
0.821654982549636 2.2002667735956
0.705451304778161 1.87003509143419
1.91752469627165 10.3449889069322
0.874637643884115 2.38294790188029
1.51463407476122 6.03663515683795
0.151964456223034 1.08602328208579
1.04920814119334 3.068717656264
-1.40228679868784 -1.8703882159323
-1.77433062051123 -4.56898824875179
-0.0807199561361247 0.950590991918523
-0.554066225390007 0.676013409967546
0.206793034906914 1.12162283249017
1.75495320295264 8.37161770096163
-0.58976643547149 0.6122332590153
-0.570265564796054 0.632088878798852
-1.95652629667675 -6.43236712937771
-1.4809655079249 -2.34962300810275
1.24656860874358 4.09436510331362
1.04805738906209 3.06644566689696
0.245998012934822 1.14616974570374
0.712250733621294 1.8985145336873
-0.178999679781882 0.937689573199967
-0.476724521976923 0.732967579431614
1.57873780306787 6.60316766960317
0.492905548214301 1.44673427787968
1.96692074715421 10.9935413794022
-0.866109375642834 0.137806550330372
-0.902579722541617 0.0533495254779157
-0.4836989537125 0.721480405653948
-1.13860622829876 -0.65812494496624
-1.68674389282483 -3.8269478035141
-0.696380684896323 0.462339397493708
-0.474718614621026 0.746846022399951
-0.55910870278528 0.649211265691872
1.82989710207585 9.26165707381093
-0.747639397853016 0.384910459455232
-1.21735597090758 -0.980617456389869
0.370162748415718 1.29014014231753
1.52374645718827 6.10901973485953
-0.973366382246718 -0.131009116090045
1.16171016526292 3.61697168798142
-0.862987828535092 0.161078907132451
0.621063996734559 1.67901339083638
1.65446335556188 7.34017160890469
1.77094426801913 8.54839534495773
-1.94388635474945 -6.28311770959015
0.761200781796524 2.0338355887882
0.0561172373005312 1.0213583364538
-1.06957567386867 -0.411001296441611
-0.0905481850787533 0.959183499127584
0.927040219844024 2.55424325573772
1.8673231985492 9.70422640604761
0.192108321913312 1.10176166062749
0.513399284106047 1.48857011176581
-1.94821049397743 -6.31319466428297
-1.1214590808955 -0.592631271826917
1.87606393659345 9.80087088545134
1.84450991390793 9.4353368625487
-0.049463327121563 0.977076978786467
-1.18160442071139 -0.836295878091666
0.304089057238058 1.20940477365912
-1.59280051013905 -3.09419187436563
-0.949890380048008 -0.0644734829889833
-1.55324124736082 -2.83858147926448
0.447202411694953 1.38807365079743
-1.34200036036863 -1.5262171878444
0.973498301084059 2.75044470769863
0.177052608250697 1.10709754591426
0.0405965700044342 1.01925157212804
-0.752339810673941 0.37289273972447
0.168360846647789 1.0924382059155
-0.147385650667953 0.935925901041053
-0.0765706138124309 0.966310494501428
-0.834073223482545 0.189431722979178
1.77400258345304 8.58815547506922
-1.10923369352111 -0.537029360030512
1.47731730787747 5.73058332341255
-1.97426591551265 -6.65678088227076
-0.46465919100746 0.743422947401185
0.255453290103878 1.18141978600141
0.401946609783347 1.33135091302226
0.7064603645273 1.87840239998449
0.195040572032665 1.11690200145903
0.492725116630803 1.45754533122342
-0.703358333592255 0.454883515873727
-1.77071940640349 -4.54100131963547
0.0263778409957616 1.0146304100638
1.5817323511687 6.64653301342941
0.188973280062564 1.12135526154422
1.00105500765951 2.85977988013338
0.665601496451699 1.78226977757198
-1.0758927035242 -0.439445577528765
-0.628467864929659 0.552521016834506
-0.462845639931512 0.743649977634857
0.875025914541497 2.37374222271564
-1.41383683667054 -1.94667718809135
-1.6697111938579 -3.68486173375494
1.09603549914051 3.28567858678585
1.62489834081238 7.03690114243045
0.416808458329184 1.33498770793675
-1.91346899374008 -5.96256237772706
0.48890415351061 1.44286865771006
-1.82702980036481 -5.05551592861472
-0.271371741436216 0.870328942223097
0.0838712538685789 1.04189708615533
1.24655440116632 4.10024175055629
-1.46183141783321 -2.22374166989591
1.28908473130799 4.36630170768102
-1.10855990085608 -0.562063462639219
0.933803589297539 2.58521487037548
0.744090662949441 1.9707947494304
-1.18542851197 -0.840491037489215
1.09662847195115 3.28482021885626
-0.891904969622721 0.0831156169344429
-0.753767622372932 0.369836018957886
-0.419699453159715 0.788081674304745
1.92783611898138 10.4584854988399
1.18434629483526 3.74481770501953
-0.850048612766949 0.176560415789933
0.308458399857949 1.21462782639147
-0.104193997970978 0.946005898372418
-0.799972736430903 0.285927701009198
0.452767708690214 1.39766719368028
0.597412080517727 1.64373993259383
0.593059513923374 1.62427200463893
-1.40250057451328 -1.8789264057244
0.798607175412876 2.13331105315066
1.93257294274525 10.5355389630815
1.99854849151513 11.4244399794226
-0.641668732087585 0.56475366966752
1.30016400132094 4.43814008524292
-1.5326104928455 -2.67493541839968
1.37502262423487 4.94986951331791
0.693800378736968 1.8515716644279
1.84517395726995 9.43599452837678
1.92798071190122 10.4753358247517
-1.40772643694339 -1.88214066370521
-1.28601120370353 -1.26647654672643
-1.26627685651145 -1.18784224817778
1.80731301646996 8.98912832188046
0.90850383276711 2.4863705914669
0.0979589637191518 1.05077299541241
-1.49663558808682 -2.44076085587208
-0.912437419464073 0.0290584337091611
0.877357979495537 2.38865486323248
0.166820627653863 1.09218585818911
-1.77159986131449 -4.5633205102713
1.30437703030859 4.47265186631175
-1.3499326931183 -1.60225917189635
-0.596182841232093 0.608163392297305
-1.87620318613111 -5.58351911620609
1.39855464279318 5.11359242740474
1.75283927641585 8.36942259144701
1.73162706612773 8.13826599634615
-0.766719221853512 0.340864384170174
-1.92279238876334 -6.02321863912057
-0.69318058382554 0.479668680603467
0.00831948249165038 0.999751590490404
1.48305548603697 5.76142487230455
0.153118246334721 1.08972946434602
0.58443191408571 1.61378691004878
0.236693925468985 1.14339170496782
0.646310871944384 1.74030084845718
-0.398390991704618 0.808366816187522
1.62011872039472 6.98953178832817
0.163170880205424 1.09101052685306
-1.05611175254839 -0.360223134163143
-0.450061164749038 0.760914089150395
1.08455697511961 3.23431779839021
1.07031030216604 3.15892859213176
-1.13369677432869 -0.640947173373567
0.954869425914125 2.6738604387378
0.190796234602346 1.09960866645781
0.34599075164786 1.25596865017729
-0.0417235161195215 0.971710379783804
1.6996522021834 7.79550916409845
-1.98068498864511 -6.68203605801189
-1.9115004631109 -5.93364414254743
1.53382768152136 6.22353879720709
1.31412729272728 4.51961993490372
-0.685326460329759 0.47418665950724
1.97256061584938 11.0547994580786
0.520690061019052 1.49111868744377
0.431498523136836 1.35789622587648
-1.54362657653036 -2.75460998020516
1.30076943936791 4.43975314877167
0.110705357748875 1.0486479677216
0.0688186092254197 1.02119784489148
1.67413450951731 7.53341721455009
-0.985207968505172 -0.160684540430415
-0.806089795979227 0.284444176437579
-1.85237667046108 -5.33233210062879
-0.920371820742431 0.0164668433527722
1.38141443382448 4.98863015282167
1.4044209338851 5.16669611818724
-0.542691600248609 0.675735850532089
1.53677101777 6.23346451860097
-1.85850522725861 -5.37882491327555
-0.444161392146742 0.763771158560042
-0.688501376551688 0.478927137327461
-0.593839572757993 0.606594772417755
-0.47441079581978 0.741345605082063
1.8831706950332 9.91285698162258
1.09169014263716 3.2517433215441
1.25389927608542 4.15534414662229
-1.59454200048405 -3.10600040525396
-0.442711305472325 0.762550665143076
1.95806353082891 10.8724704838653
-1.78460742639196 -4.68605248225971
0.40343560525519 1.3213194545941
0.652220564362389 1.7449551608856
-0.380836178645997 0.813053290144024
0.301719509164299 1.20512017543195
-1.21097150035761 -0.945192574788016
-0.029223722261007 0.986101783902391
-0.0358890414286468 0.967455226587691
-1.73775276817218 -4.2353480015459
-0.974754376510504 -0.1267811188567
1.44203622854204 5.45062670730321
-1.59890642565992 -3.14658288961961
0.795623049507605 2.12756096210532
-0.150922502332556 0.928688080958424
-1.53940041912788 -2.71723181405971
0.0487454727564445 1.01873554029746
1.96324448991474 10.9345743710465
1.66781793345191 7.45481886645438
1.73884391761524 8.21231629267577
-1.92817622064621 -6.11153730590321
-0.45376515808234 0.739633548853883
1.56230883155345 6.45409691855904
0.54828212047107 1.54415654792256
0.442437873944731 1.37678215206925
0.192524261196695 1.11685238085851
-1.18547716382506 -0.833172875923745
0.131479945441129 1.07644065677529
-1.09917176544636 -0.529014837658806
-0.470282670431948 0.730487027283281
-0.398534332506407 0.809759655887622
-0.425917514759681 0.774729820255762
-1.20315602945709 -0.901621628900615
-0.658854777674013 0.520887961468839
1.90279869634687 10.1535144403424
0.692627430585937 1.85850380394309
1.75702419283417 8.40248244830576
-0.218196717797614 0.908040475854448
-1.78916450307182 -4.7072198549223
0.129070596408469 1.07891023355875
0.605038800031316 1.65138168952477
-0.296316256804238 0.867483154147729
0.0164951209081132 0.988124447555131
1.21346412797034 3.90729131164088
0.574011131001371 1.59115392326239
0.462290671416839 1.40614477192918
0.40410819275721 1.32633180954348
-0.481595493587628 0.733282230587212
-1.31782552436269 -1.43093933687504
-1.5194837321693 -2.58315848179824
-1.62116265544403 -3.30344496903047
1.26310081801528 4.20151526945381
0.127202279000513 1.07081464024437
-1.68520250790848 -3.77793523091417
-0.603284843250533 0.581234017543621
-1.71673296848097 -4.07683475528603
-0.532651132542052 0.672261847207747
0.261358540370296 1.18512451859773
0.558581812877678 1.55736841573768
0.666199452291657 1.77515665700216
-1.00171427776995 -0.204103549526015
0.216887403452102 1.12488879547923
-1.66491700448442 -3.63645953599789
1.62371233364995 7.03647864252876
1.89298814563213 10.0075273939971
1.99696540172668 11.3985956317353
1.49857458851087 5.91810351458466
0.0259271242716275 0.99490040667827
1.25066053839307 4.12725819342969
0.741554342666891 1.97024401050884
-0.211134376905973 0.904265588112284
-0.21293296225532 0.902208372935693
1.17443099534439 3.68928772174936
-1.81670082980822 -4.96647251739284
0.236732509242013 1.13315323922089
0.643870709298875 1.73772272122517
-1.72227657824281 -4.13521216446846
1.63583921447564 7.14238562238487
-1.68186374417258 -3.78531224191223
-1.48826071885157 -2.37873968411393
-1.49310781433945 -2.39855010937928
-1.10696078276606 -0.543277142424544
1.12452649312155 3.42889377843165
0.783650456241737 2.09083865222502
-0.161911748499459 0.931137194158423
0.628007684776644 1.68442530419979
-0.243924773625153 0.893116388521999
0.834772814765674 2.2351337827488
-0.223841541824117 0.900516913726931
1.07132477350497 3.16306477307216
-1.18266111651542 -0.832822823434673
-0.428500086152374 0.778049775291733
0.90169194696583 2.46648254129569
1.09745444799678 3.28796605696869
0.523909850278403 1.50151748649013
1.73717403201051 8.19682065367904
-1.22113069512167 -0.995120760996621
-1.20955849396782 -0.953255169415617
-0.399486529556002 0.808781879310853
1.45060068843825 5.53087848795961
0.495836028886981 1.44879256164344
-1.54563926520025 -2.76119064927738
0.664636176471893 1.78163979951814
-0.84536772623423 0.195978403180779
1.84784612583578 9.444972752919
-0.102234273221019 0.960155648523784
-0.92791030503047 -0.00201279096713414
-1.58862031670824 -3.09677838620582
-1.40597534148641 -1.90921472245161
-0.0567426886504618 0.967309510794382
1.95112557211344 10.7836771757371
-1.16202891862751 -0.743642158251616
1.21626034128942 3.92421811806867
-0.80452779989253 0.26609097967275
1.92498870321335 10.4303878655003
-1.92515688031637 -6.10821307405095
-0.852504867179035 0.170600669614789
1.55705189462523 6.42445316743795
-1.51788427122346 -2.55981081658546
1.3689591098897 4.90303972175912
1.75972680655051 8.43497694100864
-1.65957374583011 -3.62665673598408
0.270920890606869 1.16949282304471
-1.16885208042951 -0.763987666325113
-1.38447964389345 -1.78579113660629
0.177339566960526 1.10697129299923
0.543295442544644 1.53305438086109
1.81238014720433 9.03883120895235
1.71910415517618 7.98951188473036
-1.38574884153756 -1.80453591712449
-0.866168705486111 0.13155224998541
0.56542842101778 1.57148386121625
-0.999358479402314 -0.195291929222586
0.517148701909557 1.48294034813878
0.385827101406312 1.29470435488792
1.5126698410979 6.00707691823652
-0.847485819932245 0.181307650065231
1.37674588490367 4.96367311113197
-0.836170270341625 0.192930597010434
0.772951559632365 2.04951697969711
1.59583203319957 6.773635455285
-1.82180810974138 -5.05045903144996
-1.69300170047109 -3.85896485817725
1.44794960970744 5.4999476493892
-1.66109627169685 -3.59464056731823
-1.7573382699851 -4.43118186833144
-1.10719019711833 -0.532950161422905
-1.51256371702266 -2.55738152032505
-1.2165312811155 -0.966154532794749
0.879251925733673 2.38729259303622
-1.87054596694108 -5.51283413585285
-0.441182691897304 0.771939195941864
-0.19493987501971 0.911954170378767
-1.37728571964468 -1.72844840209457
-1.82483177421839 -5.06024739494962
-0.702010881428292 0.451843329310979
0.326122906767451 1.23172693663433
1.16851737260205 3.66289673262974
-0.214871572242574 0.902305420744069
1.31468918333912 4.5327764118568
0.62424958546635 1.68687541631648
-0.41248770407115 0.78026988517054
0.462518280013717 1.40117426976985
-0.876558973447731 0.121068561770403
-0.445309663339691 0.765122542337629
0.286003384268398 1.19420332378833
-1.38403511544209 -1.77359364826383
1.88068779520718 9.85376193718621
-1.6185680719818 -3.27523195342785
0.152180559735233 1.07754121780664
-0.109158171926152 0.94797019633471
0.24651978919582 1.14965891306799
0.628965416458113 1.69882274744483
1.91935588472357 10.3540013605604
-0.0373830886566422 0.976154099548856
-1.68054947050608 -3.76688627554973
-1.87451027850194 -5.53804599169035
-0.555079240761019 0.66428892526408
1.43463336473141 5.38729732564786
0.296471880579543 1.21214083308364
0.596403994269193 1.63471382573155
-1.9599885782736 -6.46862056778319
-1.18377169591759 -0.835764211058809
-1.06598170603066 -0.406347917054451
-0.0744541532140537 0.970544054476839
-1.76337010414204 -4.47980815903837
-1.09667441590765 -0.497958138448634
1.62826999259492 7.09346422247264
-1.66496683308716 -3.65859501527713
1.1752654100337 3.6867880267327
-0.711313673950508 0.441159206437018
0.949181662762235 2.64913485004861
0.23509675462388 1.14631030410878
-0.132282055480576 0.951380021157754
1.15162918388688 3.56590005483656
-0.825309704656912 0.235428161007987
0.0416208378466676 1.01916426816218
1.28237397513989 4.32803419372785
1.67804663468333 7.57560936511758
1.12847037380373 3.44583730827938
-0.666677668801422 0.500535466902455
0.0771935567965247 1.0251579835824
0.382430291118693 1.30125913420781
-1.75396408351992 -4.38684640189132
-1.27713484050294 -1.23728764579277
1.26371408079828 4.20747471043581
0.631521703007308 1.71097872801138
-1.9927489392525 -6.84210703404694
0.429833334690461 1.36756426971277
0.885000450618656 2.40899164257351
-0.928739515085039 -0.00112897243548949
-1.79090100863839 -4.70858495483779
-0.190402521964136 0.920291720165465
-1.02277075142579 -0.254672518218796
-1.98986414127296 -6.81967763997728
-1.56202189656636 -2.8700002969686
-0.823305383673487 0.240385903454283
-0.16535725880631 0.924942878158168
1.60375707727229 6.83113643466909
0.473228880387381 1.41554736679962
0.265718505149918 1.17273273686391
-0.0597395478923923 0.971368513214256
0.300775466668908 1.20781495684076
-1.05886021200778 -0.380369566542633
1.57753799391147 6.60566982934968
1.2678036676601 4.24054628185889
0.408969622909501 1.33434109200796
-0.802400565314632 0.289453107342089
0.797390436275073 2.11921434311487
0.269292520108907 1.17748066233802
1.09863981296441 3.30238844951017
1.66270555828155 7.39872202244838
-0.0678474745177926 0.961163352850507
1.15694791286823 3.59120166028488
-0.811165066606464 0.267983807143246
-0.948634762579282 -0.0490322431644212
0.98120715897871 2.7593947116178
-0.228724559837756 0.900312851955522
-1.37228948746877 -1.72229450435978
1.34818409339734 4.7540661109
-0.63180800074183 0.559632310925091
0.969694764923288 2.72146353859967
-0.15563669746887 0.925302312133414
0.732961867207918 1.94489665160677
0.985039768614193 2.77722119348587
0.320869401360116 1.22146660615769
1.73657869247921 8.1709769068018
-1.48569607335105 -2.35751632668528
-1.67052307021378 -3.67684260041769
1.80986213936798 9.00751847410175
0.269901033496012 1.17497245115711
1.86964035709248 9.72053894079343
-0.618116427779094 0.57948319057898
-0.559013229623111 0.640029635332599
1.86053966331359 9.62113266205902
1.83639477816247 9.32089258884026
-0.144643275194269 0.934702694488004
-0.661368634771165 0.518005380333269
-1.49996211674519 -2.47060303978315
1.61565620264765 6.9556738208655
-0.525992789448132 0.690153670706809
0.616639082134211 1.67233716321322
1.11048851688358 3.34846455523854
0.490327890360109 1.44696062751948
-1.70630821016035 -3.97942357433033
-1.08432082336125 -0.447866032025121
-0.470247615489478 0.737040127491929
-1.12330616920672 -0.592688203625215
-1.01151194857804 -0.22531529845062
0.93537143033542 2.58899557709597
-0.780955892956971 0.32494943131485
-1.64128116553575 -3.45549750121784
1.19263378023814 3.78334139440739
1.10397943127227 3.32477236366575
0.564707687955063 1.57096864135258
-1.07973399874306 -0.462843967286802
1.20149098397145 3.84041420340939
0.695933299560299 1.84328315610021
0.495176883716653 1.45111124871371
-1.67550375125215 -3.72282087079692
-0.529803028699704 0.67391141147433
-0.445856021780258 0.746872603821456
0.433812532010605 1.36821956949014
1.54202085374851 6.27014082733516
-1.32375381765984 -1.47198532477877
-1.62134922807424 -3.32157292176848
1.51933493963951 6.08477696358128
0.471822244683532 1.42329463587884
-1.7933624880284 -4.7584821763407
0.596742359496119 1.63294109206395
0.552130519421354 1.55492174251572
1.47710227591609 5.73121027778143
0.00302994761583175 0.998517527547431
-0.236203654112023 0.890720008570176
0.970165770754015 2.72487587611247
-1.34277493490018 -1.560845574621
0.206076871985257 1.1257152390431
1.46412316541014 5.63221243541157
-1.80601682524143 -4.87953733752454
-1.391266729644 -1.80715428918603
1.88631551186505 9.94610960451085
-1.10746678059245 -0.531319186183653
1.05027379222569 3.07342814811358
-0.542183740128499 0.661768181124191
0.0271656137687568 1.00594390041955
0.298087731551128 1.20094009138946
0.434215781068596 1.37296325698884
0.853518173197716 2.29721032854808
-1.02478655719921 -0.278179128954753
0.930280869478582 2.56653063252946
1.282077171204 4.33841398892212
-0.867967962223037 0.137217640904829
-1.77505278071821 -4.57906190273618
-0.485529113216811 0.70779431482014
-0.935709399072525 -0.0297072239274747
-1.7124804972844 -4.03787929313712
1.7597432806956 8.43146031674183
0.756160878053701 2.00445575272812
-0.00706127845737115 0.981964054736906
1.8189676207847 9.12539100938857
-1.43708373555605 -2.07284379552645
-0.391867214393602 0.81005629802083
-0.125622315968765 0.948908915942806
-0.227608808523419 0.901222918334211
-0.771183404944305 0.339516244786058
1.46047134279569 5.5976889022238
0.467766545173386 1.40672102796338
1.95848779282826 10.8896618409186
0.864335094856651 2.33375476079497
-1.63506328198021 -3.41060197363597
1.71568012752807 7.94562910214681
1.39272912489522 5.08961515331874
-1.22465164856736 -0.999370845561272
-1.02616643763754 -0.281109670883781
0.442734584377311 1.3636453659406
1.22412292980746 3.97718119164815
1.97421205804188 11.0926394512552
1.48145057673646 5.75859204672228
0.819496054168538 2.19628102305374
1.4980055219723 5.90529065694513
-1.18573895693203 -0.831044404668815
-1.99278274877207 -6.8494856711783
-1.91941536557922 -6.03689343835302
0.131241931655927 1.0600975339464
-0.580657766359237 0.621323641971421
-0.888088359400202 0.0879748598093739
-0.499052602078298 0.710622739621337
1.71820673011863 7.97388064367148
-0.316414951511492 0.850172933138876
0.393727886022705 1.30739732012249
1.00589706885763 2.87467650861009
0.57688395039735 1.58492314006907
-0.426697312104449 0.778529184768187
0.375189921780094 1.28553206713643
-0.813630157701015 0.272796922860551
1.2372573491537 4.04711291749065
-1.53832901940263 -2.71516049093013
1.35062857979094 4.78353830796967
-1.33694973116921 -1.5097033149936
0.564734755108878 1.58133376269604
-1.03761236511533 -0.304364625652098
-1.7347823407846 -4.19876453668963
0.945711070922103 2.63280043844309
-1.95258942984414 -6.39542582524028
0.0826143577681773 1.04656434275939
-0.511421420547976 0.700390007367014
-0.689362310384055 0.484189815360375
0.954490563853327 2.66189127695141
1.66434830933929 7.42551155767517
1.83138414083698 9.2535906483007
-1.02932672246069 -0.302454377177136
-1.04661156110905 -0.349194696608168
0.805645153296944 2.14727751708212
1.77401916557127 8.58561787796415
-1.40235998658197 -1.85626609020103
-1.0095984473455 -0.219611478121682
-0.270919141607282 0.871927628762499
1.2570941367174 4.15531025075
-1.86122822980676 -5.41679154352957
1.84661759143133 9.47138804968631
-0.90332609244693 0.0560399363281534
1.39563280875743 5.09338234894223
-1.98651168890195 -6.75881861620785
-1.9996278349795 -6.9117747602424
1.35079226518494 4.77394550653772
-1.79705482162434 -4.77476543435062
0.733839749036802 1.95167204064777
0.825728470288538 2.20865121947308
-1.999966065817 -6.9613984030445
-0.843299266844224 0.208944452179643
1.24300068592906 4.09229023697789
-0.568510673671339 0.628967296580392
1.51908286676269 6.08838913446026
-1.87867682212743 -5.56365333032191
0.51612906772595 1.48055749728735
-0.859068978962361 0.154159105057622
1.94861230872647 10.7340006651497
0.0534360058709065 1.02654421919225
1.15519390843326 3.59722956870279
-1.74527861230335 -4.29824285096523
1.66524052084331 7.43899384484359
-0.325981587735932 0.836218281066738
-0.752689738584645 0.372297506566086
0.128978317546487 1.07630978592339
-0.5621304215494 0.641110563792218
0.955494604929388 2.66243994646497
-0.555682648203415 0.657877220211347
-1.94548297516413 -6.30660150158887
-1.43026669699117 -2.03133524617525
0.356931141155441 1.26497851786683
-1.77861129362684 -4.64027561767086
0.986070692175463 2.79238643308185
-1.12135417504645 -0.583873973582881
-0.874192151028235 0.121619859150448
1.64662777642106 7.23880275344668
-1.36031739836864 -1.63023469181487
1.43962972373616 5.43679065708496
0.452585949580314 1.38020612816226
-1.76731956908791 -4.51087184778122
0.407565146203995 1.32590473134467
1.98898051125618 11.2966858096881
-0.598156498020338 0.603843415879494
-0.331884353093339 0.843716252343694
1.9650023216185 10.9636901570038
1.47465463159076 5.70579701648684
1.5204818278323 6.09842556468083
0.382297270832394 1.29907946604983
-0.308472069208821 0.86452974033168
0.464213345558603 1.40962225619966
-0.313121928977522 0.851139529144977
0.404197074317428 1.32297254888198
0.106771342877163 1.05328070645161
1.23841881398682 4.05574132595654
-0.282973759325347 0.86927530549534
-1.73525840893654 -4.24571728575114
1.20188368354561 3.84484408833895
-0.429258807295439 0.781389744582732
-0.748427956519946 0.38035750152032
0.0804517475120305 1.0282753335618
1.97951469313928 11.1614158365608
0.124479142005626 1.07152735282127
-0.68196947072659 0.478051156977546
0.223057422568356 1.12770066107423
0.469529140163063 1.42042477119935
1.81420517068728 9.06347397745196
1.09682487909601 3.29017544589656
0.914181350597536 2.51081156794732
-1.44096664107871 -2.09895744289672
0.659760682063931 1.77014541399254
1.04063674372985 3.0245021484701
-0.410649290239818 0.772472104365948
1.18214697910478 3.73390336265526
0.672531383775982 1.7818910668283
0.833061843126512 2.23739284468989
-0.617849314672589 0.576766640827125
-1.07259635030233 -0.439464162323758
0.391562842556896 1.29947440262709
-1.32972617708255 -1.48797880023631
1.70646260742689 7.86914077390867
-1.30706265571546 -1.35497932712865
0.187434893599617 1.11386759115505
-1.73820046861836 -4.24171387296479
-0.953602384085788 -0.0634957917389822
1.29875352212481 4.43489366176754
-0.103768824953191 0.943596096010008
1.78570975428274 8.72670519279494
-1.53634762343696 -2.68365510961318
-0.845841128729985 0.179914162737134
1.96130458726239 10.9217955746843
-1.38123143847418 -1.74150095436957
-1.82357045233629 -5.05015042661105
1.17694940504624 3.70540656238348
-0.810272952936271 0.261757370921051
-0.991398036985326 -0.172896205243066
-0.961061734282657 -0.115018311465007
-0.214642191506338 0.912795507245668
-1.80569826891484 -4.85454330736145
-1.26660129022973 -1.19561802966693
1.71984917692486 7.99792161372587
0.855714417828803 2.30130420773794
-1.9857480554809 -6.77474122630028
-1.72758247087356 -4.12988884755024
0.400499487781533 1.31846776525232
0.114574489690743 1.064962908841
-0.779632733148505 0.32180448146621
1.41151444990222 5.21844113975241
-0.899575225078841 0.0588786822150742
1.24093702114785 4.07803491765101
-1.67889529411105 -3.76011061763772
0.261032694419074 1.16507301696946
0.145927319726146 1.08550642394716
1.72277198899299 8.03579300253467
-0.881853643689245 0.101790168217609
0.799560413958917 2.14568536438307
-0.0893640599908605 0.968936056859639
-1.62447012296027 -3.31089685950611
-0.818707933572382 0.235670890205997
1.41743799518074 5.25552410123793
-1.2788388253118 -1.22818174476524
1.02741670432943 2.96620221224974
-1.75287345650198 -4.40168189949109
0.638185008569665 1.72221985914351
-0.566120518358522 0.650562467976798
-0.400109521832959 0.81283133976865
1.28489563036768 4.35346886871966
0.458876720195964 1.39828053433356
1.66049162432052 7.38465069349612
1.42024196877238 5.28362108847663
1.48326877753877 5.78490256315112
-0.588650003265751 0.610370712375689
-0.927367872736143 -0.0114696038192499
-1.02199237234107 -0.286769881922343
-1.81491938247115 -4.96943751358194
0.198105630543795 1.10133534625769
1.05566735595332 3.0952690904825
0.216283757072662 1.13320891762436
0.0952747408355625 1.04483947637352
1.88356421222412 9.88305662489422
-1.50949082880058 -2.53185428408598
-1.61612601851821 -3.25947685794374
-0.355798702545588 0.822434457477091
0.389974202292196 1.3071554623711
-0.346350159364043 0.837496587477261
-0.731906551736923 0.4231017017737
0.406415533261588 1.32501739654885
0.718342248995067 1.90723924704651
0.0356803131817549 0.995954228879776
0.32991610507093 1.23539495995339
1.42183905565164 5.29867881100847
1.11622051534856 3.38301214580797
1.92780392950029 10.4727861730491
0.674874856790324 1.79951472002745
-1.23425008498926 -1.05251765071292
-1.3974131057156 -1.82541063839319
-1.43099673513162 -2.02815987997362
1.82551334714146 9.17010534096805
1.79738352841659 8.8557328134358
-1.88660253039814 -5.67393300807767
0.26405973820585 1.17900471321331
-1.65724856526241 -3.59570924152451
-0.656847438749077 0.546819432088693
-1.85985870708892 -5.41022761275473
-1.92669888790178 -6.10120496681913
1.29510024203283 4.41730776155684
1.21689615980552 3.92920219065836
-1.48920215574009 -2.38781776744757
0.169083726631659 1.10055457991979
0.0312554397253666 1.0056786571507
0.117235240035499 1.05365418816086
1.61530394923782 6.94640681714602
-1.08506533136718 -0.471932555403605
-1.50655348044929 -2.51834677153378
0.542047461016566 1.52317408442638
-1.49196218005367 -2.40584651618691
-1.62953219152099 -3.38030741319258
0.653952115193265 1.75577793997946
1.16589918784997 3.64381270220462
-1.77321322818014 -4.56469465484214
-1.44063688957111 -2.09918107331092
-0.786776926887563 0.303985331212182
1.56757347418732 6.51428252945648
-1.46068373870209 -2.20864555530647
1.51617594077454 6.04677205451449
-0.576313733093636 0.634815342485373
1.34090731590764 4.71533528035314
0.521243193049925 1.50107230710231
-1.23207534268808 -1.02999184790934
-1.19751185489956 -0.896932767351459
1.42298322693977 5.31141137619722
-1.84564276286538 -5.24061024182928
-1.31186592728896 -1.38985918746882
-0.922858524377307 0.0126962294143067
1.60549991204032 6.85970352455356
1.1792532466363 3.70185784389861
-1.45845606087949 -2.18931104790498
1.97842696473836 11.1559755107588
-0.889592956516454 0.0881258849970086
-0.159391035282863 0.93137945195953
1.41663934420473 5.27030191748982
1.72248227506626 8.03246200431787
-1.76016432377812 -4.42610150387204
-1.82777538997854 -5.0940822768977
-1.78206005009434 -4.67457738753414
1.46006475348857 5.58108255384519
-0.938027545334918 -0.0241690864101531
-1.89713331915919 -5.79558084756152
-0.569747605280394 0.63397940911139
0.124837037063116 1.06821918165989
-0.684232287818634 0.491088808639237
-1.83385619125808 -5.15769800925519
-1.26330160100964 -1.17849027537144
1.37217992254188 4.93689519001865
-0.849073677961068 0.167156637959228
-0.0294605789009212 0.979995443089457
-1.75569733601588 -4.39528696677677
-1.8642781527774 -5.42629141379155
1.00355790811262 2.85991001296316
0.215009660003008 1.12504404521865
-1.32660668313371 -1.47256222141998
1.38069310867288 5.01125720875715
1.86924894822072 9.71332902292094
-0.190122966059632 0.921551851526446
-0.152887041387415 0.925827416569911
-0.0428840413177665 0.979160071851424
1.65420267720902 7.33863979494432
-0.900901126415482 0.0560309237293428
-1.90168442933663 -5.83625996430053
-0.779415341955448 0.322351078873821
-1.45067678889729 -2.15145455279476
-0.998250686285214 -0.171300019597086
1.66790312588376 7.46235873854451
0.28681978794021 1.19691243674556
-0.484523640375409 0.722736324751854
-1.6555459953985 -3.57807298722176
-1.5455072230979 -2.77803191414267
-0.688771224907164 0.483977747667041
-1.13893090991256 -0.648782171993345
1.74139892441216 8.24162070521052
1.49893617150363 5.90031251014836
-1.08765843446177 -0.487572985944805
1.7835459703997 8.69772460480078
-1.10577444808059 -0.535344053874771
-0.905368974653909 0.0580592417348249
-1.75191942613126 -4.38392739531254
1.93356751038377 10.5522770070376
0.9575879972616 2.67385742229081
-0.796119159539778 0.275863167409546
-1.17064460231169 -0.78528781265218
0.526940860518066 1.50994238785575
-1.18559700128084 -0.831772801428173
-0.0335359284307351 0.983857790419807
-0.818512156031854 0.245098220106386
1.04792454807403 3.05831772259283
0.127512311118942 1.07848873391252
1.15092609088398 3.55360510416694
1.37906280384393 4.97418955877048
-0.0473423907324682 0.96888765579321
-1.81903180933046 -4.97334268297504
-0.67902321115456 0.493726309635186
-0.798945544129008 0.283558204400877
0.598411739509126 1.63323878745168
0.301391070356309 1.20040807529484
1.47141508416494 5.68106703834634
0.883240667186142 2.40434783390999
0.156345836899162 1.0835720610823
0.749642134354526 1.98493438217823
1.49307307902606 5.86243439468387
1.0051133217067 2.87566950658169
1.22060969375902 3.94798487387617
-0.360249201898012 0.818908380812852
0.925829811942897 2.56283549067487
-1.3265763354304 -1.47290186967351
-1.79276493501864 -4.74355242164767
0.666266542292513 1.77844030769164
-0.95305747324246 -0.088946051202931
1.92203859375713 10.3913976509573
1.60206068652581 6.82197328090426
-0.424429206266169 0.787192414010418
-1.80954533229095 -4.92551946312734
1.91009723362869 10.2448499958386
1.68732742763117 7.67365132459873
-1.63516046158355 -3.40894235039875
-0.459631207281084 0.753705430502516
-0.522873646404107 0.68248796383917
1.59857006948995 6.79465092547699
1.03205890678154 2.98085490625002
1.59776386078881 6.78937914715549
1.15190579811508 3.57596356932545
-1.84684367541041 -5.26939484200659
1.26353314362285 4.2062199699076
1.35352808434662 4.81378397228656
1.40621035681587 5.18601440856323
-0.0272575569321387 0.971616125992548
-1.94627739820394 -6.30517712192766
-0.320949431762209 0.853498206307889
-1.59390277865162 -3.10199050406274
1.03232614482696 2.98732112354215
-0.288476953259873 0.869825858368562
-0.111058309972998 0.954146397701777
1.5736115516678 6.55485856481457
-0.860588998866321 0.15642658047556
1.39521178802958 5.09339707406236
0.0423370688837927 1.01865161261728
-1.91965920542692 -6.00410035946931
1.61115595534765 6.91347479932424
1.4840555846706 5.77788686376411
1.34085225582999 4.72451646530128
0.595661324885158 1.63354063840449
-0.138836632916785 0.931557022836676
0.156291343415457 1.07807508598658
-0.199381414145481 0.918076848212177
-0.559238193346088 0.649944599417361
0.0290059939260847 1.00886399855458
-0.698820135672535 0.46917212264515
-0.87237557527917 0.123745015732272
-1.58195014208742 -3.01173280754049
-1.9689231345982 -6.53962061378453
-0.410543849357809 0.774965062938156
1.3643124949043 4.88762618880428
0.823704960749702 2.20013420747358
0.674083822044872 1.78666363937089
-0.0325492464892863 0.963388674789698
1.94758130859223 10.721755324163
0.612287849542449 1.66267016705614
-0.690234865295915 0.477630742418045
-1.48019927598632 -2.33678801276316
-1.13292069294096 -0.626449522089767
-1.16648749781285 -0.747041932511638
1.3021658674702 4.43795673604559
-1.1482368174498 -0.678728219680581
-0.328263345948502 0.836237638536323
1.07923511541064 3.20608357727173
-1.02190446706387 -0.265216820748055
1.24442771631325 4.07920038204131
-1.65796184062742 -3.57106322333758
-0.846881701827113 0.178251765787283
-0.195611341148256 0.91061814545446
-1.25153965425888 -1.11718835843621
1.09206189549521 3.25655635179998
0.787016617069217 2.09829918227663
0.638861992142554 1.71771868617151
0.399893941960438 1.31866926108832
-0.761502475471771 0.351803496175628
1.0527372198015 3.08520865018572
-0.0214878931824463 0.975799084075638
1.0568998378917 3.10330396030909
-1.46873068032241 -2.27852019339702
1.54491279479151 6.31880603915764
-1.30793049787522 -1.41517139604437
0.864324612137999 2.33487841538761
0.208219021114417 1.11889714003132
1.70114170828014 7.81190050955589
1.06753378547388 3.1404300027043
-1.11001488165708 -0.566475339272693
0.900206776703087 2.45633244911888
-0.971142814603263 -0.114440229600136
-0.588370612047532 0.622154241959811
-1.36107915505973 -1.66834163500672
-0.580659726949028 0.62587338231478
-1.71050106264944 -3.98719767008242
-0.838727982174718 0.200416642319025
-0.248301588894354 0.89533448102124
-0.989899385886345 -0.170950406653217
1.45858345806112 5.58376216154754
-0.827374754255352 0.227448429646508
-1.28680681315339 -1.26625444942153
-0.560470344749795 0.65092888444907
1.34103519777274 4.72632065303982
0.854124982911279 2.30360432367002
0.727499969443455 1.93095709065338
-0.831947031058761 0.233215952035243
1.45247471287064 5.52279789355315
-1.98284281776198 -6.72926578696236
-1.03189856907482 -0.304522152519759
-0.340538292167053 0.843966930170733
1.13223942151434 3.4621697814609
-0.25423408601707 0.874988364124284
1.15671650638182 3.59695541308704
-1.67804201036224 -3.76369997152317
-0.55737849175 0.641312139375349
1.22284522480552 3.96610427024584
0.303291470771417 1.21458446684847
-1.74160573469583 -4.2921974139226
1.51465156566754 6.06014896988245
1.85431344649103 9.53042038623335
-0.958264059974216 -0.0779086575288621
-0.181367974652247 0.921003247623564
0.354703920530524 1.2681057921221
1.5957858792884 6.77614092967193
1.88579874047252 9.93153684685712
-0.708036601941509 0.459243219841501
1.23034559038518 4.00717188621409
-0.842252021191427 0.192630036329356
-1.71655053080056 -4.07657785364739
-0.881557445068461 0.0991392922692949
0.26514318781149 1.16717793776953
0.232989434111001 1.1482522835511
-1.82715605861806 -5.06253576136783
1.1497195166075 3.54793790777054
0.00757987893164103 0.997700018868247
-0.61739697037267 0.5735170991148
1.50647848392169 5.99546984833221
1.39166169440554 5.09328381826458
-1.96514402110851 -6.52215451332572
-0.838737524950258 0.194299200340564
0.138339494100933 1.06186840384019
-0.213443547284491 0.904404608021305
-1.46046481821283 -2.2080135401958
0.609831401488206 1.66300067914825
-1.26495257795061 -1.1561069134446
-1.16687188815107 -0.778955859164267
1.27510839832142 4.28933685985629
1.14154528439393 3.50920383108221
-0.778864741792832 0.327268260447748
1.63754071403956 7.15964537774718
-0.643346611279531 0.549356031200366
1.67127472441855 7.50541451384146
0.111885802770817 1.05179206964457
0.882107696941235 2.39247794919952
0.91120952438826 2.50577477205716
0.0119238114203624 1.00140676882301
-1.38083405007942 -1.77009842808448
-1.90610125215633 -5.86952783934642
0.175143291573622 1.09829986463524
-1.61207890964624 -3.24037009230353
1.5512325813554 6.35356088789581
1.65048286681893 7.28520831121552
0.188159596204922 1.11595697837751
-0.0382019512333382 0.975155447665413
1.4540122369886 5.55071321768358
0.432059746865559 1.35824805871841
1.67199448097685 7.50303519222687
0.257309056406947 1.1677196249847
1.06166139347091 3.10876862857168
1.17483889410732 3.68968549532277
-1.26118992652891 -1.16115793532479
0.382797281524637 1.30668725064353
0.889016236667162 2.42748823070005
1.8514711371672 9.50847084255472
-0.0193655465799174 0.975163024843158
0.496029176730886 1.45262652122023
0.870030643099832 2.34669385276523
-1.76505887531598 -4.51997178880659
-1.49081964038111 -2.41529704137638
1.56004730908803 6.42839628279718
-0.566720010329685 0.649859838998548
-0.000648148169025031 0.994814329710489
-1.36290341870177 -1.65567881249887
-1.93639048032525 -6.21241344520818
-0.981602974126864 -0.149832691429882
0.853315405048908 2.30571530395305
1.82818383675168 9.24487137613128
1.97604295543575 11.1219852538549
-0.545612270990488 0.667582385127273
-0.155994484131787 0.926496782955658
-0.674605850503229 0.496796399880046
0.0708290577944006 1.03235432288971
-0.71072226618099 0.445681361775277
1.99239397624334 11.3604064288272
-1.72307522658745 -4.13217941900431
-1.3528138003879 -1.62026796776499
1.34167153626317 4.71668204199445
-1.46411382512332 -2.22783314293238
0.39152830503833 1.30147245218844
-0.908375413808081 0.0288199848965107
-1.94765578353395 -6.31316002076307
0.85951902865061 2.31840828190778
0.44754279041818 1.38058500588807
1.0067674631368 2.86274815808275
0.326349266204949 1.23939686078795
-0.540902655479463 0.660366168807756
1.60714789447584 6.88435406585097
1.28467265919759 4.33791827404069
0.418555441252085 1.34600156677161
-1.19450012731273 -0.887995026790252
-1.65377266173249 -3.54407195701156
1.58540353990366 6.6678132010446
-0.944656579769672 -0.0624543370852131
-1.07655252991842 -0.441720469500251
-1.59444925894351 -3.11427511740346
1.71387417043291 7.95638094156658
-1.26811232935142 -1.19141718266164
0.729303047022139 1.93695961928968
-0.691313419934584 0.474337501633529
-0.333711135403621 0.846686242271427
0.743107835090247 1.97670809350009
1.92591919534008 10.4377933417182
0.473148974962963 1.41366645468216
0.828791786541936 2.21895247598261
0.3345891138474 1.23955509502903
1.39616687331661 5.10666305825946
-0.107275505718615 0.950290375938965
-0.666267250289023 0.505426329319135
1.09289645024831 3.26411664524251
0.765859836476295 2.03488783891601
-0.076768188464758 0.964850464205859
1.03012343761375 2.97897332767509
0.759985684719819 2.02728103831829
0.843133608500746 2.26551528180674
0.0185066691335232 0.997453045371327
0.81233786159402 2.17018711706274
-1.65239328606655 -3.5672486883531
1.13197953472704 3.45242391677667
1.4057918676896 5.17409631571188
-1.54510507422204 -2.76209138479126
-1.39250318746843 -1.81557285837887
0.951075027305254 2.65346207742809
0.552736525770104 1.55705750348945
0.75878836535122 2.01236637732636
-1.16362024732662 -0.756914679674317
-1.11217610609222 -0.564261933287231
0.29255119146397 1.20298932043435
0.479236055210636 1.42928091158792
0.496025141932312 1.44951455884393
1.36609464527703 4.88068154557778
0.648538620336751 1.74753746436698
-0.71533225144047 0.433695281884308
-0.286770268938291 0.871823398856179
-1.95954722302836 -6.44929625957437
0.0739036573349594 1.03753704458365
1.72379301526689 8.04513583149042
0.783420714537717 2.08166388865928
0.0340859451793007 1.0132308773025
-1.74428102155042 -4.32525431791812
-0.281295155032678 0.863976874564564
0.0382462130772492 1.0180676212334
0.0983137141314621 1.05255513583602
1.20290104395388 3.85108434055182
-1.29481912896917 -1.30459771930081
-0.472474670183368 0.74509676010242
0.995401823040681 2.83115558456327
1.75319062696321 8.36567906912779
1.93603985687797 10.5723349974804
-1.26379812183442 -1.18303648649257
-1.18369692778102 -0.834816459090062
-0.152466784570433 0.920038060380442
0.720518130186124 1.9218428064194
0.894330754604709 2.43343420666996
-0.366421366664984 0.819732553669731
-1.42045248163259 -1.98412367895874
-1.91774835420265 -5.9759897017323
-1.20154198012988 -0.904325167068791
-1.33402732468644 -1.49111658514942
-1.50006564953072 -2.45603628183577
-1.66573140633805 -3.65543141562523
-1.5756079616456 -2.95917698659341
-0.392833384024106 0.800043458612075
0.240577748960603 1.16030800232116
-0.331661484133019 0.843324377323934
1.53620483944148 6.22849660991294
1.65714912226404 7.36720614017756
-0.0807157883464811 0.950704651969986
-1.77898436190461 -4.62025313227013
1.67585648664319 7.54823942403045
1.52477266006196 6.1187316608363
1.91797135475423 10.3417542070811
1.04023328929125 3.02492428773039
0.152241395502356 1.07670621900785
0.267046725387243 1.18055550785525
-1.56787123692317 -2.90639158108232
1.79178200361487 8.79783347146289
-1.23832791286711 -1.05843066198729
-0.635983682522904 0.554167833607614
1.60367031537368 6.83510040594383
-1.43481455291798 -2.05228495787823
-1.34739155945832 -1.5820111799188
1.04189068051751 3.02340727837807
0.115200091567095 1.05584975421508
-1.56120307441528 -2.85405514439073
1.48105121191228 5.76149935691055
0.603419123923464 1.64265840860011
1.48273036672582 5.78463909421347
-1.85380337466648 -5.33897504761528
-1.57355092300985 -2.95677414336917
1.47731221536187 5.7269004017349
1.07094623968228 3.1601858383006
1.31681938520243 4.55115143856544
-0.187214975522651 0.926397790816692
-1.76804267009889 -4.52164572165682
-1.77431958323804 -4.56931709534463
-0.871989619000863 0.124054020531704
-1.96752937866711 -6.54008401687717
-1.38883862653797 -1.80297219382232
-0.327448169569295 0.834872236128242
-0.690183808076499 0.480177524505144
1.34261236925349 4.72281373472197
-0.234061054505653 0.909474538769884
-1.85406584301217 -5.34226013439623
-1.11703539684548 -0.580381249126704
-1.67582335398275 -3.71999950522554
0.620588046902504 1.67816478821467
0.0395400978842821 1.01551799149946
-0.381418696904148 0.817207506981972
0.323309013266111 1.23600467523781
1.87142246633171 9.74870087342748
0.914983635785953 2.51724871692523
-1.98702218833896 -6.7279904564761
1.23668201082246 4.04673891861823
-1.61541830725903 -3.26578016768269
-0.536681404637868 0.6646827534073
1.03184549074563 2.97977288547089
-1.43017253565283 -2.04495330253704
1.17441696521193 3.68646265301762
-1.80765575659796 -4.87649258477049
0.474712397966165 1.42047152054773
-1.34420669032224 -1.56172252528927
-1.49750533480565 -2.44561101717458
0.513620072846923 1.48015025720336
-0.181096461189659 0.916883890875462
-1.40583515211248 -1.87647974602402
-1.06265554006933 -0.390387309735581
1.65282905824244 7.29493166451245
1.66613735323944 7.43914211227543
0.382870698033663 1.29193673507371
-0.590190346697981 0.616899363724993
0.674235447856505 1.80008301025492
1.10470625697213 3.32133411715915
0.298221482999279 1.20649053030239
-1.77844578663561 -4.61137131191489
-0.251845902583502 0.893052687740206
-0.526362363048593 0.677052334110852
-1.9285328629071 -6.13142398806669
0.900312461603144 2.45765483573309
-1.16068571417867 -0.725991110077104
0.788977412787463 2.10122441798541
-0.618023494494006 0.590162006628996
0.654998547681636 1.75288201335859
1.18336330269809 3.73071029040782
1.31524089609012 4.53769166491941
0.536173360459293 1.51739394892642
-1.73823963116706 -4.2535981564127
-0.937413290639119 -0.0188526248895549
0.90753282930239 2.47685833734623
0.931842474996466 2.57309961283514
0.83486413335356 2.22652023320826
-1.94351384966993 -6.29627520712077
0.490866577973689 1.4484370249403
0.0556246733436576 1.0154465907918
0.430478704730366 1.36158879573413
0.328921568783732 1.23664315263756
-0.0505001537468308 0.964305599299581
-0.701819536199354 0.449999116700887
-0.788744350627737 0.300664730132906
-0.883644406730378 0.0877547887264485
0.777204138795909 2.08098907305371
0.404023384471348 1.33036729637145
1.76799442449208 8.52326451192898
-0.5371473705932 0.668579714078399
1.43450835051504 5.39236754399515
1.59992625107588 6.80108342442818
-0.0145317781918641 0.973809921922492
-0.858960912599482 0.161889779176557
1.64230744253453 7.22339541450917
1.66998788212528 7.4942110818753
1.68929584078918 7.69716483625987
-0.543860238380735 0.666641322435849
0.256868355219774 1.16414014155318
0.376161311061836 1.27664840274201
0.921408393990831 2.53631401686555
0.325325670266654 1.21472848356356
-0.973116046489098 -0.126730500007163
-0.132667533932598 0.931236035795649
-0.810451709626184 0.267886963544335
-1.06430194672865 -0.405383782335702
-1.23502125439532 -1.0397945704866
0.616029416380815 1.67516031035656
1.78745526782503 8.76227787418742
1.27931288916847 4.29798436573562
-1.96119267589845 -6.46255441288198
0.918509441967202 2.5320154212315
-0.306940973882254 0.858326142702147
0.248852555225668 1.16555921034989
1.35316821236359 4.79508743215688
0.371451472248911 1.28008006402277
-0.0794440959535696 0.956005125018601
-0.341680082781666 0.844119094518881
-0.955971468985779 -0.0852165990259612
0.230866145284785 1.14945751241327
-0.0534421397766298 0.969914319455691
-1.43957202939354 -2.09791697219914
0.882489151516282 2.40664510197794
-0.285313963934104 0.880573605380485
0.505935499872656 1.46343379814224
0.128261451177393 1.07816194692088
-0.586429144990519 0.629530107770523
1.23866600370325 4.06559185397949
0.608518453790227 1.65703005032368
-1.35766524022517 -1.63538355867661
-0.37282628974832 0.799384448782269
-0.419963720245379 0.791005201947998
1.8799898503159 9.87557110974731
1.61158609663821 6.91480338354218
0.664953488897894 1.78236988607563
0.0381909020082576 1.0174443363817
-0.761353550774154 0.350851006317331
-0.938026164897382 -0.024747382537045
-1.92236700880279 -6.08004135770381
-1.07957429545738 -0.462580093288629
-1.34622173676857 -1.57747081911211
-1.85871116212759 -5.3772350007477
-1.22803534582444 -1.02576141340772
0.953203011928426 2.66258882259319
-0.0857827511584714 0.960057081465361
0.650724796933962 1.75055291964157
-0.689073718481559 0.479183477617647
-0.584716312051022 0.620072841706773
0.894474614097228 2.44429276115889
0.79865611645924 2.13584353133673
1.77078759706236 8.5596311596453
-1.03304944262268 -0.296886437274474
-1.08241491519637 -0.46778174459126
0.133838570248027 1.06426121203456
0.704280169796623 1.86488251242162
-1.45523923732272 -2.1786120769966
1.86704022144864 9.70734519076174
0.576394656998143 1.59523376434528
0.33815543603102 1.24583696334757
-1.70620039544918 -3.97625834036472
1.75644320812458 8.40681772519677
-1.12422479455467 -0.607728519914606
1.88394526347206 9.89404180276951
1.34119140011109 4.71394366242212
-0.866659018725533 0.13342644160029
1.42751332053673 5.33351738479513
-0.911531703377452 0.0316959680587874
0.48508057268549 1.44718587288087
-1.64027078153502 -3.46066925693651
-1.25139406637923 -1.11242585131484
-0.248185975672745 0.897321916366227
-1.28391033792466 -1.28402663811459
1.02742077214036 2.95866689031474
-1.00469821661397 -0.200025010181867
-1.4839169486916 -2.36476061584865
-0.388063615699482 0.791749832117417
-0.351330238210304 0.829615571305214
-0.919162594924984 0.0179998135616073
-1.8300296068598 -5.09523489832429
0.261126692492549 1.16381348789023
-0.434250311374255 0.760268745883126
1.71127701069478 7.91589387950685
-1.74396833588139 -4.29086220370818
1.99084425319697 11.3249840251345
1.96740657851932 11.0128349752937
0.00519723053157728 1.00314128653432
-1.0300674545336 -0.292760531150931
-0.679871421735784 0.497273222659833
0.346398392238981 1.26222464776238
1.86986837689732 9.72720250684752
0.681415810398582 1.82471633810466
-1.77451887462146 -4.5905765817216
0.407108934350533 1.33052102469306
0.234319557119511 1.1441133039476
0.680255483899439 1.82443646987879
-1.46806636854791 -2.24633672359298
1.37432472355754 4.94560301954173
-0.0444665624422655 0.967684314859427
0.946510929382118 2.62683889570353
-0.348613222191393 0.840643230744185
-1.16085167224828 -0.752985307128517
1.52131471522301 6.10602184516401
-0.651129722285528 0.522096171976065
-0.974782964327846 -0.130492050354493
1.38933505926152 5.06701358419365
-0.405141209582377 0.788636045619015
-1.07427216939099 -0.439348576157792
0.287925799214639 1.20249444515467
1.99219825484057 11.3300781732418
-1.97984787642545 -6.67440166993524
-1.16710498101462 -0.766106919849088
0.0880439799314834 1.04798139144606
1.65176684501771 7.31505205040777
-0.417609509438676 0.781733792741341
-0.383443127401129 0.804453481067922
1.01000388646387 2.88434280725849
-1.32293438896487 -1.47265021133377
-0.99370936187379 -0.173512541750945
-0.404147661596262 0.789443384056748
-0.105702670056522 0.956508392161341
0.219292884334177 1.12972481089977
0.693092568668425 1.84314500579793
-1.49873438446745 -2.42808076757983
-1.71917638967302 -4.11445980893961
1.92275290133258 10.4060905577529
0.557888885180773 1.55207039096411
0.52577897427859 1.50941467333881
1.19379941638111 3.80067611582445
-1.77051760303498 -4.55469830292259
-1.45032825014804 -2.14115707845768
-1.23795852109155 -1.04877741172978
-1.27410510105383 -1.22052728026993
-0.468704783116342 0.748660167008506
-0.886149958616648 0.108477253199736
-0.545832084491524 0.677538442559796
-0.313696873254803 0.868101603557104
1.87479353239457 9.8019035285421
-1.96168063254752 -6.46823505732823
1.96440026679769 10.9681385483752
-0.662848989822634 0.526116383312213
1.20039942205076 3.83742650014747
-1.54721858629723 -2.79417593391984
-1.63364823827773 -3.39429408125394
1.57659720521503 6.58152772647782
1.21200313281051 3.90426995352027
-1.83574674575522 -5.16861220219619
0.566320040889504 1.57474620086365
0.972707935299989 2.74195158574218
-0.134613050994941 0.938027907445401
-1.49592388715788 -2.40495280111003
0.51808299716567 1.48451456005604
0.327121432995443 1.24307092490041
-0.0117717724824216 0.98816938612114
0.746398682964023 1.9860250421071
-0.160139522955972 0.926953715518133
-0.426027531146679 0.77686273662751
-1.69477632375206 -3.89258825943529
-1.05936959994091 -0.374798803510702
-1.1240287678796 -0.591747658886704
-0.512397011144832 0.697801414844525
0.1151272306282 1.06359410266455
-0.757649192903382 0.365441869230087
1.62827887860161 7.072337004527
-0.956679625097431 -0.0770897233418417
1.48797995432694 5.81793273760318
-1.21173687511944 -0.942758423942144
0.760812354488629 2.0248690601893
1.98588523589092 11.2553963644691
-1.23865491219261 -1.06419656255286
-1.16524180671887 -0.754310282516707
-1.74637746992985 -4.33715527923672
-1.77321833494424 -4.56462033849922
-1.33398224041622 -1.52005657626804
-0.326070805576438 0.851593511677802
0.256559935709201 1.16176848434099
-0.542645927277806 0.672885495947295
-0.971484127948466 -0.110229010415843
-0.230176814593698 0.901620005359611
-0.646934580672173 0.533070758814424
0.611919937847676 1.66399803907441
1.39915511741084 5.1443280764869
-0.313547573274856 0.864464716742631
-0.524558158878761 0.689965272099956
-1.76473060004123 -4.48645796483076
1.58150824165447 6.61680769943754
1.88474461581527 9.94195904014881
-0.403718263151636 0.805022294489218
0.549520352891367 1.52742518740342
-0.535840524475981 0.676708758476658
0.292676339923188 1.2041758733986
-0.066385069686731 0.949511916891023
1.36977344495431 4.908901766673
0.80996594457818 2.15548891267912
1.99370354145932 11.3561875597128
1.18161107639942 3.71275108063636
0.754047140010336 2.00558932006914
-1.88210308314506 -5.62783733524903
0.732503302683575 1.94171425844159
-0.105199099402469 0.946278688183064
0.208559701528764 1.12431237278009
-1.02994191709767 -0.28548779700985
-0.478799425230845 0.729109789391446
0.364156171385083 1.27495206463488
0.237431538421525 1.14882985176487
1.56652522547804 6.4928423735079
1.11753751981401 3.38651113266133
0.559766732161253 1.56331223476094
-1.05253303063569 -0.3624041255737
-1.13665085908765 -0.623529463788013
1.33676157043896 4.68030204116338
-1.25855088245022 -1.15405200948732
-0.728781804869114 0.418785600098235
0.332252831390097 1.24183419559873
1.75266336276354 8.34030576346291
-0.40671405183982 0.79036816665767
0.662488101353502 1.77281326360528
-0.527760172621683 0.689230694951209
-1.98480003264102 -6.71804058404472
-1.22220688253495 -0.988741786050339
-1.32471942362105 -1.46071557583938
-1.17864005024045 -0.811013937624279
1.52887589294808 6.17035021378948
0.414821241537061 1.32832430744308
1.35356563000651 4.80017499293894
-0.449204147693734 0.750592488597213
-0.00290057132472299 0.984750224695719
0.324368361284517 1.22909634265887
-0.521848177920158 0.692713691182687
1.1533626147704 3.57618636017796
-1.10677675050792 -0.56455294394963
-0.318060382757006 0.846300804142516
1.87730983895634 9.82142181400893
-1.96746098497336 -6.54993728917854
-1.93965862086028 -6.25476790416755
-0.459787076254924 0.752114786759366
-1.31348172702031 -1.43242130060961
-1.14514328231759 -0.673155888416164
0.682021135275135 1.81691810531443
0.281605217003446 1.18687597092314
-0.918989576387054 0.0173575925587022
1.86009333852669 9.5959629724972
0.422851653026241 1.35683263155697
-1.18753310918885 -0.857527370097498
-1.05991208496293 -0.386661829729422
0.0788137458913938 1.02237543048186
1.81703979429071 9.09427555649752
-1.77523326667127 -4.60471049943597
-0.134965858344948 0.939760171760071
-0.774427290685961 0.330664922257729
1.84912287266476 9.48631256447653
-0.0473840523157882 0.975596591749443
-0.310445770973367 0.855826804192119
-1.85040726666538 -5.28123090480442
0.0195484778286565 1.00949379122544
-0.32148485732909 0.85171566029906
1.00341624502118 2.85243197422087
-0.0413769862398357 0.968545930426121
-1.50428660641862 -2.4734882991516
-0.315956120587869 0.850716428317477
0.379662199511765 1.28561243609028
1.26056841548983 4.17614570857556
0.117109705720536 1.0719757158902
-0.72026478263644 0.426526965944616
1.76009071973641 8.44462688866445
};
\addlegendentry{new\_neural\_network\_estimate(x)}
\end{axis}
\end{tikzpicture}}
	 \caption{Visueller Vergleich der Funktion $m_1$ mit ihrer Schätzung.}
        \label{fig:subfig810}
    \end{subfigure}
    \caption{Darstellung der Funktion $m_1$ und dessen Schätzung durch den Schätzer \textit{new\_neural\_network\_estimate}.} 
\label{fig:subfig1.a.4}
\end{figure}

%\begin{figure}
%\centering
%%\input{tikzfigure1.tikz}
%\begin{tikzpicture}
%\begin{axis}[x dir=reverse]
%\addplot3 [scatter, only marks]
%  table[x=x, y=y, z=f, col sep=comma] {plotpostpro_sort.csv};
%\end{axis}
%\end{tikzpicture}
%\end{figure}

Wie wir in Tabelle~\ref{tab:truthTablesm1} und Tabelle~\ref{tab:truthTablesm2} anhand des skalierten $L_2$-Fehlers sehen können, übertrifft unserer Neuronale-Netze-Regressionsschätzer in allen Fällen die Leistung der anderen Schätzer. 
\begin{table}
\centering
\begin{tabular}{ |p{5cm}||p{4cm}|p{4cm}|}
 \hline
 & \multicolumn{2}{|c|}{$m_1$}\\
 \hline
 \textit{noise}& $5\%$ & $10\%$ \\
 \hline
 $\bar{\epsilon}_{L_2,N}(avg)$& $13.4482362$ & $13.3925910$ \\
 \hline
 \textit{approach}&  Median (IQR) &  Median (IQR)   \\
 \hline
 new\_neural\_network\_estimate & $\mathbf{2.482\textbf{e-}05 (1.612\textbf{e-}05)}$   & $\mathbf{6.908\textbf{e-}05 (3.936\textbf{e-}05)}$  \\
 fc\_neural\_1\_estimate & $4.384\text{e-}04 (2.14\text{e-}03)$ &   $7.261\text{e-}04 (4.57\text{e-}03)$ \\
 nearest\_neighbor\_estimate & $2.9527\text{e-}04 (9.312\text{e-}05)$ & $9.0864\text{e-}04 (2.8952\text{e-}04)$\\
 \hline
\end{tabular}
    \caption{Median und IQR des skalierten empirischen $L_2$-Fehlers für Schätzungen für $m_1$ mit Stichprobengröße $N = 1000$.}
     \label{tab:truthTablesm1}   
\end{table}
    
    \begin{table}
\centering
\begin{tabular}{ |p{5cm}||p{4cm}|p{4cm}|}
 \hline
 & \multicolumn{2}{|c|}{$m_2$}\\
 \hline
 \textit{noise}& $5\%$ & $10\%$ \\
 \hline
 $\bar{\epsilon}_{L_2,N}(avg)$& $0.0324$ & $0.0311$ \\
 \hline
 \textit{approach}&  Median (IQR) &  Median (IQR)   \\
 \hline
 new\_neural\_network\_estimate & $\mathbf{0.003961 (0.000932)}$   & $\mathbf{0.00431 (0.000973)}$  \\
 fc\_neural\_1\_estimate & $0.0257 (0.3803)$ &   $0.0559 (0.52033)$ \\
 nearest\_neighbor\_estimate & $0.01616 (0.005906)$ &$0.01763 (0.007081)$\\
 \hline
\end{tabular}
    \caption{Median und IQR des skalierten empirischen $L_2$-Fehlers für Schätzungen $m_2$ mit Stichprobengröße $N = 1000$.}
    \label{tab:truthTablesm2}   
\end{table}

%\include{chap5}
%\include{chap6}
%\chapter*{Appendix}
\markboth{APPENDIX}{}
\addcontentsline{toc}{chapter}{Appendix}
\label{chap:app}

Der Programmcode ist wie folgt aufgebaut:
%\linespread{1.0}
\begin{itemize}
    \item[-] \texttt{main.py} ist das Hauptprogramm welches alle Schätzer aufruft und die Ouputs generiert.
    \item[-] \texttt{data\_gen.py} generiert die Daten die wir für unsere Simulation benötigen.
    \item[-] \texttt{help\_neural\_networks.py} fasst alle Hilfsfunktionen zusammen.
    \item[-] \texttt{new\_neural\_network.py} enthält unseren Neuronale-Netze-Regressionsschätzer.
    \item[-] \texttt{fc\_neural\_network.py} enthält das neuronale Netz mit einer verborgenen Schicht.
    \item[-] \texttt{nearest\_neighbor.py} enthält einen Nächste-Nachbar Schätzer.
    \item[-] \texttt{constant.py} enthält den konstanten Schätzer.
\end{itemize}
\renewcommand\lstlistingname{Listing}
\lstinputlisting[basicstyle=\scriptsize, caption={\texttt{main.py}},captionpos=t]{../01_Code/main.py}
\lstinputlisting[basicstyle=\scriptsize, caption={\texttt{data\_gen.py}},captionpos=t]{../01_Code/data_gen.py}
\lstinputlisting[basicstyle=\scriptsize, caption={\texttt{help\_neural\_networks.py}},captionpos=t]{../01_Code/help_neural_networks.py}
\lstinputlisting[basicstyle=\scriptsize, caption={\texttt{new\_neural\_network.py}},captionpos=t]{../01_Code/new_neural_network.py}
\lstinputlisting[basicstyle=\scriptsize, caption={\texttt{fc\_neural\_network.py}},captionpos=t]{../01_Code/fc_neural_network.py}
\lstinputlisting[basicstyle=\scriptsize, caption={\texttt{nearest\_neighbor.py}},captionpos=t]{../01_Code/nearest_neighbor.py}
\lstinputlisting[basicstyle=\scriptsize, caption={\texttt{constant.py}},captionpos=t]{../01_Code/constant.py}

% Literaturverzeichnis (beginnt auf einer ungeraden Seite)
%\cleardoublepage
\nocite{*}
\bibliographystyle{alpha}
%\bibliographystyle{acm}
\bibliography{Masterthesis_Adrian_Gabel}
\addcontentsline{toc}{chapter}{Literaturverzeichnis}

\end{document}
